% =====================================================================
%  THE DISTILLED PROOF  --  phase 1.9 of the augusst knot campaign.
%  A main theorem and an internal theorem, three declared inputs, the shortest complete proof.
%  Lean-ready; no Lean.
% =====================================================================
\documentclass[11pt]{article}
\usepackage[margin=1in]{geometry}
\usepackage{amsmath,amssymb,amsthm}
\usepackage{xcolor}
\usepackage[colorlinks,linkcolor=blue!50!black,citecolor=blue!50!black]{hyperref}

\newtheorem{theorem}{Theorem}[section]
\newtheorem{proposition}[theorem]{Proposition}
\newtheorem{lemma}[theorem]{Lemma}
\newtheorem{corollary}[theorem]{Corollary}
\theoremstyle{definition}
\newtheorem{definition}[theorem]{Definition}
\newtheorem{convention}[theorem]{Convention}
\newtheorem{axiom}{Axiom}
\theoremstyle{remark}
\newtheorem{remark}[theorem]{Remark}

\newcommand{\Ind}{\operatorname{Ind}}
\newcommand{\rot}{\operatorname{rot}}
\newcommand{\wind}{\operatorname{wind}}
\newcommand{\sgn}{\operatorname{sgn}}
\newcommand{\Cat}{\operatorname{Cat}}

\title{The polygon coefficient $X_1$ satisfies the wall-crossing laws:\\
a distilled proof}
\author{augusst knot campaign, phase 1.9}
\date{}

\begin{document}
\maketitle

\begin{abstract}
This document proves the main theorem that the state sum $X_1$ of
Definition~\ref{def:X1} agrees with the lawful quantity $A$ on every generic
polygon, and therefore satisfies the wall-crossing laws (S3)--(S8) on the
domains where the specification states them. It assumes exactly three
things \emph{of the campaign}, all listed in the first group of the axiom list
of Section~\ref{sec:axioms}: the law (S2) for $X_1$; the laws asserted of the
campaign's declared constant $A$ --- declared, not asserted to exist, the
declaration being what supplies it; and the project owner's normalization of
$A$ on one explicit polygon per rotation class. It also proves internally the
cyclic-classification Theorem~\ref{thm:mycyclic} used in the transport step.
It quotes five theorems from the literature, listed in the same section with
their sources and with the reason each is not proved here; an earlier abstract
said ``exactly three'' without that distinction, which read as a claim to
assume nothing else.
Every other definition it uses is stated here, and every other statement it
uses is proved here. Where the text names an evidence pack, a ledger edge or
the phase-1.5 manuscript, it does so as source history only: no conclusion
depends on any of them, every argument so credited being printed here in
full. It is written so that a formalizer can proceed line by line with no
mathematics left to invent, and the eight axiom entries are the interface
such a formalization would have to import.
\end{abstract}

\tableofcontents

% Section 1 of the distilled proof: everything the readout is made of.
% Self-containment rule: no phrase of the form ``as in the specification''
% survives here without the specification's own text printed beside it.
\section{Setup}\label{sec:setup}

\subsection{Polygons, guarded predicates, genericity}

\begin{definition}[polygon]\label{def:polygon}
Let $n\geq3$. A \emph{polygon} on $n$ vertices is a tuple
$P=(p_1,\dots,p_n)\in(\mathbb{R}^2)^n$ with $p_{i+1}\neq p_i$ for every $i$,
indices read cyclically modulo $n$ throughout. Its \emph{edges} are the closed
segments $e_i=[p_i,p_{i+1}]$ and its \emph{edge directions} are
$d_i=p_{i+1}-p_i\in\mathbb{R}^2\setminus\{0\}$. Two edges are \emph{consecutive}
if their index sets meet, and \emph{remote} otherwise. An edge is \emph{remote
to a vertex} $p_M$ when it is remote to both edges incident to $M$ --- the
stronger of the two possible readings, and the one every consumer of the phrase
uses: it excludes an edge that merely avoids $p_M$ while sharing a vertex with
one of its edges. We write
$\det(u,v)=u_1v_2-u_2v_1$.
\end{definition}

\begin{definition}[principal turns and the regular locus]\label{def:regular}
\begin{enumerate}
\item[(A)] Let $L$ be a closed polygon with corners $q_0,\dots,q_{c-1}$ and edge directions
$\delta_i=q_{i+1}-q_i$. The \emph{principal turn} at $q_i$ is the unique
$\tau_i\in(-\pi,\pi)$ with
$\cos\tau_i=\frac{\langle\delta_{i-1},\delta_i\rangle}{|\delta_{i-1}||\delta_i|}$
and $\sgn\tau_i=\sgn\det(\delta_{i-1},\delta_i)$. It exists precisely when
$\delta_i$ is not a negative multiple of $\delta_{i-1}$: if
$\det(\delta_{i-1},\delta_i)\neq0$ it is nonzero, and if $\delta_i$ is a
positive multiple of $\delta_{i-1}$ it is $0$. The excluded case
$\delta_i\in\mathbb R_{<0}\delta_{i-1}$ is the \emph{kink}, where the two
candidate values $\pm\pi$ are both excluded from the open interval.
\item[(B)] The \emph{regular locus} $\mathcal R_n\subseteq(\mathbb R^2)^n$ is the set of
$P=(p_1,\dots,p_n)$ with $d_i\neq0$ for all $i$ and $d_{i+1}\neq-\lambda d_i$
for all $i$ and all $\lambda>0$: nonzero edges and no consecutive pair that
doubles back. Equivalently, every principal turn of
clause~(A) exists.
\end{enumerate}
\end{definition}

\begin{definition}[the guarded list $\mathcal G$]\label{def:guarded}
Fix $n$. The \emph{guarded list} is a finite family of real polynomial functions
on $(\mathbb R^2)^n$, indexed once and for all by combinatorial data that does
not depend on the configuration. Write $d_i=p_{i+1}-p_i$ and, for an edge index
$e$, write $\ell_e(x)=\det(d_e,\,x-p_e)$ for the affine form vanishing on the
line of $e$, with coefficient row
$\bigl(A_e,B_e,C_e\bigr)=\bigl(-d_{e,2},\;d_{e,1},\;d_{e,2}p_{e,1}-d_{e,1}p_{e,2}\bigr)$,
so that $\ell_e(x)=A_ex_1+B_ex_2+C_e$.

\emph{Unconditional members.}
\begin{enumerate}
\item[(G1)] $\mathrm{G1}_i=\det(d_{i-1},d_i)$, for every $i\in\{1,\dots,n\}$.
\item[(G2)] $\mathrm{G2}_{e,i}=\ell_e(p_i)=\det(d_e,\,p_i-p_e)$, for every edge
index $e$ and every $i\notin\{e,e+1\}$.
\end{enumerate}

\emph{Activation.} Two remote edges $e,f$ are \emph{defined} to cross when
\[
   \mathrm{G2}_{e,f}\,\mathrm{G2}_{e,f+1}<0
   \quad\text{and}\quad
   \mathrm{G2}_{f,e}\,\mathrm{G2}_{f,e+1}<0 ,
\]
a condition on unconditional members alone. The products are \emph{strict}, and
that is a decision, not a formality: an earlier revision wrote
$\sgn\neq\sgn$, which is undecided exactly where it is consumed. At a wall
point one of the four members vanishes, and the three available readings
disagree there. The strict product makes the pair \emph{inactive} at such a
point; a three-valued $\sgn\in\{-1,0,+1\}$ would make it active, and then the
sliding bundle of the wall dictionary of
Section~\ref{sec:dictionary} would be $\{\mathrm{G2},\mathrm{G4}\}$ rather than
the singleton that dictionary's item~(iii) records;
and the shipped evaluator's test treats a zero as negative, which is asymmetric
between the two edges. The strict product is taken, and the divergence from the
evaluator is recorded here rather than left silent. What can be said about its
consequences is scoped to this document: every statement below evaluates $X_1$
at \emph{generic} polygons, where no $\mathrm{G2}$ vanishes and the two
readings agree, so no statement below depends on which reading the code takes.
What an external caller passes to that code is outside this document's claims
--- a caller that hands it a collinear configuration reaches the divergence, and
nothing here says otherwise. An earlier revision wrote that the code "never
reaches" such a configuration, which is a claim about the code's callers and not
about this document. It is the sign
condition, and not the geometric phrase ``their relative interiors meet'', that
every consumer below reads and that the shipped evaluator computes.

The two agree wherever the four members are nonzero, and there they say that
each edge has its two endpoints strictly on opposite sides of the other's line,
which for segments is exactly a transverse interior crossing. They do
\emph{not} agree everywhere in $\mathcal R_n$, and an earlier revision wrote
``exactly when'' as though they did: two collinear overlapping remote edges have
meeting relative interiors while all four members vanish, so the sign condition
fails. Such a configuration is not generic --- those vanishing members are
unconditional, hence relevant --- so the generic locus defined below and the
discriminant $\mathcal D$ of Section~\ref{sec:dictionary} are the same sets under
either reading, and no statement
below changes. What changes is that activation is now decided by a formula.

\emph{Conditional members.}
\begin{enumerate}
\item[(G5)] $\mathrm{G5}_{e,f}=\det(d_e,d_f)$, for every pair of remote edge
indices with $1\leq e<f\leq n$ as integers --- the comparison is between the
integer representatives $1,\dots,n$ of the edge indices and not a cyclic
comparison, there being no such thing; adjacency and remoteness continue to be
read modulo $n$. This fixes one \emph{ordered representative} per pair, and an
unordered pair will not do, because
$\det(d_f,d_e)=-\det(d_e,d_f)$ and the family is used where signs are read, in
the sign-change condition that defines a transversal event below; \emph{active} when $e$
and $f$ cross. Making it
conditional rather than unconditional is a fidelity point: the specification
guards the turn at a smoothing site, and a smoothing site exists only where two
edges cross. Demanding $\det(d_e,d_f)\neq0$ for a pair that does not cross
would shrink the generic locus below the specification's and exclude polygons
the laws are stated about.
\item[(G3)] $\mathrm{G3}_{e,f,g}=\det\begin{pmatrix}A_e&B_e&C_e\\
A_f&B_f&C_f\\ A_g&B_g&C_g\end{pmatrix}$, the rows being the line-coefficient
rows bound at the head of this definition, one for each of $e$, $f$ and $g$;
for every triple of pairwise remote edges $1\leq e<f<g\leq n$, ordered by their
integer representatives; \emph{active} when $e,f,g$ pairwise cross. It vanishes exactly when the
three lines share a point of the projective plane: either they are concurrent in
the plane, or all three are parallel. A parallel \emph{pair} with a transverse
third does \emph{not} make it vanish --- an earlier revision said it did, and the
configuration $p_0=(0,0)$, $p_1=(1,0)$, $p_2=(0,1)$, $p_3=(1,1)$, $p_4=(2,-1)$,
$p_5=(2,2)$ refutes it, the rows of $e_0,e_2,e_4$ being $(0,1,0)$, $(0,1,-1)$
and $(-3,0,6)$ with determinant $3$. On the active domain the three edges
pairwise cross, so no two of the lines are parallel and the vanishing is exactly
a concurrence at a finite point, which is what every consumer reads.
\item[(G4)] $\mathrm{G4}_{e;f,g}
=\det(d_f,\,p_f-p_e)\,\det(d_g,d_e)-\det(d_g,\,p_g-p_e)\,\det(d_f,d_e)$,
for every edge $e$ and every pair of edges remote to $e$ with
$1\leq f<g\leq n$ as integers --- again the comparison is between integer
representatives, and again one ordered representative per pair, since
$\mathrm{G4}_{e;g,f}=-\mathrm{G4}_{e;f,g}$; \emph{active} when $f$ and $g$ both
cross $e$.

\emph{Normalization of cyclic aliases.} Edge indices are read cyclically, so
one edge has many names: $e_0$, $e_n$ and $e_{2n}$ are the same edge. Every
comparison of indices in this document --- the $f<g$ of (G5), the $e<f<g$ of
(G3), the $f<g$ of (G4) --- is a comparison of \emph{representatives}: each
index is first reduced to the unique integer in $\{1,\dots,n\}$ naming that
edge, and the reduced integers are compared. Adjacency, remoteness and the
edges themselves are unaffected by the reduction; only the choice of ordered
representative depends on it, and only that choice carries a sign.

\emph{The two accessors.} Consumers meet a pair of edges in the order the
geometry gives, which after reduction may be the reverse of the representative
order. Let $f,g$ be distinct edges remote to $e$, presented in some order and
possibly under aliases, and let $\bar f,\bar g\in\{1,\dots,n\}$ be their
representatives. Define
\[
   \mathrm{G4}\langle e;f,g\rangle
   =\mathrm{G4}_{e;\min(\bar f,\bar g),\,\max(\bar f,\bar g)}\in\mathcal G,
   \qquad
   \epsilon(f,g)=\begin{cases}+1,&\bar f<\bar g,\\ -1,&\bar f>\bar g,\end{cases}
\]
the first a \emph{member-valued} accessor --- it returns an element of the
guarded list, the one indexed by the ordered representative pair --- and the
second an \emph{orientation sign} in $\{\pm1\}$, a constant of the index data
and not a function of the configuration. Their product
\[
   \widehat{\mathrm{G4}}_{e;f,g}
   =\epsilon(f,g)\cdot\mathrm{G4}\langle e;f,g\rangle
\]
is the \emph{oriented value}: a real polynomial function on $(\mathbb R^2)^n$,
antisymmetric in the ordered pair, equal to the member up to the orientation
sign.

The two are used in different places and the distinction is not cosmetic.
Every set of guarded predicates formed below --- in particular the zero set of
an event, defined later in this section --- is a set of \emph{members} of
$\mathcal G$, and membership there is by \emph{index}: the element is
$\mathrm{G4}\langle e;f,g\rangle$, the member indexed by the ordered
representative pair, and not the oriented value, which when $\epsilon=-1$ is
that member's negative and is indexed by nothing. The distinction is one of
indexing and not of polynomials, and the difference matters: as polynomials a
negated member can coincide with another member, since for three pairwise
remote edges $\mathrm{G4}_{e;f,g}=-\mathrm{G3}_{e,f,g}$ --- an identity of
polynomials, proved in the dictionary section where the canonical factors are
analysed --- so $-\mathrm{G4}_{e;f,g}$ is the polynomial $\mathrm{G3}_{e,f,g}$,
which \emph{is} a member. An earlier revision said the
oriented value "is not a member of $\mathcal G$ when $\epsilon=-1$", which as a
statement about polynomials is false.

Statements about vanishing, activation and relevance are statements about the
member and are insensitive to the presentation order. Statements about
\emph{sign} --- a sign change across an event, a factorization whose two sides
must agree in sign, the order of two crossings along an edge --- read the
oriented value, and there the orientation sign $\epsilon$ is carried
explicitly. Both kinds occur below, and the next display is of the first kind.

Two illustrations of the reduction, both used below. At $j=0$ the pair
$(e_{j-1},e_j)$ is $(e_{n-1},e_n)$, whose representatives are already
increasing: $\epsilon=+1$, and the oriented value is the member. At $j=1$ it is
$(e_0,e_1)$, whose representatives are $(n,1)$: $\epsilon=-1$, and the member
is $\mathrm{G4}_{e;1,n}$. An earlier revision wrote
$\mathrm{G4}_{h;e_{j-1},e_j}$ where the reduction makes that name empty and the
sign opposite, and a later one placed the wrap at $j=0$, where there is none. When active, both
$\det(d_f,d_e)$ and $\det(d_g,d_e)$ are nonzero, the crossing of $e$ with $f$
sits at the parameter $t_f=\det(d_f,p_f-p_e)/\det(d_f,d_e)$ along $e$, and, for
the ordered representative pair $\bar f<\bar g$ that indexes the member,
\[
   \mathrm{G4}_{e;\bar f,\bar g}
      =(t_{\bar f}-t_{\bar g})\det(d_{\bar f},d_e)\det(d_{\bar g},d_e),
\]
so its vanishing is exactly a tie $t_{\bar f}=t_{\bar g}$ in the order of
crossings along $e$. Presented in the other order the same identity holds for the oriented value,
both sides changing sign. Two kinds of consumer read this display, and they read
different things. Those that need only the \emph{vanishing} --- the two collision arguments
below, in the chamber-invariance proposition and in the silence lemma --- read
the member, and the presentation order is irrelevant to them. Those that
need the \emph{sign} --- the order of the two crossings along $e$, and the
bigon-bundle computation of Section~\ref{sec:dictionary} --- read the oriented
value $\widehat{\mathrm{G4}}_{e;f,g}=\epsilon(f,g)\,\mathrm{G4}\langle
e;f,g\rangle$ and carry $\epsilon$. An earlier revision said the vanishing was
"the only reading the consumers of this display take", which the sign-reading
sites contradict.
\end{enumerate}
\end{definition}

\begin{definition}[relevant members, generic polygons and chambers]\label{def:generic}
\begin{enumerate}
\item[(A)] A member of $\mathcal G$ is \emph{relevant} at a polygon $P$ if it is
unconditional, or if it is conditional and active at $P$
(Definition~\ref{def:guarded}). A polygon $P$
(Definition~\ref{def:polygon}) is \emph{generic} if every member relevant at
$P$ is nonzero at $P$; equivalently, if every unconditional member of
$\mathcal G$ is nonzero at $P$ and every conditional member that is active at
$P$ is nonzero at $P$. We write $\mathcal U_n$ for the set of generic polygons
on $n$ vertices.

\item[(B)] A \emph{chamber} is a connected component of the set $\mathcal U_n$ of generic
polygons on $n$ vertices (clause~(A)). It is \emph{not} in
general a component of the set where every member of $\mathcal G$ is nonzero:
a conditional member that is inactive may vanish, or vanish identically, at a
perfectly generic polygon, and nothing in this document asks it not to.
\end{enumerate}
\end{definition}

\begin{remark}[what genericity excludes, item by item]\label{rem:genericmeans}
$\mathrm{G1}_i\neq0$ says $d_{i-1}$ and $d_i$ are linearly independent, which
excludes \emph{both} degenerate turns at once: the \emph{flat} turn, where
$d_i$ is a positive multiple of $d_{i-1}$ and $p_i$ lies strictly between its
neighbours, and the \emph{kink}, where $d_i$ is a negative multiple of $d_{i-1}$
and the polygon reverses at $p_i$. $\mathrm{G2}_{e,i}\neq0$ says no vertex lies
on the line of a non-incident edge, hence none lies on the edge.
$\mathrm{G3}\neq0$ says no three pairwise remote edges are concurrent, so no two
crossings collide. $\mathrm{G4}\neq0$ says the crossings along each edge are
totally ordered with no ties. $\mathrm{G5}\neq0$ says two crossing edges are
never parallel, so every crossing is a transversal double point. Consequently a
generic polygon meets each of its remote edges in finitely many transversal
double points, has no triple points, and has no vertex on a non-incident edge:
it is an A'Campo divide.
\end{remark}

\begin{proposition}[active (G5) and the frozen generic locus]\label{prop:fidelity}
Fix $n\geq3$. Let $\mathcal U_n^{\flat}$ be the set of polygons at which every member of the
four families (G1)--(G4) that is relevant is nonzero --- the specification's
guarded list, which does not mention $\det(d_e,d_f)$.
\begin{enumerate}
\item[(i)] Let $P$ be a polygon on $n$ vertices with $\mathrm{G1}_i(P)\neq0$ for every $i$ and
$\mathrm{G2}_{e,i}(P)\neq0$ for every edge $e$ and every valid
$i\notin\{e,e+1\}$. Then every (G5) member active at $P$ is nonzero,
equivalently every remote pair whose relative interiors meet has nonzero direction determinant.
\item[(ii)] $\mathcal U_n=\mathcal U_n^{\flat}$.
\end{enumerate}
\end{proposition}
\begin{proof}
(i) Let $e,f$ be remote with a common interior point and suppose
$\det(d_e,d_f)=0$. Then $d_e$ and $d_f$ are proportional and the two segments,
having a common point, lie on one line and overlap in a segment of positive
length or in a point. If they overlap in a segment, an endpoint of one lies in
the other, so the corresponding $\mathrm{G2}$ vanishes, contrary to hypothesis.
If they meet in a single point, that point is interior to both by assumption,
and two collinear segments meeting in exactly one interior point of each is
impossible. So $\det(d_e,d_f)\neq0$.

(ii) $\mathcal U_n\subseteq\mathcal U_n^{\flat}$ because $\mathcal U_n$ asks for more.
Conversely let $P\in\mathcal U_n^{\flat}$; every relevant (G1), (G2), (G3) and
(G4) is nonzero there. Since (G1) and (G2) are unconditional, clause (i) gives
$\mathrm{G5}_{e,f}\neq0$ for every active pair. An inactive (G5) is not
relevant, so every member relevant at $P$ is nonzero and $P\in\mathcal U_n$.
\end{proof}

\begin{remark}[why (G5) is conditional, and what it would cost to change]
\label{rem:fidelity}
Proposition~\ref{prop:fidelity}(ii) is the fidelity statement, and it is stated
rather than remarked because it is load-bearing: it says no theorem below is
proved on a narrower domain than the specification states it for.
Definition~\ref{def:guarded} adds (G5) as a \emph{conditional} member so that
the turn at a smoothing site has a named predicate, and the proposition says
that this adds nothing.

Making (G5) unconditional would \emph{not} be harmless: the pinned bow-tie $K_0$
of Definition~\ref{def:circularstar} has two parallel remote edges that do not
meet, so it would cease to be generic and the rank-zero anchor of
Theorem~\ref{thm:main}(B) would not exist. The same bow-tie later refuted an
argument that ran general position on the unconditional zero set
(Lemma~\ref{lem:relgp}(B)); it is the configuration this whole distinction exists
to keep.
\end{remark}


\subsection{The Gauss word, supports and carriers}

Fix a \emph{diagrammatic} polygon $P$ --- the class defined below, of which
every generic polygon is a member --- and let $m=m(P)$ be the number of its
double points, labelled $1,\dots,m$. Everything in this subsection, up to and
including the residual pieces and their diagrams, is stated for that class:
genericity enters later, with the guards, and an earlier revision fixed a
generic polygon here and then claimed the wider scope downstream. Traversing $P$ once from $p_1$ in the direction of
increasing index visits each double point exactly twice; recording the labels in
the order visited gives a cyclic word of length $2m$ in which every letter occurs
twice, the \emph{Gauss word} of $P$.

\begin{definition}[diagrammatic polygon]\label{def:diagrammatic}
A polygon $P$ is \emph{diagrammatic} if its self-intersections are finitely
many, each of them an isolated transverse crossing of two edges with exactly two
preimages on the traversal circle, no two of them sharing an image or a
preimage, none of them at a corner, and no vertex of $P$ lying on an edge not
incident to it. The clause about preimages and shared images is not implied by
the others, and an earlier revision omitted it: the polygon
$((-3,0),(3,0),(0,-3),(0,3),(-2,-2),(2,2))$ has its edges $e_0,e_2,e_4$ meeting
pairwise transversally --- the three determinants are $36$, $24$, $-24$ --- at
the \emph{single} point $(0,0)$, so three strands pass through one point, and the
$2m$-letter Gauss word does not exist.

Every generic polygon is diagrammatic: the (G2) members forbid a vertex on a
remote edge, the (G5) members make each crossing transversal, and a triple point
would make a (G3) member vanish while all three pairs cross, so that member is
active, hence relevant, hence nonzero. So is the polygon at a simple positive
flat wall, whose only degeneracy is a vanishing turn with the middle vertex
\emph{inside} the segment. That last qualification matters: at an exterior
collinear vertex the two incident edges overlap and the far one carries the
middle vertex, as in $((0,0),(2,0),(1,0),(0,1))$, which is not diagrammatic ---
an earlier revision said "any polygon whose only degeneracy is a vanishing
turn".

The class is named because the Gauss word, the interlacement graph, the
independent sets, the undominated sets, the residual pieces and their diagrams
are read off a diagrammatic polygon and need nothing more; genericity is what
$X_1$ needs, through the guards, and is a strictly stronger condition. Three
sections use exactly this class: the flat law at its wall, the rounding lemma,
and the carrier floor.
\end{definition}

\begin{definition}[interlacement graph]\label{def:interlace}
The \emph{interlacement graph} $G_P$ has vertex set $[m]=\{1,\dots,m\}$, with
$c\sim c'$ if and only if exactly one of the two occurrences of $c'$ lies between
the two occurrences of $c$ in the Gauss word. (The relation is symmetric.) We
write $\Ind(G_P)$ for the set of independent sets of $G_P$, including
$\varnothing$, and $N_{G_P}(S)$ for the set of vertices adjacent to some element
of $S$.
\end{definition}

\begin{definition}[oriented smoothing]\label{def:smoothing}
For $S\in\Ind(G_P)$, the \emph{oriented smoothing of $P$ along $S$} replaces, at
each double point of $S$, the two transversally crossing arcs by the two arcs
that respect the orientation of $P$ and do not cross. The result is a disjoint
union of closed oriented curves, called the \emph{carriers} of $S$.
\end{definition}

\begin{lemma}[the carriers of a support]\label{lem:carriers}
Let $P$ be diagrammatic (Definition~\ref{def:diagrammatic}) and
$S\in\Ind(G_P)$ --- genericity is not needed here and an earlier revision
assumed it; what the argument reads is the Gauss word. Let $\Gamma$ be the
\emph{traversal circle}: an abstract oriented circle
together with the traversal map $\Gamma\to\mathbb R^2$ that runs once around
$P$, on which each double point of $P$ has exactly two preimages, so that
$\Gamma$ carries $2m$ marked points. (No coordinate is placed on $\Gamma$; in
particular the crossing-free case $m=0$ is allowed, with no marked points.) Then:
\begin{enumerate}
\item[(i)] the oriented smoothing of $P$ along $S$ has exactly $|S|+1$ carriers;
\item[(ii)] the assignment of traversal points to carriers is \emph{non-crossing}:
there are no four points $u_1,u_2,u_3,u_4$ of $\Gamma$ in this cyclic order,
none of them a preimage of an element of $S$, with $u_1,u_3$ on one carrier and
$u_2,u_4$ on a different carrier;
\item[(iii)] if $c\in[m]\setminus S$ is non-adjacent in $G_P$ to every element of
$S$, then both occurrences of $c$ lie on the same carrier;
\item[(iv)] if $H$ is a connected component of the induced graph
$G_P\bigl[[m]\setminus(S\cup N_{G_P}(S))\bigr]$, then there is exactly one
carrier on which both occurrences of every $c\in H$ lie.
\end{enumerate}
\end{lemma}
\begin{proof}
Induct on $|S|$, proving (i), (ii) and (iii) together; the induction is on the
conjunction, and (iii) is what makes the step go through.

\emph{Base $S=\varnothing$.} There is one carrier, namely $P$ itself: (i) and
(iii) are immediate and (ii) is vacuous, all four points lying on that one
carrier.

\emph{Step.} Let $|S|\geq1$, fix $c\in S$ and put $S'=S\setminus\{c\}$, so
$S'\in\Ind(G_P)$ and $|S'|=|S|-1$. Since $S$ is independent, $c$ is non-adjacent
to every element of $S'$, so the inductive hypothesis (iii) applies to $S'$ and
$c$: both occurrences of $c$ lie on one and the same carrier $L$ of $S'$. Write $x_1,x_2\in\Gamma$ for its two preimages; they cut $\Gamma$ into two
arcs $\alpha,\beta$.

(i) Oriented smoothing at $c$ replaces the two crossing arcs at that double point
by the two non-crossing arcs respecting the orientation, which severs $L$ at
$x_1$ and $x_2$ and reconnects it as the two closed curves carried by
$L\cap\alpha$ and $L\cap\beta$; every other carrier of $S'$ is untouched, since
$c$ meets none of them. So the number of carriers rises by exactly one, from
$|S'|+1$ to $|S|+1$.

(ii) The carriers of $S$ are those of $S'$ with the single block $L$ replaced by
the two blocks $L\cap\alpha$ and $L\cap\beta$. Suppose $u_1,u_2,u_3,u_4$ in cyclic
order violate (ii) for $S$, with $u_1,u_3$ on a carrier $A$ and $u_2,u_4$ on a
carrier $B\neq A$, none of the $u_k$ an occurrence of an element of $S$. If
neither $A$ nor $B$ is one of the two new blocks, the same four points violate
(ii) for $S'$, contradicting the inductive hypothesis. If exactly one of them, say
$A$, is a new block, then $A\subseteq L$ and $B$ is a carrier of $S'$ distinct
from $L$, so the four points again violate (ii) for $S'$ with $L$ in place of $A$.
If both are new blocks, then $\{A,B\}=\{L\cap\alpha,\,L\cap\beta\}$, so
$u_1,u_3\in\alpha$ and $u_2,u_4\in\beta$ (or the same with $\alpha,\beta$
exchanged). None of the $u_k$ equals $x_1$ or $x_2$, so each of the four cyclic
transitions $u_1\to u_2\to u_3\to u_4\to u_1$ passes from $\alpha$ to $\beta$ or
back, and every such passage contains one of the two points $x_1,x_2$ in its
interior. Four disjoint open arcs would then each contain a point of the
two-element set $\{x_1,x_2\}$, which is impossible

(iii) Let $c''\in[m]\setminus S$ be non-adjacent to every element of $S$, hence to
every element of $S'$. By the inductive hypothesis both occurrences of $c''$ lie
on one carrier $L''$ of $S'$. If $L''\neq L$ then $L''$ is a carrier of $S$ as
well and the claim holds. If $L''=L$, use that $c''$ is non-adjacent to $c$:
exactly one occurrence of $c$ lying between the two occurrences of $c''$ is what
adjacency means, so non-adjacency puts both occurrences of $c''$ strictly inside
the same arc, $\alpha$ or $\beta$. Hence both lie on the same one of
$L\cap\alpha$, $L\cap\beta$.

(iv) Put $U=[m]\setminus(S\cup N_{G_P}(S))$. For $c\in H$, clause~(iii) puts
both occurrences of $c$ on one carrier, which we call $L(c)$. If $c,c'\in H$
are adjacent in $G_P[U]$, exactly one occurrence of $c'$ lies between the two
occurrences of $c$. Thus their four occurrences alternate on $\Gamma$, and none
is an occurrence of an element of $S$. If $L(c)\neq L(c')$, those four points
violate clause~(ii). Hence $L(c)=L(c')$. The assignment $c\mapsto L(c)$ is
therefore constant on every edge of the connected graph $H$, hence constant on
$H$.

It is one carrier, not two, because it is a function. Uniqueness is immediate:
the carriers partition $\Gamma$, so distinct carriers are disjoint \emph{as
subsets of the traversal circle} --- not as subsets of the plane, where two
distinct carriers can meet only at \emph{dominated} crossings, those of
$N_{G_P}(S)$. At a crossing outside $S\cup N_{G_P}(S)$ both visits lie on one
carrier by clause~(iii), so distinct carriers cannot meet there, and the
crossings of $S$ are smoothed. An earlier revision said distinct carriers meet
at every shared double point \emph{outside} $S\cup N_{G_P}(S)$, which names
exactly the set where they cannot.
\end{proof}

\begin{lemma}[smoothing preserves the induced cyclic order]\label{lem:carrierword}
Let $C$ be a closed curve with a traversal circle $\Gamma_C$ and finitely many
transverse double points, and let $S$ be a set of them no two of which
interlace. Then each closed curve produced by smoothing $C$ along $S$ traverses
the marked points of $\Gamma_C$ lying on it in the cyclic order they have on
$\Gamma_C$.

In particular, taking $C=P$ generic and $S\in\Ind(G_P)$, each carrier of $S$
traverses the marked points lying on it in the order induced from $\Gamma$.
\end{lemma}
\begin{proof}
Induction on $|S|$. If $S=\varnothing$ the only curve is $C$ itself, traversing
$\Gamma_C$, and there is nothing to prove.

Let $|S|\geq1$ and pick $c\in S$, with traversal preimages $c_1,c_2$. Smoothing
$c$ alone joins the incoming branch at $c_1$ to the outgoing branch at $c_2$ and
the incoming at $c_2$ to the outgoing at $c_1$, so the two resulting closed
curves traverse exactly the two arcs $\Gamma_1=(c_1,c_2)$ and
$\Gamma_2=(c_2,c_1)$ of $\Gamma_C$, each in the order induced from $\Gamma_C$;
no marked point changes arcs and none is visited twice.

Every other $e\in S$ is nonadjacent to $c$, so its two preimages do not separate
$c_1$ from $c_2$: both lie in $\Gamma_1$ or both in $\Gamma_2$. Hence
$S\setminus\{c\}$ is partitioned into two sets, one carried by each of the two
closed curves, no two members of either interlacing; and the remaining
smoothings are performed inside those two curves separately. Each is a closed
curve with a traversal circle and finitely many transverse double points, which
is the hypothesis of this lemma, so the induction hypothesis applies to each and
gives the induced order there; composing with the previous paragraph gives it on
$\Gamma_C$.

The statement was restated at this generality because the induction applies it
to curves obtained by smoothing, which are not polygons; an earlier revision
quantified it over generic polygons only, and the peer was right that its own
proof then stepped outside its domain.
\end{proof}

\begin{definition}[corner, turn sign, and $\wind$]\label{def:wind}
Let $P$ be \emph{generic} and $S\in\Ind(G_P)$ --- the binders under which the
definition is read, and which an earlier revision left to the ambient
subsection, whose standing hypothesis is only that $P$ is diagrammatic. A
\emph{corner} of a carrier of $S$ is either a vertex $p_i$ of $P$ traversed by
it or a smoothing site of $S$ traversed by it. At each corner the carrier turns
\emph{left} or \emph{right} according to the sign of the determinant of the
incoming and outgoing directions, which under genericity is nonzero: at a
vertex by (G1), unconditional and hence relevant, at a smoothing site by (G5),
active at a crossing. On a merely diagrammatic polygon the vertex clause is
false --- $((0,0),(1,0),(2,0),(0,1))$ is diagrammatic and its turn at $p_2$ is
exactly zero --- which is why the generic binder is part of this definition and
not decoration. A carrier is \emph{uniform} if all its corners turn the
same way and \emph{mixed} otherwise. Its \emph{weight} is
\[
   \mathrm{wt}(L)=\begin{cases}
     +1, & L \text{ uniform, all turns right},\\
     (-1)^{c(L)}, & L \text{ uniform, all turns left, } c(L)=\#\text{corners},\\
     0, & L \text{ mixed},
   \end{cases}
\]
and $\wind(S)=\prod_{L}\mathrm{wt}(L)$, the product over the $|S|+1$ carriers of
$S$. In particular $\wind(S)\neq0$ forces every carrier of $S$ to be
uniform: $\wind(S)$ is a product of integers, a product of integers is
nonzero only if every factor is, and the weight of a mixed carrier is $0$.
\end{definition}

\begin{definition}[undominated set and residual pieces]\label{def:pieces}
For $S\in\Ind(G_P)$ put $U(S)=[m]\setminus\bigl(S\cup N_{G_P}(S)\bigr)$. The
\emph{residual pieces} of $S$ are the connected components
$H\in\pi_0\bigl(G_P[U(S)]\bigr)$ of the induced subgraph on $U(S)$.
\end{definition}

\subsection{The readout}

\begin{definition}[record, and record isomorphism]\label{def:record}
A \emph{record} consists of an oriented circle $\Gamma$; a finite subset
$V\subset\Gamma$ of \emph{marked points}; a partition of $V$ into two-element
sets, the \emph{double points}; for each double point a designation of one of
its two points as the \emph{over} position and the other as the \emph{under}
position; and for each double point a \emph{sign} in $\{\pm1\}$. The record of a
link diagram is the one obtained by taking $\Gamma$ to be its traversal circle,
$V$ the preimages of its crossings, the partition the crossing correspondence,
and the over/under and sign data those of the diagram.

A \emph{record isomorphism} from $(\Gamma,V,\dots)$ to $(\Gamma',V',\dots)$ is a
bijection $\varphi:V\to V'$ that
\begin{enumerate}
\item[(a)] preserves the cyclic order: for all $v_1,v_2,v_3\in V$ occurring in
that cyclic order on $\Gamma$, the points $\varphi v_1,\varphi v_2,\varphi v_3$
occur in that cyclic order on $\Gamma'$;
\item[(b)] carries double points to double points;
\item[(c)] carries over positions to over positions and under to under;
\item[(d)] preserves signs.
\end{enumerate}
Condition (a) says exactly that $\varphi$ extends to an orientation-preserving
homeomorphism $\Gamma\to\Gamma'$ carrying $V$ onto $V'$, and any two such
extensions are isotopic through such homeomorphisms: an orientation-preserving
circle homeomorphism is determined up to isotopy by the cyclic-order-preserving
bijection it induces on a finite subset, the complementary arcs being carried to
the complementary arcs in the induced order and any two orientation-preserving
homeomorphisms of a closed arc rel endpoints being isotopic. Records are
\emph{isomorphic} when such a $\varphi$ exists; the relation is an equivalence.
\end{definition}

\begin{definition}[HOMFLY--PT normalization]\label{def:homfly}
$P_H$ is normalized by the first two identities below; the third is their
consequence, obtained by applying the skein relation at a crossing between $L$
and a split unknot, and is displayed because it fixes the reader's convention
for a split component:
\begin{equation}
P(\bigcirc)=1,\qquad
aP(L_+)-a^{-1}P(L_-)=zP(L_0),\qquad
P(L\sqcup\bigcirc)=\frac{a-a^{-1}}{z}\,P(L).
\label{eq:homfly}
\end{equation}
\end{definition}

\begin{definition}[piece diagram, positivity, writhe]\label{def:piecediagram}
For a diagrammatic parent $P$, the diagram of a residual piece $H$ is the parent
curve $P$ with every double
point outside $H$ erased --- the two strands drawn as passing without
interaction --- and every double point of $H$ resolved by the divide convention:
the branch whose direction $u_{\mathrm{over}}$ satisfies
$\det(u_{\mathrm{over}},u_{\mathrm{under}})>0$ passes over. Under this convention
every crossing of the diagram is positive, so its writhe is
$w(H)=|H|$, the number of double points of $H$. We write $P_H(a,z)$ for the
HOMFLY--PT polynomial of the link so presented, normalized as in
Definition~\ref{def:homfly}.

\emph{What ``erased'' means, stated once, and where the statement comes from.}
The evaluator that defines $X_1$ builds a piece's datum as the parent Gauss word
restricted to the crossings of $H$, in the parent's cyclic order, together with
the parent's per-crossing counter-clockwise half-edge order restricted to the
surviving half-edges and re-expressed in the restricted word's arc labels
(\texttt{sub\_word\_and\_rot},
\texttt{engines/polygon\_motivic/ruling\_dp\_statesum\_fin.py:28} of the shipped
code). So the datum is a word \emph{together with a rotation system}, not a bare
word, and the definition above is to be read as naming that pair.

To erase a double point is to omit it from that datum. The combinatorial shadow
of the datum --- the word with its over/under and signs and nothing else --- is
Definition~\ref{def:record}.
\end{definition}

\begin{lemma}[a residual piece is carried by an actual closed curve]
\label{lem:piececurve}
Let $P$ be diagrammatic (Definition~\ref{def:diagrammatic}),
$S\in\Ind(G_P)$, and let $H$ be a residual piece of $S$. Let $C_H$ be the closed
plane curve that the proof's own route constructs, and nothing else: smooth
every crossing of $S$ --- of $S$ \emph{only} --- so that the curve falls into
the $|S|+1$ carriers of $S$; take the one carrier that carries $H$, which is
unique by Lemma~\ref{lem:carriers}(iv); and apply to it Step 5's iteration,
smoothing at each step one double point outside $H$ and keeping the daughter
curve that carries $H$, which Step 5 shows is well defined. Then $C_H$ is a
closed plane curve whose double points are exactly the crossings of $H$ and
whose Gauss word is the parent word restricted to $H$ in the parent's cyclic
order.

This binding is the construction and not a description of its result. The
route above smooths nothing but $S$, so $H$ lies on one carrier before the
iteration begins, and each subsequent step preserves that. In particular that restricted word is realizable, and the datum of
Definition~\ref{def:piecediagram} is the datum of $C_H$.
\end{lemma}
\begin{proof}
\emph{Step 1: a dominated crossing has its two visits on different carriers.}
Let $c\notin S$ interlace some $s\in S$. Smoothing $s$ alone splits the curve
into two closed curves traversing the two arcs $(s_1,s_2)$ and $(s_2,s_1)$ of
$\Gamma$, and $c$ interlacing $s$ means one visit of $c$ lies in each. Smoothing
the remaining elements of $S$ only ever splits: the carriers number $|S|+1$ by
Lemma~\ref{lem:carriers}(i), which is the count reached when every one of the
$|S|$ smoothings increases the component count, so none of them merges two
components. Hence the two visits of $c$ lie on different carriers.

\emph{Step 2: the double points of a carrier are exactly the undominated
crossings it carries.} A double point of a carrier $L$ is a crossing outside $S$
both of whose visits lie on $L$. Step 1 excludes the dominated ones, and
Lemma~\ref{lem:carriers}(iii) puts both visits of each undominated one on a
single carrier.

\emph{Step 3: a carrier's Gauss word is realizable and is a restriction.} A
carrier is a closed plane curve with transverse double points, so its Gauss word
is realizable, being read off an actual curve. By Step 2 its double points are
the undominated crossings it carries, and by Lemma~\ref{lem:carrierword} it
traverses them in the cyclic order they have on $\Gamma$. So the carrier's word
is the parent word restricted to those crossings.

\emph{Step 4: a residual piece interlaces nothing else in that word.} $H$ is a
connected component of $G_P[U(S)]$, so an undominated crossing interlacing a
member of $H$ lies in $H$; and $H$ lies on one carrier by
Lemma~\ref{lem:carriers}(iv).

\emph{Step 5: restricting to such a set is realizable, constructively.} Let $C$
be a closed plane curve and $K$ a nonempty set of its double points such that no
member of $K$ interlaces a double point outside $K$ \emph{and $K$ induces a
connected interlacement subgraph}. If $C$ has a double point $d$ outside $K$,
smooth it. The curve becomes two closed curves, traversing the two arcs of $d$.
Since $d$ interlaces no member of $K$, each member of $K$ has both visits in one
of those two arcs; and two members lying in different arcs have their visits in
disjoint arcs, so they do not interlace. Connectedness of $K$ therefore forces
all of $K$ into one arc, hence onto one of the two curves. Keep that one and
discard the other, which carries no member of $K$; the retained visits keep
their cyclic order by Lemma~\ref{lem:carrierword} applied to the single smoothed
chord. The hypotheses persist: $K$ is unchanged, so it is still connected, and
it still interlaces no surviving double point outside it. Each step removes at
least one double point outside $K$, so after finitely many steps the curve
$C_H$ has double points exactly $K$ and Gauss word the restriction.

The connectedness hypothesis is not decoration. Without it the step is false,
and the peer's exact witness is the word $d\,a\,a\,d\,b\,b$ with
$K=\{a,b\}$: neither $a$ nor $b$ interlaces $d$, yet $a$ lies inside the arc
$(d_1,d_2)$ and $b$ outside it, so $K$ is not on one side. A residual piece
supplies the hypothesis by Step 4, being a connected component.

Applying Step 5 to the carrier of $H$ with $K=H$, which Steps 3 and 4
licence --- Step 4 supplying both the non-interlacing and the connectedness
hypotheses --- performs exactly the statement's third operation and gives
$C_H$. Its rotation system at each surviving crossing is the carrier's,
which is the parent's, since smoothing at other points does not disturb the
half-edge order at a survivor; that is the second half of the datum.
\end{proof}

\begin{remark}[what the piece datum presents, and a reading that fails]
\label{rem:piecedatum}
An earlier revision of Definition~\ref{def:piecediagram} defined the diagram by
its record alone. That is not the evaluator's datum and it is not safe:
restricting a realizable Gauss word to an arbitrary label set can produce a word
no plane curve realizes --- the peer's witness restricts $123123$ to two of its
three crossings and gets $2323$, in which each chord interlaces exactly one
other, against the realizable-word parity used in the proof of
Corollary~\ref{cor:evendominated}. What rescues the
definition is not a reading but the lemma above: a residual piece is not an
arbitrary label set. The record shadow is what Lemma~\ref{lem:pieceintrinsic}
compares.

Definition~\ref{def:piecediagram} names a word together with a rotation system
and calls it the diagram of a piece. Lemma~\ref{lem:piececurve} is what makes
that a link diagram rather than a formal datum: the pair is the datum of an
actual closed plane curve, so $P_H$ is a classical invariant and no
realizability question is left open. The fact is recorded here, after the lemma,
rather than inside the definition, which cannot draw on a result printed after
it.
\end{remark}

\begin{corollary}[undominated crossings meet an even number of dominated ones]
\label{cor:evendominated}
Let $P$ be generic and $S\in\Ind(G_P)$. Every undominated crossing interlaces an
even number of dominated crossings.
\end{corollary}
\begin{proof}
Let $h$ be undominated, on the carrier $L$. The word of $L$ is realizable
(Step 3 above), so $h$ interlaces an even number of chords of that word, which
by Step 2 are the undominated crossings on $L$; and $h$ interlaces no
undominated crossing off $L$, by the non-crossing property
Lemma~\ref{lem:carriers}(ii). So $h$ interlaces an even number of undominated
crossings in $\Gamma$'s word. That word is realizable too, so $h$ interlaces an
even number of chords altogether, and $h$ interlaces no member of $S$, being
undominated. Subtracting leaves an even number of dominated ones.
\end{proof}
\begin{proof}[Second proof, the peer's]
This argument is not this author's and is printed with that attribution. For any
independent $T$, the set $N(T)$ is exactly the set of crossings whose two visits
lie on \emph{different} carriers of $T$, by Steps 1 and 2 above. Distinct closed
curves in the plane meet in an even number of transverse points, so summing over
the pairs of carriers, $|N(T)|$ is even. Now $h$ is undominated, so $S$,
$\{h\}$ and $S\cup\{h\}$ are all independent, and
\[
   N\bigl(S\cup\{h\}\bigr)=N(S)\cup N(h).
\]
Taking cardinalities modulo two and using that all three of
$|N(S\cup\{h\})|$, $|N(S)|$ and $|N(h)|$ are even gives $|N(S)\cap N(h)|$ even,
and that intersection is exactly the set of dominated crossings interlacing $h$.

The peer's caveat is worth keeping with the proof: this parity fact is necessary
for realizability but not equivalent to it, so it does not replace
Lemma~\ref{lem:piececurve}, whose carrier construction remains load-bearing.
\end{proof}


\begin{definition}[rotation number, computably]\label{def:rot}
Let $L$ be a closed polygon through corner points $q_0,\dots,q_{c-1}$, lying
in the regular locus $\mathcal R_c$ of Definition~\ref{def:regular}(B) ---
nonzero edges, no consecutive pair doubling back; equivalently, every
principal turn of Definition~\ref{def:regular}(A) exists --- and put
$\delta_i=q_{i+1}-q_i$. Choose $r\in\mathbb{R}^2$ with $\det(r,\delta_i)\neq0$
for all $i$; such $r$ exists because the $\delta_i$ are finitely many. Then
\[
   \rot(L)=\sum_{i}\epsilon_i,\qquad
   \epsilon_i=\begin{cases}
     +1,&\det(\delta_{i-1},\delta_i)>0,\ \det(\delta_{i-1},r)>0,\ \det(r,\delta_i)>0,\\
     -1,&\det(\delta_{i-1},\delta_i)<0,\ \det(\delta_{i-1},r)<0,\ \det(r,\delta_i)<0,\\
     0,&\text{otherwise.}
   \end{cases}
\]
All comparisons are exact sign tests of determinants: no angles and no
tolerances occur.

\emph{Why regularity is part of the definition.} The sum is independent of
the admissible $r$, and regularity is what makes it so; a vanishing turn is
harmless. First delete flat corners: if $\delta_i=\lambda\delta_{i-1}$ with
$\lambda>0$ then $\det(\delta_{i-1},\delta_i)=0$, so $\epsilon_i=0$ for every
$r$, and deleting $q_i$ merges the two edges into one of the same direction,
changing no other $\epsilon_j$, since every determinant a surviving term
reads is altered only by a positive scale. The deletion lands in
$\mathcal R_{c-1}$, and no polygon of any $\mathcal R_c$ is all-flat --- a
closed polygon with every turn flat would traverse one direction forever and
could not close --- so finitely many deletions reach a polygon with every
$\det(\delta_{i-1},\delta_i)\neq0$ and the same sum. On that polygon,
$\epsilon_i$ depends on $r$ only through $\det(\delta_{i-1},r)$ and
$\det(r,\delta_i)$, so it moves only as $r$ crosses the line of
$\delta_{i-1}$ or of $\delta_i$. Crossing the line of $\delta_j$ moves
exactly $\epsilon_j$ and $\epsilon_{j+1}$, and in each of the four sign cases
one gains what the other loses, so the sum is unchanged; if several edges
share that line they are pairwise nonconsecutive --- consecutive ones would
be flat or doubling back, and both are gone --- so their pairs are disjoint
and each is preserved separately. Regularity is not decoration: the polygon
$((-7,4),(-9,6),(9,-9),(3,3),(9,-9),(-7,-7),(-7,-6))$, whose consecutive
pair $(-6,12),(6,-12)$ doubles back, returns both $-1$ and $0$ at admissible
rays.

\emph{Direction loops and smooth curves.}
For a continuous loop $T\colon\mathbb R/\mathbb Z\to S^1$, a
\emph{tangent-angle lift} at a seam is a continuous
$\theta\colon[0,1]\to\mathbb R$ such that
$T(s)=(\cos\theta(s),\sin\theta(s))$. Put
\[
   \operatorname{tw}(T)=\frac{\theta(1)-\theta(0)}{2\pi}.
\]
For a closed $C^1$ regular oriented curve
$\gamma\colon\mathbb R/\mathbb Z\to\mathbb R^2$, put
\[
   T_\gamma=\frac{\gamma'}{|\gamma'|},
   \qquad
   \rot(\gamma)=\operatorname{tw}(T_\gamma).
\]
Lemma~\ref{lem:turnlift} constructs the lift and proves that these values do
not depend on its choices, on the seam, or on an orientation-preserving regular
reparametrisation. For either a polygon or a $C^1$ regular closed curve put
$R(L)=|\rot(L)|$.

\end{definition}

\begin{lemma}[tangent lifts and principal turns]\label{lem:turnlift}
Here $\rot$ is as in Definition~\ref{def:rot}.
\begin{enumerate}
\item[(i)] Every continuous direction loop has a tangent-angle lift.
The integer $\operatorname{tw}(T)$ is independent of the lift and seam,
is unchanged by an orientation-preserving reparametrisation, and is constant
under a continuous homotopy through direction loops. Consequently the rotation
of a closed $C^1$ regular curve is invariant under regular homotopy, and
reversing its orientation negates it.
\item[(ii)] If $L\in\mathcal R_c$ has principal turns $\tau_i$, then
\[
   2\pi\rot(L)=\sum_i\tau_i.
\]
Consequently rotation is constant along every path in $\mathcal R_c$, is
unchanged by a positive flat subdivision, and is negated by orientation
reversal. Every regular three-corner polygon has rotation $+1$ or $-1$
according to the common sign of its three turns.
\item[(iii)] Suppose a closed $C^1$ regular curve is obtained from $L$ by
replacing every corner by a regular arc whose compatible tangent-angle lift
has increment $\tau_i$. Its rotation equals $\rot(L)$.

More generally, let $b,c$ be regular oriented arcs having the same endpoints
and the same oriented tangent rays at both endpoints, and suppose replacing
$b$ by $c$ in a closed curve gives closed $C^1$ regular curves $F_b,F_c$.
For compatible tangent lifts with equal initial values, write their increments
as $\Delta_b,\Delta_c$. Then
\[
   \rot(F_c)-\rot(F_b)=\frac{\Delta_c-\Delta_b}{2\pi}.
\]
If $A_t$, $0\leq t\leq1$, is a supplied continuous path in
$\mathrm{GL}^+(2,\mathbb R)$ with $A_0=I$ and $A_1=A$, applying $A$ to both
arcs leaves $\Delta_c-\Delta_b$ unchanged.
\end{enumerate}
\end{lemma}
\begin{proof}
For~(i), uniform continuity gives
$0=s_0<\cdots<s_m=1$ such that
$T(s)\cdot T(s_{j-1})>0$ on every
$[s_{j-1},s_j]$. After choosing $\theta(0)$, continue it there by
\[
 \theta(s)=\theta(s_{j-1})+
 \operatorname{atan2}\!\left(
   \det(T(s_{j-1}),T(s)),\,
   T(s_{j-1})\cdot T(s)\right).
\]
The second argument is positive, so the added angle lies in
$(-\pi/2,\pi/2)$; the formulas match at the subdivision points and construct a
continuous lift without a covering-space theorem. Since $T(1)=T(0)$, its
endpoint increment belongs to $2\pi\mathbb Z$. Two lifts differ continuously
by a value in $2\pi\mathbb Z$, hence by one constant.

If the increment is $2\pi k$, extend the lift by
$\theta(s+n)=\theta(s)+2\pi nk$. Every interval of length one then has the same
increment, proving seam independence. An orientation-preserving
reparametrisation of the circle has an increasing representative
$\phi\colon\mathbb R\to\mathbb R$ with
$\phi(s+1)=\phi(s)+1$; the lift $\theta\circ\phi$ has the same increment.

For a homotopy $T(s,t)$ fix $t_0$. Uniform continuity gives a neighbourhood of
$t_0$ on which
$T(s,t)\cdot T(s,t_0)>0$ for every $s$. There is therefore
a unique continuous relative angle
\[
 \alpha(s,t)=\operatorname{atan2}\!\left(
 \det(T(s,t_0),T(s,t)),\,
 T(s,t_0)\cdot T(s,t)\right)
 \in(-\pi/2,\pi/2).
\]
If $\theta_0$ lifts $T(\,\cdot\,,t_0)$, then
$\theta_0+\alpha(\,\cdot\,,t)$ lifts $T(\,\cdot\,,t)$.
Periodicity gives $\alpha(1,t)=\alpha(0,t)$, so the endpoint increment is
locally, hence globally, constant in $t$. A regular homotopy supplies such a
homotopy of normalised tangents. For the oppositely oriented curve,
$-T_\gamma(1-s)$ has lift $\theta(1-s)+\pi$, whose increment is the negative
of the original one.

For~(ii), give the admissible reference vector $r$ of
Definition~\ref{def:rot} an angle $\rho$, and let
$\beta_i\in(\rho,\rho+2\pi)$ represent the direction of the $i$th edge.
The three cases in that definition give exactly
\[
   \tau_i=\beta_i-\beta_{i-1}+2\pi\epsilon_i:
\]
a positive crossing of the reference ray contributes $+2\pi$, a negative
crossing contributes $-2\pi$, and otherwise no correction occurs.
Summing cyclically cancels the $\beta$ terms and proves the displayed identity.

Along a path in $\mathcal R_c$, every $\tau_i$ varies continuously in
$(-\pi,\pi)$; hence the displayed integer is constant. A positive flat
subdivision inserts one zero turn and changes no other turn. Reversal negates
and reverses the turn list.

For the three-corner claim, if the vertices are $A,B,C$, all three cyclic turn
determinants equal $\det(B-A,C-A)$. This determinant cannot vanish: otherwise
the three nonzero collinear edge scalars sum to zero, and a cyclic sign change
is a consecutive negative-multiple pair, contrary to regularity. Thus all
three turns have one sign. Their sum is nonzero and has magnitude below
$3\pi$; being $2\pi$ times an integer, it has rotation $+1$ or $-1$ according
to that sign.

For~(iii), concatenate constant lifts on the straight pieces with the
prescribed corner lifts. Their total increment is $\sum_i\tau_i$, so~(ii)
proves the rounding assertion.

For the replacement assertion, use the same lift on the unchanged
complementary arc. Shifting that complementary lift at a join changes neither
its increment nor the computation; subtraction cancels it and leaves
$\Delta_c-\Delta_b$.

Finally follow the tangent path of $c$ and then the tangent path of $b$
backwards. This is a closed direction loop of increment
$\Delta_c-\Delta_b$. The maps
\[
   \nu_t(z)=\frac{A_tz}{|A_tz|}
\]
give a homotopy of that direction loop. Clause~(i) keeps its increment
constant, proving the last assertion without invoking degree.
\end{proof}

\begin{definition}[the slot, the factor, and $X_1$]\label{def:X1}
Let $P$ be generic and $S\in\Ind(G_P)$. For a carrier $L$ of $S$ set
\[
  P_{S,L}=\prod_{H\text{ carried by }L}P_H,\qquad
  w_{S,L}=\sum_{H\text{ carried by }L}w(H),
\]
the products and sums taken over the residual pieces assigned to $L$ by
Lemma~\ref{lem:carriers}(iv), with the empty product $P_{S,L}=1$ and the empty
sum $w_{S,L}=0$ when no piece is carried by $L$. Here $P_H$ is the HOMFLY--PT
polynomial of the link that the piece $H$ presents, in the normalization of
Definition~\ref{def:homfly}; that such a polynomial exists at all is
Axiom~\ref{ax:homfly}, and this definition is where it is first read. The \emph{slot} of $L$ is the
integer $1-w_{S,L}-R(L)$, the \emph{factor} is
\[
   \Omega_1(S,L)=\bigl[a^{\,1-w_{S,L}-R(L)}z^{0}\bigr]P_{S,L}(a,z)
\]
--- the coefficient itself, taken as a value, with no sign gate and no
zero-gate applied --- and
\[
   X_1(P)\;=\;\sum_{S\in\Ind(G_P)}\wind(S)\prod_{L}\Omega_1(S,L),
\]
the inner product running over the $|S|+1$ carriers of $S$.
\end{definition}

\begin{proposition}[chambers and chamber invariance]\label{prop:chamberinv}
Fix $n\geq3$.
\begin{enumerate}
\item[(i)] The generic locus $\mathcal U_n$ is open in $(\mathbb R^2)^n$, and
every chamber --- every connected component of it --- is open and path connected.
\item[(ii)] $X_1$ is constant on each chamber.
\end{enumerate}
\end{proposition}
\begin{proof}
For clause (i), fix $P\in\mathcal U_n$. The argument must control the members that are
\emph{not} relevant at $P$ as well as those that are: openness of an activation
region says that a member active at $P$ stays active nearby, and says nothing
about one that is inactive at $P$ becoming active. An earlier revision inferred
the second from the first.

What controls both is the unconditional layer. Every unconditional member ---
every $\mathrm{G1}_i$ and every $\mathrm{G2}_{e,i}$ --- is relevant at every
polygon, so all of them are nonzero at $P$, they are finitely many, and they are
polynomials; let $N$ be a neighbourhood of $P$ on which each of them keeps its
sign. On $N$ the crossing status of every pair of remote edges is constant,
being the four-sign condition of Definition~\ref{def:guarded} in exactly those
members; hence the activation of every conditional member is constant on $N$,
and so is the set of members relevant at a point. Shrinking $N$ so that the
finitely many members relevant at $P$ --- now the members relevant everywhere on
$N$ --- also keep their signs, every point of $N$ is generic. So
$\mathcal U_n$ is open.

A connected component of an open subset of $\mathbb R^N$ is open, the space
being locally connected, and an open connected subset of $\mathbb R^N$ is path
connected: the set of points reachable from a fixed one by a path in it is open
and its complement in the component is open, and the component is connected. An
earlier revision asserted both facts inside the proof that consumes them,
rather than proving them.

For clause (ii), let $P^{0},P^{1}$ lie in one chamber and choose a path
$t\mapsto P(t)$, $t\in[0,1]$, between them inside it, as clause (i) allows.
Every statement below is proved along that path.

\emph{The combinatorics.} Each unconditional member $g\in\mathcal G$ is a
polynomial, so the composite $t\mapsto g(P(t))$ is continuous and, by genericity,
nowhere zero along the path. Its sign is therefore constant on the connected
interval $[0,1]$. The four-sign crossing test of
Definition~\ref{def:guarded} therefore fixes the set of crossing pairs
along the path, and hence a conditional member is active at one parameter
exactly when it is active at every parameter. Each $g$ in this now-fixed active
set is a polynomial relevant at every $P(t)$, so $t\mapsto g(P(t))$ is likewise
continuous and nowhere zero by genericity. Its sign is therefore constant as
well. So: the set of double points, indexed by the pairs of edges carrying them, is
constant; and no two of them collide, for the following reason, which splits in
two cases and uses (G3) in one and (G4) in the other --- an earlier revision
announced it as "the (G4) one and not the (G3) one", which is the half that is
true when the two double points share an edge. Two double points can collide in two ways, and an earlier revision admitted
only one of them, saying that four distinct edges through a point "is not a
coincidence of any member". It is one. \emph{If the four edges are distinct}, the collision puts four edges through a
single point, and there are two sub-cases. If some three of the four are
pairwise nonparallel, they cross pairwise at that point, so their
$\mathrm{G3}$ member is active, and it vanishes, three concurrent lines sharing
a projective point; being active it is relevant, so genericity forbids it. If
no such triple exists, then two of the four are parallel, and being parallel
lines through a common point they are the same line: two distinct edges lying
on one line, each containing that point in its relative interior, hence
overlapping in a segment. Then an endpoint of one lies on the other, an
\emph{unconditional} $\mathrm{G2}$ vanishes, and genericity forbids that as
well. Either way no such collision occurs along a path of generic polygons.
\emph{Otherwise the two double points share an edge}, and there the concurrency
member says nothing --- the two carrying edges may be parallel --- so the
argument is the (G4) one: let them lie on a common edge
$e$, carried by $f$ and $g$, in the order the path presents them. The member
that decides their order along $e$ is
$\mathrm{G4}\langle e;f,g\rangle$, the member-valued accessor of
Definition~\ref{def:guarded}, and the order itself is the sign of the oriented
value $\widehat{\mathrm{G4}}_{e;f,g}=(t_f-t_g)\det(d_f,d_e)\det(d_g,d_e)$,
which is that member times the orientation sign $\epsilon(f,g)$. The member is
active exactly when both $f$ and $g$ cross $e$, which is the case, so it is
relevant and hence nonzero of constant sign along the path; and since
$\epsilon(f,g)$ is a constant of the index data, the oriented value is of
constant sign too. Only the nonvanishing is used below, so the conclusion is
the same under either reading. A collision would be $t_f=t_g$, so no collision occurs and the
order along each edge is constant. When $f$ and $g$ share a vertex $M$ the same
member does the work, its vanishing being that of
$\mathrm{G2}_{e,M}\mathrm{G1}_M$. Both factors are unconditional members, so on
this path --- along which every polygon is \emph{generic} --- both are nonzero;
in particular $\mathrm{G1}_M\neq0$, and that is genericity and not regularity.
Membership of $\mathcal R_n$ (Definition~\ref{def:regular}(B)) excludes only zero
edges and a consecutive pair that doubles back, and a flat vertex with
$\mathrm{G1}_M=0$ is regular. So the product vanishes only if
$\mathrm{G2}_{e,M}$ does, which puts $M$ on the edge $e$ and is forbidden here.
An earlier revision of this sentence said "the turn is nonzero on a regular
polygon", which is false of $\mathcal R_n$. An earlier revision credited
the constancy of the active (G3) signs for all of this; a concurrency member is
inactive when its three edges do not pairwise cross, and then it says nothing,
while two crossings on one edge can still swap. Hence the cyclic sequence of crossing visits along
the traversal --- the Gauss word --- is constant, and so are $G_P$, $\Ind(G_P)$,
each $U(S)$, each family of residual pieces, and, by
Lemma~\ref{lem:carriers}(iv), the assignment of pieces to carriers. Each
carrier is determined by the Gauss word and the smoothing sites, so the carriers
correspond canonically along the path, with corners varying continuously.

\emph{The weights.} At a vertex corner the turn sign is $\sgn\mathrm{G1}_i$,
constant; at a smoothing corner of the crossing of $e,f$ it is the sign of
$\mathrm{G5}_{e,f}$ up to the fixed sign of the smoothing convention, also
constant. So each carrier is uniform on the whole path or mixed on the whole
path, its corner count is constant, and $\mathrm{wt}(L)$ and $\wind(S)$ are
constant (Definition~\ref{def:wind}).

\emph{The rotations.} Turn signs alone do \emph{not} determine a rotation
number --- the convex regular pentagon and the regular pentagram have the same
five positive turn signs and rotations $1$ and $2$ --- so this step is a path
argument, not a sign count. Along the path each carrier $L(t)$ has corners
varying continuously and no vanishing turn determinant, so each principal turn
$\tau_k(t)$ is continuous with values in $(-\pi,\pi)$, and
$\rot(L(t))=\frac1{2\pi}\sum_k\tau_k(t)$ by Lemma~\ref{lem:turnlift}(ii) is a
continuous integer-valued function of $t$ on the connected interval, hence
constant. So each $R(L)$ is constant.

\emph{The piece polynomials.} Fix two parameters $t_0,t_1$ and a residual
piece $H$, identified at the two parameters by the common Gauss word. By
Lemma~\ref{lem:piececurve}, each piece datum is realized by a connected closed
plane curve whose double points are exactly the crossings of $H$ and whose
cyclic word is the parent word restricted to $H$. Match each marked crossing--
carrying-edge incidence $(c,e)$ at $t_0$ with $(c,e)$ at $t_1$. This bijection
preserves cyclic order, and the two incidences bearing $c$ remain its crossing
pair. It preserves signs because every piece crossing is positive, and it
preserves over/under
positions because the divide convention chooses them from the sign of the
active $\mathrm{G5}$ member of the two carrying edges, which is constant along
the chamber. Thus it is an orientation-preserving record isomorphism
(Definition~\ref{def:record}). Lemma~\ref{lem:piececurve} retains exactly
the finite, distinct transverse crossings of $H$
and creates no self-intersection when it smooths locally; hence each realizing
curve has one component and no triple points. Axiom~\ref{ax:gausscode}
therefore says that the piece diagrams present the same oriented link. Hence
$P_H$ is constant by Axiom~\ref{ax:homfly}, while $w(H)=|H|$ is constant because
$H$ is.

Hence every $\Omega_1(S,L)=[a^{1-w_{S,L}-R(L)}z^0]P_{S,L}$ is constant along the
path, and so is the finite sum defining $X_1$. In particular
$X_1(P^0)=X_1(P^1)$.
\end{proof}

\begin{definition}[event, and its zero set]\label{def:event}
An \emph{event} is a continuous path $t\mapsto P(t)$ of polygons on $n$
vertices, $t\in(-\varepsilon,\varepsilon)$, such that $P(t)$ is generic for
every $t\neq0$ and $P(0)$ is not. Its \emph{zero set} is
\[
   Z=\bigl\{g\in\mathcal G:\ g(P(0))=0\ \text{ and }\ g\ \text{is relevant at }
   P(t)\ \text{for some }t\neq0\bigr\}.
\]
The \emph{two chambers of the event} are the chambers containing
$P\bigl((-\varepsilon,0)\bigr)$ and $P\bigl((0,\varepsilon)\bigr)$; each of
those two sets is connected and consists of generic polygons, so each does lie
in one chamber. We say the event is \emph{guarded relative
to $Z$}, and every appeal to simplicity below names its $Z$.

The event is \emph{transversal} if every member of $Z$ changes sign at $t=0$.
Every named event of the convention that lists them is required to be
transversal. Without that requirement the definition admits a tangency: the
path
\[
   p_1=(0,0),\quad p_2(t)=(2,-t^2),\quad p_3=(1,0),\quad p_4=(3,2)
\]
is generic for $t\neq0$ and has exactly the cusp zero bundle at $t=0$, yet
$P(t)=P(-t)$, so its two punctured sides are the same chamber and no crossing
is created. A tangency is not a wall crossing and every wall law would be
vacuous or false there.
\end{definition}

\begin{remark}[why the two chambers give two values]\label{rem:eventvalues}
$X_1$ takes one value on each of the two chambers of an event, by
Proposition~\ref{prop:chamberinv}(ii), which is what makes the jump across the event
a well-defined number and is the reason the definition names the chambers at
all. The fact is recorded here rather than in the definition, so that the
definition depends on nothing proved after it.
\end{remark}

\begin{lemma}[the guards outside the zero set]\label{lem:guardconst}
Let $t\mapsto P(t)$, $t\in(-\varepsilon,\varepsilon)$, be an event with zero set
$Z$, and let $g\in\mathcal G$ be relevant at $P(t)$ for some $t\neq0$ and not a
member of $Z$. Then $g(P(0))\neq0$, and there is
$\varepsilon'\in(0,\varepsilon]$ such that $g$ is nonzero and of constant sign
on the whole of $(-\varepsilon',\varepsilon')$.
\end{lemma}
\begin{proof}
By Definition~\ref{def:event}, $Z$ consists of the members that vanish at
$P(0)$ \emph{and} are relevant at some $t\neq0$. The hypothesis supplies the
second clause and denies membership, so the first clause fails: $g(P(0))\neq0$.
Since $g$ is a polynomial in the coordinates and $t\mapsto P(t)$ is continuous,
$g\circ P$ is continuous and nonzero at $0$, hence nonzero on some interval
$(-\varepsilon',\varepsilon')$, on which, being continuous and nowhere zero, it
has constant sign.
\end{proof}

\begin{remark}[why the conclusion is local and not global]\label{rem:guardlocal}
An earlier revision concluded that $g$ is nonzero on the \emph{whole} interval.
That does not follow and its own proof conceded as much: $g$ may be relevant on
one side of the event and inactive on the other, and an inactive conditional
member is under no constraint --- it may vanish, or vanish identically. What is
true, and what every consumer uses, is the local statement above. Since an
event is used only as a germ at $t=0$, shrinking $\varepsilon$ costs nothing,
and every appeal below is to the shrunken interval.
\end{remark}

\begin{remark}[why the zero set is defined and not imposed]\label{rem:zeroset}
An earlier form of this definition required \emph{every} member of
$\mathcal G\setminus Z$ to be nonzero throughout. That is too strong and makes
the intended domains empty: a conditional member that is inactive on both sides
of the event may vanish identically --- for instance a crossing-order predicate
$\mathrm{G4}_{e;f,g}$ for a pair $f,g$ that never crosses $e$ --- so no event
would satisfy it. It was also too weak in the other direction, since it did not
require the two punctured sides to be generic, and without that the values
$X_1(P_\pm)$ that every wall law compares need not exist.
Definition~\ref{def:event} asks only for genericity off the wall and reads the
zero set off the path; Lemma~\ref{lem:guardconst} then \emph{proves} the
constancy that used to be assumed.

The specification's Definition~3 declares an event simple when \emph{the}
predicate defining the named event vanishes and every other guarded predicate
keeps its sign. Read with $|Z|=1$ that has empty domain at some events used
below: at a flat event three members vanish together --- the turn
$\mathrm{G1}_j$ and the two incidences it forces --- so no such event has a
singleton zero set. The set-valued form above is what every appeal here uses.
\end{remark}

\begin{convention}[the named events]\label{conv:events}
Throughout, all three named events are required to be transversal in the
sense of Definition~\ref{def:event}, and for a polygon on $n$ vertices and an
index $j$:
\begin{itemize}
\item write $\mathcal T_j=\{\mathrm{G1}_j,\ \mathrm{G2}_{(j-1,j),\,j+1},\
\mathrm{G2}_{(j,j+1),\,j-1}\}$ for the \emph{turn trio at $j$}, and
\[
   \mathcal F_j=\bigl\{\,\mathrm{G4}\langle h;\,e_{j-1},e_j\rangle\ :\ h\
   \text{an edge remote to the vertex } p_j\,\bigr\}
\]
for the order predicates that compare the two crossings a single edge makes with
the two edges at $j$. The entries of $\mathcal F_j$ are \emph{members} of
$\mathcal G$ --- the member-valued accessor of Definition~\ref{def:guarded}
returns the member indexed by the reduced representatives --- as
Definition~\ref{def:event} requires of anything that can lie in a zero set. The
factorization is a statement about signs and is therefore written on the
oriented value:
$\widehat{\mathrm{G4}}_{h;e_{j-1},e_j}=-\mathrm{G2}_{h,j}\,\mathrm{G1}_j$,
where the orientation sign $\epsilon(e_{j-1},e_j)$ is $-1$ exactly when the
pair wraps, that is exactly when $j\equiv1$. Since $\epsilon$ is a nonzero
constant, the member and the oriented value vanish together and change sign
together, so the conclusion below may be read off either. On a
neighbourhood of a configuration with $\mathrm{G2}_{h,j}\neq0$ it is a nonzero
multiple of $\mathrm{G1}_j$: it vanishes with the trio and changes sign with it.
\item a \emph{simple flat event at $j$} is an event whose zero set is exactly
$\mathcal T_j$, at which $p_j$ lies strictly between $p_{j-1}$ and $p_{j+1}$ at
$t=0$. No member of $\mathcal F_j$ can be relevant at such an event, so nothing
is lost by requiring the zero set to be the trio alone: at $t=0$ the two edges
at $j$ lie on one line $\ell$ and meet only at $p_j$, which is interior to their
union. A remote edge $h$ is then in one of two positions. If $h$ is not
contained in $\ell$ it meets $\ell$ in at most one point, so it cannot meet both
legs in their relative interiors and cannot cross both. If $h$ \emph{is}
contained in $\ell$ it is parallel to both legs, and a parallel pair does not
cross --- the four-sign condition of Definition~\ref{def:guarded} fails, its
four $\mathrm{G2}$ members vanishing --- so it crosses neither. Either way no
member of $\mathcal F_j$ is active at $t=0$.

Activation on a punctured side is a germ condition and is settled in the
limit, not by the value at $t=0$ alone. Suppose $h$ crosses both legs at
parameters $t_n\neq0$ with $t_n\to0$; let $x_n$ and $y_n$ be the crossing
points on the two legs. The vertices are continuous in $t$ and the segments
are compact, so after passing to a subsequence $x_n\to x$ and $y_n\to y$, with
$x$ and $y$ in the closed edge $h(0)$ and in the respective closed legs at
$t=0$, all three of which lie where their $t=0$ positions are. Two positions of
$h(0)$ against the flat line $\ell$ remain, and both force $p_j\in h(0)$. If
$h(0)$ is not contained in $\ell$, it meets $\ell$ in at most one point; both
$x$ and $y$ lie in $h(0)\cap\ell$, so $x=y$, and that one point lies in both
closed legs, whose intersection is $\{p_j\}$. If $h(0)$ \emph{is} contained in
$\ell$, then it is a segment of $\ell$ meeting both closed legs, which lie on
opposite sides of $p_j$ along $\ell$, $p_j$ being interior to their union; a
connected set meeting both sides contains the point between them. Either way
$p_j\in h(0)$, so the unconditional member $\mathrm{G2}_{h,j}$ vanishes at
$t=0$, joins the zero set, and the simplicity requirement excludes it. A
previous repair of this paragraph argued instead from "$h$ meets at most one
leg's relative interior at $t=0$", which the contained-line position falsifies
--- legs $(-2,0)$--$(0,0)$--$(2,0)$ with $h(0)=[(-1,0),(1,0)]$ meet in both
relative interiors --- while the conclusion survives, $h(0)$ containing $p_j$
there. So on a
punctured neighbourhood of $0$ no member of $\mathcal F_j$ is relevant, which is
the condition Definition~\ref{def:event} imposes. An earlier revision argued the
first position only, writing "a straight edge $h$ meeting that line once", which
says nothing about an $h$ lying along it;
\item a \emph{simple cusp event at $j$} is an event with a zero set of that same
form --- $\mathcal T_j$ together with the relevant members of $\mathcal F_j$ ---
at which $p_j$ lies outside the segment $[p_{j-1},p_{j+1}]$ at $t=0$. The event
is \emph{threaded} when some member of $\mathcal F_j$ is relevant, that is when
some edge crosses both collapsing legs, and \emph{unthreaded} otherwise; both
are simple, and the laws are stated for both.

This is where an earlier revision of this document narrowed the domain. It
required the zero set to be exactly $\mathcal T_j$, which excludes every
threaded cusp. The witness, printed in Remark~\ref{rem:witnesses} in its
own indexing, has the cusp at $j=1$ on seven vertices: the trio vanishes as
$\tfrac25\to0\to-\tfrac25$ across the wall and the thread's member
$\mathrm{G4}_{4;1,7}$ as $-\tfrac{16}5\to0\to\tfrac{16}5$ --- oriented value
$\tfrac{16}5\to0\to-\tfrac{16}5$, the pair at $j=1$ wrapping --- the
thread crossing both legs on both sides --- while the cusp law and the main
theorem are both stated "threaded or not".
Phase 1.5 admitted those events under \emph{relative} simplicity; the clause
above is that condition printed. The enlarged zero set is still \emph{one}
condition: every added member is a nonzero multiple of $\mathrm{G1}_j$ near the
wall, so it vanishes and changes sign exactly when the turn does, which is the
flat bundle's pattern and is what the transversality requirement asks;
\item a \emph{simple vertex-on-edge event}, with moving vertex $M$ and remote
edge $(a,b)$, is an event at which $p_M$ lies in the relative interior of
$[p_a,p_b]$ at $t=0$, and whose zero set is \emph{branch-specific}:
\[
   Z_{\mathrm{bigon}}
   =\bigl\{\mathrm{G2}_{(a,b),\,M},\
     \mathrm{G4}\langle (a,b);\,e_{M-1},e_{M}\rangle\bigr\},
   \qquad
   Z_{\mathrm{sliding}}=\bigl\{\mathrm{G2}_{(a,b),\,M}\bigr\}.
\]
In the \emph{bigon} branch the two neighbours of $M$ lie strictly on one side of
the remote line, so on the two-crossing chamber both edges incident to $M$ meet
the remote edge; the order predicate comparing their crossing parameters along
$(a,b)$ is active there, and at the event the two parameters coincide at $p_M$,
so it vanishes. It is therefore relevant at some $t\neq0$ and vanishes at
$t=0$: by Definition~\ref{def:event} it belongs to $Z$, and a singleton zero
set would have empty domain, exactly as at a flat event.

In the \emph{sliding} branch the two neighbours lie strictly on opposite sides,
so on \emph{each} punctured side exactly one incident edge meets the remote
edge. That order predicate is then active on neither side; it vanishes at the
wall but is relevant at no $t\neq0$, so Definition~\ref{def:event} does not put
it in $Z$, and imposing it would empty the branch
\end{itemize}
\end{convention}

\begin{definition}[silent event]\label{def:silent}
An event is \emph{silent} if every member of its zero set $Z$ is a predicate
$\mathrm{G2}_{e,i}$ whose vanishing at $t=0$ places $p_i$ on the \emph{line} of
$e$ but not on the segment $e$. In particular no member of $Z$ is a turn (G1),
a crossing-order predicate (G4), a concurrency (G3) or a direction determinant
(G5).
\end{definition}

\begin{remark}[two clauses that were removed because they are empty]
\label{rem:silentclauses}
An earlier revision of Definition~\ref{def:silent} admitted two further kinds of
member of $Z$: a conditional (G3) whose common point at $t=0$ misses one of the
three segments, and a (G5) whose pair of edges has relative interiors not
meeting at $t=0$. Neither can occur.

The first is Lemma~\ref{lem:dictionary}(A)(a). For the second, split on whether
the two closed segments meet at $t=0$. If they do not, it is
Lemma~\ref{lem:dictionary}(A)(b). If they do, then since they are parallel
with disjoint relative interiors they meet at a point that is an endpoint of at
least one of them, say the vertex $p_i$ of $e$ lying on $f$; then
$\mathrm{G2}_{f,i}$ vanishes at $t=0$ \emph{with the vertex on the segment}, and
$\mathrm{G2}$ is unconditional, hence relevant at every $P(t)$, so that
predicate lies in $Z$ and the event is not silent, since
Definition~\ref{def:silent} asks of every member of $Z$ that its vanishing put
the vertex on the \emph{line} and not on the segment; no clause about its (G5)
would have made it so.

Removing the clauses does not change the class of silent events. What it costs
is that the proof of Lemma~\ref{lem:silence} can no longer say ``clause~(ii)
excludes that''; it argues instead from the dichotomy the relevance filter
provides, which is the honest form of the same step. The peer raised this; the
entry records it.
\end{remark}


\begin{lemma}[$X_1$ is unchanged across a silent event]\label{lem:silence}
Let $t\mapsto P(t)$ be a silent event. Then $X_1(P_+)=X_1(P_-)$.
\end{lemma}
\begin{proof}
By the proof of Proposition~\ref{prop:chamberinv}(ii), $X_1$ is determined by the
set of double points indexed by the pairs of edges carrying them, their order
along each edge, the turn signs at vertices and at smoothing sites, the
\emph{absolute rotation} $R(L)$ of each carrier, and the polynomial $P_H$ of each
residual piece. An earlier revision omitted $R(L)$ from this list,
although the slot of Definition~\ref{def:X1} reads it. It suffices to show each is
constant on the whole interval, $t=0$ included.

A double point of the pair $(e,f)$ exists exactly when the relative interiors of
$e$ and $f$ meet, and by continuity it can appear or disappear only through a
configuration in which an endpoint of one segment lies on the other segment ---
a $\mathrm{G2}$ predicate vanishing \emph{with the vertex on the segment}.
Definition~\ref{def:silent} excludes that --- a member of $Z$ may put $p_i$ on
the line of $e$ but never on the segment --- and every other relevant
$\mathrm{G2}$ is nonzero throughout by Lemma~\ref{lem:guardconst}.
So the set of double points is constant, and with it the set of active
conditional members. Nor do two double points on \emph{disjoint} pairs of edges collide: that puts
four edges through one point, and the two sub-cases are those of
Proposition~\ref{prop:chamberinv}(ii). If some three of the four are pairwise
nonparallel they cross pairwise there, so their $\mathrm{G3}$ member is active
and vanishes, and an active member is relevant, hence lies in $Z$ if it
vanishes at $t=0$ --- which Definition~\ref{def:silent} forbids, its zero sets
containing no (G3). Otherwise two of the four are collinear and overlap at the
point, so an unconditional $\mathrm{G2}$ vanishes there with the vertex on the
\emph{segment}, which Definition~\ref{def:silent} also forbids.
The collision would leave the Gauss word alone but would destroy the
diagrammatic hypothesis the isotopy step below reads.

Two double points on a common edge exchange order only by colliding, which puts
three segments through one common point. Fix the common edge $e$ and the two edges $f,g$ carrying the colliding double
points; both are remote to $e$, a crossing requiring remoteness. The predicate
that decides the order of those two crossings along $e$ is the member
$\mathrm{G4}\langle e;f,g\rangle$ --- the accessor, since $f$ and $g$ are
presented in the order the geometry gives and their representatives need not be
increasing --- and it is the right tool: it is \emph{active} exactly when $f$
and $g$ both cross $e$, which is the situation in hand, and it vanishes exactly
when the two crossing parameters coincide, which is the collision. Being active
it is relevant, so if it vanished at $t=0$ it would lie in the zero set $Z$;
and $Z$ contains no (G4) member, by Definition~\ref{def:silent} --- the zero
set holds members, so it is membership of the accessor's value that is
excluded, and the orientation sign is irrelevant to it. So it does not
vanish, and by Lemma~\ref{lem:guardconst} it keeps its sign: no collision.

An earlier revision reached instead for the concurrency member
$\mathrm{G3}_{e,f,g}$ and argued by cases on whether $f$ and $g$ are remote or
consecutive. That misses a configuration, and the peer's chart is exact: with
$e$ from $(-2,0)$ of direction $(4,0)$, $f_t$ from $(t,-1)$ of direction
$(0,2)$ and $g_t$ from $(-t,-2)$ of direction $(0,4)$, the triple is pairwise
remote, $f$ and $g$ are \emph{parallel} so $\mathrm{G3}$ is inactive and says
nothing, yet both cross $e$ and their crossings swap at $t=0$;
$\mathrm{G3}=-64t$ and $\mathrm{G4}_{e;f,g}=64t$, and it is the (G4) that is
active and in the zero set's way.

The two edges $f,g$ may also share a vertex $M$, and then the same predicate
does the work, being indexed for any pair remote to $e$. Its vanishing is then
the vanishing of $\mathrm{G2}_{e,M}\,\mathrm{G1}_M$. The second factor is a
turn, and it is nonzero here because $\mathrm{G1}_M$ is unconditional and every
polygon of a silent event's punctured sides is generic, with $\mathrm{G1}_M$
outside the zero set by Definition~\ref{def:silent} and so nonzero at $t=0$ as
well by Lemma~\ref{lem:guardconst} --- regularity alone would not give it, a
flat vertex being regular. So a collision would need $M$ on the line of $e$, and
on the \emph{segment}, since the collision point is the common crossing
--- a $\mathrm{G2}$ vanishing with the vertex on the segment, which
Definition~\ref{def:silent} excludes from $Z$ while $\mathrm{G2}$ is
unconditional and hence relevant.

So no collision occurs, the order along each edge is constant, and the Gauss
word with it.

The turn sign at a vertex is $\sgn\mathrm{G1}_i$, which is unconditional,
hence relevant, and lies in no $Z$ of a silent event, so it is nonzero of
constant sign by Lemma~\ref{lem:guardconst}. The
turn sign at a smoothing site of the crossing of $e$ and $f$ is the sign of
$\mathrm{G5}_{e,f}$ up to the fixed sign of the smoothing convention. That
predicate is not in $Z$, so by the same dichotomy either it is nonzero on a
neighbourhood, or $e$ and $f$ carry no double point for $t\neq0$ and hence no
smoothing site whose turn could change.

The absolute rotation $R(L)$ of a carrier is next, and it needs its own
argument: turn signs alone do not determine a rotation number, and link
isotopy cannot substitute for it, plane rotation not being a knot-type
invariant. Each carrier $L(t)$ has, by the no-birth and no-collision facts
above, a fixed finite cyclic corner set, each corner varying continuously in
$t$. Each corner's turn determinant --- the $\mathrm{G1}$ at a vertex, the
$\mathrm{G5}$ at a smoothing site --- is nonzero through the wall by the two
preceding paragraphs, $t=0$ included, so each principal turn $\tau_k(t)$ is
continuous with values in $(-\pi,\pi)$, and
$\rot(L(t))=\frac1{2\pi}\sum_k\tau_k(t)$ by Lemma~\ref{lem:turnlift}(ii) is a
continuous integer-valued function on the connected event interval, hence
constant. This is the path argument of Proposition~\ref{prop:chamberinv}(ii) run
through the wall, and it is printed rather than cited because a silent wall
lies inside no chamber. So each $R(L)$ is constant.

Finally, the common Gauss word identifies every residual piece $H$ on the two
sides. The edge-pair indexing is constant as proved above: a silent zero puts its vertex
off the remote segment, so no crossing changes carrying edge at the wall.
Lemma~\ref{lem:piececurve} realizes its datum on each side by a connected
closed plane curve whose double points are exactly the crossings of $H$. Match each marked crossing--carrying-edge
incidence $(c,e)$ on the first side with the same incidence $(c,e)$
on the second. The common restricted word makes this bijection cyclic-order
preserving; the two incidences bearing $c$ remain its crossing pair. Every sign
is positive, and the $\mathrm{G5}$ argument above shows that the divide
convention designates the same visit as over on the two sides. Thus the map is
an orientation-preserving record isomorphism
(Definition~\ref{def:record}). Lemma~\ref{lem:piececurve} retains exactly
the finite, distinct transverse crossings of $H$ from the two generic parents
and creates no self-intersection when it smooths locally; hence each realizing
curve has one component and no triple points. Axiom~\ref{ax:gausscode} says that the piece diagrams present the same
oriented link, so Axiom~\ref{ax:homfly} gives the same $P_H$; and
$w(H)=|H|$ is unchanged because $H$ is.
Every input to $X_1$ is therefore constant on $(-\varepsilon,\varepsilon)$, and
$X_1(P_+)=X_1(P_-)$.
\end{proof}

\begin{remark}[why this lemma is here]\label{rem:silencenew}
The induction that proves the main theorem transports a polygon along a path
and needs $\Delta=X_1-A$ to be unchanged wherever the path is not at one of the
four named walls. For $A$ that is a declared input (Axiom~\ref{ax:lawful}). For
$X_1$ it is two facts, not one: constancy inside a chamber
(Proposition~\ref{prop:chamberinv}(ii)) and invariance across the silent walls
between chambers (Lemma~\ref{lem:silence}). A silent wall \emph{is} a wall ---
a guarded predicate vanishes on it --- so chamber constancy does not cross it.
\end{remark}

\begin{remark}[the guarded list cannot be trimmed]\label{rem:chamberguards}
The proof of Proposition~\ref{prop:chamberinv}(ii) uses all five families, and an
implementation enforcing only some of them defines a coarser partition than
Definition~\ref{def:generic}(B), on which the proposition is false. (G1) fixes the
vertex turn signs, (G2) keeps vertices off edges, (G3) keeps three edges from
concurring, (G4) fixes the crossing order along each edge, (G5) fixes the
smoothing-site turn signs. Which pairs of edges cross, and hence which
conditional members are active, is decided by the (G2) signs alone. Dropping (G3) or (G4) permits two crossings to
exchange order along an edge with no guarded predicate changing sign; that
changes the Gauss word and can change $X_1$.
\end{remark}

\input{d2_laws}
% Section 3: the carrier floor.  Printed in full, including the three
% constructions the phase-1.5 manuscript discharged by citing an evidence
% pack (the rounding, the vertical-direction choice, the negative curl).
\section{The carrier floor}\label{sec:floor}

The selector in Definition~\ref{def:X1} makes the contributing carriers
geometrically special, and that specialness bounds the degrees available to
their HOMFLY rows from below. The bound is the engine of Section~\ref{sec:s7}:
every vanishing proved there is an instance of it.

\subsection{What the selector forces}

\begin{remark}[attribution of the constructions in this section]\label{rem:floorsources}
The phase-1.5 manuscript proved Theorem~\ref{thm:carrierfloor}(C) but discharged
three of its steps by citing an evidence pack of the counterpart bench: the
rounding of the polygon corners, the choice of a vertical direction with its
tangency count, and the exact negative curl. Those three are printed here in
full --- Lemma~\ref{lem:rounding}, Theorem~\ref{thm:carrierfloor}(B) and Lemma~\ref{lem:curl} ---
and no conclusion here depends on that pack. The direct floor count in
Theorem~\ref{thm:carrierfloor}(B) replaces its former degree-theory last step;
attribution is otherwise unchanged: the rational curl model and its affine
insertion, the dual-vector argument of
Lemma~\ref{lem:uniformrot}(ii), and the computation identifying the forbidden
crossing as positive are the counterpart's; Lemma~\ref{lem:uniformrot}(i), the mirror calculation in the proof of
Theorem~\ref{thm:carrierfloor}(C), and that theorem's assembly are this author's. Printing a proof is not
claiming it; it is taking responsibility for it.
\end{remark}

\subsection{Rounding}

\begin{lemma}[rounding]\label{lem:rounding}
Here $\rot$ is as in Definition~\ref{def:rot}.
Let $L$ be a closed polygon with corners $q_1,\dots,q_c$, and let $D$ be an
oriented diagram whose underlying plane curve is $L$ --- $L$ together with an
over/under assignment at each of its double points. Assume the principal turns
of $L$ all exist and are \emph{nonzero}; that $L$ has finitely many double
points, all transversal, none of them a corner; and that no corner of $L$ lies
on an edge of $L$ other than the two incident to it. Then there is a \emph{clearance} $\varepsilon_0(L)>0$ such that for every
$\varepsilon\in\bigl(0,\varepsilon_0(L)\bigr)$ there is a $C^\infty$ regular
closed plane curve $L_\varepsilon$, and a diagram $D_\varepsilon$ carried by it, with
these properties. The clearance is part of the conclusion and is quantified
here: an earlier revision wrote "below the bound displayed in the proof", which
names an object chosen inside a proof as though the statement supplied it.

The hypothesis names a curve \emph{and} a diagram: clause~(c) speaks of an
over/under assignment, of crossing signs and of a writhe, none of which a plane
curve carries, and an earlier revision bound only $L$. It does \emph{not} name a
generic parent polygon. An earlier revision did, and then
Theorem~\ref{thm:carrierfloor}(C) applied this lemma on its own hypotheses, which
do not mention one --- the peer's polygon $((0,0),(1,0),(-2,-2),(2,0),(-2,1))$
satisfies the floor's hypotheses and is not a carrier of a support of a generic
polygon. What genericity was supplying is written out instead: finitely many
transverse double points, none at a corner, and no corner on a non-incident
edge, which is exactly what makes the finitely many distances in the proof
positive.
\begin{enumerate}
\item[(a)] $L_\varepsilon$ coincides with $L$ outside the union of the discs of
radius $\varepsilon$ about the corners.
\item[(b)] Inside the disc about $q_i$ the unit tangent moves \emph{strictly}
monotonically, in the sense of $\sgn\tau_i$, from $\delta_{i-1}/|\delta_{i-1}|$
to $\delta_i/|\delta_i|$, sweeping an arc of length exactly $|\tau_i|$ and no
more, and attaining each direction of that arc at exactly one parameter.
Moreover the unit-tangent map is an \emph{immersion} on the open junction arc:
parametrized by arclength, its angular derivative is nonzero at every interior
parameter, while at the two ends it and all its derivatives vanish, the
junction meeting the straight edges flat to infinite order.

Both clauses are exported because both are read downstream, and the second does
not follow from the first: a strictly monotone angle may still have a vanishing
derivative at a parameter, and the direct count below needs the derivative
itself to be nonzero, not merely the monotonicity. "Moves monotonically" alone
is weaker still --- it is
satisfied by a junction that pauses on one direction for a whole sub-arc, and
then a count of the parameters carrying a prescribed direction is infinite. Both
hold because the profile fixed in the proof below is strictly increasing on
$(0,1)$ with $\varphi'>0$ there, and flat to infinite order at $0$ and $1$: the
tangent angle at arclength $\sigma$ along the junction of length $\ell$ is
$\theta_u+\tau_i\varphi(\sigma/\ell)$, whose derivative is
$\tau_i\varphi'(\sigma/\ell)/\ell\neq0$ for $0<\sigma<\ell$.
\item[(c)] $L_\varepsilon$ has the same double points as $L$, with the same
strands, the same over/under assignment and hence the same crossing signs and
the same writhe.
\item[(d)] $\rot(L_\varepsilon)=\rot(L)$.
\item[(e)] \emph{The disc package is returned.} There are pairwise disjoint
closed discs $D_1,\dots,D_c$, one about each corner $q_i$, such that each $D_i$
meets no edge of $L$ that is not incident to $q_i$, contains no double point of
$L$, meets the two incident edges exactly in the two sub-segments of length
$\varepsilon$ at $q_i$, and contains the whole of the modification made at
$q_i$. This clause is a conclusion of the lemma, not a package handed over by a
paragraph of its proof; its consumer asks for a disc meeting the diagram in one
embedded arc and reads it here.
\end{enumerate}
\end{lemma}
\begin{proof}
\emph{The junction, constructed.} Since $0<|\tau_i|<\pi$, the unit directions
$u=\delta_{i-1}/|\delta_{i-1}|$ and $v=\delta_i/|\delta_i|$ are neither parallel
nor antiparallel. Write $A_0=q_i-\varepsilon u$ and $A_1=q_i+\varepsilon v$ for
the two points at distance $\varepsilon$ from $q_i$ along the incident edges, so
that the displacement to be realized is
\[
   A_1-A_0=\varepsilon\,(u+v),
\]
a nonzero vector along the bisector of the convex angle at $q_i$. Fix once and
for all the \emph{transition profile}
\[
   \varphi(t)=\frac{f(t)}{f(t)+f(1-t)},
   \qquad
   f(t)=\begin{cases} e^{-1/t}, & t>0,\\ 0,& t\leq0,\end{cases}
\]
which is exhibited rather than posited, and which has the four properties every
use below reads. It is $C^\infty$ on $[0,1]$: $f$ is $C^\infty$ on $\mathbb R$
with all derivatives vanishing at $0$, and the denominator
$f(t)+f(1-t)$ is positive on $[0,1]$, being $f(1)>0$ at $t=0$ and $f(1)>0$ at
$t=1$ and a sum of two nonnegative terms not both zero in between. It has
$\varphi(0)=0$ and $\varphi(1)=1$. It satisfies $\varphi(1-t)=1-\varphi(t)$, by
inspection of the formula. And it is \emph{strictly increasing on $(0,1)$}:
\[
   \varphi'(t)=\frac{f'(t)f(1-t)+f(t)f'(1-t)}{\bigl(f(t)+f(1-t)\bigr)^2}>0
   \qquad (0<t<1),
\]
both terms of the numerator being positive there, since $f>0$ and
$f'(t)=t^{-2}e^{-1/t}>0$ on $(0,1)$. Every derivative of $\varphi$ vanishes at
$0$ and at $1$, because every derivative of $f$ vanishes at $0$ and the
denominator is smooth and nonvanishing.

Strict increase is not a convenience. An earlier revision asked only that
$\varphi$ be nondecreasing, and a nondecreasing profile may be constant on a
subinterval; the tangent direction would then be constant on a whole sub-arc of
the junction, so a direction attained there would be attained at a continuum of
parameters rather than at isolated ones, and every count of the points carrying
a prescribed tangent direction --- the counts this section takes later --- would
be infinite. With $\varphi$ strictly increasing on $(0,1)$ the tangent angle
$\theta_u+\tau_i\varphi(t/\ell)$ is strictly monotone across the junction, and
each direction in the swept arc is attained exactly once. Let
$\theta_u$ be the angle of $u$ and, for a length $\ell>0$, let
\[
   \gamma_\ell(s)=A_0+\int_0^{s}
   \bigl(\cos(\theta_u+\tau_i\varphi(t/\ell)),\,
         \sin(\theta_u+\tau_i\varphi(t/\ell))\bigr)\,\mathrm dt,
   \qquad s\in[0,\ell].
\]
Then $\gamma_\ell$ is a unit-speed $C^\infty$ arc. Its tangent equals $u$
\emph{at} $s=0$ and $v$ \emph{at} $s=\ell$, and agrees with those constants
\emph{to infinite order} at those two parameters, every derivative of $\varphi$
vanishing there; it does not equal them on a neighbourhood, and cannot, since
the angle is strictly monotone on $(0,\ell)$ --- an earlier revision wrote
"near $s=0$", which contradicts the strictness this same construction exports.
Infinite-order agreement at the endpoint is what the junction needs: it is
exactly the condition under which replacing the corner by the arc leaves a
$C^\infty$ regular curve, since all one-sided derivatives match those of the
straight edge. The tangent angle moves strictly monotonically through exactly
$\tau_i$ and no more, and its derivative
$\tau_i\varphi'(s/\ell)/\ell$ is nonzero on $(0,\ell)$, which is (b).

Its endpoint is $A_0+\ell\,m(\varphi)$ where
$m(\varphi)=\int_0^1(\cos(\theta_u+\tau_i\varphi),\sin(\theta_u+\tau_i\varphi))\,
\mathrm dt$. The symmetry $\varphi(1-t)=1-\varphi(t)$ makes the angles
$\theta_u+\tau_i\varphi(t)$ and $\theta_u+\tau_i\varphi(1-t)$ reflections of
each other in the bisector, so $m(\varphi)$ points along $u+v$; and
$|m(\varphi)|>0$ because all the directions integrated lie in a closed angular
interval of width $|\tau_i|<\pi$. Taking
$\ell=\varepsilon\,|u+v|/|m(\varphi)|$ makes the endpoint exactly $A_1$. This is
where the circular arc of an earlier revision has been replaced: that arc is
$C^1$ and not $C^\infty$ at the two junctions, and the document then carried a
lemma to repair the loss.

\emph{The arc lies in the triangle $A_0q_iA_1$, in coordinates.} Put
$\alpha=|\tau_i|/2\in(0,\pi/2)$ and $s=\sgn\tau_i$, let $w$ be the unit vector
along $u+v$ and $z$ the unit vector with $\det(w,z)=1$, and take $q_i$ as
origin, writing a point as $(x_w,x_z)$ in that frame. Then
\[
   u=\cos\alpha\,w-s\sin\alpha\,z,
   \qquad
   v=\cos\alpha\,w+s\sin\alpha\,z,
\]
because $u$ and $v$ are unit vectors whose sum is along $w$ and whose angle is
$|\tau_i|=2\alpha$, turning from $u$ to $v$ by $\tau_i$. Hence
$A_0=-\varepsilon u$ has coordinates $(-\varepsilon\cos\alpha,\ s\varepsilon
\sin\alpha)$ and $A_1=\varepsilon v$ has $(\varepsilon\cos\alpha,\
s\varepsilon\sin\alpha)$, and the closed triangle with vertices $A_0,q_i,A_1$
is exactly
\begin{equation}
   T_i=\bigl\{\,(x_w,x_z):\ |x_w|\leq\cot\alpha\cdot(s\,x_z),
   \quad 0\leq s\,x_z\leq\varepsilon\sin\alpha\,\bigr\},
\label{eq:triangle}
\end{equation}
the first inequality being the pair of cone conditions at the apex $q_i$ and the
second the chord through $A_0$ and $A_1$: a point $\lambda(-u)+\mu v$ with
$\lambda,\mu\geq0$ has $x_w=(\mu-\lambda)\cos\alpha$ and
$s\,x_z=(\lambda+\mu)\sin\alpha$, so $|x_w|\leq(\lambda+\mu)\cos\alpha=\cot\alpha
\cdot(s\,x_z)$ with equality exactly on the two edges, and $\lambda+\mu\leq
\varepsilon$ is the chord.

The junction arc satisfies all three inequalities of \eqref{eq:triangle}. Write
its unit tangent as $T(\sigma)=\cos\beta(\sigma)\,w+s\sin\beta(\sigma)\,z$ with
$\beta(\sigma)=-\alpha+2\alpha\,\varphi(\sigma/\ell)$, which is the tangent
angle constructed above read in this frame: at $\sigma=0$ it is $u$
($\beta=-\alpha$), at $\sigma=\ell$ it is $v$ ($\beta=\alpha$), and $\beta$ is
strictly increasing with values in $[-\alpha,\alpha]$.
\emph{The two cone inequalities.} Let $n_1$ be the unit normal to $u$ with
$\langle n_1,v\rangle>0$. The line $\mathbb Ru$ contains both $q_i$ and $A_0$,
and $T_i$ lies in $\{\langle\cdot,n_1\rangle\geq0\}$. Along the arc,
$\tfrac{d}{d\sigma}\langle\gamma,n_1\rangle=\langle T,n_1\rangle\geq0$, because
$T=au+bv$ with $a,b\geq0$ --- a direction at angle $\beta\in[-\alpha,\alpha]$
from $w$ lies in the closed cone of $u$ and $v$ --- so
$\langle T,n_1\rangle=b\langle v,n_1\rangle\geq0$; and
$\langle\gamma(0),n_1\rangle=\langle-\varepsilon u,n_1\rangle=0$. Hence
$\langle\gamma,n_1\rangle\geq0$ throughout. Symmetrically, with $n_2$ the unit
normal to $v$ with $\langle n_2,A_0\rangle>0$, one has $\langle
T,n_2\rangle=a\langle u,n_2\rangle\leq0$, so $\langle\gamma,n_2\rangle$ is
nonincreasing and equals $\langle\varepsilon v,n_2\rangle=0$ at $\sigma=\ell$;
hence it is $\geq0$ throughout. Those two are the first inequality of
\eqref{eq:triangle}.
\emph{The chord inequality.} $s\,x_z(\sigma)=\varepsilon\sin\alpha+\int_0^\sigma
\sin\beta(t)\,dt$, and $\int_0^\sigma\sin\beta\leq0$ for every
$\sigma\in[0,\ell]$: $\beta$ is negative on $(0,\ell/2)$ and positive on
$(\ell/2,\ell)$, and the symmetry $\varphi(1-t)=1-\varphi(t)$ makes
$\beta(\ell-t)=-\beta(t)$, so the integral decreases to its minimum at
$\ell/2$ and increases back to $0$ at $\ell$. Hence $s\,x_z\leq\varepsilon
\sin\alpha$, which is the second. The lower bound $s\,x_z\geq0$ is implied by
the two cone inequalities, the apex being the origin.

So the arc lies in $T_i$, and $T_i$ lies in the closed disc $D_i$ of radius
$\varepsilon$ about $q_i$: the disc is convex and contains the three vertices,
$|A_0-q_i|=|A_1-q_i|=\varepsilon$. An earlier revision argued this by saying the
arc ``cannot leave the region cut off by the two tangent lines'', which names
the conclusion rather than proving it and leaves the chord side unaddressed.

For (a) and (c) put
\[
   \varepsilon_0(L)=\tfrac13\min\Bigl\{
 \min_{i}\min_{k\neq i}|q_k-q_i|,\;
 \min_{i}\operatorname{dist}\bigl(q_i,\,\textstyle\bigcup\{\text{closed edges of }L\text{ not incident to }q_i\}\bigr),\;
 \min_i|\delta_i|,\;
 \min_{i,\,x\text{ a double point}}|x-q_i|
\Bigr\},
\]
and require $\varepsilon<\varepsilon_0(L)$; the index $i$ is bound by the outer
minimum in every term, an earlier revision having left it free in the first.
Each minimum over an empty index set is $+\infty$ --- which is how the last
one is read when $L$ has no double point. Every listed quantity is positive: the
corners are finitely many and distinct; a corner is at positive distance from a
non-incident closed edge because the edges are compact and a corner on one would
contradict the hypothesis that no corner lies on a non-incident edge; the edges have positive length; and a double point lies
in the relative interiors of two edges, so it is distinct from every corner.
With this $\varepsilon$ the discs $D_i$ are pairwise disjoint, each $D_i$ meets
no edge not incident to $q_i$, each $D_i$ contains no double point, and each
$D_i$ meets the two incident edges in the two sub-segments of length
$\varepsilon$ at $q_i$, and each contains the whole modification made at its
corner, the replacement arc lying within distance $\varepsilon$ of $q_i$. That
is clause~(e). Hence the rounding changes the curve only inside those
discs, where it meets nothing else, so no double point is created or destroyed
and the strands through each surviving double point are unchanged as arcs, with
their orientations; the over/under assignment is inherited, so the crossing signs
and the writhe are unchanged. That is (a) and (c). For (d), the total turning of
$L_\varepsilon$ is the sum of the turnings of its arcs, which is
$\sum_i\tau_i$ because the tangent is constant along the straight parts; and the
rotation number of a $C^1$ regular closed curve is $\frac1{2\pi}$ times its
total turning (Lemma~\ref{lem:turnlift}(iii)), which is also the value assigned to $L$
by Definition~\ref{def:rot}.
\end{proof}

\begin{lemma}[rotation bounds for uniform and one-dissent polygons]
\label{lem:uniformrot}
Here $\rot$ is as in Definition~\ref{def:rot}.
\begin{enumerate}
\item[(i)] Let $L$ be a closed polygon all of whose principal turns exist and are
positive. Then $\rot(L)\geq1$, with equality if $L$ has three corners. If all
turns are negative, $\rot(L)\leq-1$, again with equality in absolute value for
three corners.
\item[(ii)] Let $L$ be a closed polygon whose principal turns all exist, with
exactly one negative principal turn, of magnitude $\alpha\in(0,\pi)$, and let
$\Pi$ be the sum of the positive ones. Then
$\Pi-\alpha=2\pi r$ with $r=\rot(L)$ an integer, and $r\geq1$.
\end{enumerate}
\end{lemma}
\begin{proof}
(i) By Lemma~\ref{lem:turnlift}(ii),
$2\pi\rot(L)=\sum_i\tau_i$. A polygon has at least three corners and each
$\tau_i>0$, so the sum is positive and $\rot(L)>0$; being an integer,
$\rot(L)\geq1$. With three corners the sum is also below $3\pi$, forcing
$\rot(L)=1$. Orientation reversal in the same lemma gives both negative
claims.

(ii) The identity is Lemma~\ref{lem:turnlift}(ii). Suppose $r\leq0$. Then
$\Pi\leq\alpha<\pi$. Cut the traversal of $L$ immediately after the negative
corner and follow it once around: the direction of the edge is turned by the
successive principal turns, all of which are \emph{nonnegative} --- the
hypothesis names one negative turn and admits turns equal to zero, which $\Pi$
omits, and an earlier revision said "all of which are positive", which is false
as written when a zero turn is present --- and they sum to $\Pi<\pi$, so every
edge direction of $L$ lies in a closed angular interval $I$ of width $\Pi<\pi$.
A zero turn does not widen that interval, so the argument is unchanged. An
interval of width less than $\pi$ has a nonempty open dual: there is
$u\in\mathbb R^2$ with $\langle u,\delta\rangle>0$ for every direction
$\delta\in I$, hence for every edge of $L$. But the edges of a closed polygon
sum to zero, so
$0=\langle u,\sum_i\delta_i\rangle=\sum_i\langle u,\delta_i\rangle>0$, a
contradiction. Hence $r\geq1$, and $R:=|\rot(L)|=r$.
\end{proof}

\subsection{The negative curl}

\begin{lemma}[exact negative-curl replacement]\label{lem:curl}
Here $\rot$ is as in Definition~\ref{def:rot}.
Let $F$ be a connected $C^\infty$ immersed circle in the plane --- one
component, with finitely many transverse double points and no triple points ---
given with an oriented diagram, and let $p$ be a point of $F$ at which the
tangent points in a fixed direction $u$, isolated among such points, lying in
an embedded arc of $F$ that contains no double point and along which the tangent
turns strictly positively. Then $F$ may be modified inside a disc $\Delta$
meeting the rest of the diagram only in that arc, so that the resulting
diagram $F'$
\begin{enumerate}
\item[(i)] is again such an oriented diagram, satisfies
$P_{F'}(a,z)=P_F(a,z)$, and has the same double points outside $\Delta$, with
the same signs;
\item[(ii)] has no point of $\Delta$ at which the tangent equals $u$, and exactly
one at which it equals $-u$;
\item[(iii)] has exactly one double point inside $\Delta$, and it is negative;
\item[(iv)] satisfies $\rot(F')=\rot(F)-1$ and $w(F')=w(F)-1$.
\end{enumerate}
\end{lemma}
\begin{proof}
Work first in the rational model on $-2\leq t\leq2$,
\[
b(t)=\Bigl(\tfrac{3(t^2+7)}{11},\,-3t\Bigr),\qquad
c(t)=\bigl(t^2-1,\;t-t^3\bigr).
\]
The endpoints and endpoint tangent rays agree up to positive scale:
$b(\pm2)=c(\pm2)=(3,\mp6)$ and $b'(\pm2)=\frac3{11}c'(\pm2)$. The old arc turns
strictly positively and has exactly one tangent pointing straight down, at
$t=0$:
\[
\det\bigl(b'(t),b''(t)\bigr)=\tfrac{18}{11}>0,\qquad b'(0)=(0,-3).
\]
The replacement turns strictly negatively and has exactly one tangent pointing
straight up, at $t=0$:
\[
\det\bigl(c'(t),c''(t)\bigr)=-2(1+3t^2)<0,\qquad c'(0)=(0,1).
\]
In each arc the first derivative has vanishing first component only at $t=0$, so
neither arc has any other vertical tangent; this is (ii) in the model. The
replacement runs between the same endpoint rays along the complementary,
clockwise tangent path, so its compatible lift increment is the old one minus
$2\pi$; Lemma~\ref{lem:turnlift}(iii) gives (iv) for $\rot$. It has exactly one double point,
\[
c(-1)=c(1)=(0,0),\qquad c'(-1)=(-2,-2),\qquad c'(1)=(2,-2),
\]
and putting the later branch over the earlier one makes it negative, since a
crossing sign is the sign of the determinant of the overpassing tangent followed
by the underpassing tangent and
$\det\bigl(c'(1),c'(-1)\bigr)=-8<0$; that is (iii), and with it (iv) for $w$.
Being a single Reidemeister-I monogon, the replacement does not change the knot
type, which is (i).

To insert the model at $p$ we first record the projective fact it needs, in the
form the insertion argument requires.

\emph{Ordered-ray lemma.} Let $(r_1,r_2,r_3)$ and $(s_1,s_2,s_3)$ be triples of
rays from the origin such that $r_2$ lies in the open convex cone spanned by
$r_1,r_3$ and $s_2$ lies in the open convex cone spanned by $s_1,s_3$, with
$r_1,r_3$ and $s_1,s_3$ independent, and suppose in addition that for positive
representatives
\begin{equation}
   \sgn\det(\rho_1,\rho_3)=\sgn\det(\sigma_1,\sigma_3),
\label{eq:orderedray}
\end{equation}
a condition independent of the representatives chosen. Then there is a linear
map $A$ with $\det A>0$ carrying each $r_k$ into $s_k$.
\emph{Proof:} if both determinants in \eqref{eq:orderedray} are negative,
exchange the names $1$ and $3$ \emph{in both triples simultaneously}; this
preserves the correspondence $r_k\mapsto s_k$ and makes both determinants
positive. The linear map $A_0$ with $A_0\rho_1=\sigma_1$, $A_0\rho_3=\sigma_3$
then has positive determinant and carries $r_1,r_3$ to $s_1,s_3$. It carries the
open cone of $r_1,r_3$ onto that of $s_1,s_3$, so $A_0r_2$ is a ray $\sigma$ in
the latter cone. Writing $s_2=\mathbb R_{>0}(x\sigma_1+y\sigma_3)$ and
$\sigma=\mathbb R_{>0}(x'\sigma_1+y'\sigma_3)$ with $x,y,x',y'>0$, the map
$D=\mathrm{diag}(x/x',\,y/y')$ in the basis $(\sigma_1,\sigma_3)$ has positive
determinant, fixes $s_1$ and $s_3$, and carries $\sigma$ to $s_2$. Take
$A=DA_0$. $\square$

An earlier revision of this lemma omitted \eqref{eq:orderedray} and said the
positive determinants could be arranged ``after swapping the names once if
necessary''. A swap performed in one triple only destroys the correspondence it
is supposed to produce, and the peer refuted that step; the hypothesis is now
stated and the swap is simultaneous. In the application below the condition is
not inferred from the two cones: both determinants are computed and are
positive.

\emph{A positive-turn chart, and the cuts.} Parametrise the embedded arc by
arclength as $\gamma(s)$ with $\gamma(0)=p$ and unit tangent $\gamma'(0)=u$.
Because the tangent turns strictly positively, after shrinking the chart there
is a continuous strictly increasing lift $\theta$ with
\[
   \gamma'(s)=\cos\theta(s)\,u+\sin\theta(s)\,Ju,\qquad
   \theta(0)=0,\qquad |\theta(s)|<\tfrac\pi2,
\]
where $J$ is the positive quarter turn. Put $v=-Ju$, so that $(v,u)$ is a
positively oriented orthonormal basis, and set
\[
   \xi(s)=\langle\gamma(s)-p,\,v\rangle,\qquad
   \eta(s)=\langle\gamma(s)-p,\,u\rangle,
\]
so that $\xi'(s)=-\sin\theta(s)$ and $\eta'(s)=\cos\theta(s)$. Hence $\xi$
increases strictly for $s<0$ and decreases strictly for $s>0$, with a strict
maximum $\xi(0)=0$, while $\eta$ increases throughout. For every level
sufficiently close to $\xi(0)$ from below, strict monotonicity and the
intermediate value theorem give unique cuts $s_-<0<s_+$ with
$\xi(s_-)=\xi(s_+)$, and both tend to $0$ as the level tends to $\xi(0)$. Write
$p_\pm=\gamma(s_\pm)$. Equal transverse coordinates and monotone $\eta$ give
\[
   p_+-p_-=\ell u,\qquad \ell>0,\qquad \ell\to0 .
\]
Two side facts follow at once and are what the earlier revision lacked: the
central arc $\gamma((s_-,s_+))$ lies strictly in $\xi>\xi(s_-)$, and both
retained tails lie strictly in $\xi<\xi(s_-)$.

\emph{The affine fit, quantitatively.} In the positively oriented model basis
$v_0=(-1,0)$, $u_0=(0,-1)$ one has $q=\tfrac4{11}$ and
\[
   b'(-2)\in\mathbb R_{>0}(qv_0+u_0),\qquad
   b'(0)\in\mathbb R_{>0}u_0,\qquad
   b'(2)\in\mathbb R_{>0}(-qv_0+u_0).
\]
Write the two target endpoint tangents in the basis $(v,u)$ as
$T_-=\gamma'(s_-)=av+bu$ and $T_+=\gamma'(s_+)=-cv+du$; then
$a=-\sin\theta(s_-)>0$, $c=\sin\theta(s_+)>0$ and $b,d=\cos\theta(s_\mp)>0$,
and $a,c\to0$, $b,d\to1$ as the cuts approach $p$. The two outer determinants
are $\det(qv_0+u_0,-qv_0+u_0)=2q>0$ and $\det(T_-,T_+)=ad+bc>0$, so
\eqref{eq:orderedray} holds; it is computed, not inferred. Define
\[
   H=\frac{2}{\tfrac ba+\tfrac dc}=\frac{2ac}{bc+ad},\qquad
   x=\frac Hq,\qquad
   y=\frac{\tfrac ba-\tfrac dc}{q\bigl(\tfrac ba+\tfrac dc\bigr)}
    =\frac{bc-ad}{q(bc+ad)},
\]
and let $A$ be the linear map with $A u_0=u$ and $Av_0=xv+yu$. Then
\[
   A(qv_0+u_0)=\tfrac Ha\,T_-,\qquad
   A(-qv_0+u_0)=\tfrac Hc\,T_+,\qquad
   \det A=x=\frac{2ac}{q(bc+ad)}>0,
\]
each by substitution, and
\[
   1-(qy)^2=\frac{4abcd}{(bc+ad)^2}>0,
\]
so $|y|<1/q$ uniformly, while $x\to0$ as the cuts approach $p$: the map $A$
stays bounded. This is the quantitative statement that the earlier revision
replaced by an unquantified ``fixed unit-normalized part of $A$''.

Since $b(2)-b(-2)=12u_0$, put $B=\tfrac{\ell}{12}A$ and
$\Phi(z)=p_-+B\bigl(z-b(-2)\bigr)$. Then $\Phi(b(-2))=p_-$,
$\Phi(b(2))=p_-+\ell A u_0=p_-+\ell u=p_+$, and the endpoint tangent rays match
with positive scale. Because $\ell\to0$ while $A$ stays bounded, the diameter of
$\Phi(c([-2,2]))$ tends to $0$, so the cuts may be chosen so that the whole
inserted arc lies in a preassigned disc $\Delta$ about $p$ whose intersection
with the diagram is contained in the embedded chart --- the disc is fixed first
and the cuts afterwards.

\emph{Only the model's own double point is created.} In the basis $(v_0,u_0)$
the transverse coordinate of $c(t)$ is $1-t^2$, which equals $-3$ at both
endpoints, so the model's transverse excess over its endpoint chord is
$4-t^2>0$ for $-2<t<2$. Since $Bv_0=\tfrac{\ell}{12}(xv+yu)$ with $x>0$ and
$Bu_0=\tfrac{\ell}{12}u$, the physical transverse excess of the inserted arc
over the chord $[p_-,p_+]$ is exactly
\[
   \frac{\ell x}{12}\bigl(4-t^2\bigr)>0 .
\]
So the inserted interior lies strictly in $\xi>\xi(s_-)$ while both retained
tails lie strictly in $\xi<\xi(s_-)$: they meet only at the two joins. Together
with the shrinking of the previous paragraph, which keeps the inserted arc away
from every nonlocal part of the diagram, this \emph{proves} rather than assumes
that the only double point created is the model's own.

\emph{The replacement, and why the curve stays $C^\infty$ regular.} Replace the
sub-arc of $F$ between $p_-$ and $p_+$ by $\Phi\circ c$, modified near its two
endpoints as follows. At each join the two one-sided unit tangents already
agree, $\Phi c'(\pm2)$ being a positive multiple of $\Phi b'(\pm2)$ and that of
the old tangent, so the concatenation is $C^1$; it need not be $C^\infty$, and
the collar that repairs that is constructed here rather than described. An
earlier revision said only that one may ``interpolate the tangent angle by the
same device'', which does not say what curve results, and interpolating angles
moves the endpoint.

\emph{The collar, as a graph.} Work at the join $p_-$ first; the join at
$p_+$ is constructed after property (3) below, with its own signs printed
rather than left to "roles exchanged". In the chart coordinates
$(\xi,\eta)$ of the positive-turn chart --- $\xi=\langle\gamma-p,v\rangle$ with
$v=-Ju$, $\eta=\langle\gamma-p,u\rangle$ --- the old arc is a graph
$\xi=f(\eta)$, since $\eta'=\langle\gamma',u\rangle=\cos\theta>0$ there. Its
slope carries a sign that an earlier revision got wrong, and the sign is
load-bearing below, so it is computed:
$\xi'=\langle\gamma',v\rangle=-\sin\theta$, whence
\[
   f'(\eta)=\frac{\xi'}{\eta'}=-\tan\theta .
\]
For $\eta<\eta(p)$ the angle $\theta$ is strictly negative, $\theta$ being
strictly increasing with $\theta(0)=0$, so $f'>0$ there --- not $f'<0$, which is
what the earlier revision's $f'=\tan\theta$ gave. Let
$\eta_-=\eta(p_-)<\eta(p)$. The inserted arc leaves $p_-$ with the same unit
tangent, so near its start it is also a graph $\xi=g(\eta)$ with
$g(\eta_-)=f(\eta_-)$ and $g'(\eta_-)=f'(\eta_-)$; and because the model turns
strictly negatively its tangent angle only decreases from $\theta(s_-)<0$, so by
the same formula
\[
   g'>f'>0
\]
on that stretch. Two consequences are used below: $g>f$ on
$(\eta_-,\eta_-+\lambda]$, the two agreeing at $\eta_-$ with $g'>f'$; and at
every point of the collar the smaller slope is $f'(\eta)$, still positive. The
bound is pointwise of necessity: $f'=-\tan\theta$ \emph{decreases} along the
collar, $\theta$ increasing, so $0<f'(\eta)<f'(\eta_-)$ there and an endpoint
bound from the left end is unavailable --- an earlier revision claimed both
slopes bounded below by $f'(\eta_-)$, which $f'$ itself violates. Fix
$\lambda>0$ small enough that $[\eta_-,\eta_-+\lambda]$ lies both below
$\eta(p)$ and inside the stretch on which the inserted arc is a graph, and set,
on $[\eta_-,\eta_-+\lambda]$,
\[
   h(\eta)=\bigl(1-\varphi_\lambda(\eta)\bigr)f(\eta)+\varphi_\lambda(\eta)\,
   g(\eta),
   \qquad
   \varphi_\lambda(\eta)=\varphi\Bigl(\frac{\eta-\eta_-}{\lambda}\Bigr),
\]
with $\varphi$ the transition profile of Lemma~\ref{lem:rounding}. The modified
curve runs along the old arc up to $\eta_-$, along the graph of $h$ across the
collar, and along the inserted arc from $\eta_-+\lambda$ on. Three properties,
each proved from the formula.

\emph{(1) It is $C^\infty$ and regular, with positive speed.} $h$ is a $C^\infty$
function, and every derivative of $\varphi$ vanishes at $0$ and $1$, so $h$
agrees with $f$ to infinite order at $\eta_-$ and with $g$ to infinite order at
$\eta_-+\lambda$: the two joins are $C^\infty$. Parametrized by $\eta$, the
collar has velocity $h'(\eta)\,v+u$ in the frame $(v,u)$, of norm
$\sqrt{1+h'(\eta)^2}\geq1$, so its speed is positive and it is regular; being a
graph, it is embedded.

\emph{(2) No tangent of the collar equals $u$, and no smallness condition is
needed.} The tangent is parallel to $u$ exactly where $h'=0$, and
\[
   h'=(1-\varphi_\lambda)f'+\varphi_\lambda g'
      +\frac{\varphi'\bigl((\eta-\eta_-)/\lambda\bigr)}{\lambda}\,
       \bigl(g-f\bigr).
\]
Every term is nonnegative, and the sum of the first two is positive: it is the
convex combination $(1-\varphi_\lambda)f'+\varphi_\lambda g'$ of two positive
numbers --- each \emph{term} separately can vanish, $\varphi_\lambda$ being $0$
at the left end and $1$ at the right, which an earlier revision's "the first
two are positive" overlooked --- while the third term is
$\varphi'\cdot(g-f)/\lambda\geq0$, since $\varphi'\geq0$ and $g-f\geq0$ there,
the two agreeing at $\eta_-$ with $g'>f'$. Hence, pointwise,
\[
   h'(\eta)\geq(1-\varphi_\lambda)f'(\eta)+\varphi_\lambda g'(\eta)
   \geq f'(\eta)>0 ,
\]
the middle inequality because $g'>f'$, so $h'$ never vanishes and the collar
has no tangent equal to $u$; and none equal to $-u$ either, a graph over $\eta$ having
$\eta'>0$ throughout. An earlier revision, working from the wrong slope sign,
had both $f'$ and $g'$ negative, needed a Taylor bound on $g-f$ to control the
third term, and chose $\lambda$ small to make the sum negative. With the frame
computed the third term is on the same side as the other two and no choice of
$\lambda$ is required.

\emph{(3) It creates no double point and changes no turning.} Pointwise
$f\leq h\leq g$ on the collar --- the value is a convex combination and $g\geq f$
there, which is the separation the corrected slope chain supplies --- so the
collar lies in the region between the two graphs; that region is inside the
disc $\Delta$ and inside the embedded chart, where the only strands of the
diagram are the old arc --- whose collar stretch is removed --- and the inserted
arc. The collar is embedded by~(1) and meets those two only at its endpoints,
so no double point is created in it, and the model's own double point, at
$|t|=1$, is interior to the inserted arc and untouched. For the turning: the
collar's tangent angle is continuous, stays in $(-\pi/2,0)$ --- the slope
$h'$ being positive and finite, so the tangent angle $\theta$ has
$-\tan\theta=h'>0$ with $\cos\theta>0$ --- and takes at its ends exactly the
angles of the two arcs it joins, so the lift of the tangent
along the modified curve is the concatenation of the old lifts with an
interpolating stretch of the same endpoint values; the total tangent winding is
unchanged by its insertion. An earlier revision stopped at $C^1$
and the document then carried a lemma to recover smoothness. This sentence is what the downstream consumers of this
lemma read, the front construction below taking tangent directions of the
curve, and an intermediate revision of this proof deleted it while rewriting
the tail; it is restored here.

\emph{The collar at $p_+$.} The same three properties are now proved at the
other join, with its own signs printed: here the \emph{incoming} arc is the
inserted one and the \emph{outgoing} arc is the old one, and both slopes are
negative. Let $\eta_+=\eta(p_+)>\eta(p)$. On a stretch
$[\eta_+-\lambda_+,\eta_+]$ chosen above $\eta(p)$ and inside the stretches
on which both arcs are graphs, write $\xi=g(\eta)$ for the inserted arc and
$\xi=f(\eta)$ for the old arc extended backward from $p_+$; the two agree at
$\eta_+$ with common slope $m_+=-\tan\theta_+<0$, the angle
$\theta_+=\theta(s_+)$ lying in $(0,\tfrac\pi2)$. Just left of the join the
inserted arc turns strictly negatively into its endpoint, so its angle
exceeds $\theta_+$ there, while the old arc turns positively, so its angle
lies below $\theta_+$; by $\text{slope}=-\tan\theta$, decreasing in
$\theta$,
\[
   g'(\eta)<m_+<f'(\eta)<0
   \qquad\text{on }[\eta_+-\lambda_+,\eta_+),
\]
the outer inequality $f'<0$ because $\theta>0$ above $\eta(p)$. Integrating
$g'-f'<0$ backward from the common value at $\eta_+$ gives $g>f$ on
$[\eta_+-\lambda_+,\eta_+)$. Set
\[
   h=\bigl(1-\chi_{\lambda_+}\bigr)g+\chi_{\lambda_+}f,
   \qquad
   \chi_{\lambda_+}(\eta)=\varphi\Bigl(\frac{\eta-(\eta_+-\lambda_+)}{\lambda_+}\Bigr),
\]
with the same transition profile $\varphi$: the modified curve runs along the
inserted arc up to $\eta_+-\lambda_+$, along the graph of $h$ across the
collar, and along the old arc from $\eta_+$ on, the two joins $C^\infty$ by
the flat endpoint jets of $\varphi$ exactly as in (1). For the derivative,
\[
   h'=\bigl(1-\chi_{\lambda_+}\bigr)g'+\chi_{\lambda_+}f'
      +\frac{\varphi'\bigl((\eta-\eta_++\lambda_+)/\lambda_+\bigr)}{\lambda_+}
       \,\bigl(f-g\bigr),
\]
and the sign of $f-g$ in the cutoff term is exactly what the negativity of
the two slopes alone does not control: here $f-g\leq0$ by the separation just
integrated, and $\varphi'\geq0$, so the cutoff term is nonpositive, while the
convex part is a convex combination of two negative numbers. Hence,
pointwise,
\[
   h'(\eta)\leq\bigl(1-\chi_{\lambda_+}\bigr)g'(\eta)+\chi_{\lambda_+}f'(\eta)
   \leq f'(\eta)<0 ,
\]
the middle inequality because $g'<f'$, and the bound pointwise for the
mirrored reason: $f'$ varies along the collar and no endpoint bound is
available. So $h'$ never vanishes, and this collar --- a graph over $\eta$
with $\eta'>0$ throughout --- has no tangent equal to $u$ and none equal to
$-u$; its tangent angle stays in $(0,\tfrac\pi2)$, with
$-\tan\theta=h'<0$ and $\cos\theta>0$, and takes at its ends exactly the
angles of the two arcs it joins, so the lift concatenation of (3) applies to
it verbatim and the total tangent winding is unchanged. Pointwise
$f\leq h\leq g$, the value being a convex combination and $g\geq f$ there, so
the collar lies in the region between the two graphs, inside the disc
$\Delta$ and the embedded chart, and it is embedded and meets the two arcs
only at its endpoints: no double point is created, by the argument of (3)
unchanged.

\emph{The four conclusions, separately.} (ii): $c'(t)=(2t,1-3t^2)$ is parallel to
$u_0$ only at $t=0$, where $c'(0)=-u_0$; since $A$ is invertible with
$Au_0=u$, the inserted arc has no tangent $u$ and exactly one tangent $-u$, and
the retained tails have neither, their angles lying in $(-\tfrac\pi2,\tfrac\pi2)$
away from $0$.

(iii): if $c(s)=c(t)$ then $s=\pm t$, and the second coordinate forces
$2t(1-t^2)=0$, so the only pair of distinct parameters is $\{-1,1\}$: exactly
one double point. With the later branch over the earlier one its sign is that of
$\det(c'(1),c'(-1))=-8<0$, and $\det B>0$ preserves it.

(iv) for $\rot$: Let $R$ be the rotation carrying the positive orthonormal basis
$(v_0,u_0)$ to $(v,u)$ and put $C=R^{-1}A$. Relative to $(v_0,u_0)$,
$C$ has matrix
$\left(\begin{smallmatrix}x&0\\y&1\end{smallmatrix}\right)$ with $x>0$, so
$C_t=\left(\begin{smallmatrix}1-t+tx&0\\ty&1\end{smallmatrix}\right)$
joins $I$ to $C$ inside $\mathrm{GL}^+(2,\mathbb R)$. If $R_t$ rotates
through $t$ times any fixed angle representing $R$, first $C_t$ and then
$R_tC$ give a path $A_t$ from $I$ to $A=RC$ in that group. Apply
$z\mapsto A_tz/|A_tz|$ to the closed direction loop obtained by following
the tangent path of $c$ and the tangent path of $b$ backwards.
Lemma~\ref{lem:turnlift}(iii) therefore preserves their compatible endpoint
lift-increment difference. The tangent of $b$ follows the short positive arc between the
endpoint rays; the tangent of $c$ follows the complementary clockwise arc, since
$\det(c',c'')=-2(1+3t^2)<0$ and it passes through $-u_0$ at $t=0$; so their
relative winding is $-1$. The old central arc also follows the short positive
arc, its lifted sweep $\theta(s_+)-\theta(s_-)$ lying in $(0,\pi)$. Hence Lemma~\ref{lem:turnlift}(iii) says that the replacement changes the
closed curve's rotation by exactly $-1$. For
$w$: the removed arc was embedded and carried no crossing, so the writhe changes
by the one new negative crossing, that is by $-1$.

\emph{The polynomial conclusion, without a knot-category detour.}
The construction just completed also proves that $F'$ remains in the lemma's
input class.  It replaces one oriented arc by one $C^\infty$ regular oriented
arc, leaves the component count unchanged, and creates exactly one transverse
double point in a disc meeting no other strand.  Thus its double points remain
finite and transverse and no triple point is created.

Let $F^{\circ}$ be the oriented diagram obtained from $F'$ by deleting the
unique monogon in $\Delta$ by a Reidemeister-I move.  The move is available:
the preceding separation argument shows that the monogon is disjoint from
all other strands, and the crossing calculation shows it is the unique new
double point.  It replaces that monogon by one crossing-free regular arc in
$\Delta$, so $F^{\circ}$ is again a one-component immersed-circle diagram with
finite transverse double points and no triple points.  Reidemeister-I
invariance in Axiom~\ref{ax:homfly} gives $P_{F'}=P_{F^{\circ}}$.

Outside $\Delta$, the diagrams $F^{\circ}$ and $F$ coincide.  Inside
$\Delta$, each has one crossing-free oriented arc joining the same two
boundary points in the same traversal direction.  Match every marked
preimage outside that local traversal interval with the identical marked
preimage of the other diagram.  The local intervals contain no marked
preimage, so this bijection preserves cyclic order, crossing pairs,
over/under positions and signs: it is a record isomorphism.  Both underlying
curves are one-component immersed circles with transverse double points and
no triple points, so Axiom~\ref{ax:gausscode} says that the diagrams present
the same oriented link.  Axiom~\ref{ax:homfly} therefore gives
$P_{F^{\circ}}=P_F$, hence $P_{F'}=P_F$.  Equality outside $\Delta$ proves the
remaining clause of~(i).
\end{proof}

\begin{remark}[authorship of the construction above, and what its referee
found]\label{rem:curlauthor}
The proof of Lemma~\ref{lem:curl} from the ordered-ray lemma through the
geometric insertion is \textbf{not this author's}.  Its earlier form was
refuted by the peer: it swapped names in one ray triple, left the affine
shrinking unquantified, and did not exclude additional intersections.  The
printed repair is the peer's: the positive-turn chart, monotone equal-level
cuts, explicit $H,x,y$ affine fit, the bound
$1-(qy)^2=4abcd/(bc+ad)^2$, the transverse-excess separation
$\frac{\ell x}{12}(4-t^2)>0$, and the separate tangent, crossing and
rotation conclusions.

The HOMFLY conclusion is narrower than the former knot-isotopy conclusion.
It uses only Reidemeister-I invariance from Axiom~\ref{ax:homfly} and the
record comparison from Axiom~\ref{ax:gausscode}; no standard realization,
arc-isotopy theorem or category-change theorem is involved.

Theorem~\ref{thm:carrierfloor}(C)'s turn hypothesis was tightened for the same
reason: as printed it read ``every turn positive, or exactly one negative'',
which does not exclude a zero turn, while Lemma~\ref{lem:rounding} needs every
turn nonzero. The hypothesis now requires nonzero turns explicitly and says
``exactly one is negative and every other is positive''. Every consumer
supplies them: a carrier's corners are
guarded by (G1) at a vertex and (G5) at a smoothing site
(Definition~\ref{def:wind}).
\end{remark}

\subsection{Global orientation reversal}



\subsection{The floor}

\begin{theorem}[global reversal, rounding record, vertical tangencies and the carrier floor]\label{thm:carrierfloor}
\begin{enumerate}
\item[(R)] Here $\rot$ is as in Definition~\ref{def:rot}.
Let $K$ be an oriented knot and $-K$ the same knot with its orientation
reversed. Then $P_{-K}=P_{K}$. Moreover a diagram of $-K$ obtained by reversing
the orientation of a diagram of $K$ has the same crossing signs, the same
writhe, and rotation number of the underlying plane curve negated.
\item[(A)] Let $L$ and $D$ satisfy the hypotheses of Lemma~\ref{lem:rounding}. For every
$\varepsilon\in(0,\varepsilon_0(L))$, the construction in the proof of that
lemma returns \emph{one} curve and \emph{one} diagram at those data: the
junction inserted at the corner $q_i$ is determined by $\varepsilon$, by the
two incident unit directions and by the transition profile, which is fixed once
and for all in that proof and is not chosen per corner or per curve; the arc
length $\ell$ is then determined by the endpoint condition, and the rest of the
curve is $L$ itself. Write
\[
   \mathrm{Round}(L,D,\varepsilon)=(L_\varepsilon,D_\varepsilon)
\]
for that pair, the \emph{rounding record} of $(L,D)$ at $\varepsilon$, and call
$L_\varepsilon$ the rounded curve and $D_\varepsilon$ the rounded diagram. The
record exists to be quantified over: an earlier revision, written when the
profile was a per-use choice, had every consumer say "every curve the lemma
produces", which is a quantifier over an unnamed family and reads differently
in each consumer.
\item[(B)] Here $\rot$ is as in Definition~\ref{def:rot}.
Let $L$ be a diagrammatic closed polygon (Definition~\ref{def:diagrammatic})
whose principal turns all exist and are nonzero, carrying a diagram $D$, and
assume, after reversing the orientation of $L$ if necessary, that either every
principal turn is positive, or exactly one is negative. Put $R=|\rot(L)|$. Then
there are a direction $u\in S^1$ and an $\varepsilon_1>0$ such that \emph{for
every} $\varepsilon\in(0,\varepsilon_1)$ the rounded curve $L_\varepsilon$ of
the record $\mathrm{Round}(L,D,\varepsilon)$
from clause~(A) has exactly $R$ points at which its unit
tangent equals $u$ and exactly $R$ at which it equals $-u$; at each of them the
tangent crosses that direction in the positive sense. The hypotheses are stated
here rather than imported by "as in Lemma~\ref{lem:rounding}", and the
clearance is produced here rather than borrowed from that lemma's existential
one. The object is named rather than quantified over: the record is one curve
and one diagram at each $\varepsilon$, the profile being fixed once and for all,
and an earlier revision said "every curve the lemma produces" because at the
time the profile was a per-use choice.
\item[(C)] Here $\rot$ is as in Definition~\ref{def:rot}.
Let $D$ be an oriented knot diagram all of whose crossings are positive, with
writhe $w$, whose underlying plane curve is a closed polygon $L$ with all
principal turns existing, \emph{nonzero}, and of magnitude below $\pi$, with
finitely many double points, all transversal, with no triple points, none of
them a corner of $L$, and no corner of $L$ lying on a non-incident edge; and let
$R$ be the absolute value of its Whitney rotation number.  The
finiteness/transversality and the two corner conditions are what
Lemma~\ref{lem:rounding} asks of its input.  The no-triple condition repeats the
shared-image clause of Definition~\ref{def:diagrammatic}; it is preserved by
rounding and is needed below when diagram records are compared.  These
conditions are stated here so that no generic parent polygon is assumed. Assume that, after reversing the
orientation if necessary --- which by clause~(R) changes neither
$P_D$, nor the crossings and their signs, nor $w$, nor $R$ --- either every
principal turn is positive, or exactly one is negative and every other is
positive. Then
\begin{equation}
\min\deg_a P_D(a,z)\;\geq\;1-w-R,
\label{eq:floor}
\end{equation}
and the same bound holds for $f_D(a)=[z^0]P_D(a,z)$ whenever $f_D\neq0$.
\item[(D)] Here $\rot$ is as in Definition~\ref{def:rot}.
Let $P$ be generic, let $S\in\Ind(G_P)$ and let $L$ be a carrier of $S$ carrying
no residual piece, so that $P_{S,L}=1$ and $w_{S,L}=0$ by
Definition~\ref{def:X1}'s empty conventions --- an earlier revision wrote those
two symbols with $S$ never introduced --- and suppose that all its principal turns are nonzero and that,
after reversing the orientation if necessary, either every turn is positive or
exactly one is negative. Then $R(L)\geq1$ and
$\min\deg_a P_{S,L}=0\geq 1-w_{S,L}-R(L)$.
\end{enumerate}
\end{theorem}
\begin{proof}
\emph{Clause (R).}
Reversing the orientation of all components of a link sends a skein triple
$(L_+,L_-,L_0)$ to a skein triple: the sign of a crossing is
$\sgn\det(u_{\mathrm{over}},u_{\mathrm{under}})$ and reversing both strands
replaces both vectors by their negatives, leaving the determinant unchanged, and
the oriented smoothing of the reversed diagram is the reverse of the oriented
smoothing. So $L\mapsto P_{-L}$ satisfies the defining skein of
Axiom~\ref{ax:homfly}, and it sends the unknot to $1$; by the uniqueness clause
it equals $L\mapsto P_L$. The diagram statements have the same scope. Every crossing sign is a
determinant of two reversed vectors, so the writhe is unchanged. For a
polygonal underlying curve every principal turn changes sign and
Lemma~\ref{lem:turnlift}(ii) negates rotation. For a $C^1$ regular curve,
if $\theta$ lifts its normalised tangent, then $\theta(1-s)+\pi$ lifts the
reversed tangent with the negative endpoint increment, as in
Lemma~\ref{lem:turnlift}(i).

\emph{Clause (A).} The transition profile is fixed once and for all in the proof of
Lemma~\ref{lem:rounding}.  At each corner, $\varepsilon$ and the two incident
unit directions determine that profile's junction, and the endpoint condition
determines its length $\ell$; outside the rounding discs the curve is $L$, and
the diagram data are inherited.  Thus the named construction returns one
specified pair at every allowed $\varepsilon$.  This asserts uniqueness of the
fixed construction's record, not uniqueness among all possible smoothings.

\emph{Clause (B).}
Reversing the orientation of $L$ negates every principal turn and the signed rotation and
preserves every self-crossing sign, so perform the allowed normalization once. In the
uniform case $\rot(L)=R\geq1$ by Lemma~\ref{lem:uniformrot}(i); in the one-dissent case the same holds by Lemma~\ref{lem:uniformrot}(ii).

Take $\varepsilon_1=\varepsilon_0(L)$, the clearance of Lemma~\ref{lem:rounding}.
Choose $u\in S^1$ so that neither $u$ nor $-u$ is an edge direction and, in the
one-dissent case, neither lies in the negative turn's closed swept arc $A$, of
length $\alpha:=|\tau_-|<\pi$. Such $u$ exists: in the one-dissent case $A\cup(-A)$
has length at most $2\alpha<2\pi$, while in the uniform case there is no such arc; the additional
forbidden set of edge directions and their antipodes is finite. This choice depends only on $L$. Fix an argument $\phi$ of $u$.

Fix $\varepsilon\in(0,\varepsilon_1)$, write
$(L_\varepsilon,D_\varepsilon)=\mathrm{Round}(L,D,\varepsilon)$, and let $T$ be
its unit-tangent map. Parameterize the rounded traversal by $s\in[0,1]$, with
$s=0$ in the interior of a straight part; because $u,-u$ are not edge directions, its tangent
is congruent to neither $\phi$ nor $\phi+\pi$. Concatenating straight-part arguments
with the junction lifts supplied by Lemma~\ref{lem:rounding}(b) gives a continuous
$\theta\colon[0,1]\to\mathbb R$, starting at an argument $\theta(0)$ of $T(0)$.
Since $T(1)=T(0)$, write $\theta(1)=\theta(0)+2\pi m$ for an integer $m$.

Choose a direction $r$ not parallel to any edge, with argument $\gamma$. The
change across junction $i$ of $\lfloor(\theta-\gamma)/(2\pi)\rfloor$ is $+1$,
$-1$ or $0$ in exactly the three determinant cases defining $\epsilon_i$ in
Definition~\ref{def:rot}. Indeed the same clause makes the lifted change
$\tau_i$, and $|\tau_i|<\pi$ permits at most one crossed level, in its
signed direction; there is no change on a straight part. Telescoping and
$\lfloor q+m\rfloor-\lfloor q\rfloor=m$ give
$m=\sum_i\epsilon_i=\rot(L)=R$. Thus
\[
 \theta(1)=\theta(0)+2\pi R.                                  \tag{*}
\]

Fix $\beta=\phi$ or $\phi+\pi$. Because $u,-u$ are not edge directions, no
level $\beta+2\pi k$ is met on a straight part or at a junction endpoint; by
the choice of $u$, none is met on the decreasing junction. Hence floor jumps
occur only in increasing junctions and never elsewhere. Lemma~\ref{lem:rounding}(b)
makes each occurrence unique with positive angular derivative. Put
\[
 x=\frac{\theta(0)-\phi}{2\pi},\qquad
 y=\frac{\theta(0)-(\phi+\pi)}{2\pi}.
\]
For $u$, (*) and the integer floor identity give, separately,
\[
 \#T^{-1}(u)=\lfloor x+R\rfloor-\lfloor x\rfloor=R.
\]
For $-u$, the same argument at the other levels gives
\[
 \#T^{-1}(-u)=\lfloor y+R\rfloor-\lfloor y\rfloor=R.
\]
There are finitely many junctions, so the preimages are finite; their positive derivatives prove the assertion for arbitrary $\varepsilon$.

\emph{Clause (C).} Fix $\varepsilon\in(0,\varepsilon_1)$ with $\varepsilon_1$ the clearance of
clause~(B) and pass to the rounding record
$(L_\varepsilon,D_\varepsilon)=\mathrm{Round}(L,D,\varepsilon)$ of
clause~(A) --- one curve and one diagram, not a member
of an unnamed family. By Lemma~\ref{lem:rounding}(c),(d) the diagram, its
crossings, its writhe and its rotation number are unchanged, and by
clause~(B) there is a direction $u$ met by the tangent of
$L_\varepsilon$ exactly $R$ times, each time positively, and likewise for
$-u$. Rotate coordinates so
that $u$ is the downward vertical. The rounded diagram then has exactly $R$
downward vertical tangencies and $R$ upward ones.

Switch every crossing. The result is a diagram $\overline D$ of the mirror knot,
with the same underlying plane curve --- hence the same tangencies --- and with
writhe $-w$; every crossing of $\overline D$ is negative. Apply
Lemma~\ref{lem:curl} at each of the $R$ downward vertical tangencies; they are
isolated, lie in the interiors of positively turning rounding arcs
(clause~(B)), and the rounding discs are pairwise disjoint and
meet no crossing, by the choice of $\varepsilon$ in the proof of
Lemma~\ref{lem:rounding}, so the $R$ discs are
disjoint and each replacement leaves the others intact. Call the result $T$. By
Lemma~\ref{lem:curl}, $P_T=P_{\overline D}$, $T$ has no downward vertical
tangency, and it has writhe
\[
w(T)=-w-R .
\]
Every crossing of $T$ is negative: the old ones by the switch, the $R$ new ones by
Lemma~\ref{lem:curl}(iii).

We now construct the transverse lift rather than import a
front criterion.  Write the smooth regular parametrisation of $T$ as
$s\mapsto(x(s),z(s))$, put
\[
 v(s)=\sqrt{x'(s)^2+z'(s)^2},\qquad
 y_0(s)=-\frac{x'(s)}{v(s)+z'(s)},
\]
and use the convention that, in the $xz$ projection, the branch with smaller
$y$-coordinate is drawn over.  Since $T$ has no downward vertical tangency,
$v+z'>0$ everywhere: equality would force $x'=0$ and $z'<0$.  Thus $y_0$ is
smooth and periodic, and direct algebra, with no division by $x'$, gives
\[
 z'-y_0x'=z'+\frac{x'^2}{v+z'}
 =z'+\frac{v^2-z'^2}{v+z'}=v>0.                 \tag{*}
\]

It remains only to impose the crossing order.  At a crossing branch there is
no vertical tangency, so put $m=z'/x'$.  The strict inequality in~(*) permits
$y<m$ on a rightward branch and $y>m$ on a leftward branch.  For every crossing
choose admissible constants $c_{\rm O}<c_{\rm U}$ at its over- and underpassing
preimages.  Such a choice is automatic if the overpass points right: after
choosing any admissible $c_{\rm U}$, take
$c_{\rm O}<\min\{m_{\rm O},c_{\rm U}\}$.  If both point left, choose any
$c_{\rm O}>m_{\rm O}$ and then
$c_{\rm U}>\max\{m_{\rm U},c_{\rm O}\}$.  The only remaining case has
$x'_{\rm O}<0<x'_{\rm U}$.  Every crossing of $T$ is negative, and here
\[
 0>\det(t_{\rm O},t_{\rm U})
   =x'_{\rm O}x'_{\rm U}(m_{\rm U}-m_{\rm O}).
\]
Since $x'_{\rm O}x'_{\rm U}<0$, this says $m_{\rm O}<m_{\rm U}$; choose
$m_{\rm O}<c_{\rm O}<c_{\rm U}<m_{\rm U}$.  These are exactly the two allowed
half-lines and the required over/under order.

There are finitely many crossing preimages.  Around each one $s_j$, choose
pairwise disjoint cyclic parameter intervals $(\alpha_j,\beta_j)$ so small that
$z'-c_jx'>0$ throughout.  Using the explicit flat transition profile
$\varphi$ displayed in the proof of Lemma~\ref{lem:rounding}, define a smooth
bump on the circle by
\[
 \psi_j(s)=
 \begin{cases}
 \varphi\!\left((s-\alpha_j)/(s_j-\alpha_j)\right),&
       \alpha_j\le s\le s_j,\\
 \varphi\!\left((\beta_j-s)/(\beta_j-s_j)\right),&
       s_j\le s\le\beta_j,\\
 0,&\text{otherwise}.
 \end{cases}
\]
All endpoint derivatives are zero, so this is smooth, equals one at $s_j$ and
has support in that interval.  Set
\[
 y=y_0+\sum_j\psi_j(c_j-y_0).
\]
The supports are disjoint.  On the $j$-th one,
\[
 z'-yx'=(1-\psi_j)(z'-y_0x')+\psi_j(z'-c_jx')>0,
\]
and outside them this is~(*).  Hence $y$ is smooth and periodic and the lift
$s\mapsto(x(s),y(s),z(s))$ is positively transverse to $dz-y\,dx$.

Finally the lift is regular because its $xz$ projection is immersed.  If two
lifted points agreed, their projected points would agree; the only distinct
parameters with that property are the two preimages of a listed double point,
and there $c_{\rm O}<c_{\rm U}$.  Thus the lift is injective, hence, by
compactness of the parameter circle, an embedded transverse knot.  Its front
and its smaller-$y$ over/under assignments are exactly $T$.

By Axiom~\ref{ax:etnyre}, $\operatorname{sl}(T)=w(T)=-w-R$. Since
$P_T=P_{\overline D}$, Axiom~\ref{ax:slbound} applied to $T$ gives
\[
\max\deg_a P_{\overline D}\;\leq\;-\operatorname{sl}(T)-1\;=\;w+R-1 .
\]
For every oriented-link isotopy class $K$, put
$\mathcal M(K)=P_{\overline K}(a,z)$. Mirroring exchanges $L_+$ with $L_-$
and fixes $L_0$, so the defining skein of Axiom~\ref{ax:homfly} on the
mirrored triple, multiplied by $-1$, is
\[
 a^{-1}\mathcal M(L_+)-a\mathcal M(L_-)=-z\,\mathcal M(L_0),
 \qquad \mathcal M(\bigcirc)=1.
\]
The substitution $(a,z)\mapsto(a^{-1},-z)$ is involutive, so applying it to
the uniqueness clause of Axiom~\ref{ax:homfly} gives uniqueness for these two
conditions. The assignment $K\mapsto P_K(a^{-1},-z)$ satisfies them; hence
$P_{\overline K}(a,z)=P_K(a^{-1},-z)$. As a Laurent-ring automorphism, the
substitution negates every $a$-exponent and cannot cancel a nonzero extreme
coefficient. Hence
$\max\deg_a P_{\overline D}=-\min\deg_a P_D$, and the displayed bound becomes
$\min\deg_a P_D=-\max\deg_a P_{\overline D}\geq1-w-R$, which is
\eqref{eq:floor}. Finally, the support of $f_D$ is contained in the support of
$P_D$, so the same bound holds for it when it is nonzero.

\emph{Clause (D).} In the uniform case $|\rot(L)|\geq1$ by Lemma~\ref{lem:uniformrot}(i); in the
one-dissent case $\rot(L)\geq1$ by Lemma~\ref{lem:uniformrot}(ii). Either way
$R(L)\geq1$, so $1-0-R(L)\leq0=\min\deg_a1$. In particular the hypothesis is met by every carrier of a live support
(Definition~\ref{def:wind}), and it is met by the marked halves of the
one-dissent geometry produced in the vertex-on-edge section, where that is
proved at the point of use.
\end{proof}

\begin{remark}[what the floor does and does not assert]\label{rem:floorscope}
Theorem~\ref{thm:carrierfloor}(C) is a lower bound on a minimal degree; it asserts
no sharpness. For a live carrier the slot of Definition~\ref{def:X1} is exactly
$1-w_{S,L}-R(L)$, so exactly one of two things holds: the slot is the bottom
supported degree of the row, or the slot lies strictly below the support and its
coefficient is zero. Both are consistent with \eqref{eq:floor} and the proofs of
Section~\ref{sec:s7} use only the second reading, namely that nothing lies below
the slot.
\end{remark}

% The wall dictionary.  Split out of the fixed-count transport section so that
% the sections that consume it --- the spike, the vertex-on-edge laws --- come
% after it in document order and not before.
\section{Relative general position and the wall dictionary}\label{sec:dictionary}

A path between generic polygons meets the discriminant, and every law proved
later is a statement about one kind of crossing. This section fixes the regular
locus, replaces an arbitrary path by a finite factor-simple path, and proves
that each relevant hit carries exactly one of five bundles. Nothing here
consumes a wall law; the sections that prove those laws consume this one.

\subsection{The regular locus}


\begin{lemma}[the regular locus is open, and generic polygons lie in it]
\label{lem:regularopen}
$\mathcal R_n$ is open in $(\mathbb R^2)^n$, and every generic polygon
(Definition~\ref{def:generic}(A)) lies in $\mathcal R_n$. Consequently a compact
path in $\mathcal R_n$ has a positive minimum edge length and positive distance
from the complement.
\end{lemma}
\begin{proof}
Write $B_i=\{d_i=0\}$ and
$C_i=\{d_{i+1}=-\lambda d_i\text{ for some }\lambda>0\}$. Each $B_i$ is closed.
The union $B_i\cup B_{i+1}\cup C_i$ is closed: let $P^{(k)}$ be a sequence in it
converging to $P$, and suppose $P\notin B_i\cup B_{i+1}$, so $d_i,d_{i+1}\neq0$
at $P$. Then for large $k$ the same holds at $P^{(k)}$, so $P^{(k)}\in C_i$ and
$d_{i+1}^{(k)}=-\lambda_kd_i^{(k)}$ with $\lambda_k>0$; taking norms,
$\lambda_k=|d^{(k)}_{i+1}|/|d^{(k)}_i|\to|d_{i+1}|/|d_i|>0$, and passing to the
limit gives $d_{i+1}=-\lambda d_i$ with $\lambda>0$, i.e.\ $P\in C_i$. The set
$B_i\cup C_i$ alone need \emph{not} be closed --- with $d_i=(1,0)$ fixed and
$d_{i+1}=(-1/k,0)$, every term lies in $C_i$ while the limit lies only in
$B_{i+1}$ --- which is why the triple union is what the paragraph proves
closed. Since the indices are cyclic,
$\bigcup_i(B_i\cup C_i)=\bigcup_i(B_i\cup B_{i+1}\cup C_i)$, a finite union of
the closed sets just proved closed; that union is the complement of
$\mathcal R_n$, so $\mathcal R_n$ is open. An earlier revision called
$B_i\cup C_i$ closed on the way. A generic polygon has
$\det(d_{i-1},d_i)\neq0$ for all $i$ by (G1), so no consecutive pair is
proportional at all. The last sentence is compactness.
\end{proof}




\subsection{Finite relative general position, and which bundles occur}

\begin{lemma}[canonical factors: definition and finite relative general position]\label{lem:relgp}
Fix $n\geq3$.
On $\mathbb R^{2n}$ define, for three distinct vertex indices $a,b,c$, for two
\emph{remote} edges $e\neq f$, and for three pairwise remote edges $e,f,g$,
\[
   \Delta_{abc}=\det\bigl(p_b-p_a,\ p_c-p_a\bigr),\qquad
   H_{e,f}=\det(d_e,d_f),\qquad
   T_{e,f,g}=\det\bigl(\ell_e,\ell_f,\ell_g\bigr),
\]
where $\ell$ is the line-coefficient row of Definition~\ref{def:guarded}. A
\emph{canonical factor} is a polynomial in one of these three families with
indices as restricted here. Two canonical factors are \emph{associates} if one
is a nonzero real multiple of the other; the associate classes are what the
counting below uses, and an earlier revision spoke of canonical factors without
ever defining either the list or the equivalence.

The index restrictions are not decoration. With $a$, $b$, $c$ not distinct,
$\Delta$ is the zero polynomial; likewise $\det(d_e,d_e)=0$, and
$\det(d_e,p_i-p_e)=0$ when $p_i$ is an endpoint of $e$. An earlier revision
printed an unrestricted list and called every member of it irreducible, which
is false at exactly those indices.

The turn and the vertex-on-line determinant are one family and not two:
\[
   \det\bigl(p_b-p_a,\ p_c-p_b\bigr)=\det\bigl(p_b-p_a,\ p_c-p_a\bigr)
   =\Delta_{abc},
\]
identically, so the turn at $b$ of the two edges $ab$, $bc$ \emph{is} the
determinant that puts $c$ on the line of $ab$. Every (G1) and every (G2) is a
$\Delta$, and so is every direction determinant of two \emph{consecutive}
edges; $H$ is reserved for remote pairs, where four distinct vertices occur.

\begin{enumerate}
\item[(A1)] Every member of $\mathcal G$ is, up to sign, a canonical factor
as defined above or a product of two: $\mathrm{G1}_j=\Delta_{(j-1)j(j+1)}$,
$\mathrm{G2}_{e,i}=\Delta_{e_0e_1i}$ with $e=(e_0,e_1)$ and $i$ not an endpoint
of $e$, $\mathrm{G5}_{e,f}=H_{e,f}$, and $\mathrm{G3}_{e,f,g}=T_{e,f,g}$; the
order predicate $\mathrm{G4}_{e;f,g}$ equals $-T_{e,f,g}$ when $f$ and $g$ are
remote to each other, and $-\Delta_{e_0e_1M}\Delta_{f_0Mg_1}$ when they meet at
the vertex $M$, which is a product of two canonical factors and the only
member that is not a single one.
\item[(A2)] Every canonical factor is irreducible in $\mathbb R[x_1,\dots,x_{2n}]$.
\item[(A3)] Two canonical factors are associates if and only if they have the same
variable support and one is a signed reindexing of the other; explicitly, the
associate class of $\Delta_{abc}$ is the six orderings of $\{a,b,c\}$ with the
sign of the permutation, that of $H_{e,f}$ is $\{H_{e,f},-H_{e,f}\}$, and that
of $T_{e,f,g}$ is its six row orderings with sign.
\item[(B)] Let $\mathcal C$ contain one representative of each associate class of the
canonical factors occurring in members of $\mathcal G$. Let
$\gamma\colon[0,1]\to\mathcal R_n$ be continuous with generic endpoints and
let $\delta>0$. There is a continuous piecewise-affine path
$\tilde\gamma\colon[0,1]\to\mathcal R_n$, with the same endpoints and within
$\delta$ of $\gamma$, whose non-generic parameters form a finite set. Every
such parameter $t_0$ lies in the interior of a nonconstant segment on which one
scalar vertex coordinate $x_j$ moves, and has this certificate:
\begin{enumerate}
\item[(B1)] exactly one $F\in\mathcal C$ vanishes at
$x=\tilde\gamma(t_0)$;
\item[(B2)] $(F\circ\tilde\gamma)'(t_0)\neq0$;
\item[(B3)] no two vertices coincide at $x$.
\end{enumerate}
Every junction outside the two endpoint collars has all canonical factors
nonzero and has both scalar-coordinate differences nonzero for every pair of
vertices. The endpoint collars lie in $\mathcal U_n$. Thus corners of the path
are wall-free, and every wall germ is an affine one-parameter family with one
simple canonical-factor root.
\end{enumerate}
\end{lemma}
\begin{proof}
For (A1), use Definition~\ref{def:guarded} together with the displays above and the
factorization of a concurrency determinant whose last two edges are
consecutive.

For (A2), consider $\Delta$ and $H$: the quadratic form $q(u,w)=u_1w_2-u_2w_1$ has rank
four, and a reducible quadratic form is a product of two linear forms and so has
rank at most two; hence $q$ is irreducible. Both are $q$ pulled back along a
surjective linear map of the vertex coordinates --- $(p_b-p_a,\,p_c-p_a)$ in one
case, $(d_e,d_f)$ in the other, and both maps are surjective precisely because
the index restrictions make the two arguments independent. A pullback $q\circ L$
along a surjective linear $L$ is again irreducible: a factorization $q\circ L=AB$
has both factors independent of $\ker L$, because a divisor of a polynomial of
degree zero in a variable has degree zero in it, so $A=A'\circ L$ and
$B=B'\circ L$ and $q=A'B'$.

For $T$: it has degree one in each of the six two-vector blocks
$p_e,d_e,p_f,d_f,p_g,d_g$, which are independent coordinates because the edges
are pairwise remote. Suppose $T=AB$ with both factors nonconstant. Since $T$ has
degree one in $d_g$, one factor, say $B$, is free of $d_g$. Expanding the
determinant along the row of $g$ and collecting,
\[
   T=d_{g,y}\bigl(p_{g,x}M_C-M_A\bigr)+d_{g,x}\bigl(M_B-p_{g,y}M_C\bigr),
\]
where $M_A,M_B,M_C$ are the $2\times2$ cofactors of that row and involve only
the blocks of $e$ and $f$. As $B$ is free of $d_g$ it divides both bracketed
polynomials; the first is free of $p_{g,y}$ and the second is free of
$p_{g,x}$, so $B$ is free of both, hence of every block of $g$. Comparing the
coefficients of the four distinct monomials $d_{g,y}p_{g,x}$, $d_{g,y}$,
$d_{g,x}$, $d_{g,x}p_{g,y}$ then gives $B\mid M_A$, $B\mid M_B$ and
$B\mid M_C$. But $M_C=\det(d_e,d_f)=H_{e,f}$, which is irreducible by the
previous paragraph, so $B$ is a unit or an associate of $H_{e,f}$; and
$H_{e,f}$ does not divide $T$, since setting $d_e=d_f=(1,0)$ makes $H_{e,f}$
vanish while $T$ becomes $-d_{g,y}(p_{e,y}-p_{f,y})$, not the zero polynomial.
So $B$ is a unit, contrary to assumption.
For (A3), first observe that associates have the same zero set and the same variable support, so it is
enough to consider two canonical factors of the same family on the same
support. For $\Delta_{abc}$ the six orderings of $a,b,c$ give the same polynomial up to
the sign $\sgn\sigma$ of the permutation $\sigma$ relating them, and nothing
else: $\Delta_{\sigma(a)\sigma(b)\sigma(c)}=\sgn(\sigma)\,\Delta_{abc}$,
a determinant being alternating in its rows. For $T_{e,f,g}$ the same holds of
the six orderings of the three line-coefficient rows,
$T_{\sigma(e)\sigma(f)\sigma(g)}=\sgn(\sigma)\,T_{e,f,g}$. So two orderings of
one triple are associates by a unit $\pm1$, which is what this clause needs; an
earlier revision said "the sign of the permutation" without binding $\sigma$ or
displaying either identity. Two canonical
factors on \emph{different} supports are never associates: a vertex coordinate
occurring in one and not in the other cannot be matched by a scalar.

Two same-support pairs remain, and they are the only ones. The support of a
remote edge pair $\{e,f\}$ is $\{e,e+1\}\sqcup\{f,f+1\}$, two disjoint blocks of
two cyclically adjacent indices. If a second remote pair has the same support,
that four-element set splits into two adjacent blocks in a second way, so some
index is blocked with its \emph{other} neighbour: $f=e+2$ and, applying the same
to the other block, $e=f+2$. Then $e=e+4$ in $\mathbb Z_n$, so $n$ divides $4$,
and $n\geq3$ leaves $n=4$. The same count settles triples: the support of a
pairwise remote triple is three disjoint adjacent blocks, a second such
decomposition pairs every index with its other neighbour, and following the
chain $a\to a+1\to a+2\to\cdots$ around returns to $a$ after six steps, so
$n=6$.

At $n=4$ the two pairs are
$H_{0,2}=\det(p_1-p_0,p_3-p_2)$ and $H_{1,3}=\det(p_2-p_1,p_0-p_3)$: the
monomial $x_0y_2$ occurs in the first with coefficient $1$ and in the second
with coefficient $-1$, while $x_0y_1$ occurs in the second and not in the first,
so no scalar carries one to the other. At $n=6$ the two triples are the
alternating ones, and the hexagon
$p_0=(1,0)$, $p_1=(2,0)$, $p_2=(0,1)$, $p_3=(0,2)$, $p_4=(1,1)$, $p_5=(2,2)$
gives $T_{0,2,4}=0$ and $T_{1,3,5}=2$, so those two are not associates either.


For (B), put $N=2n$ and write the scalar coordinates as $x_1,\ldots,x_N$. The set
$\mathcal C$ is finite by (A1). We first print the
one-variable algebra used to choose the path.

Fix $F\in\mathcal C$ depending on $x_j$, and put
$K_j=\operatorname{Frac}(\mathbb R[x_1,\ldots,\widehat{x_j},\ldots,x_N])$.
Irreducibility of $F$ in the full polynomial ring
(A2) and Gauss's lemma make it irreducible in
$K_j[x_j]$. Characteristic zero and dependence on $x_j$ give
$\gcd(F,\partial_jF)=1$ there. A Bezout identity, after clearing its finitely
many denominators, therefore has the form
\[
 A_{F,j}F+B_{F,j}\partial_jF=D_{F,j},
 \qquad 0\neq D_{F,j}\in
 \mathbb R[x_1,\ldots,\widehat{x_j},\ldots,x_N].
\]
If nonassociate $F,G\in\mathcal C$ both depend on $x_j$, they are coprime in
$K_j[x_j]$: otherwise their primitive irreducible representatives would become
associates over $K_j$, and clearing the scalar would make them associates in
the full ring. Clearing a second Bezout identity gives
\[
 A_{FG,j}F+B_{FG,j}G=R_{FG,j},
 \qquad 0\neq R_{FG,j}\in
 \mathbb R[x_1,\ldots,\widehat{x_j},\ldots,x_N].
\]
At a specialization where $D_{F,j}\neq0$, the univariate specialization of
$F$ has no repeated root; where $R_{FG,j}\neq0$, the specializations of $F$
and $G$ have no common root. A factor independent of $x_j$ is constant on that
coordinate segment and is kept nonzero at its initial junction. This includes
the two exceptional same-support comparisons: the $n=4$ pair of $H$ factors
and the $n=6$ alternating pair of $T$ factors are nonassociate by the explicit
monomial and hexagon tests in (A3), so their Bezout
certificates occur in the same finite list.

We now choose the whole path at once, including its parametrization. Use the
same norm as in the asserted $\delta$-bound. Compactness of $\gamma([0,1])$
and openness of $\mathcal R_n$ (Lemma~\ref{lem:regularopen}) give an open
cover of $\gamma([0,1])$ by axis-aligned cubes whose closures lie in
$\mathcal R_n$ and whose diameters are less than $\delta$. At the two
endpoints choose the cover cubes inside $\mathcal U_n$, using openness of
$\mathcal U_n$ (Proposition~\ref{prop:chamberinv}(i)). Uniform continuity of $\gamma$
(equivalently, a Lebesgue-number subdivision of the pulled-back cover) gives
a partition
\[
 0=t_0<t_1<\cdots<t_m=1
\]
and cubes $Q_0,\ldots,Q_{m-1}$ from this cover such that
$\gamma([t_i,t_{i+1}])\subset Q_i$; refine at the ends so that
$Q_0,Q_{m-1}\subset\mathcal U_n$. In particular
$\gamma(t_i)\in Q_{i-1}\cap Q_i$, so every consecutive overlap is nonempty.
Put $q_0=\gamma(0)$ and $q_m=\gamma(1)$, and choose each internal waypoint
$q_i$ in $Q_{i-1}\cap Q_i$. Then $q_i,q_{i+1}\in Q_i$, so every
coordinate-order hybrid between them remains in $Q_i$.

Call the intervals with $1\leq i\leq m-2$ central. Their waypoints are chosen
jointly in the nonempty open product of the overlaps. Impose the following
finite list of nonvanishings for every central
coordinate segment and every one of its hybrids: $F$ at the hybrid for each
$F\in\mathcal C$; the corresponding $D_{F,j}$ and $R_{FG,j}$ on the fixed
coordinates of the segment; the coordinate step
$q_{i+1,j}-q_{i,j}$; and, at every hybrid, both coordinate differences for
every pair of vertices. Each condition pulls back to a nonzero polynomial on
the waypoint product. Indeed, a hybrid or fixed-coordinate map selects distinct
scalar coordinates from the two adjacent variable waypoints, so its image
contains an open coordinate box; a nonzero polynomial cannot vanish on that
box. The step and vertex-difference conditions are visibly nonzero linear
polynomials. Their finite product is nonzero because the polynomial ring is an
integral domain, and a nonzero polynomial cannot vanish on a nonempty open set
(all Taylor coefficients at an interior point would then vanish). Hence one
joint waypoint choice satisfies every condition. Sequential choice is not used.

On each original interval $[t_i,t_{i+1}]$, join $q_i$ to $q_{i+1}$ by
changing the $N$ scalar coordinates in the fixed order and parametrizing the
resulting coordinate legs inside that same interval (constant legs are simply
omitted). The first and last such pieces lie in
$Q_0,Q_{m-1}\subset\mathcal U_n$ and therefore contain no wall, even if an
inactive canonical factor vanishes there. For every $t\in[t_i,t_{i+1}]$, both
$\tilde\gamma(t)$ and $\gamma(t)$ lie in $Q_i$. Consequently
\[
 \|\tilde\gamma(t)-\gamma(t)\|<\operatorname{diam}(Q_i)<\delta,
\]
which is the required pointwise, hence uniform, bound. The central path lies
in $\mathcal R_n$; its junctions are factor-free, and every leg is nonconstant.
Along a leg, the Bezout certificates make each canonical-factor root simple
and prevent roots of two nonassociate classes from coinciding. Each
restriction is a nonzero univariate polynomial, so the union of their root
sets over finitely many legs is finite.

Finally, vertices cannot collide. On a leg moving an $x$-coordinate, every
relevant $y$-difference is the fixed nonzero value imposed at its two hybrid
ends; on a leg moving a $y$-coordinate, the corresponding $x$-difference is
fixed and nonzero; pairs not involving the moving vertex do not change. A
non-generic parameter is a zero of a relevant member, hence of one of its one
or two canonical factors (A1); the root audit gives
the unique $F$, its simple derivative, and the collision exclusion claimed.
\end{proof}

\begin{remark}[bundles, not single predicates]\label{rem:bundles}
Two mistakes in an earlier general-position route are worth keeping visible.
It treated any two simultaneously vanishing members as independent conditions,
which is false for every bundle below; and it asserted that every active (G4)
tie is a three-remote-edge concurrence, missing the incident-edge tie of the
bigon. The exact witness is the remote edge $(0,0)\to(4,0)$ with
$M=(2,t)$, $p=(1,1)$ and $q=(3,1)$: there
$\mathrm{G2}_{e,M}=4t$ and $\mathrm{G4}_{e;f,g}=-8t(1-t)$.
\end{remark}

\begin{remark}[where bundle transversality is consumed]\label{rem:bundlescope}
At a Reidemeister-III hit the induction verifies event transversality directly
from the common simple factor supplied by Lemma~\ref{lem:dictionary}; at a
silent hit no sign change is consumed. Thus the fixed-count transport step in the induction
consumes the dictionary's simple-zero conclusion only for flat, bigon and
sliding bundles; the induction treats Reidemeister III separately.
\end{remark}

\begin{lemma}[inactive-zero exclusions and the exhaustive wall dictionary]\label{lem:dictionary}
\begin{enumerate}
\item[(A)] Let $t\mapsto P(t)$ be an event and let $Z$ be its zero set.
\begin{enumerate}
\item[(a)] If $\mathrm{G3}_{e,f,g}$ vanishes at $t=0$, the three lines are not
all equal, and their common point fails to lie on at least one of the three
segments, then $\mathrm{G3}_{e,f,g}\notin Z$. (When the three lines coincide
there is no common point to speak of and the hypothesis says nothing; that case
is excluded here and is not used below.)
\item[(b)] If $\mathrm{G5}_{e,f}$ vanishes at $t=0$ and the two closed segments
$e(0)$ and $f(0)$ are disjoint, then $\mathrm{G5}_{e,f}\notin Z$.
\end{enumerate}
\item[(B)] Let $\tilde\gamma$ be a path supplied by Lemma~\ref{lem:relgp}(B), let $t_0$ be
one of its finitely many non-generic parameters, and put
$x=\tilde\gamma(t_0)$. Define
\[
 \mathcal B(x)=\{g\in\mathcal G:g(x)=0\text{ and }g\text{ is relevant at }
 \tilde\gamma(t)\text{ for some }t\neq t_0\text{ in every neighbourhood of }t_0\}.
\]
On an interval containing no other non-generic parameter this is exactly the
zero set of the event of Definition~\ref{def:event}. There is one canonical
factor $F$ with $F(x)=0$ and $\mathrm dF(x)\neq0$ such that every
$g\in\mathcal B(x)$ has $g=F B_g$ with $B_g(x)\neq0$. In particular all
members share one smooth branch and their restrictions to the path have simple
zeros, so every member changes sign at $t_0$. The bundle is exactly one of:
\begin{enumerate}
\item[(i)] \emph{Flat}: $\{\mathrm{G1}_j,
\mathrm{G2}_{(j-1,j),j+1},\mathrm{G2}_{(j,j+1),j-1}\}$ with $p_j$ strictly
between its neighbours; wall (S3).
\item[(ii)] \emph{Vertex-on-edge, bigon}:
$\{\mathrm{G2}_{(a,b),M},\mathrm{G4}\langle(a,b);e_{M-1},e_M\rangle\}$,
the foot in the segment interior and the two neighbours of $M$ strictly on one
side of the remote line; wall (S7).
\item[(iii)] \emph{Vertex-on-edge, sliding}: the singleton
$\{\mathrm{G2}_{(a,b),M}\}$, the foot in the segment interior and the two
neighbours strictly on opposite sides; wall (S7).
\item[(iv)] \emph{Reidemeister III}: $\{\mathrm{G3}_{e,f,g}\}$ together with
the three (G4) members whose ties it forces, for three pairwise remote edges
concurrent in the interior of all three; wall (S2).
\item[(v)] \emph{Silent}: a (G2) vanishing with its point off the segment,
together with whatever members it forces, silent in the sense of
Definition~\ref{def:silent}.
\end{enumerate}
The exterior consecutive-collinearity branch is the cusp and lies outside
$\mathcal R_n$. Remote-edge overlap, vertex coincidence, four-edge
concurrence and two independent bundles cannot occur at a certified hit.
\end{enumerate}
\end{lemma}
\begin{proof}
\emph{Clause (A).} \emph{Clause (a).} Say the common point $q$ misses
the segment $g$. Activation asks that all
three pairs cross (Definition~\ref{def:guarded}), so it is enough to exhibit
\emph{one} pair that does not cross at $t=0$ and does not cross for $|t|$ small,
and the pair must be chosen with care: an earlier revision took $(e,g)$ and said
"the lines of $e$ and $g$ meet only at $q$", which is false when those two lines
coincide.

Choose the pair as follows. The three lines are not all equal, so at least one
of $e,f$ has a line different from $g$'s; call it $h$, so that the lines of $h$
and $g$ are distinct and meet exactly at $q$. Since $q$ misses the segment $g$,
the edges $h$ and $g$ do not cross at $t=0$. For nearby $t$: if the lines of $h$
and $g$ are not parallel their intersection point is a continuous function of
the vertices, and at $t=0$ it lies at positive distance from the compact segment
$g$, so it misses $g$ for all $|t|$ small; if they are parallel --- which at
$t=0$ they are not, distinct lines through $q$ --- they can become so only by
first ceasing to meet near $g$, and in any case two parallel edges meet at most
in a collinear overlap and do not cross transversally. Either way $h$ and $g$ do
not cross for $|t|$ small, so $\mathrm{G3}_{e,f,g}$ is inactive there and
relevant at no $t\neq0$ near $0$. By Definition~\ref{def:event} it is not in
$Z$.

\emph{Clause (b).} Disjoint compact sets stay disjoint under a small
perturbation: the
distance $\min\{|u-v|:u\in e(t),v\in f(t)\}$ is a continuous function of $t$,
positive at $t=0$, hence positive for $|t|$ small. So $e$ and $f$ do not meet,
in particular do not cross, for those $t$; $\mathrm{G5}_{e,f}$ is inactive
there and relevant at no such $t\neq0$, so it is not in $Z$.

\emph{Clause (B).} Choose an interval $I$ about $t_0$ containing no other non-generic parameter.
Every point of $I\setminus\{t_0\}$ is generic, so the restricted path is an
event. Activation is constant on each punctured side because it is decided by
strict sign conditions in unconditional (G2) members. This proves the asserted
equality between $\mathcal B(x)$ and the event zero set.

\emph{Event-cell check.} Fix a vertex $M$ and a remote edge $e$. Put
$s_{e,M}=\langle p_M-p_e,d_e\rangle$. Its three strict incidence regimes are
$s_{e,M}<0$, $0<s_{e,M}<\lVert d_e\rVert^2$ and
$s_{e,M}>\lVert d_e\rVert^2$, according as the perpendicular foot lies before,
inside or after the segment; the outer two are the silent regime. The two
strict side regimes are distinguished by the sign of
$\mathrm{G2}_{e,M-1}\mathrm{G2}_{e,M+1}$: positive puts the neighbours on one
side of the remote line and negative puts them on opposite sides. Finally, for
each $g\in\mathcal G$, its activation locus is the finite conjunction of
strict guarded inequalities of Definition~\ref{def:guarded} (all of
$\mathbb R^{2n}$ when $g$ is unconditional), and its two sign regimes are
$g>0$ and $g<0$. All these regimes are semialgebraic, being cut out by finitely
many polynomial equalities and strict inequalities. This is their only use.

By Lemma~\ref{lem:relgp}(A1), every member of $\mathcal G$ is, up to
sign, one canonical factor or a product of two. Lemma~\ref{lem:relgp}(B1) says that
exactly one associate class vanishes at $x$. Thus a vanishing single-factor
member is a unit multiple of its representative $F$; for a product exactly one
factor is associate to $F$ and the other is nonzero at $x$. Hence every member
of $\mathcal B(x)$ is $F B_g$ with $B_g(x)\neq0$. The simple derivative of
$F\circ\tilde\gamma$ implies $\mathrm dF(x)\neq0$ and makes every restricted
member change sign with the same branch.

It remains to classify the possible defining family. A (G1) zero inside
$\mathcal R_n$ is positive collinearity, so the middle vertex lies between its
neighbours; its two (G2) companions are identically the same determinant. This
is (i). Negative collinearity is outside $\mathcal R_n$ and is the cusp.

For a (G2) zero, first remove the adjacent case: if $M=a-1$ or $M=b+1$, the
member is the corresponding (G1), hence (i). Otherwise the perpendicular foot
cannot be an endpoint of the remote segment, because at a zero that would make
$p_M$ equal to a remote-edge endpoint, a vertex collision excluded by
Lemma~\ref{lem:relgp}(B3). Nor can a neighbour of $M$ lie on the remote line: its
vertex-on-line factor is nonassociate to the defining factor outside the
already separated adjacent case, so two classes would vanish. The incidence
and side regimes in the event-cell check above are therefore strict. With
the foot inside the segment, neighbours on one side give the bigon bundle (ii):
the order member is active on the two-crossing side and its consecutive-edge
factorization has the incidence as its unique vanishing factor. Neighbours on
opposite sides give the sliding singleton (iii), since the order member is
active on neither side. With the foot off the segment, clause~(A)
leaves the silent case (v).

For three pairwise remote edges, the concurrency determinant and the three
order predicates satisfy, for a fixed ordering,
\[
 \mathrm{G4}_{e;f,g}=-\mathrm{G3}_{e,f,g},\qquad
 \mathrm{G4}_{f;e,g}=+\mathrm{G3}_{e,f,g},\qquad
 \mathrm{G4}_{g;e,f}=-\mathrm{G3}_{e,f,g}.
\]
When the three crossings are active their common point is interior to all three
segments, and this is (iv). If the last two edges of a (G4) are consecutive at
$M$, its polynomial instead factors as the vertex-on-line factor
$\Delta_{e_0e_1M}$ times the turn at $M$; the unique-factor certificate assigns
the hit to the corresponding (G2) or (G1) case. A punctured-side-relevant (G5) contributes no new case. Suppose it
vanishes at $x$. If some (G1) or (G2) member also vanishes there, the preceding
flat or vertex-on-edge case owns the hit, and the unique-factor certificate
leaves no independent (G5) family. Otherwise all (G1) and (G2) members are
nonzero at $x$. Their strict signs persist to $x$; an endpoint limit would force
a (G2) zero, so the pair is active at $x$. Proposition~\ref{prop:fidelity}(i)
then contradicts the (G5) vanishing.

Finally, each excluded coincidence violates a printed part of the certificate.
Remote-edge overlap forces two nonassociate vertex-on-line factors. Four-edge
concurrence forces two nonassociate concurrency factors; the possible
same-support $n=6$ pair is nonassociate by Lemma~\ref{lem:relgp}(A3).
Two independent bundles likewise force two nonassociate classes. A vertex
coincidence is excluded separately by Lemma~\ref{lem:relgp}(B3). These
exhaust the guarded families and the event-cell boundaries, proving the list.
\end{proof}


\input{d4_flat}
\input{d5_spike}
% Section 6: the vertex-on-edge law.
\section{The vertex-on-edge law}\label{sec:s7}

\begin{convention}[the standing hypotheses of this section]
\label{conv:s7standing}
Every statement of this section that speaks of a half, of a half's carriers,
rows or selectors, or of the jump those are compared against, is made under the
following three hypotheses, which are named here once and cited by name rather
than repeated or inherited silently:
\begin{enumerate}
\item[(H1)] the event is a \emph{simple vertex-on-edge event} in the sense of
Convention~\ref{conv:events}, with moving vertex $M$ and remote edge $(a,b)$,
occurring on a path inside $\mathcal R_n$ whose polygons have $n$ vertices;
\item[(H2)] the polygon $P_\bullet$ of whichever chamber is under discussion is
\emph{generic} (Definition~\ref{def:generic}(A));
\item[(H3)] both halves $\lambda_1,\lambda_2$ of Definition~\ref{def:halves},
taken at the event, are \emph{generic polygons}, each with at least three
vertices.
\end{enumerate}
When the event is of bigon type, on the side of the event where both edges
incident to $M$ cross the remote edge $(a,b)$, the two crossings so created are
the \emph{newborn} crossings: $x$, where the edge $M-1\to M$ meets $(a,b)$,
and $y$, where $M\to M+1$ meets it.
These hypotheses are in force for the whole section \emph{by this declaration},
not by each statement repeating them: a statement of this section is to be read
with (H1)--(H3) available whether or not it names them, in the way
Convention~\ref{conv:lawful} binds the letter $A$ everywhere below without
being cited at each use. Several statements do name the convention --- those
where a reader would otherwise wonder which polygons are meant --- and that is
a courtesy, not the mechanism. What the declaration is \emph{not} is an
inheritance from a neighbouring proof: the hypotheses are written out here, in
one numbered place, and a formalizer reading any statement of the section has
one address to look at.

The convention exists because the objects of this section do not exist without
it. The half quantities --- a half's carriers, its grouped rows $Q_i$, its
selectors $W_i$, its sector decomposition, and the products
$X_1(\lambda_1)X_1(\lambda_2)$ that the jump is compared against --- are read
off polygons through $X_1$, and $X_1$ is defined on \emph{generic} polygons
only (Definition~\ref{def:X1}). An earlier arrangement stated the hypothesis at
the two endpoint theorems and, in the middle of the section, said that "every
consumer below inherits" it, which is not a quantifier: a reader of the sector
definition or of the sliding restart could not see from the page what was
assumed. (H3) is discharged, at every event the induction meets, by the
certificate proved with the main theorem; nothing in this section proves it.
\end{convention}

\subsection{Grouped carriers are connected sums}

The factor $\Omega_1(S,L)$ of Definition~\ref{def:X1} reads one coefficient of
the \emph{product} $P_{S,L}=\prod_H P_H$ over the pieces carried by $L$. The
carrier floor of Theorem~\ref{thm:carrierfloor}(C), on the other hand, is a
statement about one knot diagram. The bridge is that the product is the
polynomial of a single knot.

\begin{lemma}[carrier restriction is the intrinsic piece diagram]
\label{lem:pieceintrinsic}
Let $P$ be generic, $S\in\Ind(G_P)$, and let $H$ be a residual piece of $S$,
carried by the carrier $L$ (Lemma~\ref{lem:carriers}(iv)). Let $D_L(H)$ be the
diagram obtained from $L$ by retaining exactly the crossings of $H$, each
resolved by the divide convention of Definition~\ref{def:piecediagram}, and
erasing all others; and let $D_P(H)$ be the intrinsic piece diagram of that
definition, obtained from the parent polygon $P$ by retaining exactly the same
crossings under the same convention. Throughout, that a crossing \emph{lies on}
a carrier means that both of its traversal preimages belong to that carrier,
which is the condition Lemma~\ref{lem:carriers}(iv) supplies.
Then the identity map on the visits of $H$ is a record isomorphism
(Definition~\ref{def:record}) from the record of $D_L(H)$ to the record of
$D_P(H)$. Consequently they present the same oriented link
(Axiom~\ref{ax:gausscode}), so
\[
   P_{D_L(H)}=P_H,\qquad w\bigl(D_L(H)\bigr)=|H| .
\]
\end{lemma}
\begin{proof}
\emph{Same double points.} Smoothing at a point $c\in S$ replaces the two
transversally crossing branches by two disjoint arcs inside a disc containing no
other double point, so it destroys the double point $c$ and creates none. Hence
the double points of $L$ are exactly the double points of $P$ that lie on $L$
and are not in $S$. Every label of $H$ has both of its visits on $L$ and
$H\cap S=\varnothing$, so each is a double point of $L$; and both diagrams
retain exactly those and no others.

\emph{Same over/under and signs.} The divide convention
(Definition~\ref{def:piecediagram}) reads the over/under at a double point from
the two branch directions there, and the sign of a crossing is read from the
same two directions with the orientation. Neither is changed by a smoothing
performed at a different point, which alters the curve only inside a disc
containing no other double point. Every crossing is positive in both diagrams,
so both writhes equal $|H|$.

\emph{Same cyclic order.} The traversal of $D_P(H)$ is the traversal circle
$\Gamma$ itself, so the cyclic order of the $2|H|$ visits is their order on
$\Gamma$. By Lemma~\ref{lem:carrierword} the marked points lying on $L$ are
traversed by $L$ in the order induced from $\Gamma$, and the visits of $H$ are
among them; restricting to those gives the same cyclic word.

\emph{The isomorphism.} The two records are not equal: their traversal circles
are different circles, one running around $P$ and one around $L$. What the three
paragraphs give is the map. Let $\varphi$ send each visit of $H$ on the
traversal circle of $D_L(H)$ to the same visit on the traversal circle of
$D_P(H)$ --- the same point of the plane, reached by the same branch --- which
is a bijection of marked-point sets by the first paragraph. It preserves the
cyclic order by the third, carries double points to double points because the
pairing is by label in both, and preserves the over/under designation and the
sign by the second. Those are conditions (a) to (d) of
Definition~\ref{def:record}, so $\varphi$ is a record isomorphism, and
Axiom~\ref{ax:gausscode} identifies the links.

An earlier revision of this lemma said the two records were \emph{equal}. They
are not: equality of records presupposes
one traversal circle, which is exactly what these two diagrams do not share.
Nothing in the three paragraphs changes; what changes is that the map they
produce is now named, and the definition it satisfies is printed.
\end{proof}

\begin{lemma}[HOMFLY products and the two-component $z^{-1}$ row]
\label{lem:homflyrows}
Put $R=\mathbb Z[a^{\pm1},z^{\pm1}]$ and
$\delta=(a-a^{-1})/z\in R$.
\begin{enumerate}
\item[(i)] For oriented knots $K,J$,
\[
 P_{K\# J}=P_KP_J;
\]
and for an oriented knot $K$ and an oriented link $J$,
\[
 P_{K\sqcup J}=\delta P_KP_J.
\]
\item[(ii)] Let $D$ be an oriented link diagram with ordered components
$D_1,D_2$ and linking number $\lambda$. Then
\[
 [z^{-1}]P_D=(a-a^{-1})a^{-2\lambda}
       [z^0]\bigl(P_{D_1}P_{D_2}\bigr).
\]
\item[(iii)] Consequently, for oriented knots $K_1,\dots,K_n$ ($n\geq1$),
\[
 P_{K_1\sqcup\dots\sqcup K_n}=\delta^{\,n-1}P_{K_1}\cdots P_{K_n}.
\]
\end{enumerate}
\end{lemma}
\begin{proof}
A \emph{based oriented link} is an oriented link with one distinguished
nonsingular point $p$ on one distinguished component; isotopies carry $p$ and
that component. Fix an oriented knot $K$. For a based link $(J,p)$ define
\[
   F_K(J,p)=P_{K\#_pJ},\qquad G_K(J,p)=P_KP_J.
\]
Here $K\#_pJ$ is made in a ball at $p$, disjoint from the rest of $J$, by
deleting a short oriented arc through $p$, cutting a copy of $K$ at a short
oriented arc, and using the orientation-compatible cross-pairing of incoming
and outgoing endpoints. Two sufficiently small choices of ball and arc are
related by an isotopy supported near $p$, and a based ambient isotopy carries
the ball and insertion with it. Thus $F_K$ is well defined on based isotopy
classes. Both $F_K$ and $G_K$ satisfy the HOMFLY--PT skein relation in every
disc disjoint from $p$ (Axiom~\ref{ax:homfly}): the insertion at $p$ is outside
the disc and therefore commutes with the local switch and oriented smoothing.

We use the following relative uniqueness argument, printed because unbased
uniqueness alone does not apply to a marked operation. Suppose two $R$-valued
invariants of based oriented links satisfy that skein relation away from the
mark and agree on every based unlink. Draw a based link diagram. Using $p$ as
the first traversal basepoint, choose auxiliary basepoints on the other
components and order those components after the marked one. Traverse the
oriented components in that order. Call a crossing \emph{bad} when its first
visit is on the underpassing branch. We prove equality by lexicographic
induction on
\[
   (\text{number of crossings},\ \text{number of bad crossings}).
\]
If a bad crossing is present, switch it. Its first visit becomes over, so the
crossing count is unchanged and the bad count drops by one; no other crossing
or first-visit order changes. The oriented smoothing has one fewer crossing.
Its disc misses $p$, so the marked point survives and the unique resulting
component containing it is distinguished; give all remaining components fresh
auxiliary basepoints and an order. For either invariant $Q$, if the original
crossing is $L_+$, solve the skein relation as
\[
   Q(L_+)=a^{-2}Q(L_-)+a^{-1}zQ(L_0);
\]
if it is $L_-$, solve it as
\[
   Q(L_-)=a^2Q(L_+)-azQ(L_0).
\]
Only the units $a^{\pm1}$ are used. The induction hypotheses for the switched
and smoothed diagrams give equality at the original diagram.

If there is no bad crossing, the ordered based diagram is descending. Keep a
small arc at the current component's basepoint fixed. Starting just after it
and following the orientation, pull each successive already-traversed arc
upward and across the still-untraversed diagram. At every crossing that moving
arc is the overpassing branch, so the pull crosses no strand. Successively
remove all self-crossings of that component; because its visits precede those
of every later component, it also passes over every later component and can be
slid into a disjoint ball as a round circle. For the first component the fixed
arc contains $p$, so this is rel the mark. Repeat in the chosen component
order. The result is the based unlink with the same number of components.
Thus agreement on based unlinks closes the induction.

Apply relative uniqueness first to $F_K,G_K$. On the based $r$-component
unlink $U_r$, with $U_0$ denoting no extra split component,
Definition~\ref{def:homfly} gives
\[
 F_K(U_r,p)=P_{K\sqcup U_{r-1}}=\delta^{r-1}P_K
 =P_KP_{U_r}=G_K(U_r,p).
\]
Hence $P_{K\#_pJ}=P_KP_J$. For the split operation put $K$ in a ball
disjoint from $J$ and define
\[
 F_K^{\sqcup}(J,p)=P_{K\sqcup J},\qquad
 G_K^{\sqcup}(J,p)=\delta P_KP_J.
\]
These again satisfy the skein relation away from $p$, and repeated use of the
split-circle identity in Definition~\ref{def:homfly} gives, on $U_r$,
\[
 F_K^{\sqcup}(U_r,p)=P_{K\sqcup U_r}=\delta^rP_K
 =\delta P_KP_{U_r}=G_K^{\sqcup}(U_r,p).
\]
Relative uniqueness proves $P_{K\sqcup J}=\delta P_KP_J$ for every based
oriented link $J$; this is the split identity of~(i). Clause~(iii) follows by
induction on $n$: the case $n=1$ is trivial, and
$K_1\sqcup(K_2\sqcup\dots\sqcup K_n)$ applies~(i)'s split identity with knot
$K_1$ and link $K_2\sqcup\dots\sqcup K_n$.

For~(ii), put $h(E)=[z^{-1}]P_E$. At a mixed crossing of a two-component
diagram, the oriented smoothing $E_0$ is a knot. Extracting the $z^{-1}$
coefficient from the skein relation gives
\[
 a h(E_+)-a^{-1}h(E_-)=[z^{-2}]P_{E_0}=0.
\]
The last equality uses the full retained clause
$P_{E_0}\in\mathbb Z[a^{\pm1},z^2]$: this is a polynomial ring in $z^2$, so
its $z$-powers are nonnegative as well as even. Thus
$h(E_+)=a^{-2}h(E_-)$. With the convention
$\lambda_+=\lambda_-+1$ at a positive versus negative mixed crossing,
$a^{2\lambda(E)}h(E)$ is invariant under every mixed crossing switch.
No self-crossing is switched: its oriented smoothing has three components,
not a knot, so the displayed vanishing is unavailable.

Switch precisely the mixed crossings at which component $D_1$ passes under
$D_2$. In the resulting diagram $D^\uparrow$, $D_1$ passes over $D_2$ at every
mixed projected intersection. Realize its self-crossings by strict local height
differences. Translate the whole first component upward and the whole second
component downward. Every mixed height gap strictly increases, the
self-crossing gaps are unchanged, and no other projected pair can collide.
This is an isotopy through embeddings; for a large translation a horizontal
plane separates the two components. Their knot types are still $D_1,D_2$ and
the endpoint is split, so its linking number is zero. Mixed-switch invariance
and~(i) now give
\[
\begin{aligned}
 h(D)
 &=a^{-2\lambda}h(D^\uparrow)
  =a^{-2\lambda}[z^{-1}]
       \bigl(\delta P_{D_1}P_{D_2}\bigr)\\
 &=a^{-2\lambda}(a-a^{-1})
       [z^0]\bigl(P_{D_1}P_{D_2}\bigr),
\end{aligned}
\]
which is~(ii).
\end{proof}

\begin{corollary}[interlacement connected sums and the grouped carrier]\label{cor:groupedknot}
\begin{enumerate}
\item[(A)] Let $S\in\Ind(G_P)$, let $L$ be a carrier of $S$ bearing at least one residual
piece, and let
\[
   W=\bigcup\{\,H: H\text{ a residual piece of }S\text{ carried by }L\,\}
\]
be the \emph{union of their label sets} --- a set of crossing labels, not a set
of pieces. By Lemma~\ref{lem:piececurve} Step 2, $W$ is exactly the set of
self-crossings of $L$, so the diagram $D(W)$ obtained from $L$ by retaining
exactly the labels of $W$, each resolved by the divide convention, and erasing
all others, is $L$ itself carrying all of its own crossings; no realization
argument is needed. Let $W=W_1\sqcup\dots\sqcup W_k$, $k\geq1$, be the partition
of $W$ into the connected components of the interlacement graph induced on $W$;
these are the label sets of the residual pieces. The components can be indexed
and equipped with a fixed ordered binary parenthesization $\mathcal T$ such that
\[
D(W)\;\cong\;\#_{\mathcal T}
   \bigl(D(W_1),\dots,D(W_k)\bigr),
\qquad
P_{D(W)}=\prod_{i=1}^{k}P_{D(W_i)} .
\]
Here $\#_{\mathcal T}$ is evaluated recursively: a leaf has value $D(W_i)$
and an internal node has the connected sum of its two ordered child values
with the displayed parenthesization. In particular, when $k=1$,
$\#_{\mathcal T}(D(W_1))=D(W_1)$.
\item[(B)] Let $S\in\Ind(G_P)$ and let $L$ be a carrier of $S$ bearing the residual pieces
$H_1,\dots,H_k$ with $k\geq1$. (If $L$ bears no piece the assertion is
Theorem~\ref{thm:carrierfloor}(D): $P_{S,L}=1$, $w_{S,L}=0$, and there is nothing to
decompose. Clause~(A) is stated for a nonempty label set and
is not applied in that case.) Then the diagram obtained from $L$ by retaining exactly the
crossings of $H_1\cup\dots\cup H_k$ is a knot diagram whose HOMFLY polynomial is
$P_{S,L}=\prod_iP_{H_i}$, whose writhe is $w_{S,L}=\sum_i|H_i|$, and whose
underlying plane curve is $L$, which has no triple points.
\end{enumerate}
\end{corollary}
\begin{proof}
\emph{Clause (A).}
Read the carrier as a cyclic word in the labels of $W$, each label occurring
twice. Call a \emph{gap} a position between two consecutive visits.

\emph{Gap lemma.} Fix a connected component $W_i$. A gap has no intrinsic pair
of sides, so the coordinate is taken per chord: for a gap $g$ and a chord
$x\in W_i$, record which of the two components of the traversal circle minus the
two visits of $x$ contains $g$. Ranging over $x\in W_i$ this is a binary vector
indexed by the chords of $W_i$, the \emph{side vector} of $g$. Distinct gaps of $W_i$ have
distinct side vectors. Suppose two distinct gaps had the same vector. The
oriented interval between the two cuts then contains either both or neither
endpoint of every chord of $W_i$; since the gaps are distinct, both the
interval and its complement contain visits of $W_i$. So the labels of $W_i$
split into two nonempty sets, one with both endpoints inside and one with both
endpoints outside, and no chord of one set interlaces a chord of the other.
That disconnects $W_i$, contrary to hypothesis.

\emph{Every other component sits inside one gap of $W_i$.} A chord outside
$W_i$ fails to interlace every chord of $W_i$ exactly when its two endpoint
gaps have the same side vector, hence, by the gap lemma, exactly when both its
endpoints lie in one and the same gap of $W_i$. If some component $W_j$,
$j\neq i$, occupied two different gaps of $W_i$, its chords in the two gaps
could not interlace each other, contradicting connectedness of $W_j$.

\emph{Closed-block tree.} Apply the preceding two paragraphs inside any
cyclic node word containing at least two of the components. They give a
nonempty proper label-closed linear block $\alpha$; its complementary linear
word $\beta$ is label-closed as well. Record the two closed cyclic child words
in the order $(\alpha,\beta)$. In each child, identifying its two cut ends
creates one distinguished \emph{join gap}. Both children have fewer labels,
so iteration gives a finite ordered binary tree $\mathcal T$ whose leaves are
exactly the $W_i$. Index the leaves in its order. Every non-root node also
remembers the one gap in its cyclic record prescribed when its parent was cut;
that output gap is distinct data from the two child join gaps used when the
node itself is assembled.

We induct upward with the following invariant. A node $V$ has a connected
generic oriented one-circle diagram $C_V$ and an orientation-preserving record
isomorphism from its record to the cyclic record of $V$; if $V$ is not the
root, it also has a clean open subarc in the exact output gap prescribed by its
parent.

For a leaf $V=W_i$, take the actual curve $C_{W_i}$ constructed by
Lemma~\ref{lem:piececurve}, with the inherited divide crossing data. It has
exactly the distinct inherited transverse double points of $W_i$, its
restricted cyclic word, and no triple point: the construction only smooths in
pairwise disjoint crossing discs and creates no intersection.
Lemma~\ref{lem:pieceintrinsic} supplies the record comparison with the carrier
factor $D(W_i)$. If the leaf is not the root, the extension clause in
Definition~\ref{def:record} carries its parent-prescribed output gap to a
complementary arc of its traversal circle. Choose an interior regular point
there and a sufficiently small subarc about it; a small disc meets the curve
in that arc alone, so it is clean. If the leaf is the root, the invariant
requires no output gap.

For an internal node with ordered children $A,B$, use their clean join arcs.
Take $C_A$ as the based host $J$, with $p$ in its join arc, and take $C_B$ as
$K$. At an interior point of $K$'s join arc choose an
orientation-preserving smooth chart of the sphere sending that point to
infinity. A sufficiently small spherical cap about infinity meets $K$ in one
embedded arc, so its complement is a compact long $K$-tangle with two oriented
ends. Shrink that tangle into the clean host disc, delete the short host arc
through $p$, place the tangle ends next to their prescribed host ends, and
cross-pair outgoing $A$ to incoming $B$ and outgoing $B$ to incoming $A$ by
two disjoint short arcs. Round the two joins inside disjoint smaller discs.
This is literally the orientation-compatible local insertion $K\#_pJ$ of
Lemma~\ref{lem:homflyrows}(i), not an unnamed planar splice.

The insertion disc contains no other point of $J$ and the inserted tangle and
connectors lie inside it. Thus it creates no crossing. The sphere chart is
orientation preserving, so all child crossing signs are preserved, and the
over/under designations are transported. Starting immediately before $A$, the
one traversal circle reads the linear word $\alpha$ and then $\beta$.
Pairings stay within a child. These are precisely conditions (a)--(d) of
Definition~\ref{def:record}. The child crossings remain transverse and
distinct, the connectors are disjoint, and the local rounding creates no
intersection, so the result is a generic one-circle diagram with no triple
point.

It remains to mark the node's output gap when the node is not the root. If that
gap is not one of the two split boundaries, it survives in the appropriate
child and is carried across by the record map. If it is a split boundary, use
the corresponding new connector subarc. In either case a smaller disc meeting
only that arc makes it clean. This proves the invariant.

Finally the actual diagram $D(W)$ has underlying curve $L$. The diagrammatic
parent has isolated transverse double points with no two sharing an image, so
the smoothing sites admit pairwise disjoint discs, each meeting the curve in
the two crossing arcs alone. Each local smoothing replaces them by two
noncrossing arcs and outside those discs the curve is unchanged. Hence $L$ is
a connected generic one-circle curve with no triple point. The constructed
root has the same properties by the invariant, and its record is isomorphic to
the record of $D(W)$: the traversal order is the tree's ordered concatenation,
labels and pairings never cross between children, and every over/under role and
sign is inherited. Axiom~\ref{ax:gausscode} identifies the two oriented links.
At the leaves Lemma~\ref{lem:pieceintrinsic} identifies $C_{W_i}$ with
$D(W_i)$. Thus the root is exactly the displayed
$\mathcal T$-parenthesized connected sum, without an associativity assumption.

The polynomial identity follows by applying
Lemma~\ref{lem:homflyrows}(i) at each internal node of the fixed tree
$\mathcal T$; every leaf is a knot diagram. Thus
$P_{D(W)}=\prod_iP_{D(W_i)}$.

\emph{Clause (B).}
Each $H_i$ has both visits of each of its labels on $L$ by
Lemma~\ref{lem:carriers}(iv), and the $H_i$ are exactly the connected
components of the interlacement graph induced on $H_1\cup\dots\cup H_k$, since
they are components of $G_P[U(S)]$ and interlacement is read from the same
word. Clause~(A) applies. The diagram has one component
because $L$ is one closed curve, the writhe is the crossing count because every
crossing is positive (Definition~\ref{def:piecediagram}), and the underlying
plane curve is $L$ because erasing a double point omits it from the record and
performs no operation on the curve (Definition~\ref{def:piecediagram}).

The generic parent polygon is diagrammatic by
Definition~\ref{def:diagrammatic}, so it has no triple points.  The parent's
self-intersections are isolated and no two share an image
(Definition~\ref{def:diagrammatic}), so they admit pairwise disjoint discs,
each meeting the curve in the two crossing arcs alone.  Inside each disc the
oriented smoothing of Definition~\ref{def:smoothing} replaces those two arcs
by two arcs that do not cross, and outside the discs the curve is unchanged;
so no intersection is created and $L$ has no triple points either.

The identification of the factors is Lemma~\ref{lem:pieceintrinsic}: the factor
$D(H_i)$ produced by clause~(A) is the diagram of $L$ with
only the crossings of $H_i$ retained, and the identity on the visits of $H_i$ is
a record \emph{isomorphism} from that diagram's record to the record of the
intrinsic piece diagram of $H_i$ on the parent polygon --- not an equality, the
two traversal circles being different --- so by Axiom~\ref{ax:gausscode} the two
present the same oriented link, whose polynomial is $P_{H_i}$ and whose writhe
is $|H_i|$. So the product in the display is a
product of the $P_{H_i}$ and not of some other family. An earlier revision
asserted this identification in one sentence, on the ground that both
constructions ``retain the crossings of $H_i$''; the peer refuted that ground
--- smoothing the support changes the traversal successor, so containment of
both visits does not identify the cyclic records --- and
Lemma~\ref{lem:pieceintrinsic} now supplies the three components of the record
separately.
\end{proof}

\begin{remark}[why clause (A) is not stated for an arbitrary label set]
\label{rem:connectedsumscope}
An earlier revision quantified clause~(A) over any set $W$ of labels with both
visits on the carrier. That is not well posed. Restricting a realizable Gauss
word to an arbitrary label set can produce a word no plane curve realizes: the
peer's witness restricts $123123$ to two of its three crossings and obtains
$2323$, in which each chord interlaces exactly one other, while every chord of a
realizable word interlaces an even number. For such a $W$ the diagram $D(W)$
does not exist and neither side of the display means anything.

For the full self-crossing set $W$ of a carrier --- the hypothesis of clause~(A) ---
no realization theorem is needed at all: by Lemma~\ref{lem:piececurve} Step 2 the
undominated crossings carried by $L$ are exactly the self-crossings of $L$, so
$D(W)$ is the carrier itself. An intermediate revision instead allowed $W$ to be
any nonempty \emph{union} of residual pieces and cited
Lemma~\ref{lem:piececurve} to realize it; the peer refused that citation and was
right --- that lemma realizes one \emph{connected} piece, its Step~5 requiring
connectedness, and a disconnected union is outside it. The citation is
withdrawn; what replaces it is the observation that the set
$W=H_1\cup\dots\cup H_k$ supplied by clause~(B) is the whole self-crossing
set of an actual curve.

The witness is excluded for a separate and still valid reason: two interlacing
crossings that interlace nothing else cannot be a residual piece, since by
Corollary~\ref{cor:evendominated} an undominated crossing meets an even number
of dominated crossings, so its degree inside its own piece is even, and in a
two-element piece it would be one.
\end{remark}

\begin{remark}[a generalization that was tried and withdrawn]\label{rem:groupedgeneral}
An intermediate revision quantified clause~(B) over an arbitrary closed
curve rather than a carrier, so that it could be applied to a sub-loop cut out
of a carrier at one of its crossings, and discharged the resulting comparison
of diagrams by asserting that the complementary excursion, carrying only erased
crossings, could be contracted. The peer rejected that twice and was right
both times: the general statement does not bind the parent curve or the map to
the intrinsic piece diagrams of Definition~\ref{def:piecediagram}, and an
embedded arc whose interior is disjoint from a diagram need not be
boundary-parallel, so ``separated in space'' does not produce the disc the
contraction needs. Clause~(B) is therefore back in its carrier form, and
Theorem~\ref{thm:s7universal}(D)(i), which wanted the general form,
is reproved below by passing to a larger support --- whose carriers \emph{are}
the two loops --- so that only carrier instances of clause~(B) are used.
\end{remark}

\begin{remark}[why this is needed]\label{rem:whygrouped}
Theorem~\ref{thm:carrierfloor}(C) bounds $\min\deg_a$ of one knot diagram in terms
of \emph{its} writhe and rotation. Corollary~\ref{cor:groupedknot}(B) says the
grouped object $P_{S,L}$, with its grouped writhe $w_{S,L}$ and the carrier's
own rotation $R(L)$, is exactly such a diagram. Without it the floor would have
to be proved for each piece separately and then convolved, which is false in
general: the sum of the pieces' floors is not the floor of the product.
\end{remark}

\subsection{The contact-turn geometry}

\begin{lemma}[the contact turns, both branches]\label{lem:contactturnboth}
Under Convention~\ref{conv:s7standing}, let a simple vertex-on-edge event of
bigon type have moving vertex $M$, remote edge $(a,b)$, and neighbours
$p=p_{M-1}$ and $q=p_{M+1}$.
Let $\epsilon=1$ if the two newborn crossings
(Convention~\ref{conv:s7standing}) interlace, and $\epsilon=0$
otherwise. Let $s_{\mathrm{low}}$ denote the sign \eqref{eq:S6} evaluated in the low
(zero-newborn) chamber. First take an orientation-preserving affine
chart sending the contact to the origin and the direction of $a\to b$ to
$+x$, and write $p=(u,h)$ and $q=(v,k)$. Then $hk>0$. If $h,k<0$,
compose with reflection in the $x$-axis; thereafter work in the resulting
upper picture, where $h,k>0$ and $0<\arg p,\arg q<\pi$. This reflection
reverses all three turns below and $s_{\mathrm{low}}$, so the sign comparisons are
independent of which side was reflected. In the upper picture:
\begin{enumerate}
\item[(i)] $\epsilon=0$ exactly when $\arg p<\arg q$, and $\epsilon=1$
exactly when $\arg q<\arg p$.
\item[(ii)] The two incident directions $-p,q$ are nonzero and transverse to
the remote direction. The contact turns of $\lambda_1,\lambda_2$ are,
respectively, $\arg q$ and $\pi-\arg p$; both lie in $(0,\pi)$ and
have the sign of $s_{\mathrm{low}}$.
\item[(iii)] The full turn at $M$ is the principal representative of
$\arg q-\arg p+\pi$. It is
$\arg q-\arg p-\pi\in(-\pi,0)$ when $\epsilon=0$, with sign
$-s_{\mathrm{low}}$, and $\arg q-\arg p+\pi\in(0,\pi)$ when
$\epsilon=1$, with sign $s_{\mathrm{low}}$.
\item[(iv)] Each half polygon inherits exactly one of these contact turns:
$\lambda_1$ closes along the remote direction and leaves along $q$, while
$\lambda_2$ enters along $-p$ and leaves along the remote direction.
Consequently, when $\epsilon=0$, any half contact carrier whose remaining
turns are inherited from a uniformly turning full contact carrier has exactly
one turn opposite to the full turn, and that dissenting turn has magnitude
strictly below $\pi$.
\end{enumerate}
\end{lemma}
\begin{proof}
For a sufficiently small displacement into the two-newborn chamber, normalize the
remote line to $y=0$ and write
$M_t=(m_t,-c_t)$ with $c_t>0$, and the two neighbours as
$p_t=(u_t,h_t)$ and $q_t=(v_t,k_t)$, tending to $p=(u,h)$ and $q=(v,k)$,
respectively. The incident-edge intersections with the remote line have abscissae
$x_p=m_t+c_t(u_t-m_t)/(h_t+c_t)$ and
$x_q=m_t+c_t(v_t-m_t)/(k_t+c_t)$, so
$(x_p-x_q)/c_t\longrightarrow u/h-v/k=(uk-hv)/(hk)$.
The vertex turn is nonzero, hence $uk-hv\neq0$, so this sign is fixed for small
$t$. Along the polygon the incoming crossing precedes the outgoing crossing;
along the oriented remote edge their order is increasing abscissa. They are
therefore noninterlaced exactly when $x_p>x_q$, and interlaced when $x_p<x_q$.
In the upper half-plane $x_p>x_q$ is exactly $\arg p<\arg q$. This proves (i).

The turn at $M$ is from incoming direction $-p$ to outgoing direction $q$,
hence is represented by $\arg q-\arg p+\pi$. If $\epsilon=0$, clause (i)
gives $\arg p<\arg q$, so its principal value is
$\arg q-\arg p-\pi\in(-\pi,0)$. If $\epsilon=1$, clause (i) gives
$\arg q<\arg p$, so the displayed representative lies in $(0,\pi)$.
In the upper low chamber $M$ is on the neighbours' side of the remote line,
so $s_{\mathrm{low}}=+1$ by \eqref{eq:S6}; this proves (iii),
and reflection negates both sides of every claimed sign comparison.

By Definition~\ref{def:halves}, $\lambda_1$ closes in direction $+x$ and
then leaves along $q$, while $\lambda_2$ arrives along $-p$ and then leaves
in direction $+x$. Their turns are $\arg q$ and $\pi-\arg p$, proving (ii)
and the inheritance assertion in (iv). In the noninterlacing case (iii)
makes the full contact turn opposite to both half contact turns. If the full
carrier is uniform (Definition~\ref{def:wind}), every remaining turn inherited by
either half has the full turn's sign; its one contact turn has the opposite
sign and magnitude in $(0,\pi)$. This proves (iv).
\end{proof}

\subsection{The two sectors of the bigon jump}\label{sec:sectors}

Fix a simple vertex-on-edge event of bigon type: the two neighbours of the
moving vertex $M$ lie strictly on the same side of the line of the remote edge
$(a,b)$, so that crossing the wall changes the crossing number by two. Orient
the parameter from the chamber with no newborn crossing to the chamber with
two; write $P_0$ for the first and $P_2$ for the second, and let
$\delta=+1$ in that direction and $-1$ in the reverse. By
Remark~\ref{rem:s7wellposed} nothing depends on the choice.

\begin{definition}[newborns, sectors, and the interlacement bit]\label{def:sectors}
Under Convention~\ref{conv:s7standing} --- in particular with both halves
generic, without which the products $X_1(\lambda_1)X_1(\lambda_2)$ compared
below are not defined --- the newborn crossings $x,y$ of
Convention~\ref{conv:s7standing} are realized in $P_2$. The directed jump is
$F=\delta\bigl(X_1(P_2)-X_1(P_0)\bigr)$. Split the supports of $P_2$ by how
many newborns they contain and write
\[
   F=R+\mathcal B,
\]
where $\mathcal B$ is $\delta$ times the total contribution of the supports of
$P_2$ containing \emph{both} newborns and $R$ is everything else, the
\emph{returned} sector. Let $\epsilon=1$ if $x$ and $y$ interlace in $P_2$ and
$\epsilon=0$ otherwise, and let $J=s\,X_1(\lambda_1)X_1(\lambda_2)$ be the
required jump, with $s$ the sign \eqref{eq:S6} taken in the initial chamber.
\end{definition}

\begin{lemma}[the contact word]\label{lem:contactword}
The cyclic Gauss word of $P_2$, rotated to begin at the first visit of the
newborn $x$ --- the rotation is named here because consumers speak of the visit
that \emph{opens} the word, which a cyclic word does not have until one is
fixed --- reads
\begin{equation}
   W_2=\begin{cases}
     x_0\,y_0\,A\,y_1\,x_1\,B,&\epsilon=0,\\
     x_0\,y_0\,A\,x_1\,y_1\,B,&\epsilon=1,
   \end{cases}
   \qquad\text{and in both cases}\qquad
   W_0=A\,B
   \label{eq:contactword}
\end{equation}
where $x_0,x_1$ are the two visits of $x$, $y_0,y_1$ those of $y$, $A$ is the
subword of old labels traversed from $M+1$ to $a$ together with the part of the
remote edge from $a$ to \emph{the first newborn visit on that edge} --- which is
$y$ when $\epsilon=0$ and $x$ when $\epsilon=1$, as the two displayed words show
--- and $B$ is the subword traversed from $b$ to $M-1$. An earlier revision said
``to $y$'' in both branches, which is false in the interlacing one; the proof
always used the order-neutral phrase. Moreover an old label $q$ has
\begin{enumerate}
\item[(i)] both visits in $A$, or both in $B$, or exactly one in each;
\item[(ii)] and it is adjacent in $G_{P_2}$ to $x$, or to $y$, exactly when it
has one visit in each --- so the labels adjacent to $x$ and those adjacent to
$y$ coincide, and form the \emph{common neighbourhood} $\mathcal N$.
\end{enumerate}
\end{lemma}
\begin{proof}
In $P_2$ the two edges incident to $M$ each meet the remote edge once, at $x$
and at $y$ respectively; the contact disc contains no other vertex or crossing,
because a further incidence at the event would be a vanishing guarded predicate
outside the zero set that Convention~\ref{conv:events} allows, which is exactly
$\{\mathrm{G2}_{(a,b),M},\ \mathrm{G4}\langle (a,b);e_{M-1},e_M\rangle\}$ ---
the member-valued accessor, a zero set holding members
(Definition~\ref{def:guarded}) --- and contains no other incidence. Traversing from just before $M$: the
edge $M-1\to M$ carries the visit $x_0$, then comes the vertex $M$, then the
edge $M\to M+1$ carries $y_0$; no crossing lies between them. Continuing, the
traversal runs $M+1,\dots,a$ and then along the remote edge, where it meets the
two newborns again, in one of the two orders, before continuing $b,\dots,M-1$.
Writing $A$ for the stretch from just after $y_0$ to just before the first
remote-edge newborn visit and $B$ for the stretch from just after the second to
just before $x_0$, the word is $x_0y_0Ay_1x_1B$ or $x_0y_0Ax_1y_1B$ according
to that order. In the first case the interval from $x_0$ to $x_1$ is
$y_0Ay_1$, which contains both visits of $y$, so $x$ and $y$ do not interlace;
in the second it is $y_0A$, which contains exactly one, so they do. The two
cases are therefore $\epsilon=0$ and $\epsilon=1$, which is
\eqref{eq:contactword}. Deleting
the four newborn visits leaves $AB$, and $P_0$ has exactly the old crossings of
$P_2$ with the same order, since the two chambers' crossing sets differ by
exactly the newborn pair (the argument of Lemma~\ref{lem:cuspcases}(iv), with the
vertex-on-edge zero in place of the cusp one). That is $W_0=AB$.

For (i) and (ii), take $\epsilon=0$ first: the interval from $x_0$ to $x_1$ is
$y_0Ay_1$ and the interval from $y_0$ to $y_1$ is $A$. An old label with both
visits in $A$ has two visits in each interval, so it interlaces neither
newborn; with both visits in $B$ it has none in either, so again neither; with
one in $A$ and one in $B$ it has exactly one in $y_0Ay_1$ and exactly one in
$A$, so it interlaces both. For $\epsilon=1$ the two intervals are $y_0A$ and
$Ax_1$, and the same three cases give two, zero and one visits respectively, so
the conclusion is identical.
\end{proof}

\begin{definition}[eligible]\label{def:eligible}
An old support $T\in\Ind(G_{P_0})$ is \emph{eligible} when $T\cap\mathcal
N=\varnothing$, i.e.\ when no element of $T$ is adjacent to a newborn.
\end{definition}

\begin{remark}[why eligibility is stated this way and not the other]
\label{rem:eligibility}
The tempting definition, and the one an earlier revision of this section used,
is "$T\cup\{x,y\}$ is independent". In the branch $\epsilon=0$ the two agree:
$x$ and $y$ are then non-adjacent to each other, and $T$ is independent, so
$T\cup\{x,y\}$ is independent exactly when $T$ misses the neighbours of both.
In the branch $\epsilon=1$ they do not agree at all --- there $x\sim y$, so
$T\cup\{x,y\}$ is \emph{never} independent and that definition has empty
domain. Every statement below about the returned sector would then be
vacuous at $\epsilon=1$ and would prove nothing about the interlacing bigon.
Definition~\ref{def:eligible} is the condition the arguments actually use: it
is what makes the old labels split between the two intervals, and it is
non-vacuous in both branches.
\end{remark}

\begin{lemma}[the eligible old supports are an ordered product]\label{lem:orderedproduct}
Here eligibility is as in Definition~\ref{def:eligible}.
Let $T$ be eligible. Then, writing $V_A$ and $V_B$ for the old labels with both
visits in $A$ and in $B$,
\[
   G_{P_2}\bigl[V_A\sqcup V_B\bigr]=G_A\sqcup G_B
\]
is a disjoint union with no edge between the parts, and
$T\mapsto(T\cap V_A,\;T\cap V_B)$ is a bijection from the eligible supports
onto $\Ind(G_A)\times\Ind(G_B)$.
\end{lemma}
\begin{proof}
An eligible $T$ contains no element of $\mathcal N$ by
Definition~\ref{def:eligible}, and by Lemma~\ref{lem:contactword}(ii)
$\mathcal N$ is exactly the set of old labels with one visit in each interval,
so by Lemma~\ref{lem:contactword}(i) every element of $T$ lies in $V_A$ or in
$V_B$. A chord of $V_A$ has both endpoints in the cyclic interval $A$ and a
chord of $V_B$ both endpoints in the disjoint interval $B$, so their endpoints
cannot alternate: there is no edge between the parts. Independent sets of a
disjoint union are exactly pairs of independent sets of the parts, and the map
displayed is the inverse of $(T_A,T_B)\mapsto T_A\sqcup T_B$. The order of the
pair is the order of the halves of Definition~\ref{def:halves}, not a quotient:
$A$ is the word of $\lambda_1$ and $B$ that of $\lambda_2$, because $\lambda_1$
runs $M\to M+1\to\dots\to a$ and closes along the remote edge, which is the
stretch $A$, and $\lambda_2$ opens along the remote edge and runs
$b\to\dots\to M-1$, which is $B$.
\end{proof}

\subsection{The noninterlacing both-newborn support calculation}\label{sec:niboth}

\begin{proposition}[the noninterlacing carrier factorization and sector target]\label{prop:sectortarget}
For clauses~(A)--(C), let $\epsilon=0$, let $T$ be an eligible old support
with ordered parts $T_A,T_B$ as in Lemma~\ref{lem:orderedproduct}, and put
$S=T\cup\{x,y\}$. Write $\mathcal N$ for the common neighbourhood of
Lemma~\ref{lem:contactword}, and for a word $W$ write
$\mathrm{Del}_{\mathcal N}(W)$ for $W$ with every visit of a label of
$\mathcal N$ deleted; put $A^{\circ}=\mathrm{Del}_{\mathcal N}(A)$ and
$B^{\circ}=\mathrm{Del}_{\mathcal N}(B)$. Let $s_{\mathrm{low}}$ be the sign
\eqref{eq:S6} computed in the low chamber $P_0$. For clauses~(D)--(F), let
$F,R,\mathcal B,\epsilon,J$ be as in Definition~\ref{def:sectors}. Then:
\begin{enumerate}
\item[(A)] The carriers of $S$ in $P_2$ are, canonically and without repetition, the
carriers of $T_A$ in $\lambda_1$, the carriers of $T_B$ in $\lambda_2$, and one
further carrier $U_{\mathrm{loc}}$ lying in the contact disc.
\item[(B)] Under the correspondence of clause~(A), each carrier of $S$
other than $U_{\mathrm{loc}}$ has the same residual pieces, the same signed piece words,
the same piece polynomials, the same grouped polynomial $P_{S,L}$ and the same
grouped writhe $w_{S,L}$ as its mate in the half. The local carrier bears no
piece: $P_{S,U_{\mathrm{loc}}}=1$ and $w_{S,U_{\mathrm{loc}}}=0$.
\item[(C)] Here $\rot$ is as in Definition~\ref{def:rot}.
Under the correspondence of clause~(A), each carrier of $S$ other than $U_{\mathrm{loc}}$ has the same cyclic
corner list with the same turn signs, hence the same weight, signed rotation,
$R$, and turn-uniformity as its half mate. The local carrier has
\[
   P_{S,U_{\mathrm{loc}}}=1,\quad w_{S,U_{\mathrm{loc}}}=0,\quad \rot(U_{\mathrm{loc}})=-s_{\mathrm{low}},\quad R(U_{\mathrm{loc}})=1,\quad
   \Omega_1(S,U_{\mathrm{loc}})=1,\quad \mathrm{wt}(U_{\mathrm{loc}})=s_{\mathrm{low}} .
\]
\item[(D)] if $\epsilon=1$, then $\mathcal B=0$;
\item[(E)] if $\epsilon=0$, then $\mathcal B=J$;
\item[(F)] in both cases $\mathcal B=(1-\epsilon)J$, and therefore
\[
   F=J\quad\Longleftrightarrow\quad R=\epsilon J .
\]
\end{enumerate}
In particular the bigon branch of (S7) holds if and only if the returned sector
is $J$ when the newborns interlace and $0$ when they do not.
\end{proposition}
\begin{proof}
\emph{Clause (A).}
Smoothing reconnects the traversal by the successor rule: at a smoothed
crossing with visits $c_0,c_1$, the arc arriving at $c_0$ continues with the
arc leaving $c_1$, and conversely. Applied to \eqref{eq:contactword} at the two
newborns this gives
\[
   \sigma(x_0)=B_0,\qquad
   \sigma(x_1)=y_0,\qquad
   \sigma(y_0)=x_1,\qquad
   \sigma(y_1)=A_0,
\]
where $A_0$ and $B_0$ are the first positions of $A$ and of $B$ --- and, when
one of those words is empty, the adjacent marker, namely $y_1$ for an empty
$A$ and $x_0$ for an empty $B$. The middle two are a $2$-cycle
$\{y_0,x_1\}$: that is the local carrier $U_{\mathrm{loc}}$, the closed curve made of the
two smoothing arcs at $x$ and $y$ together with the stretch of the remote edge
between them and the two incident stretches through $M$. The other two
successors leave $x_0$ into $B$ and $y_1$ into $A$, and by
Lemma~\ref{lem:orderedproduct} the selected old labels of $A$ and of $B$ are
$T_A$ and $T_B$, whose own smoothing is internal to their intervals. So the
remaining cycles are one family confined to $A$ and one confined to $B$.

To match them with the halves, delete from each the designated newborn boundary
marker and every visit of a label of $\mathcal N$. Such a label crosses between
the two halves, so it appears in neither half's self-crossing word; and it is
not smoothed, since an eligible support misses $\mathcal N$, so its visits are
through-points and not corners --- deleting them removes no corner. A step that
lands on the deleted boundary marker becomes the cyclic wraparound of the
cleaned word. The projected cycles are therefore exactly the smoothing cycles
of $T_A$ on $A^{\circ}$ and of $T_B$ on $B^{\circ}$, which are the carrier
books of $\lambda_1$ and $\lambda_2$ by Lemma~\ref{lem:carriers} applied inside
each half. If a cleaned word is empty its single boundary cycle projects to the
one crossing-free carrier of that half, which is that half itself.


\emph{Clause (B).}
A label of $\mathcal N$ is adjacent to both newborns
(Lemma~\ref{lem:contactword}(ii)), so in $P_2$ it lies in
$N_{G_{P_2}}(S)$ and is absent from $U(S)$; geometrically it is a crossing between
the two halves and appears in neither half's diagram. Every old label outside
$\mathcal N$ has both visits in one interval, so its interlacements with other
such labels and its cyclic visit order are those of that half. The over/under
role of each persistent old crossing is the sign of its $\mathrm{G5}$
determinant, which is unchanged across the event: the crossing persists, so that
member is \emph{active} and hence relevant, and a relevant member outside the
zero set keeps its sign by Lemma~\ref{lem:guardconst}. An earlier revision
called $\mathrm{G5}$ unconditional, which it is not. By Lemma~\ref{lem:carriers}(iv) each
residual component sits on the unique carrier holding all its visits, and
clause~(A) identifies that carrier with its half mate. Hence the signed residual words agree; equal records give equal links by
Axiom~\ref{ax:gausscode}, and hence equal polynomials, and the writhes agree
because they are crossing counts. An earlier revision cited
Axiom~\ref{ax:homfly} for the passage from equal records to equal polynomials,
which is the axiom that computes a polynomial, not the one that identifies the
link a record presents. The cycle
$\{y_0,x_1\}$ contains only newborn visits, so it carries no residual label,
giving the last sentence.


\emph{Clause (C).}
Choose a contact disc containing no old vertex, no old crossing and no
smoothing site, which exists because the contact is isolated. Inside it, deform
each smoothed arc of $P_2$ to the corresponding closure of the half, keeping
the disc boundary fixed. Along the deformation every corner of the carrier is a
polygon vertex or a smoothing site of a crossing of $S$, and no corner
determinant vanishes: the vertex turns are guarded by (G1), the smoothing-site
turns by the active (G5) of their crossings, and the two directions incident to
the contact are nonzero and transverse to the remote direction by
Lemma~\ref{lem:contactturnboth}(ii). Ordinary crossing visits are not corners, so
nothing inside the disc passes through a degenerate turn. The cyclic corner list is fixed along the deformation. Every incident
edge is nonzero and no consecutive pair is a negative multiple, because the
corresponding corner determinant was just shown nonzero. Thus the carrier is a
path in its fixed $\mathcal R_c$, and Lemma~\ref{lem:turnlift}(ii) keeps its
signed rotation constant; the corner signs outside the disc and the matched contact turn keep their signs;
their products, hence the weights of Definition~\ref{def:wind}, agree. This gives the first sentence.

The local carrier is the closed triangle with corners at $x$, at $M$ and at
$y$. If $M$ lies to the left of the oriented remote edge in the low chamber,
i.e.\ $s_{\mathrm{low}}=+1$, then in the high chamber the triangle is traversed clockwise,
with three right turns; Lemma~\ref{lem:turnlift}(ii) gives signed rotation
$-1$. If $s_{\mathrm{low}}=-1$ it is traversed counterclockwise with three left turns, and
the same lemma gives signed rotation $+1$. In
both cases $\rot(U_{\mathrm{loc}})=-s_{\mathrm{low}}$ and $R(U_{\mathrm{loc}})=1$. Its weight is $+1$ in the first case
and $(-1)^3=-1$ in the second (Definition~\ref{def:wind}), i.e.\
$\mathrm{wt}(U_{\mathrm{loc}})=s_{\mathrm{low}}$. It bears no piece by clause~(B), so
$P_{S,U_{\mathrm{loc}}}=1$, $w_{S,U_{\mathrm{loc}}}=0$, its slot is $1-0-1=0$ and
$\Omega_1(S,U_{\mathrm{loc}})=[a^0z^0]1=1$.



\emph{Clause (D).} If $\epsilon=1$, then $x$ and $y$ interlace and hence are adjacent in
$G_{P_2}$ (Definition~\ref{def:interlace}). No independent set contains both
(Definition~\ref{def:smoothing}), so the both-newborn sector has no supports
and its empty sum is $\mathcal B=0$.

\emph{Clause (E).} Assume $\epsilon=0$. Fix an arbitrary eligible old support $T$, let
$(T_A,T_B)$ be its ordered parts from Lemma~\ref{lem:orderedproduct}, and put
$S=T\cup\{x,y\}$. Let $s_{\mathrm{low}}$ be the sign \eqref{eq:S6} in $P_0$.
By clauses~(A)--(C), the term of $S$ factors as
\[
\wind(S)\prod_{L}\Omega_1(S,L)
=\Bigl(\mathrm{wt}(U_{\mathrm{loc}})\!\!\prod_{L\neq U_{\mathrm{loc}}}\!\!\mathrm{wt}(L)\Bigr)
 \Bigl(\Omega_1(S,U_{\mathrm{loc}})\!\!\prod_{L\neq U_{\mathrm{loc}}}\!\!\Omega_1(S,L)\Bigr)
=s_{\mathrm{low}}\,t_{\lambda_1}(T_A)t_{\lambda_2}(T_B),
\]
where $t_{\lambda_i}$ denotes the corresponding term of the half state sum,
because every factor other than the local one matches a factor of exactly one
half. The both-newborn supports are exactly the $T\cup\{x,y\}$ with $T$
eligible: independence follows from eligibility and $x\not\sim y$, while the
old part of any both-newborn support misses the common neighbourhood.
Lemma~\ref{lem:orderedproduct} indexes these supports by
$\Ind(G_A)\times\Ind(G_B)$, so finite distributivity gives
\[
 \sum_T s_{\mathrm{low}}t_{\lambda_1}(T_A)t_{\lambda_2}(T_B)
 =s_{\mathrm{low}}\Bigl(\sum_{T_A}t_{\lambda_1}(T_A)\Bigr)
  \Bigl(\sum_{T_B}t_{\lambda_2}(T_B)\Bigr)
 =s_{\mathrm{low}}X_1(\lambda_1)X_1(\lambda_2),
\]
including the empty supports in both complete half state sums. Let
$s_{\mathrm{init}}$ be the sign in the initial chamber used by
Definition~\ref{def:sectors}. The two chambers have opposite side signs; the
initial chamber is $P_0$ for $\delta=+1$ and $P_2$ for $\delta=-1$. Hence
$s_{\mathrm{init}}=\delta s_{\mathrm{low}}$, and the directed sector is
\[
 \mathcal B=\delta s_{\mathrm{low}}X_1(\lambda_1)X_1(\lambda_2)
 =s_{\mathrm{init}}X_1(\lambda_1)X_1(\lambda_2)=J.
\]

\emph{Clause (F).} Clauses (D) and (E), with $\epsilon\in\{0,1\}$, give
$\mathcal B=(1-\epsilon)J$. By Definition~\ref{def:sectors},
$F-J=R+\mathcal B-J=R-\epsilon J$, proving the equivalence.
\end{proof}

\subsection{The returned sector: framing and the spanning recurrence}
\label{sec:spanning}

By Proposition~\ref{prop:sectortarget}(F) the bigon branch of (S7) is exactly the
statement $R=\epsilon J$. This subsection puts the returned sector in a form in
which the interactions between the two contact intervals drop out.

\begin{definition}[the four returned branches]\label{def:branches}
Old labels have the same interlacements in $P_0$ and in $P_2$, since those
depend only on the cyclic order of their four visits, which the newborn visits
do not disturb. So $\Ind(G_{P_0})$ is exactly the set of newborn-free supports
of $P_2$. For $T\in\Ind(G_{P_0})$ write $L(T)$ for its term in $X_1(P_0)$,
$N(T)$ for its term in $X_1(P_2)$, and $X(T)$, $Y(T)$ for the terms of
$T\cup\{x\}$ and $T\cup\{y\}$ in $X_1(P_2)$, each read as $0$ when the set is
not independent. Then
\[
   R=\delta\sum_{T\in\Ind(G_{P_0})}\bigl(N(T)-L(T)+X(T)+Y(T)\bigr).
\]
\end{definition}

\subsection{The selector identities}\label{sec:selectors}



\begin{lemma}[the contact carrier splits]\label{lem:contactsplit}
Here eligibility is as in Definition~\ref{def:eligible}.
Under Convention~\ref{conv:s7standing}, let $T$ be eligible, with parts
$T_A,T_B$. The genericity of the halves is stated here,
where the half objects are first formed, and not only at the endpoint theorems
that consume them: the conclusion speaks of the carriers of $T_A$ in
$\lambda_1$ and of $T_B$ in $\lambda_2$, and a carrier of a support is read off
a polygon's Gauss word and its smoothing sites, which a nongeneric half need
not possess. The hypothesis is (H3) of Convention~\ref{conv:s7standing}, in force for this
section by that convention's declaration; it is named here because this is
where the half objects are first formed, not because naming is the mechanism. In the high-neither state the
carriers of $T$ in $P_2$ are: one distinguished carrier $L^{*}$ containing the
corner $M$ and the remote-edge stretch through the contact, and a family of
others, each of which lies entirely in $A$ or entirely in $B$. Cutting at the
contact carries $L^{*}$ to one carrier $L_1$ of $T_A$ in $\lambda_1$ and one
carrier $L_2$ of $T_B$ in $\lambda_2$, each containing that half's contact
corner, and carries the others bijectively to the remaining carriers of the two
halves.

Five further data are part of the conclusion, because they are used later
and a statement's consumers must find their inputs in its clauses
rather than in a picture:
\begin{enumerate}
\item[(a)] \emph{Corner partition.} Every corner of every carrier of $T$ in
$P_2$ other than the contact corner $M$ of $L^{*}$ --- each being a vertex of
the polygon or a smoothing site, lying strictly inside $A$ or strictly inside
$B$ --- is a corner of exactly one carrier of exactly one half, namely the
half of its interval, under the correspondence above.
\item[(b)] \emph{Ordered direction pairs.} At each such corner the ordered
pair (incoming direction, outgoing direction) read in the half equals the pair
read in $P_2$: the cut alters the curve only at the two contact stretches, and
a corner strictly inside an interval keeps both its incident arcs.
\item[(c)] \emph{Turn signs.} Consequently the principal turn at each such
corner, a function of that ordered pair, is equal in the half and in $P_2$ ---
in magnitude and in sign.
\item[(d)] \emph{The contact corner.} The single corner $M$ of $L^{*}$ is
replaced by exactly two corners: the contact corner of $L_1$ in $\lambda_1$
and the contact corner of $L_2$ in $\lambda_2$, the ones
Lemma~\ref{lem:contactturnboth} computes. In particular
$c(L_1)+c(L_2)=c(L^{*})+1$.
\item[(e)] \emph{Compatibility.} The bijection of carriers matches each
non-distinguished carrier with a half carrier having the \emph{same} corners
with the same ordered direction pairs and turn signs --- (a)--(c) restricted to
it --- so matched carriers have equal corner counts, equal turn-sign lists, and
equal uniformity status; and $L^{*}$ corresponds to the pair $(L_1,L_2)$ with
(a)--(d).
\end{enumerate}
\end{lemma}
\begin{proof}
$T$ is eligible, so $T\subseteq V_A\sqcup V_B$
(Lemma~\ref{lem:orderedproduct}) and every smoothing site lies strictly inside
$A$ or strictly inside $B$. The newborns are not smoothed, so the traversal
runs uninterrupted through $x_0y_0$ --- the stretch carrying the corner $M$ ---
and through the two remote-edge visits. Write the word as in \eqref{eq:contactword}. Because no newborn is smoothed, the
successor map moves through the four newborn visits without reconnecting: from
$x_0$ it continues to $y_0$ and from $y_0$ into $A$; along the remote stretch it
passes the two remote-edge newborn visits and continues into $B$; and from the
end of $B$ it returns to $x_0$. So the arcs $x_0y_0$ and the remote stretch lie
on one and the same successor cycle, together with whichever arcs of $A$ and of
$B$ the smoothing of $T_A$ and $T_B$ leaves attached to them. \emph{Interval subclaim.} Let an open interval of a cyclic traversal carry finitely many chords,
each with both visits inside it and with no two visit pairs interlacing. Smoothing all of them
leaves exactly one open successor strand joining the two boundary half-edges; every other
successor cycle is internal, with all its arcs strictly inside the interval.

Induct on the number of chords. With none, the interval itself is the one boundary strand and
there is no internal cycle. For the step, order the chords by first visit and choose an
\emph{innermost} one, with no other chord having both visits between its two; one exists because
finiteness and noncrossing leave only nesting. Smoothing it first replaces the sub-arc between
its visits by a closed internal cycle --- the sub-arc contains no other chord visit --- and leaves
the ambient strand entering and leaving where it did, with the two visits spliced out. The
remaining chords lie on that strand or in internal cycles. Those on the ambient strand form a
smaller noncrossing family on an interval with the same boundary half-edges, so induction
applies; a chord in an internal cycle splits it into two internal cycles and does not touch the
boundary strand. Thus exactly one boundary strand remains and every new cycle is internal.
Nested chords give the innermost-first recursion; the empty interval is the base case.
Reconnections at distinct chord pairs act on disjoint visit half-edges and therefore commute,
so simultaneous smoothing may be evaluated in this innermost-first order.

Apply the subclaim first to $A$. The chords of $T_A$ lie wholly inside $A$ and are pairwise
noncrossing there: eligibility makes the support independent, and interlacement of two chords
whose four visits lie in one interval is exactly their crossing there. Hence smoothing $T_A$
leaves one open strand joining the two boundary half-edges of $A$ and otherwise only internal
cycles. The same argument for $B$ and $T_B$ gives its unique boundary strand and internal
cycles. These two boundary strands are exactly the arcs attached to the $x_0y_0$ stretch and
the remote stretch, so they belong to the distinguished cycle, while every remaining cycle is
internal to its interval. Hence the cycles are exactly those inside $A$, those inside $B$, and
the single cycle carrying both stretches, namely $L^{*}$. An earlier revision asserted the
internal-cycle claim in one sentence; the induction above carries it.

Cutting the polygon at the contact replaces the two stretches by the two
closing sub-segments of Definition~\ref{def:halves}: the arcs of $L^{*}$ lying
in $A$ close up into a carrier $L_1$ of $\lambda_1$, and those in $B$ into a
carrier $L_2$ of $\lambda_2$, because the cut separates exactly the two
stretches and the smoothing of $T_A$, $T_B$ is unchanged. The other cycles are
untouched and become carriers of the half they lie in.

For the five exported clauses. (a): every non-contact corner is a vertex or a
smoothing site strictly inside one interval, the newborns being unsmoothed and
$M$ the only corner on the contact stretches; the cut leaves the interior of
each interval pointwise fixed, so the corner lies on the one half carrier that
its cycle becomes. (b): both arcs incident to such a corner lie strictly
inside its interval --- a corner is not a boundary half-edge --- and the cut
alters the curve only at the two contact stretches, so the incoming and
outgoing directions are unchanged as an ordered pair. (c): the principal turn
is a function of that ordered pair, so it is unchanged in magnitude and sign.
(d): the corner $M$ of $L^{*}$ lies on the cut itself; the two closing
sub-segments of Definition~\ref{def:halves} give $L_1$ and $L_2$ one new
corner each at the contact, the ones Lemma~\ref{lem:contactturnboth} computes,
and no other corner is created, the closing sub-segments being straight. Hence
$c(L_1)+c(L_2)=(c(L^{*})-1)+2$. (e): a non-distinguished carrier is untouched
by the cut --- its cycle lies inside one interval --- so its matched half
carrier has literally the same corner set with the same pairs and turns, by
(a)--(c); equal turn-sign lists give equal uniformity status
(Definition~\ref{def:wind}); and $L^{*}$'s data pass to $(L_1,L_2)$ by
(a)--(d).
\end{proof}

\begin{lemma}[minimum carrier corners and selector identities]\label{lem:selectorid}
\begin{enumerate}
\item[(A)] Under the guards of Definition~\ref{def:generic}(A), every carrier of every
support has at least three corners.
\item[(B)] Here eligibility is as in Definition~\ref{def:eligible}.
Let $T$ be eligible, with parts $T_A,T_B$. Write $P_{\mathrm L}$ and
$P_{\mathrm H}$ for the low and high chambers of the event --- upright
subscripts, since the italic $P_H$ already names the polynomial of a residual
piece $H$ (Corollary~\ref{cor:groupedknot}(B)), and an earlier revision wrote
$\wind_{P_H}$ with the two collided. Write $W_T=\wind_{P_{\mathrm H}}(T)$,
$W_1=\wind_{\lambda_1}(T_A)$ and $W_2=\wind_{\lambda_2}(T_B)$, each with its
argument and its ambient polygon. Let $s_{\mathrm{low}}$ be the sign \eqref{eq:S6} evaluated with $P_-$ taken to
be the \emph{low} chamber of the event --- the sign of
$\det(p_b-p_a,\,p_M-p_a)$ there --- which is a datum of the event and of the
chamber, bound here and not borrowed from a lemma printed nine hundred lines
later. Then
\[
\epsilon=1\ \Longrightarrow\ W_T=-s_{\mathrm{low}}\,W_1W_2,
\qquad
\epsilon=0\ \Longrightarrow\ W_T\,W_1W_2=0 .
\]
Moreover, at $\epsilon=1$, the contact carrier $L^{*}$ is uniform if and only if
both half contact carriers $L_1,L_2$ are, and then all three are uniform with
the same turn sign. That clause is proved below with the identities and is
stated here because a consumer reads it as a statement --- an earlier revision
left it inside this proof and cited the lemma for it.

The sign here is $s_{\mathrm{low}}$, fixed by the chamber, and not the
\eqref{eq:S6} sign $s$, which depends on which chamber the parametrisation calls
$P_-$; an earlier revision wrote a bare $s$.
\end{enumerate}
\end{lemma}
\begin{proof}
\emph{Clause (A).} A carrier is a closed curve made of straight segments joined at corners of
nonzero turn. With two corners it would consist of two straight segments
joining the same two points, so the two segments would overlap, which puts an
endpoint of one on the other and vanishes a $\mathrm{G2}$. With one corner it
would contain a straight return of zero length, which is not a segment of a
generic polygon. With no corner at all it would be a closed curve made of a
single straight segment, whose two endpoints coincide, so that segment has zero
length --- excluded for the same reason, and by
Definition~\ref{def:regular}(B)'s requirement that every edge be nonzero. An
earlier revision omitted this case. So there are at least three. In particular a zero selector
means mixed turns, and no one- or two-corner convention is needed.

\emph{Clause (B).} By Lemma~\ref{lem:contactsplit} the carriers other than $L^{*},L_1,L_2$ match
one for one and lie in one half each, with the same corners and hence the same
weights. Write $Q$ for the product of those weights, so that
\[
   W_T=Q\cdot\mathrm{wt}(L^{*}),
   \qquad
   W_1W_2=Q\cdot\mathrm{wt}(L_1)\mathrm{wt}(L_2).
\]
$Q$ is \emph{not} cancelled --- it may be zero, and an earlier revision
cancelled it --- so what is proved below is the local identity between
$\mathrm{wt}(L^{*})$ and $\mathrm{wt}(L_1)\mathrm{wt}(L_2)$, and both displayed
identities follow by multiplying it by $Q$, which is legitimate whatever $Q$ is.
It remains to compare $\mathrm{wt}(L^{*})$ with
$\mathrm{wt}(L_1)\mathrm{wt}(L_2)$. Each of $L_1,L_2$ has at least three
corners (clause~(A), now stated above rather than quoted from
inside another lemma's proof), hence at least two besides its
contact corner.

\emph{Case $\epsilon=1$.} By Lemma~\ref{lem:contactturnboth} all three contact
turns --- the full one and the two half ones --- have the sign of
$s_{\mathrm{low}}$. The non-contact corners of $L^{*}$ partition between $L_1$
and $L_2$ with their signs. So $L^{*}$ is uniform if and only if all its
non-contact corners have the sign of $s_{\mathrm{low}}$, if and only if $L_1$
and $L_2$ are both uniform, and then all
three are uniform of the same sign. If they are mixed, both sides of the
local identity are $0$. If the common sign is right, then $s_{\mathrm{low}}=-1$
and all three weights are $+1$, so
$\mathrm{wt}(L^{*})=1=-s_{\mathrm{low}}\mathrm{wt}(L_1)\mathrm{wt}(L_2)$. If it
is left, then $s_{\mathrm{low}}=+1$, $\mathrm{wt}(L^{*})=(-1)^{c(L^{*})}$ and
$\mathrm{wt}(L_1)\mathrm{wt}(L_2)=(-1)^{c(L_1)+c(L_2)}=(-1)^{c(L^{*})+1}$, so
$\mathrm{wt}(L^{*})=-\mathrm{wt}(L_1)\mathrm{wt}(L_2)
=-s_{\mathrm{low}}\mathrm{wt}(L_1)\mathrm{wt}(L_2)$. Multiplying by $Q$ gives
$W_T=-s_{\mathrm{low}}W_1W_2$.

\emph{Case $\epsilon=0$.} Now the two half contact turns have the sign of
$s_{\mathrm{low}}$ and the full contact turn the sign of $-s_{\mathrm{low}}$.
Suppose $\mathrm{wt}(L^{*})\neq0$; then $L^{*}$ is uniform, necessarily of the
sign of its contact turn, i.e.\ of $-s_{\mathrm{low}}$, so all its non-contact
corners have that sign. Each half carrier inherits at least one such corner and
also carries its own contact corner, of sign $s_{\mathrm{low}}$; so each is
mixed and $\mathrm{wt}(L_1)\mathrm{wt}(L_2)=0$. Symmetrically, if that product
is nonzero then both halves are uniform of sign $s_{\mathrm{low}}$, so every
non-contact corner has that sign, and $L^{*}$ carries those together with its
contact corner of sign $-s_{\mathrm{low}}$, hence is mixed and
$\mathrm{wt}(L^{*})=0$. Either way the local product vanishes, and multiplying
by $Q^2$ gives $W_TW_1W_2=0$.
\end{proof}

\begin{remark}[what the noninterlacing identity replaces]\label{rem:noscalarfact}
Lemma~\ref{lem:selectorid}, clause~(B) does \emph{not} say that in the noninterlacing
branch the full selector factors through the half selectors with a sign. It
says the opposite: a live full selector and a live product of half selectors
are mutually exclusive. Any argument that replaces the full selector by the
half-selector product in that branch is refuted by that clause, not supported by
it.
\end{remark}

\subsection{The interlacing endpoint}\label{sec:sr12}

\begin{lemma}[what the two chambers of a bigon event share]\label{lem:chambercommon}
Here eligibility is as in Definition~\ref{def:eligible} and $\rot$ is as in
Definition~\ref{def:rot}.
Let $T$ be eligible at a simple bigon event. Write $C_{\mathrm L}$ for its
contact carrier in the low chamber and $C_{\mathrm H}$ for the same carrier in
the two-crossing chamber. Then, with no hypothesis on any selector:
\begin{enumerate}
\item[(i)] every carrier of $T$ other than the contact carrier is the same
closed curve, with the same corners, the same turn signs and the same residual
pieces, in the two chambers; write $U_T$ for the product of the reads
$\Omega_1$ over those carriers, one number for both chambers;
\item[(ii)] $C_{\mathrm L}$ and $C_{\mathrm H}$ have the same corners and the
same turn signs, hence $\wind(T)$ takes one value $W_T$ in both chambers, and
$\rot(C_{\mathrm L})=\rot(C_{\mathrm H})$, hence
$R(C_{\mathrm L})=R(C_{\mathrm H})$.
\end{enumerate}
\end{lemma}
\begin{proof}
(i) The two chambers differ by moving the single vertex $M$ across the remote
edge, and the newborn crossings lie in the bigon disc that move sweeps. A
carrier not through the contact contains neither $M$ nor either newborn, so
neither its corner list nor its crossing list changes; the reads over them
therefore agree, and their product is one number $U_T$.

(ii) The two copies of the contact carrier differ by a Reidemeister-II finger.
Neither newborn is a polygon vertex and neither lies in $T$, so neither is a
corner; the corner list is therefore the same list of polygon vertices and
smoothing sites, and each corner's turn is a determinant of two edge directions
whose sign is guarded by (G1) or (G5) and constant on the event's interval
(Lemma~\ref{lem:guardconst}). Hence the weights, and $\wind(T)$, agree. During a standard finger
deformation this fixed cyclic corner list has nonzero edges and no consecutive
negative-multiple pair: at every corner its two directions have the guarded
nonzero determinant just identified. It is therefore a path in the corresponding
$\mathcal R_c$, so Lemma~\ref{lem:turnlift}(ii) keeps $\rot$ unchanged.
\end{proof}

\begin{definition}[the marked data of an eligible support]\label{def:markeddata}
Under Convention~\ref{conv:s7standing}, let $T$ be eligible with parts
$T_A,T_B$, and let
$L_1,L_2$ be the contact carriers of $T_A$ in $\lambda_1$ and of $T_B$ in
$\lambda_2$ (Lemma~\ref{lem:contactsplit}). Write
\[
   f_i(a)=\bigl[z^0\bigr]Q_{i}(a,z),\qquad
   d_i=1-w_i-R_i,\qquad
   \omega_i=\bigl[a^{d_i}\bigr]f_i,\qquad
   D=d_1+d_2,
\]
where $Q_i$ and $w_i$ are the grouped polynomial and grouped writhe of $L_i$,
$\rot_i=\rot(L_i)$ its signed rotation and $R_i=|\rot_i|$, so that $\Omega_1$ on
$L_i$ is $\omega_i$. The signed rotation carries one name in this document,
$\rot$ with a subscript --- numerals for halves, upright $\mathrm L,\mathrm H$
for chambers; earlier revisions also wrote $\rho_i$ and $r_i$ for the same
quantity, which is three names for one thing. The chambers and the half selectors are bound here as
well, since both are used outside the lemma that first introduced them:
$P_{\mathrm L}$ and $P_{\mathrm H}$ are the low and the high chamber of the
event --- upright subscripts, the italic $P_H$ being the polynomial of a
residual piece $H$ --- and
\[
   W_1=\wind_{\lambda_1}(T_A),\qquad W_2=\wind_{\lambda_2}(T_B)
\]
are the two half selectors, each with its argument and its ambient half. An
earlier revision left them bound only inside Lemma~\ref{lem:selectorid}, clause~(B) and used
them in several statements after it. The letter $Q$ is used for a half's grouped polynomial and
$F_\nu$ for a chamber's, $\nu\in\{\mathrm L,\mathrm H,\mathrm A\}$. For the
symbols bound in this definition --- $f_i$, $d_i$, $\omega_i$, $Q_i$, $w_i$,
$\rot_i$, $R_i$, $W_i$, $L_i$ --- a numerical subscript means a half and an
upright $\mathrm L$ or $\mathrm H$ means a chamber. The rule is stated for those
symbols and not as a general convention, because chambers elsewhere do carry
numerical names: $P_0$ and $P_2$ are the no-crossing and two-crossing chambers
of a vertex-on-edge event, and $d_0,d_1$ are the incident wall tangents of the
sliding branch. An earlier revision wrote "numerical subscripts mean halves and
never chambers", which those three names contradict.

\emph{The chamber data.} For $\nu\in\{\mathrm L,\mathrm H\}$ write $C_\nu$ for
the contact carrier of $T$ in that chamber and $D_\nu$ for the grouped diagram
that carrier bears (Corollary~\ref{cor:groupedknot}(B)), and put
\[
   w_\nu=w(D_\nu),\quad R_\nu=R(C_\nu),\quad \rot_\nu=\rot(C_\nu),\quad
   d_\nu=1-w_\nu-R_\nu,\quad F_\nu=P_{D_\nu},\quad
   \Omega_\nu=\bigl[a^{d_\nu}z^0\bigr]F_\nu,
\]
the writhe of that diagram, the absolute and signed rotations of its curve, the
slot, the diagram's polynomial and its read. The quantities are defined
\emph{from the diagrams} $D_{\mathrm L},D_{\mathrm H}$, which is what the
full-twist lemma quantifies over; an earlier revision defined them from the
chambers and left the diagrams to be inferred. These six symbols are used throughout the sections
below and were bound only in running prose, which is not a binder.

Two factors are bound \emph{separately}, an earlier revision having named only
their product:
\[
   W_T=\wind_{P_{\mathrm H}}(T)=\wind_{P_{\mathrm L}}(T)
   \quad\text{(the two agree by Lemma~\ref{lem:chambercommon}(ii))},
\]
and $U_T$ the product of the reads $\Omega_1$ over the carriers of $T$ other
than the contact carrier --- the untouched ones, common to the two chambers by
Lemma~\ref{lem:chambercommon}(i), which is the clause about those carriers;
clause (ii) is about the contact carrier and the selector, and an earlier
revision cited it here. The term of $T$ is then $W_TU_T$ times the read on the contact
carrier.
\end{definition}

\begin{remark}[why the hypothesis is the selector alone]\label{rem:sr12hyp}
The hypothesis of Lemma~\ref{lem:triplebridge}(vii) is $W_T\neq0$ and not $W_TU_T\neq0$. The proof below uses only
that the full carrier through the contact is live, which is what $W_T\neq0$
gives; the spectator product $U_T$ appears nowhere in it. The recorded defect was actual: an earlier revision's statement asked for
$W_TU_T\neq0$, while its consumer had only $W_T\neq0$; the invocation therefore
exceeded the stated hypothesis --- the peer's point. The repair weakens that
statement to $W_T\neq0$, which its proof consumes, and the live full-twist
calculation below consumes exactly that repaired interface.

\end{remark}

\subsection{The sliding branch}\label{sec:sliding}

Now let the event be of sliding type: the two neighbours $p=p_{M-1}$ and
$q=p_{M+1}$ lie strictly on \emph{opposite} sides of the line of the remote
edge. Convention~\ref{conv:s7standing} is in force throughout this subsection: the
restart changes the branch, not the standing hypotheses, and (H3) is what makes
the halves of this branch polygons that $X_1$ is defined on. Write $r=p_b-p_a$,
$d_0=p_M-p$, $d_1=q-p_M$ for the three directed wall tangents at the event.

\begin{lemma}[signed-word transport, the containing-sector Cartesian product, and the chamber-to-half carrier map]
\label{lem:slidingcartesian}
\begin{enumerate}
\item[(A)] For this statement, the \emph{signed cyclic Gauss word} of a generic polygon
means its Gauss word (Section~\ref{sec:setup}) together with, at each double
point, the record of which visit is the overpass under the divide convention of
Definition~\ref{def:piecediagram}. Every divide crossing is positive, so this
is the whole decoration; there is no separate crossing-sign decoration.

Let $x_-$ be the crossing present in $P_-$ and $x_+$ the one present in $P_+$,
and let $\phi$ relabel $x_-$ as $x_+$ and fix every other label. Then $\phi$ is
an isomorphism of those signed cyclic Gauss words
$W(P_-)\to W(P_+)$. Consequently it
is an isomorphism of interlacement graphs, a bijection of supports, and it
matches smoothing cycles, undominated sets, residual components, their signed
piece diagrams, their polynomials, their writhes, and the carrier owning each
component.
\item[(B)] Let the event be a simple sliding vertex-on-edge event, let
$P_\bullet\in\{P_-,P_+\}$ be either chamber, let $x_\bullet$ be the relocated
crossing in it --- $x_-$ in $P_-$ and its image $x_+$ in $P_+$ --- and write
$G=G_{P_\bullet}$ and $G_{\lambda_1},G_{\lambda_2}$ for the interlacement graphs
of that chamber and of the two halves at the wall. Nothing below distinguishes
the two chambers: the statement is proved once for either, and an earlier
revision stated it for $P_-$ alone while its consumers quantified over both.
Cutting the traversal circle of $P_\bullet$ at the two visits of $x_\bullet$
separates the old labels into two intervals;
the induced subgraph of $G$ on the labels of one interval \emph{is} the
interlacement graph of the corresponding half, since interlacement of two chords
with all four visits inside one interval is read from their order inside it,
which the half inherits. With that identification,
\[
   \{\,S\in\Ind(G):x_\bullet\in S\,\}
   \;\longleftrightarrow\;
   \Ind(G_{\lambda_1})\times\Ind(G_{\lambda_2}),
   \qquad
   S\longmapsto\bigl(S_1,S_2\bigr),
\]
is a bijection, where $S_i=(S\setminus\{x_\bullet\})\cap V(G_{\lambda_i})$.
\item[(C)] Fix a simple sliding event, a chamber $P_\bullet\in\{P_-,P_+\}$, and a support
$S\in\Ind(G_{P_\bullet})$ containing the relocated crossing. Let $S_1,S_2$ be its
halves (Lemma~\ref{lem:slidingcartesian}(B)). Let $\chi_\bullet$ send a carrier $C$
of $S$ in $P_\bullet$ to the successor cycle of the half $\lambda_i$ carrying
the same retained visits, $i$ being the interval $C$ meets. Then $\chi_\bullet$
is a well-defined map from the carriers of $S$ in $P_\bullet$ to the carriers of
$S_1$ in $\lambda_1$ and of $S_2$ in $\lambda_2$, and it has the following
properties.
\begin{enumerate}
\item[(i)] \emph{Correspondence.} Every carrier of $S$ in $P_\bullet$ meets
exactly one of the two intervals, and $\chi_\bullet$ is a bijection onto the
union of the two halves' carrier sets. Exactly two carriers pass through the
contact, one in each interval, and they are sent to the two halves' contact
carriers, one each. Only one of the two carries the short collapsing edge, and
it is the one in the interval containing that edge.
\item[(ii)] \emph{Records.} For each carrier $C$, the retained-visit map from
$C$'s marked points to those of $\chi_\bullet(C)$ is an orientation-preserving
record isomorphism (Definition~\ref{def:record}).
\item[(iii)] \emph{Owners and grouped data.} The residual pieces of $S$ carried
by $C$ are exactly the residual pieces of the corresponding half support carried
by $\chi_\bullet(C)$; each has the same label set, hence the same piece diagram,
polynomial and writhe. Thus the matched carriers have the same grouped polynomial
and grouped writhe.
\item[(iv)] \emph{Corners.} Away from the contact, $C$ and $\chi_\bullet(C)$
have the same corner list with the same signs. Of the two carriers through the
contact, the one \emph{not} carrying the short collapsing edge also has the same
corner list as its image; the one carrying it has, at the contact, one direction
more in its local direction list than its image has, and nothing is claimed here
about the two lists beyond that --- the refinement is computed where it is
consumed.
\item[(v)] \emph{Compatibility with the chamber map.} If $S$ contains the
relocated crossing of $P_-$ and $\phi$ is the chamber-to-chamber relabelling,
then $\chi_+\circ\phi=\chi_-$ as maps from the carriers of $S$ in $P_-$ to the
halves' carriers.
\end{enumerate}
\end{enumerate}
\end{lemma}
\begin{proof}
\emph{Clause (A).}
\emph{Source.} This proof is a transcription of the peer's pack \texttt{evidence/codex/phase15-sliding-universal-attack/REPORT.md}, SHA-256 \texttt{4f6da162\dots9318}, §``Signed word and support transport'',
lines 118--150. The argument is the peer's and not this author's; the display
numbering $\mathrm{U}$ is that document's.

The two visits of the relocated crossing occupy two local intervals of the
cyclic traversal word, and event isolation keeps every old visit out of those
intervals. Sliding replaces the visit on one incident edge by the visit on the
other, so deleting the relocated label leaves the same old cyclic word.

The opposite-side condition supplies the signed refinement: both directed
incident tangents point into the same open half-plane bounded by the directed
remote tangent, which is
\begin{equation}
   \sgn\det(r,d_0)=\sgn\det(r,d_1)\neq0 .
   \tag{U8}
\end{equation}
So the determinants against the remote tangent have the same sign, and under
the positive-crossing convention the over/under role and the oriented sign of
the gained visit are those of the lost visit. Hence $\phi$ is an isomorphism of
the \emph{signed} cyclic word and not only of its unsigned chord diagram; it
induces an isomorphism of interlacement graphs and a bijection of all
independent supports.

The smoothing successor permutation is determined by the signed traversal word
together with the selected labels, so corresponding supports have corresponding
successor cycles. The undominated residual graph is preserved with them, hence
so are its connected components. For the piece data the map is restricted
rather than waved at: for a residual piece $H$, $\phi$ carries the visits of $H$
bijectively to those of $\phi H$, preserving their cyclic order --- $\phi$ being
an isomorphism of the cyclic word --- and preserving the over/under designation
and the sign at each, by the positive-divide clause above. Those are the four
conditions of Definition~\ref{def:record}, so $\phi|_H$ is a record isomorphism
between the two piece diagrams' records, and Axiom~\ref{ax:gausscode}
identifies the links: $P_H=P_{\phi H}$ and $|H|=|\phi H|$. A residual component
is owned by the successor cycle containing its visits, so carrier ownership is
preserved as well. Every combinatorial and polynomial assertion of the statement follows. Rotation
is not among them: this lemma claims nothing about it, and the carrier rotations
are the subject of a separate lemma below.

\emph{Clause (B).}
\emph{Source.} A transcription of the peer's pack \texttt{evidence/codex/phase15-sliding-universal-attack/REPORT.md}, SHA-256 \texttt{4f6da162\dots9318}, §``Carrier rotation'', lines 165--170, where it is
Display~U6.

A support chord is independent from the relocated chord $x_\bullet$ exactly when
both of its ends lie in one half interval, since a chord with one end in each
interval interlaces $x_\bullet$. Two chords whose ends lie in \emph{different}
half intervals cannot interlace each other, their visits being in disjoint arcs.
Hence, \emph{for a set of old labels each of which is nonadjacent to
$x_\bullet$}, independence in the old graph is independence of its two halves in
the respective half graphs.

That hypothesis is not decoration and the statement is false without it: in the
word $123123$ the chords $1,2,3$ pairwise interlace, so a set may fail to be
independent while its two halves --- one of them empty --- are independent in
their own graphs. The bijection of this lemma is unaffected, because the sets it
maps are supports \emph{containing} $x_\bullet$, and independence of such a
support already forces every other member to be nonadjacent to $x_\bullet$,
which is the hypothesis. An earlier revision stated the equivalence for every
set of old labels. Conversely any pair of independent supports, one from
each half, combines with the relocated crossing to an independent support
containing it. The two assignments are mutually inverse.

\emph{Clause (C).}
Cut the traversal circle of $P_\bullet$ at the two visits of the relocated
crossing $x_\bullet$. By Lemma~\ref{lem:slidingcartesian}(B) the old labels split
into the two half intervals and the induced subgraphs are the half interlacement
graphs, so $S_i$ is a support of $\lambda_i$.

\emph{The two intervals close separately.} $S$ contains $x_\bullet$, so
$x_\bullet$ \emph{is} smoothed. An oriented smoothing at a self-crossing of one
closed curve joins each outgoing branch to the incoming branch that continues
the orientation, and so cuts the traversal circle into exactly two cycles, one
carrying the labels of each interval. Every other element of $S$ is independent
from $x_\bullet$ and therefore has both visits inside one interval
(Lemma~\ref{lem:slidingcartesian}(B)), so no later smoothing joins the two
families. Hence every carrier meets exactly one interval, and smoothing $S$
inside an interval performs exactly the reconnections that smoothing $S_i$
performs in $\lambda_i$; the successor cycles correspond one for one, which is
the bijection of (i). Two of them run through the contact, one per interval,
and the short collapsing edge lies in exactly one interval, hence on exactly one
of those two.

An earlier revision said instead that the contact carrier was ``the one cycle
that meets both intervals''. That is the picture of a \emph{bigon} event, where
the newborns are not smoothed and one cycle does meet both; here the relocated
crossing is smoothed and it is never true. The minimal support $S=\{x_\bullet\}$
is the sharpest case: nothing else is smoothed, and the two intervals close into
two carriers, each mapping to the contact carrier of its half. Those carriers
need not be pieceless --- a minimal support forbids further \emph{smoothings},
not residual pieces, and the half may carry pieces of its own --- and an earlier
revision called them pieceless. The
selector identity below already reads the two-carrier picture --- it multiplies
one factor $c$ for the chamber's short-edge carrier by one factor $h$ for the
other --- so the correction removes an inconsistency rather than creating one.

For (ii): the correspondence is the identity on the retained visits, it
preserves their cyclic order by Lemma~\ref{lem:carrierword} applied inside each
interval, and it preserves the over/under designation and the sign at each,
those being read from the two branch directions, which the cut does not move.
For (iii) the label sets must first be shown to agree, and an earlier revision
assumed it here and proved it later from this very clause. The argument is
three lines and belongs here. Remove the relocated crossing from $S$ and
intersect what remains with the two half label sets, which is the definition of
$S_1,S_2$; every label interlacing the relocated crossing is dominated and so
absent from $U(S)$; every remaining label has both visits in one interval, and
there is no interlacement edge between the two intervals, a chord with ends in
different intervals interlacing the relocated chord and so not being present. So
\[
   G_{P_\bullet}\bigl[U(S)\bigr]=G_{\lambda_1}\bigl[U(S_1)\bigr]\sqcup
                     G_{\lambda_2}\bigl[U(S_2)\bigr],
\]
and the residual components on the two sides are the same sets of labels. Then
Lemma~\ref{lem:carriers}(iv), applied in $\lambda_i$, puts each component on
one carrier, and Definition~\ref{def:piecediagram} gives it the same piece
diagram, the label set and its signs being unchanged; the passage from the equal
records to equal polynomials is Axiom~\ref{ax:gausscode}, which is what licenses
reading a polynomial off a record at all. For (iv): a corner of a carrier is a polygon vertex or a smoothing site
(Definition~\ref{def:wind}), and away from the contact each is present in
$C$ and in $\chi_\bullet(C)$ with the same incoming and outgoing directions,
since the wall moves only the sliding vertex; so the corner lists and their
signs agree there. The carrier through the contact that does not carry the short
edge meets the contact in a single corner, present in both with the same
directions, so its whole list agrees too. The remaining carrier traverses the
short edge, which the wall contracts to a point: its local direction list at the
contact has the two directions of the collapsing corner where the image has
one, which is the one extra direction claimed, and no more is asserted here.

For (v): $\phi$ is the identity on the old labels and carries $x_-$ to $x_+$, so
the two chambers' traversal circles have the same cyclic word on the labels of
$S$ and $\phi S$, the same cut, and the same reconnections. The successor cycles
therefore correspond under $\phi$, and each pair is sent to the same half
carrier, since the half data are read from the same interval word in both. That
is $\chi_+\circ\phi=\chi_-$.
\end{proof}

\begin{remark}[what this map replaces]\label{rem:chamberhalfwhy}
Three arguments below --- rotation transport, coefficient-product transport and
selector transport --- previously each reached for the halves in its own way: by
``restricting'' the chamber-to-chamber relabelling $\phi$, which is a map
between the two \emph{chambers} and not between a chamber and a half, or by
citing Proposition~\ref{prop:sectortarget}(A), which is stated for the both-newborn support
of a \emph{bigon} event and not for a sliding one. The missing object is one and the
same map, and it is printed here.
\end{remark}

\begin{lemma}[rotation transport, sliding incidence, the collapsing half, coefficient products, selector transport and the sliding jump]\label{lem:slidingcoeff}
\begin{enumerate}
\item[(A)] Here $\rot$ is as in Definition~\ref{def:rot}.
Under $\phi$, corresponding carriers have the same signed rotation and the same
$R$. More precisely:
\begin{enumerate}
\item[(i)] for a support \emph{avoiding} the relocated crossing, every carrier
of one chamber is a regular homotopy of its mate, and they have the same signed
rotation, the same $R$, \emph{and the same corner-sign list};
\item[(ii)] for a support \emph{containing} it, every carrier $C$ of either
chamber has the signed rotation and $R$ of its wall-half mate
$\chi_\bullet(C)$; hence $\chi_+\circ\phi=\chi_-$ gives the two chambers the same
signed and absolute rotations, carrier by carrier. No claim is made about
corner-sign lists in this sector: the carrier carrying the short collapsing edge
has one corner more than its wall-half mate, so the lists have different lengths.
\end{enumerate}
\item[(B)] Under Convention~\ref{conv:s7standing}, let the event be of sliding type, write
$r=p_b-p_a$ for the direction of the remote edge and $d_0=p_M-p_{M-1}$,
$d_1=p_{M+1}-p_M$ for the two incident edge directions at the moving vertex. These are the three directed wall tangents, restated
here rather than imported by citation.
\begin{enumerate}
\item[(i)] At a simple sliding event, $\sgn\det(r,d_0)=\sgn\det(r,d_1)\neq0$. Moreover
there is an interval of the event path around the wall on which both statements
below hold, and they are about different objects. First, each incident
\emph{line} meets the \emph{line} of the remote edge at a point whose parameter
along the remote \emph{segment} lies strictly inside $(0,1)$ --- that is a
condition on the remote segment, and it holds for both incident lines. Second,
exactly one incident \emph{edge} meets the remote edge in the relative interiors
of both, in each of the two chambers, and the two chambers exchange which
incident edge that is; the other incident edge has its line meeting the remote
segment at a point beyond its own far endpoint. No other crossing changes on the
interval. An earlier revision compressed the two into one sentence that read
"each incident line meets the remote line at a parameter strictly inside $(0,1)$
along both segments", which asks the intersection to lie inside the incident
segment as well and so contradicts the "exactly one" it was supporting.

The interval restriction is not a formality. Opposite sides of the remote
\emph{line} do not by themselves put the intersection inside the remote
\emph{segment}: the segments $(0,0)$--$(1,0)$ and $(2,1)$--$(2,-1)$ have
endpoints strictly on opposite sides of the first one's line and meet that line
at $(2,0)$, outside the first segment. The peer's example is that one. Since the
incidence parameter is continuous and equals the contact point's at the wall,
which lies strictly inside the remote segment by
Convention~\ref{conv:events}, shrinking the interval puts it inside for every
parameter of the interval. The statement is phrased symmetrically in the two
chambers rather than assigning a direction, which the event's parametrisation
does not fix.
\item[(ii)] Let $s$ be the sign
\eqref{eq:S6} taken in the initial chamber. Put
\[
   \eta_1=\sgn\det(r,d_1),\qquad \eta_2=\sgn\det(d_0,r).
\]
Then $\eta_2=-\eta_1$; in each chamber exactly one half carries the short
collapsing edge; and that half has contact-corner sign $s$ in the initial
chamber and $-s$ in the other.

The last clause is \emph{true} and is proved below from the smoothing successor.
What an earlier revision got wrong was the \emph{reason}, not the conclusion: it
inferred the owner from which side of the remote line the incident edge lies on,
and the peer's chart refutes that inference --- with the remote edge
$(-\tfrac12,0)\to(\tfrac12,0)$, $M=(0,0)$, $p=(0,-1)$, $q=(1,1)$ and
$M_\mp=(0,\pm\tfrac1{10})$, the half with $\eta_1=s$ is $q$'s while $q$ lies on
the \emph{same} side as $M_-$. A later revision then over-corrected and printed
that the conclusion itself is false, which it is not. Ownership is read from the
smoothing successor. Smoothing the relocated crossing splits the traversal at its two
visits into two cyclic words, and the short edge --- the stretch of the incident
edge between the moving vertex and the contact --- lies in exactly one of them,
namely the one whose arc contains the moving vertex; that arc is the half whose
contact corner is the turn at the contact on that side. There are exactly two
cases, according to which incident edge carries the crossing in the chamber
under consideration, and clause~(B)(i) says the two chambers
exchange them. The caveat recorded there applies here too: opposite sides of the
remote \emph{line} do not place an intersection inside the remote
\emph{segment}, so the event interval is shrunk as clause~(B)(i) prescribes before
any of this is read off.
\end{enumerate}
\item[(C)] Let $S\in\Ind(G_{P_-})$ be arbitrary, and let $\phi S$ be its image under the
relabelling of Lemma~\ref{lem:slidingcartesian}(A). Write $A_-(S)$ for the product
of the coefficient factors $\Omega_1$ over the carriers of $S$ in the initial
chamber and $A_+(\phi S)$ for the same product over the carriers of $\phi S$ in
the other chamber.
\begin{enumerate}
\item[(i)] One always has $A_-(S)=A_+(\phi S)$. If $x_-\in S$, put
$S_i=(S\setminus\{x_-\})\cap V(G_{\lambda_i})$ and write $A_i(S_i)$ for the
product of the coefficient factors over the carriers of $S_i$ in $\lambda_i$.
Every argument is bound: an earlier revision wrote $A_i$ with none. Under that condition,
Lemmas~\ref{lem:slidingcartesian}(C)(i)--(iii) and~\ref{lem:slidingcoeff}(A)(ii) identify the two
chamber products carrier by carrier with the two half products, and
\[
   A_-(S)=A_+(\phi S)=A_1(S_1)A_2(S_2).
\]
\item[(ii)] If $x_-\in S$, then, under Convention~\ref{conv:s7standing}, put
$S_i=(S\setminus\{x_-\})\cap V(G_{\lambda_i})$, the two halves of
Lemma~\ref{lem:slidingcartesian}(B), and write
$W_-=\wind(S)$, $W_+=\wind(\phi S)$ and $W_i=\wind(S_i)$, each with its
argument. With $r,d_0,d_1$ the three directed wall tangents as above, put
\[
 \eta_1=\sgn\det(r,d_1),\qquad \eta_2=\sgn\det(d_0,r),\qquad
 \tau=\sgn\det(d_0,d_1).
\]
Here $\eta_2=-\eta_1$, and the collapsing half has contact-corner sign $s$ in
the initial chamber and $-s$ in the other, where $s$ is the sign~\eqref{eq:S6}
in the initial chamber. For $\eta\in\{\pm1\}$, write $h_\eta$ for the weight of
the distinguished carrier of the half with contact-corner sign $\eta$, and
$c_\eta$ for the weight of the corresponding chamber carrier carrying the short
edge, whose contact-corner sign there is $\eta$. Write $Q$ for the product of
the weights of all carriers not distinguished at the contact. It is the same
number in both chambers and both half terms: Lemma~\ref{lem:slidingcartesian}(C)(i)
matches those carriers bijectively, and clauses (iv) and (v) give the matched
corner-sign lists. An earlier revision wrote $W_-,W_+,W_1,W_2$ and $Q$ with no
arguments and no bijection. Then
\[
   W_1W_2=Q\,h_s h_{-s},\qquad
   W_-=Q\,c_s h_{-s},\qquad
   W_+=Q\,h_s c_{-s},\qquad
   c_\eta=\begin{cases}-\eta\,h_\eta,&\tau=\eta,\\ 0,&\tau=-\eta,\end{cases}
\]
and consequently $W_+-W_-=s\,W_1W_2$.
\end{enumerate}

\item[(D)] At every simple sliding vertex-on-edge event, in the sense of
Convention~\ref{conv:events}, whose two halves are generic polygons,
\[
   X_1(P_+)-X_1(P_-)=s\,X_1(\lambda_1)X_1(\lambda_2),
\]
with $s$ the sign \eqref{eq:S6} in the initial chamber and $\lambda_1,\lambda_2$
the halves of Definition~\ref{def:halves}.
\end{enumerate}
\end{lemma}
\begin{proof}
\emph{Clause (A).}
\emph{Source.} This proof is a transcription of the peer's pack \texttt{evidence/codex/phase15-sliding-universal-attack/REPORT.md}, SHA-256 \texttt{4f6da162\dots9318}, §``Carrier rotation'', lines 152--218.

\emph{(i), the avoiding sector.} Every carrier corner is either a polygon
vertex or the smoothing site of a persistent transverse crossing. Its point and
its incoming and outgoing directions vary continuously through the wall, and all
corner determinants remain nonzero by the guards: the vertex turn by (G1), the
smoothing-site turn by the active (G5) of the crossing, which persists. The
contact itself is only a transverse coincidence of parameter values, the
relocated crossing not being smoothed here. Thus each successor cycle has a fixed cyclic corner list. The nonzero
corner determinants make every incident edge nonzero and exclude consecutive
negative multiples, so it is a path in the corresponding $\mathcal R_c$.
Lemma~\ref{lem:turnlift}(ii) keeps its signed rotation constant; taking
magnitudes gives $R$. The same nonzero
determinants preserve the entire corner-sign list, which is also the selector
statement in this sector.

\emph{Three corners.} Every smoothed carrier has at least three corners under
the generic straight-edge guards: a carrier with two corners would consist of
two straight segments joining the same two points, which would overlap, and one
with a single corner would contain a zero-length straight return. This matters
because a zero selector then means mixed turns rather than a special one- or
two-corner convention.

\emph{(ii), the containing sector.} The correspondence with the halves is
Lemma~\ref{lem:slidingcartesian}(C)(i), and clause~(iv) of that lemma says the corner
lists agree away from the contact. There are \emph{two} carriers through the
contact, one in each interval, and clause (iv) gives the one not carrying the
short collapsing edge the same corner list as its mate as well; so every carrier
except the single one carrying that edge is a regular homotopy of its mate, by
the argument of (i). An earlier revision spoke of ``the one through the
contact'', which is the picture Lemma~\ref{lem:slidingcartesian}(C)(i) withdraws. The
distinguished carrier has one short edge that contracts
at the wall, and its local direction list is one of the two refinements
\begin{equation}
   (r,d_0,d_1)\longrightarrow(r,d_1),
   \qquad
   (d_0,d_1,r)\longrightarrow(d_0,r),
   \tag{U9}
\end{equation}
the two-direction lists on the right being the corresponding wall-half corners.
Let $\theta$ denote the principal oriented turn between two nonparallel
directions. Display~(U8) places both incident directions in one open angular
interval of length $\pi$ bounded by the remote direction, so no branch cut is
crossed and the two local angle additions
\begin{equation}
   \theta(r,d_0)+\theta(d_0,d_1)=\theta(r,d_1),
   \qquad
   \theta(d_0,d_1)+\theta(d_1,r)=\theta(d_0,r)
   \tag{U10}
\end{equation}
are exact, with no hidden full turn. All other turns agree by the polygon
path of (i). Lemma~\ref{lem:turnlift}(ii) therefore gives every chamber
carrier the same signed rotation and,
after taking absolute values, the same $R$ as its wall-half mate; clause~(v)
then makes the two chambers agree in both values carrier by carrier.

The pack also gives an equivalent form that avoids the angle bookkeeping, and it
is worth printing because it is the one a formalizer would prefer: choose one
reference ray avoiding the finite set of wall carrier directions and shrink the
chamber interval until that same ray avoids all corresponding directions, then
compute every rotation with that common ray. On either refinement in~(U9) a ray
in a backtracked subarc is crossed once with each sign and cancels, and every
other ray gets the direct-corner count.

\emph{Clause (B).} (i) Both incident tangents are transverse to $r$: a parallel one would put its far
endpoint on the remote line, an incidence outside the zero set
Convention~\ref{conv:events} allows, which contains no incidence but the
contact one. The neighbours lie on opposite sides of
the line, so the traversal $p\to M\to q$ crosses the line at $M$ and the two
incident directions point into the same open half-plane bounded by $r$, i.e.\
the two determinants have one sign. The statement is invariant: for each of the two incident edges, that edge meets
the remote edge exactly when its two endpoints lie strictly on opposite sides of
the remote line, and $p$, $q$ lie on opposite sides of it while $M$ passes from
one side to the other. So in each chamber exactly one incident edge meets the
remote edge, and the two chambers exchange which one it is --- no direction of
travel is named. No other crossing changes: any other change would vanish a
member of $\mathcal G$ outside the zero set, and by
Lemma~\ref{lem:guardconst} every relevant member outside $Z$ is nonzero of
constant sign on a possibly smaller interval, on which the whole argument is
read.

(ii) By clause~(B)(i) the two incident directions satisfy
$\sgn\det(r,d_0)=\sgn\det(r,d_1)$, and $\det(d_0,r)=-\det(r,d_0)$, so
$\eta_2=-\eta_1$: the two wall-half contact turns are opposite.

For the assignment the reason is the \emph{cross-pairing} of the oriented
smoothing, and not a side of a line. An earlier revision of this proof concluded
by the far-side inference that the statement above already withdraws --- the
statement was corrected and the proof under it was not, which is a way of
leaving a refuted argument in print.

At a slide exactly one incident edge meets the remote segment
(clause~(B)(i)), so in the chamber under consideration the
relocated crossing sits on one of them; write $X$ for the crossing point. Take
first the incoming edge $pM$. The traversal visits $X$ twice, once running
$p\to M$ and once running along the remote edge $a\to b$. The oriented smoothing
joins the branch arriving along $pM$ to the branch leaving along $X\to b$, and
the branch arriving along $a\to X$ to the branch leaving along $X\to M$. So the
short stretch $X\to M$ continues into $M\to q$: it lies on the successor cycle
carrying $q$'s arc, which is the half $\lambda_1$, whose contact corner is
$\eta_1$.

Take instead the outgoing edge $Mq$, which is the case in the other chamber by
clause~(B)(i). Now the two visits of the crossing $Y$ run $M\to q$
and $a\to b$, the smoothing joins the branch arriving along $M\to Y$ to
$Y\to b$ and the branch arriving along $a\to Y$ to $Y\to q$, and the short
stretch $M\to Y$ is entered from $p$: it lies on the cycle carrying $p$'s arc,
the half $\lambda_2$, whose contact corner is $\eta_2=-\eta_1$.

Both orientations occur and both are printed, an earlier revision of this proof
having treated only the first. In the initial chamber the crossing sits on the
incoming edge exactly when the moving vertex has passed to the side where $pM$
meets the remote segment; then the short stretch lands on $\lambda_1$, whose
contact corner is $\eta_1$, and $\eta_1=s$ there. In the reversed configuration
--- the peer's chart, with $M_-=(0,-\tfrac1{10})$ and $s=-1$, where $M_-q$
crosses the remote edge at $(\tfrac1{11},0)$ --- the crossing sits on the
outgoing edge instead, the short stretch lands on $\lambda_2$, and $\eta_2=s$.
The two cases are exchanged by the chamber, by clause~(B)(i), so
in either configuration the collapsing half has contact-corner sign $s$ in the
initial chamber and $-s$ in the other, which is the assignment the consumers
read.

What carries the argument in both cases is which branch the smoothing attaches
the short stretch to. That is a statement about the successor, and not about
which side of the remote line anything lies on.

\emph{Clause (C).} (i) Lemma~\ref{lem:slidingcartesian}(A) preserves the grouped polynomials and the
grouped writhes carrier by carrier, and Lemma~\ref{lem:slidingcoeff}(A) preserves
the absolute rotations, so each factor
$\Omega_1=[a^{1-w-R}z^0]P$ is selected at the same index from the same
polynomial: the first identity follows factor by factor. When $x_-\in S$, the
Lemmas~\ref{lem:slidingcartesian}(C)(i)--(iii) and~\ref{lem:slidingcoeff}(A)(ii) identify
the carriers with the disjoint union of the two sets of half carriers and
preserve each owner, grouped polynomial, grouped writhe and $R$, so the product
over the mapped carriers splits into the two half products. No coefficient
extraction is distributed over a polynomial product: each carrier lies in one
half.

(ii) By clause~(B)(ii) the collapsing half is the one of
contact-corner sign $s$ in the initial chamber and the one of sign $-s$ in the
other; that is what makes the three displayed products the $\pm$ pair they are.
In the chamber in which a given half carries the collapsing short edge, its
single wall corner of sign $\eta$ is refined into two corners, of signs $\eta$
and $\tau$. If the wall carrier is mixed, both $h_\eta$ and $c_\eta$ are $0$,
since adding a corner cannot remove an existing opposite turn. If it turns
right only and $\tau$ agrees, the refined carrier still turns right only and
its weight is $1=-\eta h_\eta$ with $\eta=-1$ and $h_\eta=1$. If it turns left
only and $\tau$ agrees, the refined carrier has one more corner, so its weight
is negated: $c_\eta=-h_\eta=-\eta h_\eta$ with $\eta=+1$. If $\tau$ disagrees
the refined carrier is mixed and $c_\eta=0$. Those are all cases, proving the
displayed rule. The three product formulas are the definition of $\wind$ with
the distinguished factor named. The single $Q$ is legitimate because
Lemma~\ref{lem:slidingcartesian}(C)(i), (iv) and (v) identifies the other carriers and
their corner-sign lists in both chambers and both half terms. Finally, if
$\tau=s$ then $c_s=-s\,h_s$ and $c_{-s}=0$, so
$W_+-W_-=Q(0+s\,h_sh_{-s})=s\,W_1W_2$; if $\tau=-s$ then $c_s=0$ and
$c_{-s}=s\,h_{-s}$, so $W_+-W_-=Q(s\,h_sh_{-s}-0)=s\,W_1W_2$.

\emph{Clause (D).} Split the supports of each chamber into those avoiding the relocated crossing
and those containing it. On the avoiding sector the coefficient products agree by
clause~(C)(i), and the selectors agree by
Lemma~\ref{lem:slidingcoeff}(A)(i), whose corner-sign list is exactly the selector
datum: no carrier acquires or loses a corner there. So that sector cancels term
by term. Clause~(C)(ii) is \emph{not} invoked here; it is
stated for a support containing the relocated crossing and an earlier revision
of this proof cited it in the avoiding sector, outside its own quantifier. On the containing sector, write the difference as the sum over the supports
themselves before any reindexing:
\[
   \sum_{\substack{S\in\Ind(G_{P_-})\\ x_-\in S}}
   \bigl(W_+(\phi S)-W_-(S)\bigr)\,A_-(S),
\]
which is legitimate because $\phi$ is a bijection of the containing supports of
the two chambers (Lemma~\ref{lem:slidingcartesian}(A)) and $A_-(S)=A_+(\phi S)$
(clause~(C)(i)), so the two chambers' containing sectors are
summed over one index set. An earlier revision passed straight from the jump to
the Cartesian sum, omitting this line, which is the first equality of the
source's final display. Now
clause~(C)(i) makes the non-selector part of the term the product
$A_1(S_1)A_2(S_2)$ of the two half terms' non-selector parts. The reindexing of
the sum by $(S_1,S_2)\in\Ind(G_{\lambda_1})\times\Ind(G_{\lambda_2})$ is
Lemma~\ref{lem:slidingcartesian}(B); the carrier, owner and piece-data
identification at a fixed $S$ is Lemma~\ref{lem:slidingcartesian}(C)(i)--(iii), and
Lemma~\ref{lem:slidingcoeff}(A)(ii) supplies its rotation data. Finally
clause~(C)(ii) makes the selector difference $s\,W_1W_2$.
Hence
\[
X_1(P_+)-X_1(P_-)
=\sum_{S_1,S_2}s\,W_1(S_1)A_1(S_1)\,W_2(S_2)A_2(S_2)
=s\,X_1(\lambda_1)X_1(\lambda_2),
\]
where $A_i$ denotes the coefficient product of the half term.
\end{proof}

\begin{remark}[printed from the source, and what it replaces]\label{rem:slidingsource}
This is the rotation argument of the counterpart's sliding pack, printed in its
own terms: regular homotopy for the avoiding sector, and exact angle addition
under the opposite-side condition for the contracting edge. An earlier revision
of this document argued it instead by an axiom about embedded arcs in a disc,
which the peer broke --- for the empty support the carrier meets a contact disc
in two arcs, not one, and they cross. That axiom is not used here, and the
displayed additions are what make the absence of a hidden full turn a
computation rather than an assurance.
\end{remark}

\subsection{Universal (S7)}\label{sec:s7universal}


\begin{remark}[the two sectors of the noninterlacing case]\label{rem:nistatus}
Four structural points of Section~\ref{sec:nijoin} are recorded here so that
no reader has to reconstruct them. The orientation factor:
Lemma~\ref{lem:nitransport}, clause~(A) names $s$
and says that the $s=-1$ configuration is the reflection of the $s=+1$ one,
with every determinant negating together, instead of printing $+1$ globally and
consuming $s$ later. The $U_T=0$ branch: the assembly multiplies by $U_T$
whatever its value, so a vanishing spectator product vanishes the term. The
same-piece hypothesis of the transport lemma: it is the source's own and it is
a hypothesis of that lemma. The
\emph{different-piece sector}: its common-neighbourhood input is
Lemma~\ref{lem:contactword}(ii); its discharge is printed in
Lemma~\ref{lem:chambertransport}, clauses~(A) and~(B),
Theorem~\ref{thm:s7universal}(D)(i), Theorem~\ref{thm:s7universal}(D)(iii) and
Theorem~\ref{thm:s7universal}(D)(v); Theorem~\ref{thm:s7universal}, clause~(E) does not exclude it.

\emph{The two sectors, named.} For an eligible $T$ at a bigon event, the
undominated set $U(T)$ carries the induced interlacement graph
$G_{P_2}[U(T)]$, and the two newborns $x,y$ lie in it. The \emph{same-piece}
sector is the set of eligible $T$ for which $x$ and $y$ lie in one connected
component of that graph, and the \emph{different-piece} sector is the set for
which they lie in two; the names are the residual pieces those components
become, and every eligible $T$ is in exactly one sector. An earlier revision of
this remark used the two names without defining them.

\emph{The different-piece discharge.} All four of its branches are proved outright, and
the branch statements are read at $\epsilon=0$ --- the noninterlacing branch,
which is where the different-piece sector arises, the interlacement bit of
Definition~\ref{def:sectors} being $0$ exactly when the two newborn chords do
not interlace. With $\epsilon=0$ fixed: the two newborns are singleton positive
kinks (Lemma~\ref{lem:chambertransport}, clause~(A), whose hypothesis that is); the row
$f_\nu=[z^0]F_\nu$ of the contact carrier in chamber $\nu$, and its prefactor,
transport between the chambers with no selector hypothesis
(Lemma~\ref{lem:chambercommon} and Lemma~\ref{lem:chambertransport}, clause~(B)); the
one-newborn terms vanish by cutting the affected carrier at its surviving
singleton (Theorem~\ref{thm:s7universal}(D)(i) and Theorem~\ref{thm:s7universal}(D)(iii)); and the
high-neither and low terms read the degrees $D-4$ and $D-2$, in the writhes
$w_{\mathrm L},w_{\mathrm H}$ of Definition~\ref{def:markeddata}, which the
floor on the two halves kills --- the same two degrees, and the same floor, as
Theorem~\ref{thm:s7universal}, clause~(C) uses in the same-piece sector
(Theorem~\ref{thm:s7universal}(D)(v)).
\end{remark}

\subsection{The returned aggregate in bracket form}\label{sec:bracket}

This subsection derives, for the interlacing branch, the bracket-form
identification of the returned aggregate. It closes two of its three terms;
what remains after it is one displayed identity, stated at the end.

\begin{lemma}[rotation across the contact, both branches]\label{lem:rotshift}
Here eligibility is as in Definition~\ref{def:eligible} and $\rot$ is as in
Definition~\ref{def:rot}.
Let $L^{*}$ be the contact carrier of an eligible $T$, let $L_1,L_2$ be its two
halves, and let
\[
   s_{\mathrm{low}}=\sgn\det\bigl(p_b-p_a,\;p_M-p_a\bigr)\ \text{ in the
   \emph{low} chamber }P_0
\]
be the side of the remote line on which the moving vertex sits before the
crossing is born. Then
\[
   \rot(L_1)+\rot(L_2)=\rot(L^{*})+\begin{cases}0,&\epsilon=1,\\
   s_{\mathrm{low}},&\epsilon=0.\end{cases}
\]
The sign here is $s_{\mathrm{low}}$ and not the sign $s$ of \eqref{eq:S6},
which is read in whichever chamber the parametrisation calls $P_-$. The two are
related by the orientation of the path: $s=\delta\,s_{\mathrm{low}}$, where
$\delta=+1$ if the parametrisation calls the low chamber $P_-$ and $\delta=-1$
if it calls the high chamber $P_-$ --- a definition, not a claim, since $s$ and
$s_{\mathrm{low}}$ are the same determinant read in the two chambers and that
determinant changes sign across the wall. The statement above is invariant
under
reparametrisation because $s_{\mathrm{low}}$ is defined by the chamber and not
by the direction of travel.
\end{lemma}
\begin{proof}
The three carriers have the same non-contact corners with the same turns
(Lemma~\ref{lem:contactsplit}), so by Lemma~\ref{lem:turnlift}(ii)
\[
   2\pi\bigl(\rot(L_1)+\rot(L_2)-\rot(L^{*})\bigr)=\eta_1+\eta_2-\tau,
\]
where $\eta_1,\eta_2$ are the two half contact turns and $\tau$ the full one.
The chart of Lemma~\ref{lem:contactturnboth} is fixed by placing the moving
vertex on a chosen side of the remote line, which is the choice
$s_{\mathrm{low}}=+1$; in that chart $\eta_1=\arg q$ and $\eta_2=\pi-\arg p$, while $\tau$ is the principal
value of $\arg q-\arg p+\pi$. At $\epsilon=1$ that value lies in $(0,\pi)$ and
equals $\arg q-\arg p+\pi$, so $\eta_1+\eta_2-\tau=0$. At $\epsilon=0$ it lies
in $(-\pi,0)$ and equals $\arg q-\arg p-\pi$, so $\eta_1+\eta_2-\tau=2\pi$ and
the shift is $+1$.

The configuration with $s_{\mathrm{low}}=-1$ is the image of the one with
$s_{\mathrm{low}}=+1$ under a reflection of the plane, which reverses every signed rotation and every turning
angle. So at $\epsilon=1$ the shift $0$ is unchanged, while at $\epsilon=0$ the
shift $+1$ becomes $-1$. Writing the two cases as the single symbol
$s_{\mathrm{low}}$ gives the display.
\end{proof}

\begin{remark}[the shift carries the side sign]\label{rem:rotshiftsign}
An earlier revision printed the constant $+1$ at $\epsilon=0$. That is the
$s=+1$ half of the statement, and the peer refuted the constant form on the
exact guarded $24$-gon of Remark~\ref{rem:diffcensus}, whose low-chamber side
sign is $s_{\mathrm{low}}=-1$
and whose signed rotations are $\rot(L_1)=\rot(L_2)=1$ and $\rot(L^{*})=3$, so
that $\rot(L_1)+\rot(L_2)-\rot(L^{*})=-1$. A second revision bound the shift to the
\eqref{eq:S6} sign $s$, which is read in whichever chamber is called $P_-$ and
therefore flips with the parametrisation while the geometry does not; the
binder is now the chamber-defined
$s_{\mathrm{low}}$. Lemma~\ref{lem:nitransport}, clause~(B) consumes the corrected form; no
conclusion drawn from it changes.
\end{remark}

\begin{remark}[this is the missing full turn]\label{rem:missingturn}
Lemma~\ref{lem:rotshift} is the reason naive addition of principal angles
across the contact is wrong in the noninterlacing branch, which the phase-1.5
record warns of without displaying the discrepancy: there the two halves carry
one full turn more than the carrier they came from, and at $\epsilon=1$ they
carry exactly as much.
\end{remark}

\begin{remark}[a withdrawn route]\label{rem:withdrawnbracket}
An earlier revision of this document printed here a lemma asserting that after
polarization the contact carrier is, in the low chamber, the connected sum of
the two half groupings, with framed row $\Gamma_L=\Gamma_1\Gamma_2$; a
proposition deriving the high-neither minus low difference from it by the
framed self-skein; and a corollary reducing the interlacing branch to an
aggregate over selector-dead supports. \textbf{All three are withdrawn.} The first
is false, and its falsity is visible inside this document: the writhe ledger
$w(D_-)=w_1+w_2+2\ell-1$ of Lemma~\ref{lem:triplebridge}(v) at $\epsilon=1$,
for a diagram whose
two parts are joined by $\ell$ between-interval crossings, so
$\Gamma_L=\Gamma_1\Gamma_2$ would force $2\ell=1$. The error was to think that
polarizing the between-interval crossings into a descending order lets
Reidemeister~II separate the two parts. It does not: in a divide every crossing
is positive, so between-interval crossings all have one sign and form a twist
region, which no Reidemeister~II move removes. The correct identity carries the
framing factor, $\Gamma_L=a^{2\ell-1}\Gamma_1\Gamma_2$, and the surviving
source factors the two parts only \emph{after} a contact crossing is smoothed,
retaining the single carrier's framing datum before that.

What this costs is the whole of my route to the interlacing endpoint. What
remains from it is stated where it belongs: Lemma~\ref{lem:rotshift}, the
rotation shift, which is independent of the framing question; and
Theorem~\ref{thm:s7universal}, clause~(B)(i) below, which is a determinant argument in
the contact chart and uses none of the withdrawn material.
\end{remark}

\begin{remark}[two lemmas retired, and one debt discharged]
\label{rem:excretired}
This subsection printed, in earlier revisions, two lemmas about the one-newborn
branches at $\epsilon=1$ --- that at most one of them is live, and that a
uniform contact carrier kills both --- and a remark calling the surviving
branch's coefficient read ``the remaining debt of the whole bigon branch''. All
three are withdrawn, and none is replaced by a weaker claim.

The peer refuted the two lemmas as printed: both treated liveness as one sign on
every raw corner of $A$ and one sign on every raw corner of $B$, whereas after a
smoothing the corners belong to \emph{carriers} and uniformity is carrierwise.
The statement that makes them unnecessary is the peer's.
Theorem~\ref{thm:s7universal}, clause~(B)(i) kills \emph{both} one-newborn branches for
\emph{every} eligible support, uniform contact carrier or not, by an exact
determinant argument in the printed chart. With both branches always dead,
``at most one is live'' is vacuous, and the debt is discharged rather than
outstanding: there is no surviving branch whose coefficient read anyone must
compute.

The consumers are rewired to Theorem~\ref{thm:s7universal}, clause~(B)(i).
\end{remark}

\subsection{The full-twist calculation}\label{sec:fulltwist}

This subsection prints the interlacing join as one live-support calculation.
The statement below carries its own $\epsilon=1$ and eligible-support binder and
imports all chamber and half data from Definition~\ref{def:markeddata}. Its
proof chooses $q=x$ and binds $D_{\mathrm A},F_{\mathrm A}$ before invoking
Lemma~\ref{lem:triplebridge}(ii); it binds $\ell$ only after that clause supplies
the ordered two components. No rotation, target or $\Omega$ is attached to the
two-component auxiliary link. Lemma~\ref{lem:fulltwist} remains branch-neutral.

\begin{lemma}[the abstract full-twist triple]\label{lem:fulltwist}
Let $D_{\mathrm L},D_{\mathrm H},D_{\mathrm A}$ be oriented diagrams and let $q$
be a positive crossing of $D_{\mathrm H}$ such that
\begin{enumerate}
\item[(T1)] the oriented smoothing of $D_{\mathrm H}$ at $q$ is
$D_{\mathrm A}$;
\item[(T2)] switching $q$ in $D_{\mathrm H}$ gives a diagram carried to
$D_{\mathrm L}$ by oriented Reidemeister-II moves.
\end{enumerate}
The symbols this lemma reads off a diagram are bound here, from the diagrams it
quantifies over, and are not imported from any definition attached to a polygon
or a chamber. They are bound in two groups, because they have two domains. For
\emph{any} oriented diagram $D$ --- any number of components --- put
\[
   F_D=P_D ,
\]
the HOMFLY--PT polynomial of the link it presents (Axiom~\ref{ax:homfly}). For
a diagram $D$ carried by a \emph{single} closed plane curve $\Gamma_D$, and
only for such a diagram, put
\[
   d(D)=1-w(D)-R(\Gamma_D),\qquad
   \Omega(D)=\bigl[a^{d(D)}z^0\bigr]F_D ,
\]
with $w(D)$ the writhe and $R(\Gamma_D)=|\rot(\Gamma_D)|$ the absolute
rotation of that curve. An earlier revision bound all three on one domain,
which typed $d$ and $\Omega$ at a two-component diagram, where the rotation
number of Definition~\ref{def:rot} is not defined. Abbreviate $F_\nu=F_{D_\nu}$ for $\nu\in\{\mathrm L,\mathrm H,\mathrm A\}$,
and $d_\nu=d(D_\nu)$, $\Omega_\nu=\Omega(D_\nu)$ for
$\nu\in\{\mathrm L,\mathrm H\}$ \emph{only}. The restriction is not tidiness:
$D_{\mathrm A}$ is the oriented smoothing of $D_{\mathrm H}$ and is carried by a
\emph{two-component} link in the application, while $R(\Gamma_D)$ is the
absolute rotation of one closed curve (Definition~\ref{def:rot}), so $d$ and
$\Omega$ are not defined at $\mathrm A$ at all. Nothing below asks for them:
$D_{\mathrm A}$ enters only through $F_{\mathrm A}$ and through the coefficient
$[a^{d_{\mathrm L}-1}z^{-1}]F_{\mathrm A}$, which is a coefficient of that
polynomial and not a read at a slot of its own. An
earlier revision used $d_\nu$ and $\Omega_\nu$ in this statement while binding
them only in Definition~\ref{def:markeddata}, which is about the contact
carriers of an eligible support; the abstract triple has no support, and the
lemma was not readable on its own hypotheses. The two bindings agree wherever
both apply, that definition computing $d_\nu$ and $\Omega_\nu$ from these same
grouped diagrams. Then
\[
   F_{\mathrm H}=a^{-2}F_{\mathrm L}+a^{-1}zF_{\mathrm A}.
\]
If moreover $d_{\mathrm H}=d_{\mathrm L}-2$, then
\[
   \Omega_{\mathrm H}-\Omega_{\mathrm L}
   =\bigl[a^{d_{\mathrm L}-1}z^{-1}\bigr]F_{\mathrm A}.
\]
\end{lemma}
\begin{proof}
Apply the skein of Axiom~\ref{ax:homfly} at $q$: its $L_+$ is $D_{\mathrm H}$,
its $L_-$ is the switched diagram, whose polynomial is $F_{\mathrm L}$ by (T2)
and the Reidemeister-II invariance of the polynomial, and its $L_0$ is
$D_{\mathrm A}$ by (T1). So $aF_{\mathrm H}-a^{-1}F_{\mathrm L}=zF_{\mathrm A}$,
which is the first display.

For the second, $\Omega_{\mathrm H}=[a^{d_{\mathrm H}}z^0]F_{\mathrm H}
=[a^{d_{\mathrm L}-2}z^0]F_{\mathrm H}$. Taking $[a^{d_{\mathrm L}-2}z^0]$ of
the first display, the first term contributes
$[a^{d_{\mathrm L}-2}z^0](a^{-2}F_{\mathrm L})
=[a^{d_{\mathrm L}}z^0]F_{\mathrm L}=\Omega_{\mathrm L}$ and the second
contributes $[a^{d_{\mathrm L}-1}z^{-1}]F_{\mathrm A}$.
\end{proof}

\begin{remark}[why this lemma is branch-neutral]\label{rem:branchneutral}
Nothing in the statement mentions $\epsilon$, a contact, or a polygon: it is an
identity about three oriented diagrams related by one crossing. That is
deliberate. An earlier revision proved the $\epsilon=1$ case and then had the
$\epsilon=0$ section import it ``verbatim'', which is not an argument --- the
two branches have different contact words and different auxiliary components.
Here each branch proves (T1) and (T2) for itself, in
Lemma~\ref{lem:triplebridge}, and the algebra is quoted once. The peer's
blueprint is the source of this arrangement.
\end{remark}

\begin{lemma}[the residual sets split, and who owns them]\label{lem:ownermap}
Here eligibility is as in Definition~\ref{def:eligible}.
Let $T$ be eligible with parts $T_A,T_B$, let $A$ and $B$ be the two intervals
of the contact word of Lemma~\ref{lem:contactword}, let $V_A,V_B$ be the old
labels with both visits in $A$ and in $B$ respectively, and let $\mathcal N$ be
the old labels with one visit in each. Write $U(T)$ for the undominated set of $T$ in the high-neither
chamber and $U_{\lambda_i}(T_i)$ for the undominated set of $T_i$ in the half
$\lambda_i$. Then
\[
   U(T)\cap V_A=U_{\lambda_1}(T_A),\qquad
   U(T)\cap V_B=U_{\lambda_2}(T_B),\qquad
   U(T)\cap\mathcal N=\mathcal N\setminus N_{G_{P_2}}(T),
\]
and every label of the third set has one visit in each interval.

\end{lemma}
\begin{proof}
A chord with both visits in one interval and a chord with both visits in the
other cannot interlace, their visits lying in disjoint arcs
(Lemma~\ref{lem:orderedproduct}); and by Lemma~\ref{lem:contactword}(i) every
old label lies in $V_A$, in $V_B$ or in $\mathcal N$. Let $c\in V_A$. Then $c$
interlaces no element of $T_B$, so $c\in U(T)$ if and only if $c\notin T_A$ and
$c$ interlaces no element of $T_A$; interlacement of two chords with all four
visits inside $A$ is read from their order inside $A$, which the half inherits,
so that is exactly $c\in U_{\lambda_1}(T_A)$. For $c\in V_B$ interchange the
two intervals throughout the previous three sentences: $c$ interlaces no element
of $T_A$, so $c\in U(T)$ if and only if $c\notin T_B$ and $c$ interlaces no
element of $T_B$, which is $c\in U_{\lambda_2}(T_B)$. For $c\in\mathcal N$ the condition $c\in U(T)$ is $c\notin T$ and
$c\notin N_{G_{P_2}}(T)$; the first holds for every eligible $T$ by
Definition~\ref{def:eligible}, so it reduces to $c\notin N_{G_{P_2}}(T)$, which
is the third identity. The subscript is not decoration: $N(T)$ without one is
bound in Definition~\ref{def:branches} as the \emph{scalar} high-neither term of
$T$, and an earlier revision wrote the graph neighbourhood with the same
symbol. The last clause is the definition of $\mathcal N$.

The third identity is the one that must be read carefully, and an earlier
revision did not: a mixed label that is \emph{not selected} may still be
\emph{dominated}, and a dominated label is erased from the diagram rather than
retained. Membership of $U(T)$, not failure of membership of $T$, is what
retention means (Definition~\ref{def:pieces}).
\end{proof}

\begin{remark}[the newborn-free statement against the both-newborn one]\label{rem:ownervssliding}
Lemma~\ref{lem:ownermap} concerns the newborn-free support, not
Proposition~\ref{prop:sectortarget}(B), which concerns $S=T\cup\{x,y\}$ at a noninterlacing bigon. There the support
contains both newborns, every label of $\mathcal N$ is adjacent to them and is
dominated outright, and the residual sets carry no mixed label at all. Here the
support contains no newborn, so a mixed label survives unless an old selection
dominates it, and which ones survive is the content of the third identity.
\end{remark}

\begin{lemma}[different-piece singleton kinks and chamber transport of an eligible row]\label{lem:chambertransport}
\begin{enumerate}
\item[(A)] Here eligibility is as in Definition~\ref{def:eligible}.
Let $\epsilon=0$ and let $T$ be eligible. If the two newborns lie in different
residual pieces of $T$, then those pieces are the singletons $\{x\}$ and
$\{y\}$; each singleton's piece diagram is an unknot, so $P_{\{x\}}=P_{\{y\}}=1$
and $w(\{x\})=w(\{y\})=1$.
\item[(B)] Here eligibility is as in Definition~\ref{def:eligible}.
Under Convention~\ref{conv:s7standing}, let $T$ be eligible at a simple bigon
event, let $C_{\mathrm L}$ and $C_{\mathrm H}$ be its contact carriers in the
low and the high chamber, let $W_T$ be its chamber selector and $U_T$ the
product of the reads over its untouched carriers --- all four as bound in
Definition~\ref{def:markeddata}, which is where they are defined, rather than
imported from the lemma that first computed with them. Then, with no hypothesis on
any selector:
\begin{enumerate}
\item[(iii)] if $T$ is different-piece, the residual pieces of $C_{\mathrm H}$ are those
of $C_{\mathrm L}$ together with the two singletons $\{x\}$ and $\{y\}$, so the grouped
rows agree, $f_{C_{\mathrm H}}=f_{C_{\mathrm L}}$, while the writhes differ by two,
$w_{T,C_{\mathrm H}}=w_{T,C_{\mathrm L}}+2$.
\end{enumerate}
Consequently --- and this conclusion is stated only under the
different-piece hypothesis of subclause~(iii), since its rows are derived from it ---
writing $U_T$ for the product of the reads on the carriers of
Lemma~\ref{lem:chambercommon}(i) --- clause (B) has no subclause (i) since the
split, and an earlier revision still pointed at one --- and
$d_0=1-w_{T,C_{\mathrm L}}-R(C_{\mathrm L})$,
\[
   L(T)=W_TU_T\bigl[a^{d_0}\bigr]f_{C_{\mathrm L}},
   \qquad
   N(T)=W_TU_T\bigl[a^{d_0-2}\bigr]f_{C_{\mathrm L}},
\]
both expressions being valid whether or not $W_TU_T$ vanishes.
\end{enumerate}
\end{lemma}
\begin{proof}
\emph{Clause (A).} At $\epsilon=0$ the newborns do not interlace, so they are not adjacent. Suppose
$x$ had an undominated old neighbour $c$. By Lemma~\ref{lem:contactword}(ii) $c$ is a
neighbour of $y$ as well, so $x,c,y$ is a path in the residual graph
$G_{P_2}[U(T)]$ and the two newborns share a piece, contrary to hypothesis.
Hence $x$ has no residual neighbour at all and $\{x\}$ is a component of the
residual graph. Every hypothesis used --- eligibility of $T$, $\epsilon=0$, and
Lemma~\ref{lem:contactword}(ii), whose conclusion is an equivalence between the two
newborns --- is invariant under interchanging the names $x$ and $y$, and so is
the conclusion that the two lie in different pieces. Interchanging them
therefore yields the identical derivation with $y$ in place of $x$, and
$\{y\}$ is a component as well.

For the diagram: by Definition~\ref{def:piecediagram} the diagram of the piece
$\{x\}$ is the parent curve with \emph{every} double point other than $x$
erased and $x$ resolved by the divide convention. A closed curve with exactly
one crossing is an unknot diagram, so its normalized polynomial is $1$
(Definition~\ref{def:homfly}), and its writhe is $w(\{x\})=|\{x\}|=1$ because
the divide convention makes every crossing positive.

\emph{Clause (B).} (iii) By clause~(A) the two newborns are singleton
residual pieces of $T$ in the high chamber; every other residual piece consists
of old crossings, and its diagram erases $x$ and $y$, so the two strands there
pass without interaction and the finger is undone by a planar isotopy, leaving
the same diagram as in the low chamber. Each singleton contributes the factor
$1$ to the grouped polynomial and the summand $1$ to the grouped writhe
(clause~(A)), which is the displayed statement.

The two displays are then Definition~\ref{def:X1} read in each chamber, the
slot in the high chamber being
$1-w_{T,C_{\mathrm H}}-R(C_{\mathrm H})=1-(w_{T,C_{\mathrm L}}+2)-R(C_{\mathrm L})=d_0-2$. Nothing above divides by a
selector or a spectator read, so both displays hold as identities of integers
whatever those factors are.
\end{proof}

\begin{lemma}[the vertex-on-edge triple, endpoint bracket, auxiliary components and live full-twist jump]
\label{lem:triplebridge}
Here eligibility is as in Definition~\ref{def:eligible}.
Let $T$ be eligible at a simple bigon event with interlacement bit $\epsilon\in\{0,1\}$.
Clauses (i)--(v) also require $\epsilon=1$, or $\epsilon=0$ with $T$ same-piece. This statement
carries its own data and inherits none: let $D_{\mathrm L}$ and $D_{\mathrm H}$
be the grouped diagrams of the contact carrier in the two chambers
(Corollary~\ref{cor:groupedknot}(B)); let $q:=x$ be the newborn crossing whose
first visit opens the contact word of Lemma~\ref{lem:contactword}; let
$D_{\mathrm A}$ be the oriented smoothing of $D_{\mathrm H}$ at $q$; and let
$F_\nu=P_{D_\nu}$ for $\nu\in\{\mathrm L,\mathrm H,\mathrm A\}$. The symbol
$\ell$ is \emph{not} bound here: its binding needs the two components that
clause~(ii) below establishes, and a conclusion cannot retroactively type a
prior let --- an earlier revision bound "the linking number of the two
components, when it has two", a partial let discharged only by its own lemma.
An earlier revision also took $F_{\mathrm A}$ and $\ell$ from a preamble that
fixes $\epsilon=1$, while the statement covers both branches. The choice of newborn is a \emph{binding} and not a
convenience: the assignment of writhes to components in (iii) is not equivariant
under exchanging the two newborns, and an earlier revision quantified $q$ over
both while proving the statement for this one. Then
\begin{enumerate}
\item[(i)] $q$ is a positive crossing of $D_{\mathrm H}$; the oriented smoothing
of $D_{\mathrm H}$ at $q$ is $D_{\mathrm A}$; and switching $q$ in
$D_{\mathrm H}$ gives a diagram carried to $D_{\mathrm L}$ by oriented
Reidemeister-II moves. These are the three hypotheses that
Lemma~\ref{lem:fulltwist} asks of an abstract triple, restated here rather than
named, so that this clause can be read without opening that lemma;
\item[(ii)] $D_{\mathrm A}$ has exactly two components, ordered by the
half-label convention --- the first is the one containing the first half's
labels, the second the other; their grouped polynomials are $Q_1$ and $Q_2$,
the two half polynomials of Definition~\ref{def:markeddata}.

\emph{With the two components of (ii) so ordered, let $\ell$ be their linking
number; clauses (iii)--(v) read it.}

\item[(iii)] the second component has writhe $w_2$, and the first has writhe
$w_1+(1-\epsilon)$;
\item[(iv)] $[z^{-1}]F_{\mathrm A}=a^{-2\ell}(a-a^{-1})f_1f_2$;
\item[(v)] $w_{\mathrm L}=w_1+w_2+2\ell-\epsilon$.
\item[(vi)] without the extra hypothesis imposed on clauses (i)--(v),
\[
   w_{\mathrm H}=w_{\mathrm L}+2,\qquad R_{\mathrm H}=R_{\mathrm L},\qquad
   d_{\mathrm H}=d_{\mathrm L}-2.
\]

\item[(vii)] With the marked-data symbols as in Definition~\ref{def:markeddata},
if $\epsilon=1$ and $W_T\neq0$, then
\[
   \bigl[a^{D-2}\bigr]f_1f_2-\bigl[a^{D}\bigr]f_1f_2+\omega_1\omega_2=0 .
\]
A summand with $W_T=0$ contributes zero to the weighted aggregate whatever its
bracket, so the weighted SR.12 aggregate vanishes.

\item[(viii)] Under Convention~\ref{conv:s7standing}, if $\epsilon=1$, bind the parts $T_A,T_B$ and the following chamber and
half data exactly as in Definition~\ref{def:markeddata}:
$C_\nu,D_\nu,F_\nu,w_\nu,\rot_\nu,R_\nu,d_\nu,\Omega_\nu$ for
$\nu\in\{\mathrm L,\mathrm H\}$; $L_i,Q_i,f_i,w_i,\rot_i,R_i,d_i,\omega_i$
for $i\in\{1,2\}$; and $D,W_T$. If $W_T\neq0$, then
\[
   \Omega_{\mathrm H}-\Omega_{\mathrm L}=-\omega_1\omega_2.
\]
\end{enumerate}
\end{lemma}
\begin{proof}
\emph{(i).} (T1) holds by the definition of $D_{\mathrm A}$. For (T2), switch
$q$: the two newborns then carry opposite signs and, the contact disc being
empty (Convention~\ref{conv:events}), they bound an empty bigon face, so an
oriented Reidemeister-II move deletes both and leaves the low grouped diagram
$D_{\mathrm L}$ --- the old pieces and their crossings being untouched by the
move, which is supported in that disc.

\emph{(ii) and (iii).} Read the contact word of Lemma~\ref{lem:contactword}. At
$\epsilon=1$ it is $x_0y_0Ax_1y_1B$; smoothing $q=x$ cuts the traversal at
$x_0,x_1$ and reconnects, so the two visits of $y$ land on \emph{different}
components. Then $y$ is a \emph{mixed} crossing of the two-component link
$D_{\mathrm A}$, and the disposition of a mixed crossing is not that of a
self-crossing: it is not deleted, and Axiom~\ref{ax:homfly} is not what applies
to it. An earlier revision wrote that it was deleted with no change of
polynomial by that axiom, which types it as a self-crossing of one component.
What is true, and what is used below, is that a mixed crossing belongs to
neither component's own diagram, so it enters neither $f_1$ nor $f_2$; it
belongs to the linking data, and enters through the exponent $-2\ell$ that
Lemma~\ref{lem:homflyrows}(ii) consumes in~(iv). At $\epsilon=0$ the word is
$x_0y_0Ay_1x_1B$; smoothing $x$ puts both visits of $y$ on one component, where
it is a positive self Reidemeister-I curl --- there the crossing \emph{is} a
self-crossing, its deletion is a Reidemeister-I move, and the polynomial is
unchanged. In both branches the two components carry, respectively, the
labels of $A$ together with the half $\lambda_1$'s and those of $B$ together
with $\lambda_2$'s.

That the two components' grouped polynomials are $Q_1,Q_2$ is the record
argument and not a picture. The residual-owner identities are
Lemma~\ref{lem:ownermap}'s --- Lemma~\ref{lem:orderedproduct} is only a support
bijection and Lemma~\ref{lem:contactsplit} only a carrier split, and neither
supplies the owner map, which an earlier sentence here attributed to them.
The identities must be read at the level of \emph{labels}, not of pieces. It is \emph{false} that each
old residual piece of $T$ on the contact carrier lies wholly in $A$ or wholly in
$B$, and an earlier revision asserted it; the peer refuted it on the label set
\texttt{contact-0001}, where the piece
$\{(0,2),(0,3),(1,3),(2,7),(3,6)\}$ has one crossing among the first half's
labels, none among the second's, and four whose two visits fall in different
intervals. A piece is a connected component of an interlacement graph and
nothing makes it respect the cut.

What the record isomorphism needs, and what those two lemmas give, is a
statement about crossings. Call a crossing of $D_{\mathrm H}$ \emph{mixed} when
its two visits lie in different intervals of the contact word; by
Lemma~\ref{lem:contactword}(ii) the mixed old labels are exactly $\mathcal N$.
Of those, the ones \emph{retained} in the diagram are
$\mathcal N\setminus N_{G_{P_2}}(T)$
and not all of $\mathcal N$: a mixed label dominated by a selected old crossing
is erased --- dominated meaning adjacent in $G_{P_2}$ to a selected label ---
and Lemma~\ref{lem:ownermap} is what says which survive. An
earlier revision wrote that an eligible $T$ contains no mixed label and
concluded that every mixed crossing is retained, which confuses not-selected
with undominated.

The retained mixed crossings belong to neither component's own diagram of
$D_{\mathrm A}$, each having one visit in each interval: they are the crossings
whose signed count is $2\ell$. Delete them from the two component records and
every remaining crossing has both visits in one interval, hence lies on one
component; by Lemma~\ref{lem:ownermap} the remaining labels of the
component through $A$ are exactly $U_{\lambda_1}(T_A)$, the residual labels of
the first half, and those of the other component exactly
$U_{\lambda_2}(T_B)$. That is the owner map, and it is a statement about
residual \emph{sets}: Lemma~\ref{lem:orderedproduct} is a bijection of supports
and Lemma~\ref{lem:contactsplit} a splitting of carriers, and neither of them
says this. Restricting the traversal
of the smoothed diagram to the retained visits of one component gives a
bijection onto the visits of the corresponding half carrier, and by
Lemma~\ref{lem:carrierword} it preserves their cyclic order; the over/under
designation and the sign at each are those of the parent, by
Definition~\ref{def:piecediagram}, since a smoothing elsewhere alters neither.
So the retained-visit map is an orientation-preserving record isomorphism in the
sense of Definition~\ref{def:record}, and Axiom~\ref{ax:gausscode}, invoked
through Corollary~\ref{cor:groupedknot}(B) on each component, identifies the two
grouped diagrams with the half grouped knots. Their polynomials are $Q_1,Q_2$.
A half carrier bearing no piece is outside the $k\geq1$ scope of
Corollary~\ref{cor:groupedknot}(B) and takes one line of its own: the component has
no retained crossing at all, so its diagram is an embedded circle, its link type
is the unknot, and its grouped polynomial is $1$ with writhe $0$, which is
Theorem~\ref{thm:carrierfloor}(D) and is the value $Q_i$ takes there.

The writhes follow by counting positive crossings: the second component carries
exactly the half $\lambda_2$'s residual crossings, so $w_2$; the first carries
$\lambda_1$'s together with the surviving newborn when $\epsilon=0$, and none
when $\epsilon=1$, so $w_1+(1-\epsilon)$.

\emph{(iv).} $D_{\mathrm A}$ is a two-component link, and it is \emph{not}
claimed to be a split union --- its two components link, and the factor
$a^{-2\ell}$ records that. Lemma~\ref{lem:homflyrows}(ii) applied to it, with the two
component rows $f_1,f_2$ from (ii), gives
$[z^{-1}]F_{\mathrm A}=a^{-2\ell}(a-a^{-1})f_1f_2$.

\emph{(v).} Count writhes, branch by branch, and count them in
$D_{\mathrm H}$, whose crossings are the residual crossings of the two halves,
the mixed ones, and the two newborns $x$ and $y$; all are positive
(Definition~\ref{def:piecediagram}). That $w_{\mathrm L}=w_{\mathrm H}-2$ is
part~(i) read for writhes and not a forward citation: switching $q$ makes its
sign $-1$ and lowers the writhe by two, and the oriented Reidemeister-II move
that follows removes $q$ and $y$, whose signs $-1$ and $+1$ cancel; what is
left is $D_{\mathrm L}$. At $\epsilon=0$ the crossing $y$ is a self-crossing
of the first component, so the mixed crossings of $D_{\mathrm A}$ are the old
mixed ones and their signed count is $2\ell$; then
$w_{\mathrm H}=w_1+w_2+2\ell+2$ and $w_{\mathrm L}=w_1+w_2+2\ell$. At
$\epsilon=1$ the crossing $y$ is itself mixed, so the old mixed crossings count
$2\ell-1$; then $w_{\mathrm H}=w_1+w_2+(2\ell-1)+2$ and
$w_{\mathrm L}=w_1+w_2+2\ell-1$. Both branches are
$w_{\mathrm L}=w_1+w_2+2\ell-\epsilon$. No crossing is \emph{deleted} in this
count, and an earlier revision said the mixed ones were; they are retained in
$D_{\mathrm A}$ and enter through $\ell$.


\emph{(vi).} First suppose $\epsilon=1$, or $\epsilon=0$ with $T$ same-piece.
By~(i), switching the positive $q$ lowers the writhe by two; the subsequent
oriented Reidemeister-II deletion removes crossings of signs $-1$ and $+1$,
so its net writhe change is zero. Thus $w_{\mathrm H}=w_{\mathrm L}+2$.
If $\epsilon=0$ and $T$ is different-piece, Lemma~\ref{lem:chambertransport}, clause~(B)(iii)
gives the same identity. Corollary~\ref{cor:groupedknot}(B), including its stated
empty-carrier case, identifies each grouped carrier writhe $w_{T,C_\nu}$ with
the writhe $w(D_\nu)$ of its grouped diagram, while
Definition~\ref{def:markeddata} binds $w_\nu=w(D_\nu)$. Thus
$w_{T,C_{\mathrm H}}=w_{\mathrm H}$ and $w_{T,C_{\mathrm L}}=w_{\mathrm L}$.
In every case Lemma~\ref{lem:chambercommon}(ii) gives $\rot_{\mathrm H}=\rot_{\mathrm L}$, hence $R_{\mathrm H}=R_{\mathrm L}$. Finally,
\begin{align*}
 d_{\mathrm H}&=1-w_{\mathrm H}-R_{\mathrm H}\\
 &=1-(w_{\mathrm L}+2)-R_{\mathrm L}\\
 &=(1-w_{\mathrm L}-R_{\mathrm L})-2\\
 &=d_{\mathrm L}-2.
\end{align*}

\emph{(vii).} $W_T\neq0$ and Lemma~\ref{lem:selectorid}, clause~(B) at $\epsilon=1$ give
$W_T=-s_{\mathrm{low}}W_1W_2$, with $s_{\mathrm{low}}$ as bound in
Lemma~\ref{lem:selectorid}, clause~(B), so both half selectors are nonzero and both half
carriers $L_1,L_2$ are uniform (Definition~\ref{def:wind}). That is the route
the hypothesis actually supports: it is a statement about the full selector, and
the halves' uniformity is deduced from it. The spectator product $U_T$ is not
used anywhere in this proof, which is why the statement asks only for
$W_T\neq0$. Each is a closed polygon whose principal turns are all
nonzero and, being turns of a generic polygon, of magnitude below $\pi$; by
Corollary~\ref{cor:groupedknot}(B) the grouped object on $L_i$ is a single knot
diagram with all crossings positive, writhe $w_i$, underlying plane curve
$L_i$ and HOMFLY polynomial $Q_i$. Theorem~\ref{thm:carrierfloor}(C) applies to it
in its uniform case --- or Theorem~\ref{thm:carrierfloor}(D), if the half bears no
residual piece at all --- and gives
\[
   \min\deg_a Q_i\;\geq\;1-w_i-R_i=d_i,
\]
hence $\min\deg_af_i\geq d_i$ as well, $f_i$ being a coefficient row of $Q_i$.

Therefore $f_1f_2$ has no support below $D=d_1+d_2$, so
$[a^{D-2}]f_1f_2=0$; and at degree $D$ the only way to write $D$ as a sum of an
exponent of $f_1$ at least $d_1$ and one of $f_2$ at least $d_2$ is
$(d_1,d_2)$, so $[a^{D}]f_1f_2=\omega_1\omega_2$. The bracket is
$0-\omega_1\omega_2+\omega_1\omega_2=0$. If one of the rows is zero then $f_1f_2=0$ and $\omega_1\omega_2=0$, so all
three terms of the bracket are zero.

The restriction to $W_T\neq0$ is not cosmetic. A summand with $W_T=0$ carries no
uniform halves for Theorem~\ref{thm:carrierfloor}(C) to be applied to, and nothing
is asserted about its bracket; it contributes zero to the weighted sum because
its weight is zero. An earlier revision wrote $W_TU_T\neq0$ here and in the
opening line while stating the proposition for $W_T\neq0$; the two stale
occurrences are corrected, and the hypothesis is the weaker one throughout.

\emph{(viii).} Put $L^*=C_{\mathrm H}$ and $q:=x$. Let $D_{\mathrm A}$ be the oriented
smoothing of $D_{\mathrm H}$ at $q$ and put $F_{\mathrm A}=P_{D_{\mathrm A}}$.
By clause~(ii), $D_{\mathrm A}$ has exactly two ordered
components with rows $f_1,f_2$; let $\ell$ be their linking number.

Clause~(i) makes $q$ positive and supplies (T1) and (T2);
clause (vi) gives $d_{\mathrm H}=d_{\mathrm L}-2$ for
Lemma~\ref{lem:fulltwist}, and clause (iv) supplies the auxiliary row.
Corollary~\ref{cor:groupedknot}(B), including its stated empty-carrier case,
identifies the grouped chamber diagrams' underlying curves as
$C_{\mathrm L},C_{\mathrm H}$, so their marked and abstract $d_\nu,F_\nu,\Omega_\nu$ bindings agree. Therefore
\begin{align}
 \Omega_{\mathrm H}-\Omega_{\mathrm L}
 &=\bigl[a^{d_{\mathrm L}-1}z^{-1}\bigr]F_{\mathrm A}\notag\\
 &=\bigl[a^{d_{\mathrm L}+2\ell-2}\bigr]f_1f_2
   -\bigl[a^{d_{\mathrm L}+2\ell}\bigr]f_1f_2. \tag{FT1}
\end{align}
At $\epsilon=1$, Lemmas~\ref{lem:rotshift} and~\ref{lem:chambercommon}(ii),
together with clause~(v), give
\begin{align}
 \rot_{\mathrm L}&=\rot(C_{\mathrm L})=\rot(C_{\mathrm H})
 =\rot(L^*)=\rot_1+\rot_2,\notag\\
 2\ell&=w_{\mathrm L}+1-w_1-w_2. \tag{FT2}
\end{align}
Put $\Delta_R=R_1+R_2-R_{\mathrm L}$. The first line of (FT2) and the triangle
inequality give $\Delta_R\geq0$, while the definitions and its second line give
\begin{align}
 D&=2-w_1-w_2-R_1-R_2\notag\\
  &=2-(w_{\mathrm L}+1-2\ell)-(R_{\mathrm L}+\Delta_R)\notag\\
  &=d_{\mathrm L}+2\ell-\Delta_R. \tag{FT3}
\end{align}
Thus (FT1) becomes
\begin{equation}
 \Omega_{\mathrm H}-\Omega_{\mathrm L}
 =\bigl[a^{D+\Delta_R-2}\bigr]f_1f_2
  -\bigl[a^{D+\Delta_R}\bigr]f_1f_2. \tag{FT4}
\end{equation}

Since $W_T\neq0$, Definition~\ref{def:wind} makes the contact carrier
$L^*=C_{\mathrm H}$ uniform. The uniformity clause of
Lemma~\ref{lem:selectorid}, clause~(B) at $\epsilon=1$ then makes both half carriers
uniform with the same turn sign. Lemma~\ref{lem:uniformrot}(i) gives
$\rot_1,\rot_2$ that same sign, so
$|\rot_1|+|\rot_2|=|\rot_1+\rot_2|$. By (FT2) this says
\begin{equation}
   \Delta_R=R_1+R_2-R_{\mathrm L}=0. \tag{FT5}
\end{equation}
With (FT5), the right side of (FT4) is exactly the first two terms of
clause~(vii)'s displayed bracket. That exported bracket gives
$(\Omega_{\mathrm H}-\Omega_{\mathrm L})+\omega_1\omega_2=0$, which proves
the claim.
\end{proof}

\begin{remark}[attribution]\label{rem:fulltwistsource}
The counterpart's full-twist calculation is retained in the proof of the live
jump: (FT1) is the abstract skein and auxiliary two-component row; (FT2) is the
signed-rotation and writhe ledger; (FT3)--(FT4) are the rotation-defect shift;
and (FT5) is the live-support collapse. The final sentence uses the exported
combined bracket of Lemma~\ref{lem:triplebridge}(vii). The writhe identity
$2\ell=w_{\mathrm L}+1-w_1-w_2$ is also the ledger that refutes the withdrawn
route of Remark~\ref{rem:withdrawnbracket}.
\end{remark}

\subsection{The noninterlacing join}\label{sec:nijoin}

The calculation of §\ref{sec:fulltwist} does transfer to $\epsilon=0$, but not
by copying: one step changes and it changes two of the bookkeeping identities
with it. Throughout, $F_{\mathrm H}$ is the grouped polynomial of the contact
carrier in the two-crossing chamber and $F_{\mathrm L}$ in the zero-crossing
chamber. The symbols $F_{\mathrm A}$ and $\ell$ are \emph{not} bound here:
"the link obtained by smoothing one newborn" does not choose the newborn, and
Lemma~\ref{lem:triplebridge}'s clause~(iii) is not equivariant under exchanging
the two, so the choice is a binding and it is made inside the lemma that reads
it. The chamber data are
$w_{\mathrm L}$, $R_{\mathrm L}$ and $d_{\mathrm L}=1-w_{\mathrm L}-R_{\mathrm
L}$ for the low chamber and likewise with $\mathrm H$; the half data are
$w_i,R_i,d_i=1-w_i-R_i$ and $D=d_1+d_2$, as in
Definition~\ref{def:markeddata}. An earlier revision wrote
$d_{\mathrm L}=1-w_{\mathrm F}-R_{\mathrm L}$, mixing a symbol from a naming
scheme this section no longer uses.

\begin{lemma}[the noninterlacing twist, its ledgers and live rotation defect]\label{lem:nitransport}
\begin{enumerate}
\item[(A)] Here eligibility is as in Definition~\ref{def:eligible}.
Let $\epsilon=0$ and let $T$ be eligible \emph{and same-piece}: after smoothing
$T$ the two newborns lie in one residual piece and on one carrier. Bind here, not by reference: let $q:=x$ be the newborn whose first visit opens
the contact word of Lemma~\ref{lem:contactword} --- the smoothed crossing is
this one and not the other newborn --- let $D_{\mathrm A}$ be the oriented
smoothing of $D_{\mathrm H}$ at $q$ and $F_{\mathrm A}=P_{D_{\mathrm A}}$; by
Lemma~\ref{lem:triplebridge}(ii) at $\epsilon=0$, $D_{\mathrm A}$ has exactly
two components, ordered by the half-label convention, and $\ell$ is their
linking number. The ledgers below are the ones clauses (ii)--(iv) of that
lemma supply. Let $s\in\{\pm1\}$ be the sign
\eqref{eq:S6} in the initial chamber; the
configuration with $s=-1$ is the reflection of the one with $s=+1$, and every
determinant below negates with it, so the statement is written for $s=+1$ and
read for $s=-1$ through that reflection. Then:
\[
   F_{\mathrm H}=a^{-2}F_{\mathrm L}+a^{-1}zF_{\mathrm A},
   \qquad
   \Omega_{\mathrm H}-\Omega_{\mathrm L}=\bigl[a^{d_{\mathrm L}-1}z^{-1}\bigr]F_{\mathrm A},
\]
and, with $f_i,d_i,D$ as in Definition~\ref{def:markeddata} and
$\Delta_R=R_1+R_2-R_{\mathrm L}$,
\[
   2\ell=w_{\mathrm L}-w_1-w_2,
   \qquad
   D=d_{\mathrm L}+2\ell+1-\Delta_R,
   \qquad
   \Omega_{\mathrm H}-\Omega_{\mathrm L}
   =\bigl[a^{D+\Delta_R-3}\bigr]f_1f_2
    -\bigl[a^{D+\Delta_R-1}\bigr]f_1f_2 .
\]
The display is the $\epsilon=0$ specialization, this lemma's own hypothesis ---
an earlier revision printed the unified two-branch exponents
$D+\Delta_R+\epsilon-3$, $D+\Delta_R+\epsilon-1$ over the ledger
$D=d_{\mathrm L}+2\ell+1-\epsilon-\Delta_R$ and called them "valid at both
branches" inside a statement that fixes $\epsilon=0$; the unified form is the
peer's blueprint's and lives in the withdrawn-obstruction remark at the end of
this section, where both branches are in scope. The value of
$\Delta_R$ on a live noninterlacing support is supplied separately, and this
statement carries the exponents in $\Delta_R$ alone rather than the numbers
they become.
\item[(B)] Here eligibility is as in Definition~\ref{def:eligible}.
Let $T$ be eligible at a simple bigon event with $\epsilon=0$, and suppose
$W_T\neq0$ --- the selector of Definition~\ref{def:markeddata}, one number for
both chambers; an earlier revision wrote "the contact carrier's selector",
which names no bound quantity, and left $T$ unbound. Write
\[
   \Delta_R=R_1+R_2-R_{\mathrm L}
\]
for the rotation defect, in the absolute rotations $R_1,R_2$ of the two half
contact carriers and $R_{\mathrm L}$ of the contact carrier in the low chamber, all bound in
Definition~\ref{def:markeddata}. Put $D=d_1+d_2$, with the half slots $d_i$
bound there. Then $\Delta_R=-1$; substituted into the two exponents of clause~(A),
that gives $D-4$ and $D-2$. The symbol $\Delta_R$ is rebound here: clause~(B)
has its own hypotheses and does not inherit clause~(A)'s local binding.
\end{enumerate}
\end{lemma}
\begin{proof}
\emph{Clause (A).} The first two displays are the two clauses of Lemma~\ref{lem:fulltwist}. Its
hypotheses (T1) and (T2) hold here by Lemma~\ref{lem:triplebridge}(i) --- that
lemma proves them at \emph{both} branches, so nothing is imported from
$\epsilon=1$ and the word ``verbatim'', which an earlier revision used here, is
gone. The \emph{second} clause carries a third antecedent,
$d_{\mathrm H}=d_{\mathrm L}-2$, which an earlier revision of this proof did not
name: it is Lemma~\ref{lem:triplebridge}(vi), stated for an eligible $T$ at a simple
bigon event and so available at this branch as well. Without it the coefficient
display is not licensed, the first display being unconditional and the second
not.

The step that differs between the branches is the identification of the two
components of the auxiliary link, and it too is
Lemma~\ref{lem:triplebridge}(ii)--(iii). At $\epsilon=1$ the chords interlace,
smoothing $x$ puts the two visits of $y$ on different components, and $y$ is a
mixed crossing --- \emph{retained}, in neither component's own diagram, and
counted in the linking exponent, which is the typing of
Lemma~\ref{lem:triplebridge}(ii); an earlier revision here said it was to be
deleted. At $\epsilon=0$ the chords are nested: smoothing
$x$ cuts the
traversal into $y_0Ay_1$ and $B$, so both visits of $y$ land on the first
component and $y$ survives there --- as a \emph{positive self Reidemeister-I
curl}. Deleting that curl changes no polynomial (Axiom~\ref{ax:homfly}) and
leaves exactly the two half groupings, so Lemma~\ref{lem:homflyrows}(ii) applies with rows
$f_1,f_2$ and gives $[z^{-1}]F_{\mathrm A}=a^{-2\ell}(a-a^{-1})f_1f_2$.

That deletion is also what changes the writhe ledger. Smoothing one positive
newborn removes one crossing from the two-crossing diagram and deleting the
curl removes one more, so $2\ell=w_{\mathrm L}-w_1-w_2$ --- without the $+1$
that Lemma~\ref{lem:triplebridge}(v) carries at $\epsilon=1$. Then
\[
   D=2-w_1-w_2-R_1-R_2
    =2-(w_{\mathrm L}-2\ell)-(R_{\mathrm L}+\Delta_R)
    =d_{\mathrm L}+2\ell+1-\Delta_R,
\]
and substituting into
$\bigl[a^{d_{\mathrm L}-1}\bigr]\bigl(a^{-2\ell}(a-a^{-1})f_1f_2\bigr)
=[a^{d_{\mathrm L}+2\ell-2}]f_1f_2-[a^{d_{\mathrm L}+2\ell}]f_1f_2$
with $d_{\mathrm L}+2\ell=D+\Delta_R-1$ gives the last display.

\emph{Clause (B).} The signed rotations $\rot_1=\rot(L_1)$ and $\rot_2=\rot(L_2)$ of the two half
contact carriers are bound in Definition~\ref{def:markeddata}, together with
$R_i=|\rot_i|$; write $\rot_{\mathrm L}$ for the signed rotation of the contact
carrier in the low chamber, so that $R_{\mathrm L}=|\rot_{\mathrm L}|$. An
earlier revision rebound these as $r_1,r_2$ here and as $\rho_i$ in the
definition, three names for one quantity. Let $s_{\mathrm{low}}$ be
the low-chamber side sign of Lemma~\ref{lem:rotshift} and $\tau$ the sign of the
turn at $M$ on the contact carrier. In the chart of Lemma~\ref{lem:contactturnboth}, normalised by
$s_{\mathrm{low}}=+1$, the two half contact turns are $\arg q$ and
$\pi-\arg p$, both in $(0,\pi)$; so the common sign of the two half contact
turns is $s_{\mathrm{low}}$, both sides negating together under the reflection
that sends $s_{\mathrm{low}}$ to $-1$. By Lemma~\ref{lem:contactturnboth} at
$\epsilon=0$ that sign and $\tau$ are opposite, $\tau=-s_{\mathrm{low}}$. All
noncontact corners partition between the halves, and
Lemma~\ref{lem:rotshift} at $\epsilon=0$ says the two halves carry one full
turn more than the carrier, of sign $s_{\mathrm{low}}$:
$\rot_1+\rot_2-\rot_{\mathrm L}=s_{\mathrm{low}}=-\tau$.

Suppose the selector is nonzero. Then the contact carrier is uniform
(Definition~\ref{def:wind}), so Lemma~\ref{lem:uniformrot}(i) gives
$\tau\rot_{\mathrm L}\geq1$. Each half inherits only noncontact turns of sign
$\tau$, together with exactly one contact turn of the opposite sign and of
magnitude strictly below $\pi$ (Lemma~\ref{lem:contactturnboth}(iv)). Apply
Lemma~\ref{lem:uniformrot}(ii) directly when $\tau=+1$ and after orientation
reversal when $\tau=-1$; it gives
$\tau\rot_1\geq1$ and $\tau\rot_2\geq1$. The absolute values may therefore all be
removed with the common sign $\tau$:
\[
   \Delta_R=|\rot_1|+|\rot_2|-|\rot_{\mathrm L}|
   =\tau\bigl(\rot_1+\rot_2-\rot_{\mathrm L}\bigr)=\tau(-\tau)=-1 .
\]
\end{proof}

\begin{theorem}[target factorization, dominated and newborn rows, both bigon branches and universal (S7)]\label{thm:s7universal}
If $T\in\Ind(G_{P_0})$ meets $\mathcal N$ then
$N(T)-L(T)+X(T)+Y(T)=0$.
\begin{enumerate}
\item[(A)] Here eligibility is as in Definition~\ref{def:eligible}.
Let $T$ be eligible with parts $T_A,T_B$, and let $W_1,W_2,U_T,\omega_1,\omega_2$
be as in Definition~\ref{def:markeddata}. Then
\[
   W_1W_2\,U_T\,\omega_1\omega_2
   \;=\;
   \Bigl(\wind_{\lambda_1}(T_A)\prod_{L}\Omega_1(T_A,L)\Bigr)
   \Bigl(\wind_{\lambda_2}(T_B)\prod_{L'}\Omega_1(T_B,L')\Bigr),
\]
the first product running over the carriers of $T_A$ in $\lambda_1$ and the
second over the carriers of $T_B$ in $\lambda_2$, as in
Definition~\ref{def:X1};
so that the two factors on the right are the terms of $T_A$ in
$X_1(\lambda_1)$ and of $T_B$ in $X_1(\lambda_2)$.
\item[(B)] Here eligibility is as in Definition~\ref{def:eligible}, and all
marked-data symbols are as in Definition~\ref{def:markeddata}.
Let the event be a simple vertex-on-edge event of bigon type, in the sense of
Convention~\ref{conv:events}, whose two newborn crossings interlace; hence
$\epsilon=1$. Then:
\begin{enumerate}
\item[(i)] For every eligible $T$, write
\[
   W_x=\wind_{P_2}\bigl(T\cup\{x\}\bigr),\qquad
   W_y=\wind_{P_2}\bigl(T\cup\{y\}\bigr)
\]
for the two one-newborn selectors. Then $W_x=W_y=0$, so
$X(T)=Y(T)=0$, whether or not the contact carrier of $T$ is uniform.
\item[(ii)] For every eligible $T$, define its returned contribution and
target share by
\[
   R_T=\delta\bigl(N(T)-L(T)+X(T)+Y(T)\bigr),\qquad
   J_T=s_{\mathrm{init}}\,W_1W_2\,U_T\,\omega_1\omega_2,
\]
where $\delta=\pm1$ is the wall orientation of Definition~\ref{def:branches},
$J$ is the sector target of Definition~\ref{def:sectors} --- an earlier revision
attributed it to Definition~\ref{def:branches} --- and
$s_{\mathrm{init}}=\delta\,s_{\mathrm{low}}$ is the \eqref{eq:S6} sign in the
initial chamber. Then $R_T=J_T$ when $W_T\neq0$; when
$W_T=0$, one has $J_T=0$ and $R_T=\delta(X(T)+Y(T))$; and
\[
   F=J\ \text{ at the interlacing bigon}
   \quad\Longleftrightarrow\quad
   \sum_{\substack{T\text{ eligible}\\ W_T=0}}
   \bigl(X(T)+Y(T)\bigr)=0 .
\]
\item[(iii)] If the two halves are generic polygons --- the hypothesis (S7)
carried in Definition~\ref{def:laws}, discharged at every such event on a
path in $\mathcal R_n$ by the certificate of the induction --- then $X_1$
satisfies (S7) there, with the sign \eqref{eq:S6} and the halves of
Definition~\ref{def:halves}.
\end{enumerate}
\item[(C)] Here eligibility is as in Definition~\ref{def:eligible} and the marked-data symbols are
as in Definition~\ref{def:markeddata}.
Let $\epsilon=0$. For every eligible $T$ with $W_T\neq0$ --- the selector alone,
the spectator product $U_T$ playing no part ---
\[
   \bigl[a^{D-4}\bigr]f_1f_2=\bigl[a^{D-2}\bigr]f_1f_2=0 .
\]
\item[(D)] The following local statements hold. In clauses~(ii)--(v),
eligibility is as in Definition~\ref{def:eligible}; let $\epsilon=0$ and let
$T$ be eligible.
\begin{enumerate}
\item[(i)] Let $S\in\Ind(G_P)$, let $A$ be a uniform carrier of $S$, and let $\{c\}$ be a
singleton residual piece of $S$ carried by $A$. Then
\[
   \min\deg_af_A\;\geq\;\bigl(1-w_{S,A}-R(A)\bigr)+2,
\]
so the factor $\Omega_1(S,A)$ of Definition~\ref{def:X1} is zero.
\item[(ii)] The set $T\cup\{x\}$ is independent and $\{y\}$ is a singleton
residual piece of it; symmetrically, $T\cup\{y\}$ is independent and $\{x\}$
is a singleton residual piece of it. No hypothesis on the pieces of $T$ is
used.
\item[(iii)] $X(T)=Y(T)=0$.
\item[(iv)] If $T$ is different-piece, put $S=T\cup\{x,y\}$. Then:
\begin{enumerate}
\item[(a)] $S$ is independent and $U(S)=U(T)\setminus\{x,y\}$;
\item[(b)] the residual pieces of $S$ are exactly the residual pieces of $T$
other than the two singletons $\{x\}$ and $\{y\}$ --- the same label sets, hence
the same piece diagrams, polynomials and writhes;
\item[(c)] with $L_1,L_2$ the two half contact carriers of
Lemma~\ref{lem:contactsplit}, rows $f_1,f_2$ and grouped writhes $w_1,w_2$, the
high-neither contact carrier satisfies $f_{C_{\mathrm H}}=f_1f_2$ and
$w_{T,C_{\mathrm H}}=w_1+w_2+2$, where $C_{\mathrm H}$ denotes that carrier;
\item[(d)] consequently $f_{C_{\mathrm L}}=f_1f_2$ and $w_{T,C_{\mathrm L}}=w_1+w_2$ in the low
chamber.
\end{enumerate}
\item[(v)] If $T$ is different-piece, the entire row vanishes:
\[
   N(T)=L(T)=X(T)=Y(T)=0 .
\]
\end{enumerate}
\item[(E)] Let the event be a simple vertex-on-edge event of bigon type, in the sense of
Convention~\ref{conv:events}, whose two newborn crossings do not interlace, and
whose two halves are generic polygons. Then $X_1$ satisfies (S7) there, with the
sign \eqref{eq:S6} and the halves of Definition~\ref{def:halves}.
\item[(F)] Let $W$ be a simple vertex-on-edge event in the sense of
Convention~\ref{conv:events}, of either the bigon or the sliding type, with any
number of exterior strands threading the region of the event, and suppose the
two halves of Definition~\ref{def:halves} are generic polygons with fewer
vertices than the parent. Then $X_1$ satisfies (S7) at $W$ with the sign (S6)
and the halves (S1) of Definition~\ref{def:halves}.

The hypothesis on the halves is not a restriction: it holds at every simple
vertex-on-edge event on a path inside $\mathcal R_n$, which is the certificate
proved in the induction's section. The genericity and vertex-count hypotheses are stated here for different reasons, and an earlier
revision gave one reason for both. The genericity is \emph{consumed} by this
proof, which evaluates $X_1$ on each half and cannot do so otherwise. The
vertex count is \emph{exported}: no step below reads it --- this theorem is not
an induction and applies no inductive hypothesis --- and it is carried because
the consumer is, the induction of the last section applying its hypothesis to
the halves and needing them smaller than the parent.
\end{enumerate}
\end{theorem}
\begin{proof}
\emph{The opening assertion.} Let $q\in T\cap\mathcal N$. Then $q$ is adjacent to both newborns
(Lemma~\ref{lem:contactword}(ii)), so neither $T\cup\{x\}$ nor $T\cup\{y\}$ is
independent and $X(T)=Y(T)=0$. Also $x,y\in N_{G_{P_2}}(T)$, so neither is in
$U(T)$: the newborns are dominated and are not residual crossings. Hence the
residual sets, the pieces, and their diagrams are computed from the old labels
alone and coincide in the two chambers, the crossing data being constant across
the event (Lemma~\ref{lem:guardconst}); the carriers are the cycles of the same
smoothing of the same word, with the same corner signs; and the rotations agree
because the two chambers are joined by a path along which no carrier turn
vanishes. So every factor of the term agrees and $N(T)=L(T)$.

\emph{Clause (A).} By Definition~\ref{def:X1} the term of $T_A$ is its selector times the product
of the reads over all of its carriers, and $\wind_{\lambda_1}(T_A)=W_1$ by
Definition~\ref{def:markeddata}; likewise for $T_B$. So the right-hand side is
$W_1W_2$ times the product of the reads over all carriers of $T_A$ in
$\lambda_1$ and of $T_B$ in $\lambda_2$. It remains to identify that product
with $U_T\omega_1\omega_2$, carrier by carrier.

\emph{The contact carriers.} $L_1$ and $L_2$ are the two halves' contact
carriers, and $\Omega_1(T_A,L_1)=\omega_1$, $\Omega_1(T_B,L_2)=\omega_2$ by
Definition~\ref{def:markeddata}, which binds $\omega_i$ as exactly that read.

\emph{The others.} Lemma~\ref{lem:contactsplit} carries the carriers of $T$ in
the two-crossing chamber other than $L^{*}$ bijectively to the carriers of
$T_A$ and $T_B$ other than $L_1$ and $L_2$. A carrier in that correspondence
bears the same residual pieces on both sides: no retained \emph{mixed} label
lies on it, since a mixed label has one visit in each interval and so its
carrier meets both intervals, which by Lemma~\ref{lem:contactsplit} is $L^{*}$
alone; and for the unmixed labels Lemma~\ref{lem:ownermap} identifies the
residual sets, $U(T)\cap V_A=U_{\lambda_1}(T_A)$ and
$U(T)\cap V_B=U_{\lambda_2}(T_B)$, while the induced interlacement graph on the
labels of one interval is the half's own. Equal label sets give equal piece \emph{records}: the map is the identity on the
retained visits, it preserves their cyclic order because the two carriers are
the same closed curve traversed the same way, and it preserves the over/under
designation and the sign at each visit, those being read from the two branch
directions (Definition~\ref{def:piecediagram}). That is a record isomorphism in
the sense of Definition~\ref{def:record}, so Axiom~\ref{ax:gausscode} gives
equal links and hence equal polynomials, and the writhes agree as crossing
counts. An earlier revision passed from equal label sets to equal polynomials
with no isomorphism and no axiom. The carrier being the same closed curve with
the same corners, the rotation agrees too; so the two reads agree, and
their product over the matched carriers is $U_T$, again by
Definition~\ref{def:markeddata}.

Multiplying the three displays gives the identity. An earlier revision asserted
it as part of the ordered-product bijection, which matches \emph{supports} and
says nothing about owners, rows, writhes or rotations.

\emph{Clause (B).} \emph{Clause (i).}
\emph{The collar invariant.} By eligibility and Lemma~\ref{lem:contactword},
both visits of every selected old label lie wholly in $A$ or wholly in $B$. So
no smoothing of $T$ cuts the boundary collar between a newborn smoothing corner
and $M$: in the $x$-branch the smoothing corner at the remote visit $x_1$
reconnects to the position immediately after the incident visit $x_0$, and that
uninterrupted arc reaches $M$ before any selected visit; in the $y$-branch the
uninterrupted incident arc from $M$ reaches $y_0$ before any selected visit and
then closes at the remote visit $y_1$. Hence, for every eligible $T$ including
nonempty ones, the relevant newborn smoothing corner and the turn at $M$ lie on
one and the same daughter carrier.

\emph{The signs, computed in the high chamber.} Work where the two newborns
actually exist. Put $M=(0,0)$, let the remote direction be $r=(1,0)$ on the line
$y=c$ with $c>0$, and write $p-M=(u,h)$ and $q-M=(v,k)$ with $h,k>c$, so that
both neighbours lie beyond the remote line. Interlacing is $u/h<v/k$, i.e.\
$hv-uk>0$. The three determinants that matter are then
\[
   \det\bigl((-u,-h),(v,k)\bigr)=hv-uk>0,
   \qquad
   \det\bigl(r,(-u,-h)\bigr)=-h<0,
   \qquad
   \det\bigl((v,k),r\bigr)=-k<0 .
\]
The first is the turn at $M$, from the incoming direction $-(u,h)$ to the
outgoing direction $(v,k)$: a \emph{left} turn. The second is the corner of the
daughter carrier produced by smoothing $x$, which arrives along $r$ and leaves
along $-(u,h)$: a \emph{right} turn. The third is the corner of the daughter
produced by smoothing $y$, which arrives along $(v,k)$ and leaves along $r$:
again a right turn. By the collar invariant each of those corners shares its
carrier with the turn at $M$, so each of those carriers is mixed, its weight is
$0$ by Definition~\ref{def:wind}, and $W_x=W_y=0$. Reflecting the configuration
negates all three determinants together and leaves every carrier mixed.


\emph{Clause (ii).}
$J_T$ has a second form, used throughout below:
$W_T=-s_{\mathrm{low}}W_1W_2$ by Lemma~\ref{lem:selectorid}, clause~(B), so
$J_T=-\delta\,W_TU_T\,\omega_1\omega_2$. The statement carries the first form
only, that being the one built from the two half terms; an earlier revision put
the derivation inside the statement.

Both quantities are directed: each carries the wall orientation $\delta$, so the
equivalence compares like with like. An earlier revision wrote an undirected row
$N-L-J(T)$ with $J(T)$ never defined, against Definition~\ref{def:branches}'s
directed $R=\delta\sum(N-L+X+Y)$; the display above is the typed replacement.

\emph{If $W_T\neq0$.} By clause~(B)(i) both one-newborn
branches are dead, so $X(T)=Y(T)=0$ and $R_T=\delta(N(T)-L(T))$. Every carrier
other than the contact carrier is common to the two chambers with the same data,
and the selector is the same in both, by
Lemma~\ref{lem:chambercommon}, clauses (i) and (ii); with $W_T$ and $U_T$ bound separately as in
Definition~\ref{def:markeddata},
\[
   N(T)-L(T)=W_TU_T\bigl(\Omega_{\mathrm H}-\Omega_{\mathrm L}\bigr)
   =-W_TU_T\,\omega_1\omega_2
\]
by Lemma~\ref{lem:triplebridge}(viii). So
$R_T=-\delta W_TU_T\omega_1\omega_2=J_T$.

\emph{If $W_T=0$.} Then $J_T=0$ by its own display, and the high-neither and low
terms vanish with their common prefactor
(Lemma~\ref{lem:chambercommon}, clauses (i) and (ii), which are what the
prefactor is made of), leaving $R_T=\delta(X(T)+Y(T))$.

\emph{Summation.} By the opening assertion, supports meeting $\mathcal N$
contribute nothing. That $\sum_TJ_T$ is the global target $J$ needs two
statements and an earlier revision cited only the first.
Lemma~\ref{lem:orderedproduct} matches the eligible supports bijectively with
the pairs $(T_A,T_B)$; clause~(A) matches each product
$W_1W_2U_T\omega_1\omega_2$ with the product of the two half terms. Summing the
second over the bijection of the first gives
$\sum_TW_1W_2U_T\omega_1\omega_2=X_1(\lambda_1)X_1(\lambda_2)$, and
$J=s_{\mathrm{init}}X_1(\lambda_1)X_1(\lambda_2)$ by
Definition~\ref{def:sectors}, which is where the sector target is defined --- the
statement above says so and this proof said otherwise, the statement having been
repaired in an earlier revision and the proof left alone. Only after that identification is the common
nonzero $\delta$ cancelled from $\sum R_T-\sum J_T$, which gives the
equivalence.


\emph{Clause (iii).}
By Proposition~\ref{prop:sectortarget}(F) it suffices that $R=\epsilon J=J$, and by clause~(B)(ii) that is
$\sum_{T\ \mathrm{eligible},\,W_T=0}\bigl(X(T)+Y(T)\bigr)=0$. By clause~(B)(i) every term of that sum is zero.

\emph{Clause (C).} \emph{The two forbidden degrees.} $W_T\neq0$ means the full carrier through
the contact is live, so by Definition~\ref{def:wind} it turns uniformly. By
Lemma~\ref{lem:contactturnboth} in the branch $\epsilon=0$ its contact turn has
the sign of $-s_{\mathrm{low}}$ while each half's contact turn has the sign of
$s_{\mathrm{low}}$ --- the low-chamber sign of Lemma~\ref{lem:selectorid}, clause~(B), which
is what Lemma~\ref{lem:contactturnboth} states and what an earlier revision here
wrote as a bare $s$; and by
Lemma~\ref{lem:contactsplit} each half carrier carries the full carrier's
non-contact corners together with its own contact turn. So each half carrier
has all its turns of one sign except exactly one, of the opposite sign and of
magnitude strictly below $\pi$ (Lemma~\ref{lem:contactturnboth}(iv)). That is the
one-dissent hypothesis of Theorem~\ref{thm:carrierfloor}(C), which applies by
Corollary~\ref{cor:groupedknot}(B) to the grouped knot on each half, or
Theorem~\ref{thm:carrierfloor}(D) if that half bears no piece, and gives
$\min\deg_af_i\geq d_i$. Hence $f_1f_2$ has no support below $D$, and in
particular none at $D-4$ or $D-2$.


\emph{Clause (D).} \emph{Clause (i).} \emph{Pass to the larger support.} Put $S'=S\cup\{c\}$. It is independent: $c$
is residual, so it is adjacent to no member of $S$. Its undominated set is
\[
   U(S')=U(S)\setminus\bigl(\{c\}\cup N_{G_P}(c)\bigr)=U(S)\setminus\{c\},
\]
because $\{c\}$ is a component of $G_P[U(S)]$, so $c$ has no neighbour inside
$U(S)$. Deleting an isolated vertex leaves every other component of a graph
unchanged, so the residual pieces of $S'$ are exactly the residual pieces of $S$
other than $\{c\}$ --- the same label sets, hence by
Definition~\ref{def:piecediagram} the same piece diagrams, the same $P_H$ and
the same $|H|$. Nothing is being compared across a change of diagram here: the
pieces are literally the same objects.

\emph{The two loops are carriers.} Smoothing $c$ cuts the traversal of $A$ at
the two visits of $c$ and reconnects, so the carriers of $S'$ are those of $S$
with $A$ replaced by the two closed curves $\Lambda_1,\Lambda_2$ carrying the
two arcs of $A$ between the visits of $c$; every other carrier of $S$ is
untouched, $c$ lying on $A$. Each residual piece of $S'$ formerly on $A$
therefore lies on $\Lambda_1$ or on $\Lambda_2$, by
Lemma~\ref{lem:carriers}(iv) applied to $S'$; no separate no-straddle
argument is needed and no sub-loop is compared with its parent. Hence
\[
   P_{S,A}=P_{S',\Lambda_1}\,P_{S',\Lambda_2}\cdot P_{\{c\}}
   =P_{S',\Lambda_1}\,P_{S',\Lambda_2},
   \qquad
   w_{S,A}=w_{S',\Lambda_1}+w_{S',\Lambda_2}+1,
\]
the singleton contributing $P_{\{c\}}=1$ and one unit of writhe
(Definition~\ref{def:piecediagram}: a closed curve with one crossing is an
unknot diagram). Each factor is a knot polynomial, hence lies in
$\mathbb Z[a^{\pm1},z^2]$ and has no negative power of $z$ ---
Axiom~\ref{ax:homfly}, whose knot-parity consequence this step reads and which
an earlier revision used without naming --- so the zeroth rows multiply:
$f_A=f_{\Lambda_1}f_{\Lambda_2}$.

\emph{Turns and rotations.} $A$ has no corner at $c$, which is neither a
polygon vertex nor a member of $S$; the two smoothing corners of $\Lambda_1$ and
$\Lambda_2$ at the site of $c$ receive turning angles that are negatives of
each other, so by Lemma~\ref{lem:turnlift}(ii)
$\rot(A)=\rot(\Lambda_1)+\rot(\Lambda_2)$. Their determinants are $\det(u,v)$
and $\det(v,u)$ for the two branch directions at $c$, hence opposite and
nonzero by (G5). Let $\sigma$ be the common sign of $A$'s turns; every corner of
either loop other than the two new ones is a corner of $A$ and has sign
$\sigma$. So one loop is uniform and the other has exactly one dissenting
corner, of magnitude below $\pi$. Lemma~\ref{lem:uniformrot}(i) applies to the
uniform loop. Apply Lemma~\ref{lem:uniformrot}(ii) to the dissenting loop
directly when $\sigma=+1$ and after orientation reversal when $\sigma=-1$.
Thus $\sigma\rot(\Lambda_i)\geq1$ for both $i$, so the two rotations carry the
sign $\sigma$ and $R(A)=R(\Lambda_1)+R(\Lambda_2)$.

\emph{The bound.} Theorem~\ref{thm:carrierfloor}(C) applies to each loop as a
carrier of $S'$ --- through Corollary~\ref{cor:groupedknot}(B) if it bears a piece
and Theorem~\ref{thm:carrierfloor}(D) if it does not, one loop uniform and the other
one-dissent --- so
$\min\deg_af_{\Lambda_i}\geq1-w_{S',\Lambda_i}-R(\Lambda_i)$ when the row is
nonzero, and if either row vanishes then $f_A=0$ and the claim is trivial.
Adding the two bounds and substituting the two displays above,
\[
   \min\deg_af_A\;\geq\;2-\bigl(w_{S,A}-1\bigr)-R(A)
   =\bigl(1-w_{S,A}-R(A)\bigr)+2 .
\]
The factor $\Omega_1(S,A)$ reads the coefficient at $1-w_{S,A}-R(A)$, two
degrees below the row's lowest nonvanishing degree, so it is zero.

\emph{Clause (ii).} $x$ is nonadjacent to $y$ because $\epsilon=0$, and nonadjacent to every element
of $T$ because $T$ is eligible; so $T\cup\{x\}$ is independent. For the same two
reasons $y$ is dominated by no element of $T\cup\{x\}$, so it is undominated and
lies in some residual piece. Let $c$ be an old crossing adjacent to $y$. By
Lemma~\ref{lem:contactword}(ii) $c$ is adjacent to $x$ as well, and $x$ now lies in
the support, so $c$ is dominated and is not residual. Hence $y$ has no residual
neighbour and $\{y\}$ is a component of the residual graph. The statement and
every step of this proof are invariant under interchanging the names $x$ and
$y$, which gives the symmetric assertion.

\emph{Clause (iii).} Take the branch that smooths $x$; the statement and the argument are invariant
under interchanging the two names, which gives the other branch. If
$\wind(T\cup\{x\})=0$ the term is zero by Definition~\ref{def:X1}. Otherwise
every carrier of $T\cup\{x\}$ is uniform. By clause~(D)(ii) $\{y\}$ is a singleton residual piece of
$T\cup\{x\}$; let $A$ be the carrier bearing it
(Lemma~\ref{lem:carriers}(iv)). Then $A$ is uniform and carries a singleton,
so $\Omega_1(T\cup\{x\},A)=0$ by clause~(D)(i), and the term, a
product with that factor, is zero.

\emph{Clause (iv).}
For clause~(D)(iv)(a), at $\epsilon=0$ the newborns do not interlace, so they are nonadjacent, and
each is nonadjacent to every member of $T$ because $T$ is eligible; hence $S$ is
independent. For the undominated set, let $c$ be an old crossing lying in
$U(T)$, and suppose $c\in N_{G_{P_2}}(x)$. By Lemma~\ref{lem:contactword}(ii) $c$ is
then a neighbour of $y$ as well, so $x,c,y$ is a path inside $G_{P_2}[U(T)]$
and the two newborns lie in one residual component of $T$, contradicting the
different-piece hypothesis. So $U(T)$ meets neither $N_{G_{P_2}}(x)$ nor $N_{G_{P_2}}(y)$, and
removing $x,y$ from $U(T)$ removes nothing else:
$U(S)=U(T)\setminus\{x,y\}$.

For clause~(D)(iv)(b), by Lemma~\ref{lem:chambertransport}, clause~(A) $\{x\}$ and $\{y\}$ are components of
$G_{P_2}[U(T)]$, that is, isolated vertices. Deleting isolated vertices leaves
every other component of a graph unchanged, so by clause~(D)(iv)(a) the components of
$G_{P_2}[U(S)]$ are exactly the components of $G_{P_2}[U(T)]$ other than those
two. A residual piece is a label set, and
Definition~\ref{def:piecediagram} builds its diagram from that label set and the
parent polygon alone, so the diagrams, the $P_H$ and the $|H|$ are the same
objects, not corresponding ones.

For clause~(D)(iv)(c), let $(T_A,T_B)$ be the ordered parts of $T$ from
Lemma~\ref{lem:orderedproduct}. Smoothing $x$ and $y$ affects only the carrier
of $T$ that contains them,
which is $C_{\mathrm H}$; every other carrier of $T$ is a carrier of $S$
unchanged, with unchanged pieces. By Proposition~\ref{prop:sectortarget}(A) the carriers of
$S$ in $P_2$ are the carriers of $T_A$ in $\lambda_1$, those of $T_B$ in
$\lambda_2$, and one local carrier $U_{\mathrm{loc}}$ in the contact disc; combining with
Lemma~\ref{lem:contactsplit}, the three of these into which $C_{\mathrm H}$
splits are the two half contact carriers $L_1,L_2$ and $U_{\mathrm{loc}}$. By
Proposition~\ref{prop:sectortarget}(B) $U_{\mathrm{loc}}$ bears no piece, and each of $L_1,L_2$ has
the same pieces, polynomials and grouped writhe as its mate in the half.

Now every residual piece of $T$ carried by $C_{\mathrm H}$ other than $\{x\}$
and $\{y\}$ is, by clause~(D)(iv)(b), a residual piece of $S$, and by
Lemma~\ref{lem:carriers}(iv) applied to $S$ it lies on exactly one carrier of
$S$; that carrier is one of the three just named, and not $U_{\mathrm{loc}}$. So those pieces
distribute between $L_1$ and $L_2$, and
\[
   P_{T,C_{\mathrm H}}
   =P_{\{x\}}\,P_{\{y\}}\prod_{H\text{ on }L_1}P_H\prod_{H\text{ on }L_2}P_H
   =P_{S,L_1}P_{S,L_2},
   \qquad
   w_{T,C_{\mathrm H}}=1+1+w_1+w_2,
\]
using $P_{\{x\}}=P_{\{y\}}=1$ and $|\{x\}|=|\{y\}|=1$ from
Lemma~\ref{lem:chambertransport}, clause~(A). Each factor is a knot polynomial, hence in
$\mathbb Z[a^{\pm1},z^2]$ with no negative power of $z$ by
Axiom~\ref{ax:homfly}, so the zeroth rows multiply.

Clause~(D)(iv)(d) follows from clause~(D)(iv)(c) together with Lemma~\ref{lem:chambertransport}, clause~(B)(iii), which gives
$f_{C_{\mathrm L}}=f_{C_{\mathrm H}}$ and $w_{T,C_{\mathrm L}}=w_{T,C_{\mathrm H}}-2$.


\emph{Clause (v).}
The one-newborn terms are clause~(D)(iii). For the other two,
Lemma~\ref{lem:chambertransport}, clause~(B) gives
$L(T)=W_TU_T[a^{d_0}]f_{C_{\mathrm L}}$ and $N(T)=W_TU_T[a^{d_0-2}]f_{C_{\mathrm L}}$ with no
hypothesis attached, so it suffices to treat the two cases of the common
prefactor.

If $W_TU_T=0$ both terms are zero as products with a zero factor.

So suppose $W_TU_T\neq0$. Then $W_T\neq0$, so every carrier of $T$ is uniform,
in particular the contact carrier. Let $L_1,L_2$ be the two half contact carriers of
Lemma~\ref{lem:contactsplit} --- the names that lemma and
Definition~\ref{def:markeddata} use; an earlier revision wrote
$L_2^{\mathrm h}$ here and then never used the symbol --- with rows $f_1,f_2$, writhes
$w_1,w_2$, rotations $R_1,R_2$, slots $d_i=1-w_i-R_i$ and $D=d_1+d_2$ as in
Definition~\ref{def:markeddata}. Three facts about them are already proved and
are used unchanged:
\begin{itemize}
\item $f_{C_{\mathrm L}}=f_1f_2$ and $w_{T,C_{\mathrm L}}=w_1+w_2$, by
clause~(D)(iv)(d), which is where the passage from the support
$T$ to the support $T\cup\{x,y\}$ --- the support that
clauses~(A) and~(B) of Proposition~\ref{prop:sectortarget} are about --- is
carried out and the old residual set is shown to survive it intact;
\item $R_1+R_2-R(C_{\mathrm L})=-1$, by Proposition~\ref{prop:sectortarget}(C) and
Lemma~\ref{lem:rotshift} at $\epsilon=0$, whose shift is the low-chamber side
sign $s_{\mathrm{low}}$ and whose absolute values may be taken with the common
sign because the contact carrier is live (the computation is that of
Lemma~\ref{lem:nitransport}, clause~(B));
\item $\min\deg_af_i\geq d_i$ for $i=1,2$: each half carrier keeps the uniform
non-contact turns of the live full carrier and acquires exactly one contact
turn of the opposite sign and of magnitude below $\pi$
(Lemma~\ref{lem:contactturnboth}(iv)), which is
the one-dissent hypothesis of Theorem~\ref{thm:carrierfloor}(C), applied through
Corollary~\ref{cor:groupedknot}(B) to the grouped knot on each half or through
Theorem~\ref{thm:carrierfloor}(D) if a half bears no piece.
\end{itemize}
Hence $f_{C_{\mathrm L}}=f_1f_2$ has no support below $D$, while
\[
   d_0=1-w_{T,C_{\mathrm L}}-R(C_{\mathrm L})=1-(w_1+w_2)-(R_1+R_2+1)=d_1+d_2-2=D-2 ,
\]
so $[a^{d_0}]f_{C_{\mathrm L}}=0$ and, two degrees lower still,
$[a^{d_0-2}]f_{C_{\mathrm L}}=[a^{D-4}]f_{C_{\mathrm L}}=0$. Both terms vanish.

\emph{Clause (E).} By Proposition~\ref{prop:sectortarget}(F) it suffices that $R=\epsilon J=0$, i.e.\
that $\sum_T\bigl(N(T)-L(T)+X(T)+Y(T)\bigr)=0$ in the notation of
Definition~\ref{def:branches}, the sum running over eligible $T$ by the
opening assertion. Fix such a $T$ and let $W$ be the
common selector of its high-neither and low supports, common by
Lemma~\ref{lem:chambercommon}(ii).

Split the eligible supports into the two sectors of
Lemma~\ref{lem:chambertransport}, clause~(A).

\emph{Different-piece supports.} Every one of the four terms of such a row is
zero by clause~(D)(v), with no hypothesis on the selector.

\emph{Same-piece supports.} These are the remaining ones, and the rest of this
proof treats them.

The one-newborn terms vanish in either case, by
clause~(D)(iii), which needs no hypothesis on any selector; so
only $N(T)-L(T)$ is at issue, and the split is on the prefactor.

\emph{If $U_T=0$.} Both $N(T)$ and $L(T)$ carry $W_TU_T$ as a factor, so both
vanish and the row is zero. The factorization is read off
Definition~\ref{def:X1} and clauses (i) and (ii) of
Lemma~\ref{lem:chambercommon} --- an earlier revision attributed those two
clauses to Lemma~\ref{lem:chambertransport}, clause~(B), which kept only (iii) when the two
were split --- and not off the latter's closing display: by
Definition~\ref{def:X1} the term of $T$ in either chamber is its selector times
the product of the reads over all of its carriers; by (i) the carriers other
than the contact carrier are the same closed curves with the same pieces in both
chambers, so their reads multiply to the same $U_T$ in both; and by (ii) the
selector is the same $W_T$ in both. What is left in each chamber is the read on
its own contact carrier. The closing display of
Lemma~\ref{lem:chambertransport}, clause~(B) is \emph{not} available here: it is stated
under the different-piece hypothesis of clause (iii), and this sector is the
same-piece one. An earlier revision of this proof cited it anyway, in both
same-piece branches. The case
is written out because the result invoked below is a statement about
coefficients, and multiplying a coefficient conclusion by a vanishing spectator
product is not the same as deducing one.

\emph{If $U_T\neq0$ and $W\neq0$}, by
Lemma~\ref{lem:nitransport}, clauses~(A) and~(B)
\[
   N(T)-L(T)=WU_T\Bigl([a^{D-4}]f_1f_2-[a^{D-2}]f_1f_2\Bigr),
\]
both of whose coefficients are zero by clause~(C), whose
hypothesis is $W\neq0$ and which is therefore invoked inside its own quantifier.

\emph{If $W=0$}, the high-neither and low terms both vanish with their common
prefactor, by the same reading of Definition~\ref{def:X1} with
Lemma~\ref{lem:chambercommon}(ii): the selector is one number $W=W_T$ in both
chambers and multiplies the whole term in each.

Every term is zero, so $R=0$ and $F=J$.

\emph{Clause (F).} The sliding type is Lemma~\ref{lem:slidingcoeff}, clause~(D). The bigon type splits by the
interlacement bit of Definition~\ref{def:sectors}: the interlacing case is
clause~(B)(iii), and the noninterlacing case is
clause~(E). Those three cases are exhaustive and mutually
exclusive: a simple vertex-on-edge event is of exactly one of the two types by
Lemma~\ref{lem:dictionary}, and a bigon event has $\epsilon\in\{0,1\}$.

No bound on the number of crossings, supports or residual pieces occurs in any
of the three arguments, and none of them inspects the region away from the
isolated contact, so arbitrary exterior crossings and arbitrary threading are
included.
\end{proof}

\begin{remark}[the superseded deletion argument]\label{rem:kinkdeletion}
Two earlier revisions --- the one-newborn half of
Theorem~\ref{thm:s7universal}, clause~(C) as first printed, and the different-piece argument
now in Theorem~\ref{thm:s7universal}(D)(v) --- argued by deleting the
surviving positive Reidemeister-I singleton from the affected carrier $A$ and
asserting that the cleaned carrier $C$ satisfies $R_A=R_C+1$, one unit less
absolute rotation; the second said explicitly that the argument did not use the
selector. Both halves of that are wrong. A crossing's sign does not determine
the sign of the change in Whitney index.

Theorem~\ref{thm:s7universal}(D)(i) replaces the deletion. It does not delete
anything: it cuts the carrier at the singleton and applies the floor to each of
the two loops, and it uses the selector, which is what supplies the uniform
turning that makes the two rotations add with one sign. Because
Theorem~\ref{thm:s7universal}(D)(ii) holds in both sectors --- the surviving
newborn is always isolated once the other is in the support, its old
neighbourhood being the other's --- the repair covers the same-piece sector as well, and
Theorem~\ref{thm:s7universal}, clause~(C) no longer carries a one-newborn half at all: it
states the two vanishing degrees, and the one-newborn terms are
Theorem~\ref{thm:s7universal}(D)(iii) wherever they are needed.
\end{remark}

\begin{remark}[what the different-piece sector cost, and what it did not]
\label{rem:diffcensus}
The proof of clause~(D)(v) above establishes vanishing in exactly the degrees $D-2$ and $D-4$
that Theorem~\ref{thm:s7universal}, clause~(C) kills in the same-piece sector, and by
the same floor on the same two halves. The different-piece sector is therefore
not a separate mechanism: what it needed was the typing --- that the two
newborns are singleton kinks (Lemma~\ref{lem:chambertransport}, clause~(A)), that the
chamber transport of the row and the prefactor holds with no selector
hypothesis (Lemma~\ref{lem:chambertransport}, clause~(B)), and that the one-newborn branch
dies by cutting at its surviving singleton rather than by deleting it
(Theorem~\ref{thm:s7universal}(D)(i)).

Two withdrawn attempts are recorded, each failing on an explicit
configuration printed here. The first asserted, as an axiom, that a
different-piece eligible support is selector-dead;
the peer refuted it with a guarded $24$-gon whose contact carrier is live,
with its data --- low word $[0,0,1,1]$, high word
$[0,2,1,1,3,3,2,0]$, edgeless interlacement, $\rot=3$ in both chambers, low slot
$-4$, high slot $-6$. The second proved the low read to vanish by
a three-loop decomposition of the contact carrier; that route is correct but it
needed a diagram identification between a proper sub-loop and its parent that
was not printed formally, and the half split above needs no such bridge, so the
route was withdrawn in favour of machinery already in the document.
\end{remark}

\begin{remark}[what this proof replaces, and whose the repair is]\label{rem:d61bridge}
An earlier revision argued these signs in the chart \emph{at the wall}, where
the two newborns coincide at $M$, and then applied the conclusion to the high
chamber where they do not. The limiting signs were right and the argument was
in the wrong chamber; and it assumed rather than printed the collar invariant
above. Both repairs --- the eligible-$T$ collar invariant and the exact
high-chamber determinants --- are the counterpart's.
\end{remark}

\begin{remark}[why the aggregate turned out to be termwise]\label{rem:aggregateclosed}
An earlier revision recorded this sum as an \emph{aggregate} identity over
selector-dead supports. It is not an aggregate identity: every term vanishes separately, because at an
interlacing bigon the turn at the moving vertex and the smoothing corner of
either newborn have \emph{opposite} signs, so whichever newborn is smoothed,
the carrier that inherits the moving vertex is mixed. The computation is three
lines in the chart of Lemma~\ref{lem:contactturnboth} and it needs neither a
pairing on supports nor a coefficient argument.

The same computation explains why the noninterlacing branch is different and
harder. There $\tau_M=\arg q-\arg p-\pi$ is a \emph{right} turn, of the same
sign as both smoothing corners, so the carriers can be uniform and the
one-newborn branches can be live --- which is exactly why they must be killed by
Theorem~\ref{thm:s7universal}(D)(iii), cutting the affected carrier at the newborn that
survives in the support, and not by any degree argument.
\end{remark}

\begin{remark}[a withdrawn obstruction]\label{rem:whynotni}
An earlier revision of this document printed a remark asserting that the
full-twist calculation \emph{cannot} transfer to $\epsilon=0$, on the ground
that smoothing one newborn leaves the other as a self-crossing so that the two
components are not the half groupings. The observation is right and the
conclusion drawn from it was wrong: the surviving crossing is a positive
Reidemeister-I curl, deleting it changes no polynomial, and what is left is
exactly the two half groupings. What the difference actually costs is one unit
in the writhe ledger --- $2\ell=w_{\mathrm L}-w_1-w_2$ here against
$2\ell=w_0+1-w_1-w_2$ there --- and, with the rotation defect, that is the
whole of why the two branches read the diagonals $D-4,D-2$ and $D-2,D$. In one
formula, over both branches: the ledger is
$D=d_{\mathrm L}+2\ell+1-\epsilon-\Delta_R$ and the jump reads
\[
   \Omega_{\mathrm H}-\Omega_{\mathrm L}
   =\bigl[a^{D+\Delta_R+\epsilon-3}\bigr]f_1f_2
    -\bigl[a^{D+\Delta_R+\epsilon-1}\bigr]f_1f_2 ,
\]
the peer's blueprint's unified form. It is printed here, where both branches
are in scope, and not inside either endpoint lemma, each of which fixes its own
$\epsilon$ and displays its own specialization.
\end{remark}

\input{d7_anchors}
\input{d8_transport}

\input{d9_induction}
% Section 10: the AXIOM LIST.
% RULE FOR THIS SECTION: an axiom entry appears here if and only if some proof
% in this document uses it and does not prove it. Every axiom entry is a closed
% statement with its hypotheses written out and is referenced by number from
% each consumer. The internal theorem stands separately from
% the imported interface.
\section{The axiom list}\label{sec:axioms}

This section contains the complete eight-entry interface of the document: the
statements it uses and does not prove. A formalization of the main theorem must
supply each of them, either by importing a library result or by proving it.
The eight axiom entries fall into two groups --- the three inputs declared by the
campaign, and five theorems taken from the literature. The document also proves
one internal cyclic-classification theorem, Theorem~\ref{thm:mycyclic}, proved in Section~\ref{sec:transport}, which stands separately from that interface.

\subsection{Declared inputs of the campaign}

\begin{axiom}[Hypothesis R for $X_1$]\label{ax:R}
$X_1(P_+)=X_1(P_-)$ across every simple Reidemeister III event, i.e.\ every
event whose zero set is the forced bundle
$Z=\{\mathrm{G3}_{e,f,g},\mathrm{G4}_{e;f,g},\mathrm{G4}_{f;e,g},
\mathrm{G4}_{g;e,f}\}$ for three pairwise remote edges with
$1\leq e<f<g\leq n$, concurrent at $t=0$ at a point interior to all three, the
event being transversal in the sense of Definition~\ref{def:event}.

\emph{Why it is not proved here.} It is not a theorem of the literature and not
a claim of this document: it is a \emph{declared input} of the campaign, the
hypothesis under which the main theorem is stated. Proving it is the open
problem the campaign exists for, and a document that proved it would not need
the entry.
\emph{Used in:} Theorem~\ref{thm:main}(B), and there only --- in the Reidemeister
III case of its induction. No other statement of this document reads it.
\end{axiom}

\begin{remark}[the domain of Hypothesis R, and what the entry said before]
\label{rem:axRhistory}
The index order and the transversality in Axiom~\ref{ax:R} are the antecedent
of (S2) in Definition~\ref{def:laws}, and that entry now carries the same one: a
declared input and the law it declares must have one domain, and an earlier
revision repaired the law's antecedent and left the entry's alone.

\emph{Fidelity: this domain is the specification's own.} The specification asks
$C(P_+)=C(P_-)$ across every Reidemeister III wall, and restricts every one of
its laws to a \emph{simple} event: only the predicate defining the named event
vanishes, all other listed vertex, edge-concurrency and crossing-order
predicates keeping their signs. Write $\mathcal E_{\mathrm{spec}}$ for the
events its (S2) thereby covers and $\mathcal E_{\mathrm{ax}}$ for those this
entry covers. Then $\mathcal E_{\mathrm{ax}}=\mathcal E_{\mathrm{spec}}$.
Containment one way is immediate, an event admitted here being a
Reidemeister-III wall at which nothing outside the forced bundle vanishes. The
other way is the reading of "the predicate defining the event": at a transverse
concurrency the two crossings on each of the three edges collide, so the three
order predicates vanish with it necessarily --- they are dependents of the
concurrency, not independent coincidences --- and a specification that asked
them to keep their signs there would have an empty (S2). The defining condition
is the bundle, exactly as the flat wall's three vanishing members are one
condition; and with that reading the two classes are the same events.

An earlier revision of this paragraph called the entry's domain a strict
subset and defended the narrowing as the safe direction for an assumption. That
was wrong in the direction that flatters: it claimed less than the entry
delivers, and it misread the specification, whose simplicity clause was already
doing the restricting. The correction is the peer's.

Where the law is \emph{read} is narrower than where it is asserted, and that is
a fact about the proof rather than about the entry: its one consumer,
Theorem~\ref{thm:main}(B), meets Reidemeister-III walls only at certified simple-factor hits of the
path supplied by Lemma~\ref{lem:relgp}(B), and it proves there that the zero set is
the forced bundle rather than assuming it here. A sentence in this position previously said that "nothing is lost
downstream by the restriction" and spoke of "the excluded events"; with the
equality above there are no excluded events, and the sentence was left over from
the framing the paragraph before it withdraws.
\emph{Wording changed, hypothesis unchanged.} An earlier revision of this entry
read "a concurrency of three pairwise remote edges together with the
crossing-order predicates it exchanges, all other members of $\mathcal G$
keeping their signs". That domain is \emph{empty}: at a transverse concurrency
the two crossings on each of the three edges collide, so all three order
predicates vanish with the concurrency and cannot be required to keep their
signs. \emph{And the change is not neutral.} A universal law whose domain is empty is
vacuously true, so the earlier print asserted nothing and the present entry,
whose domain is the actual Reidemeister-III zero set, is formally
\textbf{stronger}. What is unchanged is the \emph{intended} input: the
specification's own words, ``$C(P_+)=C(P_-)$ across every Reidemeister III
wall''. That intention was formalized vacuously in the earlier revision and is
formalized on the real domain here. An earlier revision of this paragraph
claimed the entry was ``neither strengthened nor weakened''; that was wrong.
\end{remark}



\begin{remark}[why the parameter is a convention and not a phrase]
\label{rem:lawfulbinder}
An earlier revision wrote the first entry as a closed existential, "there is a
function $A$ which \ldots", and then used the letter $A$ in the second entry and
in the theorem with nothing binding it. Two closed existentials have two
witnesses: the peer's two-object model satisfies each separately and neither
jointly, and the document contains no common binder for them.

A second revision --- also mine --- read $A$ as \emph{universally} quantified
and described that in print as strictly stronger than the existential form. It
is not: the two are \emph{incomparable}. With $X_1=0$ and no function satisfying
both properties, the universal statement is vacuously true and the existential
false; with two such functions, one agreeing with $X_1$ and one not, the
existential-and-agreement is true and the universal false. What the campaign
assumes is neither: its owner \emph{supplies} the quantity, so $A$ is a declared
constant (Convention~\ref{conv:lawful}) of which the two entries below assert
properties. The content assumed is unchanged by either repair; what moved is
where the quantifier sits, and it now has a numbered home a formalizer can point
at.

\emph{What the axiom entries of this section are.} Each axiom entry is an
assumption about objects already declared: about $X_1$, which
Definition~\ref{def:X1} constructs; about the declared constant $A$; or about
the mathematical literature, in which case it is a theorem quoted and not
proved here. No axiom entry declares an object of its own, and none asserts a
bare existence.
Theorem~\ref{thm:mycyclic} stands separately as a result proved in Section~\ref{sec:transport}.
\end{remark}

\begin{axiom}[the lawful quantity]\label{ax:lawful}
With $A$ as in Convention~\ref{conv:lawful}, $A$ is constant on chambers,
unchanged across silent events, and satisfies (S2), (S3), (S4), (S5), (S7) with
the sign (S6), (S8), and (S9) of Definition~\ref{def:laws}, each on the domain
where Definition~\ref{def:laws} states it. Write $\mathrm{Lawful}(A)$ for this
conjunction. In particular
\begin{equation}
  A(\sigma^{k}P)=A(P)\qquad\text{for every generic }P\text{ and every }k\geq0,
  \label{eq:lawful-cyclic}
\end{equation}
by induction on $k$ from \eqref{eq:S9}.

\emph{Why it is not proved here.} It is a declared input, of the same kind as
the rest of the entry: the properties the campaign's owner asserts of the
quantity $A$ it supplies. Nothing in this document computes $A$.

\emph{Stability of the remaining laws under \eqref{eq:sigma}.} The relabeling
permutes the labels and fixes the point set, the plane and its orientation. The
census is stated as equalities evaluated at $\sigma P$, in the direction of
\eqref{eq:pullback}.

\begin{itemize}
\item[(chambers)] The relabeling is a homeomorphism of the generic locus with
  inverse $\sigma^{n-1}$, by Definition~\ref{def:cycshift}, so it carries
  chambers to chambers and the constancy clause transports.
\item[(silence)] Definition~\ref{def:silent} asks of \emph{every} member of the
  zero set $Z$ that it be a predicate $\mathrm{G2}_{e,i}$ whose vanishing at the
  event places $p_i$ on the \emph{line} of $e$ but not on the segment $e$; in
  particular no member of $Z$ is a turn (G1), a concurrency (G3), a
  crossing-order predicate (G4) or a direction determinant (G5). Both halves
  transport. By \eqref{eq:pullback} the $\mathrm{G2}$ expressions satisfy
  $\mathrm{G2}_{e,i}(\sigma P)=\mathrm{G2}_{e^{+},i^{+}}(P)$, and the
  correspondence of index tuples is a bijection of each guard family with
  itself, so $Z(\sigma P)$ consists of $\mathrm{G2}$ members exactly when $Z(P)$
  does and the exclusion of the other four families is preserved. For the
  remaining half, the line of an edge and the segment spanned by it are
  determined by the edge as a point set; $\sigma$ fixes every edge as a point
  set and moves no vertex, so the corresponding member's vanishing places the
  same vertex on the same line and off the same segment. The condition holds at
  $\sigma P$ for every member of $Z(\sigma P)$ exactly when it holds at $P$ for
  every member of $Z(P)$, and silent events correspond to silent events.
\item[(genericity)] The guards are indexed by labels, and evaluation at
  $\sigma P$ reads the guard at the successor tuple:
  $\mathrm{G1}_i(\sigma P)=\mathrm{G1}_{i^{+}}(P)$ and
  $\mathrm{G2}_{e,i}(\sigma P)=\mathrm{G2}_{e^{+},i^{+}}(P)$, and likewise for
  $\mathrm{G3}$, $\mathrm{G4}$ and $\mathrm{G5}$ on their index tuples.
  Remoteness, incidence and the crossing of two segments are conditions on the
  edges and vertices rather than on their labels, and $\sigma$ fixes those point
  sets, so the activation conditions of $\mathrm{G3}$, $\mathrm{G4}$ and
  $\mathrm{G5}$ hold at $\sigma P$ exactly at the tuples corresponding under
  \eqref{eq:pullback}. The generic locus is carried to itself.
% UNSIGNED TYPESETTING (D5-ACC-3): item-heading repair; rival re-verification pending.
\item[(cut)]\emph{Guard values at the cut.} The conditional members carry an
  \emph{ordered representative}: (G5) is indexed by $e<f$, (G3) by $e<f<g$ and
  (G4) by a distinguished edge together with $f<g$, the comparisons being
  between the integer representatives in $\{1,\dots,n\}$, as
  Definition~\ref{def:guarded} prescribes. The relabeling sends the edge with
  representative $i$ at $\sigma P$ to the edge with representative $i^{+}$ at
  $P$, and $i\mapsto i^{+}$ preserves the order of every pair except one: the
  top representative $n$ wraps to $1$. The three identities, each an equality of
  values \emph{evaluated at} $\sigma P$, are therefore
  \begin{align}
    \mathrm{G5}_{e,f}(\sigma P)&=
      \begin{cases}
        \mathrm{G5}_{e^{+},f^{+}}(P), & f<n,\\[2pt]
        -\,\mathrm{G5}_{1,e^{+}}(P), & f=n,
      \end{cases}
    \label{eq:g5-wrap}\\[4pt]
    \mathrm{G4}_{e;f,g}(\sigma P)&=
      \begin{cases}
        \mathrm{G4}_{e^{+};f^{+},g^{+}}(P), & g<n,\\[2pt]
        -\,\mathrm{G4}_{e^{+};1,f^{+}}(P), & g=n,
      \end{cases}
    \label{eq:g4-wrap}\\[4pt]
    \mathrm{G3}_{e,f,g}(\sigma P)&=
      \begin{cases}
        \mathrm{G3}_{e^{+},f^{+},g^{+}}(P), & g<n,\\[2pt]
        \phantom{-\,}\mathrm{G3}_{1,e^{+},f^{+}}(P), & g=n.
      \end{cases}
    \label{eq:g3-wrap}
  \end{align}
  Off the wrap the representatives increase with the labels and the member is
  read unchanged. At the wrap the two determinantal families behave differently,
  and the difference is the parity of the permutation that restores the sorted
  order. In \eqref{eq:g5-wrap} with $f=n$ the pair $(e^{+},f^{+})=(e^{+},1)$ is
  \emph{decreasing}, and $\det(d_f,d_e)=-\det(d_e,d_f)$ gives the sign. In
  \eqref{eq:g4-wrap} with $g=n$ the ordered pair $(f^{+},1)$ is likewise
  decreasing, and the antisymmetry
  $\mathrm{G4}_{e;g,f}=-\mathrm{G4}_{e;f,g}$ of Definition~\ref{def:guarded}
  gives the sign; the distinguished edge carries no order constraint and simply
  becomes $e^{+}$. In \eqref{eq:g3-wrap} with $g=n$ the three representatives
  present themselves as $(e^{+},f^{+},1)$ and the sorted triple is
  $(1,e^{+},f^{+})$, reached by a cyclic permutation of three rows; that
  permutation is \emph{even}, so the $3\times3$ determinant is unchanged and no
  sign appears. Genericity asks that the guard values be nonvanishing, which
  every case above preserves; the two sign changes are changes of the
  representative, not of the geometry.
\item[(S2)] The domain asks that the zero set be the forced bundle
  $\{\mathrm{G3}_{e,f,g},\mathrm{G4}_{e;f,g},\mathrm{G4}_{f;e,g},
  \mathrm{G4}_{g;e,f}\}$ of three pairwise remote edges, concurrent and
  transversal. All four members correspond under \eqref{eq:pullback} to the
  members of the successor triple, remoteness and concurrency are conditions on
  the edges, and transversality is the simultaneous sign change of the bundle,
  which the permutation preserves. The asserted equality $C(P_+)=C(P_-)$ carries
  no label.
\item[(S3), (S4)] The events and the genericity of the wall deletion are
  conditions on the polygon. The statements evaluate $C$ on the deletion, and
  \begin{equation}
    (\sigma P)\setminus j^{-}=
    \begin{cases}
      P\setminus j, & j=1,\\[2pt]
      \sigma\,(P\setminus j), & j\neq1,
    \end{cases}
    \label{eq:deletion-shift}
  \end{equation}
  as configurations. In the case $j=1$ the deleted label crosses the cut and the
  child is the deletion itself; otherwise the child is a relabeling of it, and
  $C$ agrees on the two by \eqref{eq:S9} applied at arity $n-1$. The left and
  right labels of the wall, which belong to (S3), and the loop and no-loop
  labels together with the change $\kappa$ of the Whitney rotation number, which
  belong to (S4), are read from the event's geometry and are carried with it,
  $\sigma$ fixing that geometry. Equation \eqref{eq:deletion-shift} is an
  identity of configurations, and it is at that level that it is used.
\item[(S5)] The empty-spike condition is a condition on the disc bounded by the
  loop-side arc, and the asserted value $0$ carries no label.
\item[(S6), (S7)] The sign $s$ of \eqref{eq:S6} is the sign of
  $\det(p_b-p_a,\,p_M-p_a)$, formed from the oriented remote edge $(a,b)$ and
  the moving vertex $M$ and read in the orientation of the plane; $\sigma$ moves
  the labels $a,b,M$ and fixes the three points and the orientation, so $s$ is
  unchanged. The halves of Definition~\ref{def:halves} are based at $M$ and
  correspond under \eqref{eq:pullback} to the halves of the relabeled polygon,
  so
  \begin{equation}
    C(P_+)-C(P_-)=s\,C(\lambda_1)C(\lambda_2)
    \label{eq:lawful-s7}
  \end{equation}
  transports with the same sign, the values $C(\lambda_i)$ agreeing with those
  of the relabeled halves by \eqref{eq:S9} at the halves' arities.
\item[(S8)] Strict convexity is stated in the cyclic order of the hull boundary,
  and equivalently by the turns having one sign with no three vertices collinear
  and the curve embedded; each condition is invariant under \eqref{eq:sigma}.
  The two asserted values, $1$ for the clockwise boundary and $(-1)^{n}$ for the
  counterclockwise one, are attached to the two orientations of the boundary,
  which $\sigma$ preserves, so each convex polygon keeps the value of its own
  orientation.
\end{itemize}

\emph{Used in:} the consumers of the clauses already present are unchanged ---
Theorem~\ref{thm:zeroanchor}(B) reads (S5), and Theorem~\ref{thm:main}(B) reads
all of them, additionally reading \eqref{eq:S9} at the relabeled endpoint of its
transport step. Those two consume it. Four
remarks --- Remarks~\ref{rem:noncircular}, \ref{rem:notassumed},
\ref{rem:s7wellposed} and \ref{rem:silencenew} --- name it while discussing what
is assumed and what is proved; naming is not consuming.
\end{axiom}

\begin{remark}[the transport read of \eqref{eq:S9}]
\label{rem:lawful-cyclic-use}
The transport step of Theorem~\ref{thm:main}(B) terminates at
$\sigma^{k}Z_{n,r}$, where $Z_{n,r}$ is the anchor selected there and $k$ is the
relabeling supplied by Theorem~\ref{thm:mycyclic}. The step reads
\begin{equation}
  A\bigl(\sigma^{k}Z_{n,r}\bigr)=A\bigl(Z_{n,r}\bigr)
  \label{eq:anchor-shift}
\end{equation}
for that $k$ and that anchor. Equation \eqref{eq:anchor-shift} is read for every
iterate returned by Theorem~\ref{thm:mycyclic} and at every selected anchor,
which is the generality \eqref{eq:lawful-cyclic} supplies.

The declared anchor values depend on the rank alone and \eqref{eq:sigma}
preserves the rank, so anchors related by \eqref{eq:S9} carry equal declared
values. At $(n,r)=(4,0)$ the anchor is $K_0$ and \eqref{eq:anchor-shift} carries
the value $A(K_0)=-1$ of Axiom~\ref{ax:ownerlist} to each $\sigma^{k}K_0$.
\end{remark}

\begin{axiom}[owner normalization on the anchors]\label{ax:ownerlist}
With $A$ as in Convention~\ref{conv:lawful}, $\Cat$ the Catalan numbers of
Definition~\ref{def:catalan} --- the display below reads $\Cat_{|r|}$, the
subscript bound there by $r$ --- and $K_r$ the circular star of
Definition~\ref{def:circularstar} for an integer $r$ with $|r|\geq1$, together
with $K_0$ the four-vertex bow-tie of that definition:
\[
   A(K_0)=-1,\qquad
   A(K_r)=\sgn(r)\,(-1)^{|r|}\,\Cat_{|r|}\quad\text{for every integer }r\neq0 .
\]
The quantifier is written out --- $r$ ranging over the nonzero integers ---
because an earlier revision left it to be read off the display; the Catalan
numbers' own quantifier lives with their definition.
Write $\mathrm{Owner}(A)$ for this pair of identities.

\emph{Why it is not proved here.} A declared input again: the two normalization
values the owner supplies on the anchors. This document computes $X_1$ on those
anchors --- that is Proposition~\ref{prop:anchorvalues}(D) --- but it does not and
cannot compute $A$ there.
\emph{Used in:} Theorem~\ref{thm:main}(B), which reads the anchor values
\emph{directly} --- at the base case $n=3$, where $\Delta(K_{\pm1})=0$ needs
$A(K_{\pm1})$, and at the minimal anchors of the induction step. An earlier
revision routed the use "through the rank-zero anchor of
Theorem~\ref{thm:zeroanchor}(B)", whose proof reads Axiom~\ref{ax:lawful}'s (S5)
clause and not this entry. That is the one consumer. Remarks
\ref{rem:notassumed}, \ref{rem:s7wellposed} and \ref{rem:whichanchors} name it
in discussion, which is not the same as reading it.
\end{axiom}

\begin{remark}[an entry that was added and then struck]\label{rem:diffdeadremoved}
An intermediate revision of this document listed a fourth entry in this group,
asserting that at a noninterlacing bigon wall an eligible support whose two
newborns lie in different residual pieces has a mixed contact carrier. It was
added because the different-piece row had been reduced to a single term whose
prefactor is that selector, with nothing yet fixing that selector's value.

It was false. The peer produced an exact guarded $24$-gon with a live
selector on a different-piece support, recorded with its data in
Remark~\ref{rem:diffcensus}. The entry is struck, and the sector it was
covering is now proved without it in Theorem~\ref{thm:s7universal}(D)(v): the row dies
by the coefficient, not by the selector. The episode is left in the text
because an axiom introduced to cover a gap in a proof is exactly the move
this document is supposed to make visible.
\end{remark}

\subsection{Literature axioms}

\begin{axiom}[HOMFLY--PT]\label{ax:homfly}
There is a unique map $L\mapsto P_L(a,z)\in\mathbb Z[a^{\pm1},z^{\pm1}]$ from
isotopy classes of oriented links in $S^3$ satisfying
$P_{\bigcirc}=1$ and $aP_{L_+}-a^{-1}P_{L_-}=zP_{L_0}$ for every skein triple.
In particular $P_L$ is invariant under the Reidemeister moves. One consequence
is consumed as part of this entry: $P_K\in\mathbb Z[a^{\pm1},z^2]$ for a
\emph{knot} $K$, so that the $z^0$ coefficient of a product of knot
polynomials is the product of their $z^0$ coefficients.

\emph{Why it is not proved here.} The existence of the invariant is a theorem
about all links at once, proved by induction over the full skein tree with a
descending complexity. This document does not construct that invariant; the relative induction in
Lemma~\ref{lem:homflyrows} starts with $P$ already supplied and compares
based-link expressions only.
\emph{Used in:} Definition~\ref{def:X1} (to have $P_H$ at all);
Proposition~\ref{prop:chamberinv}(ii), Lemma~\ref{lem:flatdata} and
Lemma~\ref{lem:silence}; Lemma~\ref{lem:curl} and Theorem~\ref{thm:carrierfloor},
clauses~(C) and~(R); Lemma~\ref{lem:homflyrows},
Proposition~\ref{prop:sectortarget}(B), Lemma~\ref{lem:fulltwist},
Lemma~\ref{lem:triplebridge} and Lemma~\ref{lem:nitransport}, clause~(A); and, through the
knot-parity clause, Theorem~\ref{thm:s7universal}(D)(i) and
Theorem~\ref{thm:s7universal}(D)(iv)(c).
\end{axiom}

\begin{remark}[the trimmed clauses of Axiom~\ref{ax:homfly}]
\label{rem:homflywhole}
This remark is about one entry and two clauses, and it makes no claim about
any other.

\emph{Axiom~\ref{ax:homfly}, clause dropped.} The disjoint-union value
$P_{L\sqcup\bigcirc}=\frac{a-a^{-1}}{z}P_L$ left the entry as a duplicate and
not as a narrowing: it is a consequence of the skein relation and it is
displayed in Definition~\ref{def:homfly} where the normalization is fixed.

\emph{Axiom~\ref{ax:homfly}, clause still dropped.} The general
$z$-order bound for a multi-component link is not restored.
Lemma~\ref{lem:homflyrows}(ii) uses only the retained knot clause: smoothing a
mixed crossing gives a knot, and
$[z^{-2}]P_K=0$ because $P_K\in\mathbb Z[a^{\pm1},z^2]$ has nonnegative even
$z$-powers. It then proves the two-component coefficient identity directly by
mixed switches and split separation. No order bound for a multi-component
polynomial is assumed.

\end{remark}

\begin{axiom}[Etnyre; self-linking from a front]\label{ax:etnyre}
Let $T$ be a front diagram of a transverse knot with no downward vertical
tangency. Then $\operatorname{sl}(T)$ equals the writhe of $T$.

\emph{Why it is not proved here.} The identity is a computation of the
self-linking number from a front, which presupposes the contact-geometric
definition of $\operatorname{sl}$; this document never defines it, and reads
the identity only as a bridge from the diagram's writhe to the transverse
HOMFLY--PT bound of Axiom~\ref{ax:slbound}.
\emph{Used in:} Theorem~\ref{thm:carrierfloor}(C).
\end{axiom}

\begin{axiom}[the HOMFLY--PT bound on the self-linking number]\label{ax:slbound}
Let $T$ be a transverse knot in the standard contact $\mathbb R^3$ and let
$P_T(a,z)$ be the HOMFLY--PT polynomial of the knot it presents, in the
normalization of Axiom~\ref{ax:homfly}. Then
\[
   \operatorname{sl}(T)\;\leq\;-\max\deg_a P_T(a,z)-1 .
\]

\emph{Why it is not proved here.} The bound is a theorem about the whole
invariant, proved in the source by a skein template argument this document has
no machinery for. Source fidelity: the source states the bound for the maximal
self-linking number of a topological knot type and proves it representative by
representative, each front's self-linking number being that of the usual
positive transverse pushoff; the form assumed here, at a single transverse
representative $T$, is that maximum read at one of its values, and that
reading is quoted rather than renamed: every transverse knot is the positive
transverse pushoff of a Legendrian knot, by the Legendrian-pushoff
construction of Etnyre's survey, Section~2.9, and the positive pushoff's
self-linking number is $\operatorname{tb}-r$ of that Legendrian by Lemma~2.22
of the same section.  So $\operatorname{sl}(T)$ is one of the values the
maximum is taken over. The normalization of the source is that of Axiom~\ref{ax:homfly}
verbatim, in the same variable $a$; no variable conversion is assumed.
\emph{Used in:} Theorem~\ref{thm:carrierfloor}(C).
\end{axiom}

\begin{axiom}[Brini--Eynard--Mari\~no; one coefficient of a torus knot]
\label{ax:bem}
For every integer $r\geq2$, with $T(r,2r+1)$ the torus knot of those parameters,
$P$ the HOMFLY--PT polynomial in the normalization of
Definition~\ref{def:homfly}, and $\Cat_r$ the Catalan number of Definition~\ref{def:catalan},
\[
   \bigl[a^{\,1-(2r+1)(r-1)-r}\,z^{0}\bigr]\,P_{T(r,2r+1)}(a,z)
   =(-1)^{r-1}\,\Cat_r .
\]

\emph{Used in:} Proposition~\ref{prop:anchorvalues}(D), and there only.
\emph{Why it is not proved here.} The coefficient is read off a closed
evaluation of the torus-knot polynomials obtained from topological recursion on
a spectral curve; nothing in this document computes a torus-knot polynomial.
Remark~\ref{rem:bemfidelity} carries that evaluation, the variable dictionary
and the derivation of the identity above from it, so the step from the
literature to this entry is on the page.
\end{axiom}

\begin{remark}[the torus-knot entry and its source formula]
\label{rem:bemfidelity}
\emph{The source's statement, in full.} Let $P,Q$ be coprime with
$P\geq Q\geq1$ --- coprime so that $K_{Q,P}$ is a knot, and $P\geq Q$ so that
every factorial below has a nonnegative argument at every $k\leq d$ --- and put
$d=\min(P-1,Q-1)=Q-1$, reading $[0]!=1$. With $q$ the quantum variable,
\[
   [n]=q^{n/2}-q^{-n/2},\qquad [n]!=[n][n-1]\cdots[1],\qquad
   z=[1]=q^{1/2}-q^{-1/2},
\]
and $c=e^{t/2}$ the exponentiated 't~Hooft coupling, the HOMFLY--PT polynomial
of the torus knot $K_{Q,P}$ in the normalization of
Definition~\ref{def:homfly}, written in the campaign variables $a=c^{-1}$ and
$z=q^{1/2}-q^{-1/2}$, satisfies
\begin{equation}
   \frac{P_{K_{Q,P}}(a,z)}{z}
   =c^{(P-1)(Q-1)}\,\frac{1}{[P][Q]}
   \sum_{k=0}^{d}(-1)^k c^{2k}\,
   \frac{[P+Q-k-1]!}{[P-k-1]!\,[Q-k-1]!\,[k]!} .
\label{eq:bemsum}
\end{equation}
That display is the source's own, read at \texttt{arxiv.org/abs/1105.2012}:
quantum integers as displayed, $c=e^{t/2}$, the sum to $d=\min(P-1,Q-1)$, and
the left side normalized by $q^{1/2}-q^{-1/2}$. Axiom~\ref{ax:bem} assumes one
coefficient of it and not the display, and the derivation of that coefficient
from this one is below, so the step from this display to what the entry
assumes is on the page, next to its source.

\emph{Why the entry is the coefficient and not the sum.} Only one number of
\eqref{eq:bemsum} is ever read, and the step that used the rest --- that the
terminal summand has the least $a$-degree in the row --- is needed only to
extract that number \emph{from} the sum. Stating the number directly removes
the step; an earlier revision kept the sum as the entry and defended it by that
step.

For the row the consumer extracts, the bookkeeping of $z$-powers has to be
exact, and an earlier revision of this paragraph was off by one. Since
$[n]=nz+O(z^3)$, a quantum factorial $[m]!$ has $z$-order $m$, so the summand
$[P+Q-k-1]!/\bigl([P-k-1]!\,[Q-k-1]!\,[k]!\bigr)$ has $z$-order
$(P+Q-k-1)-(P-k-1)-(Q-k-1)-k=1$, while $[P][Q]$ has $z$-order $2$: the right
side, and so the left, has $z$-order $-1$. The row the consumer wants is
therefore the coefficient of $z^{-1}$ in $P_{K_{Q,P}}/z$, which is the
coefficient of $z^0$ in $P_{K_{Q,P}}$ itself:
\[
   \bigl[z^{-1}\bigr]\bigl(P_{K_{Q,P}}/z\bigr)
   =\bigl[z^{0}\bigr]P_{K_{Q,P}}
   =\sum_k(-1)^k c^{(P-1)(Q-1)+2k}\,C_k/(PQ),
   \qquad
   C_k=\frac{(P+Q-k-1)!}{(P-k-1)!\,(Q-k-1)!\,k!} .
\]
Writing $[z^0]$ of the divided left side, as that revision did, names a
coefficient one power too high. For $(Q,P)=(2,3)$ this gives
$2c^{2}-c^{4}=2a^{-2}-a^{-4}$, the $z^0$ row of the positive trefoil; for
$(Q,P)=(2,5)$ it gives $3c^{4}-2c^{6}=3a^{-4}-2a^{-6}$, the $z^0$ row of the
$(2,5)$ torus knot. \emph{The derivation of Axiom~\ref{ax:bem} from \eqref{eq:bemsum}, in five
steps.} (1) The $z^0$ row of the polynomial is the displayed sum's constant
term, by the bookkeeping just given. (2) The $k$-th summand carries the campaign
exponent $-\bigl((P-1)(Q-1)+2k\bigr)$ in $a$, since $c=a^{-1}$, and that
exponent is strictly decreasing in $k$; so the summands sit at distinct
$a$-degrees and the terminal one, $k=Q-1$, sits at the least. (3) Its quantum
ratio has classical limit
$C_{Q-1}=\frac{P!}{(P-Q)!\,(Q-1)!}$, so its coefficient is
$(-1)^{Q-1}\frac1Q\binom{P-1}{Q-1}$ at $a^{1-P(Q-1)-Q}$. (4) Specialize to the
pairs the circular stars present, $(Q,P)=(r,2r+1)$ with $r\geq2$: the exponent
is $1-(2r+1)(r-1)-r$ and the coefficient is
$(-1)^{r-1}\frac1r\binom{2r}{r-1}$. (5) That number is $(-1)^{r-1}\Cat_r$,
since $\frac1r\binom{2r}{r-1}$ and $\frac1{r+1}\binom{2r}{r}$ are both
$\frac{(2r)!}{r!\,(r+1)!}$. Steps (1)--(5) are exactly the identity
Axiom~\ref{ax:bem} states, and step (2) --- the least-degree step --- lives
here, where it is used to read one number off a formula, and nowhere in the
document's own chain.
\end{remark}

\begin{remark}[an entry removed]\label{rem:prremoved}
The phase-1.5 manuscript proved the genericity of the circular star by
importing Poonen and Rubinstein's theorem on the diagonals of a regular
$n$-gon. Lemma~\ref{lem:stargeneric}(i) proves it instead from the fact that the
star's edges are tangent to one concentric circle and that at most two tangents
to a circle pass through a point, so that import is gone from this list. Its
removal is the kind of change the axiom list exists to make visible.
\end{remark}

\begin{axiom}[isomorphic records present the same link]\label{ax:gausscode}
Let $D$ and $D'$ be oriented link diagrams, each of whose underlying plane
curves is a connected generic immersed circle in the sphere --- one component,
transverse double points, no triple points --- and suppose there is an
orientation-preserving record isomorphism from the record of $D$ to the record
of $D'$ (Definition~\ref{def:record}). Then $D$ and $D'$ present the same
oriented link.
\emph{Used in:} Lemma~\ref{lem:pieceintrinsic}, Lemma~\ref{lem:curl},
Proposition~\ref{prop:chamberinv}(ii), Lemma~\ref{lem:flatdata},
Lemma~\ref{lem:silence}, Corollary~\ref{cor:groupedknot}(A) and~(B),
Proposition~\ref{prop:sectortarget}(B),
Lemma~\ref{lem:slidingcartesian},
clauses~(A) and~(C), Theorem~\ref{thm:s7universal}, clause~(A) and
Lemma~\ref{lem:triplebridge}. An earlier revision named the first and then said
``every consumer of a grouped carrier row'', which is a description and not a
list.
\emph{Why it is not proved here.} This is the classical statement that a
realizable classical signed Gauss code determines its oriented link. Its content
is Carter's classification of generic immersed curves by their Gauss data: two
such curves in a surface with isomorphic data differ by a homeomorphism of the
surface, \emph{provided the curves fill it}, that is, provided every
complementary region is a disc. The hypothesis above is the case in which that
proviso is automatic and is stated for that reason rather than dropped: the
underlying curve of each diagram is a connected four-valent graph embedded in
the sphere, and every face of a connected graph embedded in the sphere is a
disc, the crossing-free case being a circle with two disc faces. The resulting
homeomorphism carries diagram to diagram and so link to link, a priori only up
to reflection if it reverses orientation; the crossing signs, which the record
isomorphism preserves, exclude that. This document does not develop the
realizability theory, and the consumers need the conclusion for diagrams built
on different curves, whose traversal circles are compared by an explicit record
isomorphism rather than identified.
\emph{Two earlier revisions.} The first did not list this entry at all: the
identification it provides was asserted inside the proof of
Corollary~\ref{cor:groupedknot}(B) in one sentence, and the peer refuted the
sentence --- smoothing the support changes the traversal successor, so retaining
the same labels does not identify the cyclic signed records. The second listed
the entry but hypothesized that the two records were \emph{equal}, which
presupposes a shared traversal circle that these two diagrams do not have; the
peer refused that typing too. The combinatorial content is proved ---
Lemma~\ref{lem:carrierword} for the order, Lemma~\ref{lem:pieceintrinsic} for
the map --- and what remains, the passage from a record isomorphism to an
equality of links, is this entry.
\end{axiom}

\begin{remark}[an entry removed, the second]\label{rem:arcremoved}
An earlier revision carried an axiom asserting that two embedded regular arcs
in a disc with the same endpoints and endpoint tangents are regularly
homotopic. It was used at two places and the peer broke the first: for the
empty support the carrier meets the contact disc in two arcs, not one, and they
cross. Both places now argue instead that the deformation inside the contact
disc is a regular homotopy of a polygonal immersion, because no corner
determinant vanishes along it --- which is how the counterpart's own packs
argue it. The entry is gone from this list.
\end{remark}

\subsection{Citations}\label{sec:axcitations}

The five axiom entries in this subsection are statements taken from the
literature. Their sources are:
\begin{itemize}
\item Axiom~\ref{ax:homfly}: W.~B.~R. Lickorish and K.~C. Millett,
\emph{A polynomial invariant of oriented links}, Topology \textbf{26}
(1987), 107--141, Section~1, especially p.~120, for the complete proof of
existence and uniqueness; and S.~Chmutov and M.~Polyak, \emph{Elementary
combinatorics of the HOMFLYPT polynomial}, Int. Math. Res. Not. IMRN 2010,
no.~3, 480--495, Remarks~1--2 on p.~3 and the proof in Section~5, for the
retained knot-parity clause in the exact normalization of
Axiom~\ref{ax:homfly}.
\item Axiom~\ref{ax:etnyre}: H.~Geiges, \emph{An Introduction to Contact
Topology}, Cambridge Studies in Advanced Mathematics \textbf{109}, Cambridge
University Press (2008), Proposition~3.5.32, p.~127, with printed proof; see
also J.~B. Etnyre, \emph{Legendrian and transversal knots},
\texttt{arXiv:math/0306256}, in \emph{Handbook of Knot Theory}, Elsevier
(2005), 105--185, Section~2.6.4, Equation~(9).
\item Axiom~\ref{ax:slbound}: L. Ng, \emph{A skein approach to Bennequin type
inequalities}, Int. Math. Res. Not. (text at
\texttt{sites.math.duke.edu/\~{}ng/math/papers/bound.pdf}), stated in its list
of Bennequin type inequalities as the HOMFLY--PT bound and proved by its
Theorem~1 (skein template), extended to oriented links by its Corollary~1;
attributed there to J. Franks, R. F. Williams, \emph{Braids and the Jones
polynomial}, Trans. Amer. Math. Soc. \textbf{303} (1987), 97--108, and
H. R. Morton, \emph{Seifert circles and knot polynomials}, Math. Proc.
Cambridge Philos. Soc. \textbf{99} (1986), 107--109; the reading of that bound
at a single transverse representative uses J.~B.~Etnyre, \emph{Legendrian and
transversal knots}, preprint \texttt{arXiv:math/0306256}, Section~2.9 --- the
Legendrian pushoff of a transverse knot, and Lemma~2.22,
$\operatorname{sl}(T_\pm(L))=\operatorname{tb}(L)\mp r(L)$. For the classical
origin of the inequality, rather than the HOMFLY--PT refinement, see
D.~Bennequin, \emph{Entrelacements et \'equations de Pfaff}, Ast\'erisque
\textbf{107--108} (1983), 87--161, Th\'eor\`eme~3, p.~96.
\item Axiom~\ref{ax:bem}: A. Brini, B. Eynard, M. Mari\~no, \emph{Torus knots
and mirror symmetry}, Ann. Henri Poincar\'e \textbf{13} (2012), 1873--1910;
preprint \texttt{arXiv:1105.2012}, the finite closed formula of Section~3 with
the quantum integers of its Equation~(3.40). Read at the preprint on
2026-08-23. The proof source for the coefficient is E.~Gorsky,
\emph{$q,t$-Catalan numbers and knot homology},
\texttt{arXiv:1003.0916v3}, Theorem~3.1 and its Appendix. Its torus-knot
HOMFLY input is V.~F.~R. Jones, \emph{Hecke algebra representations of braid
groups and link polynomials}, Ann. of Math. \textbf{126} (1987), 335--388,
Theorem~9.7, p.~359, whose proof is printed. The rewriting uses N.~M.
Dunfield, S.~Gukov and J.~Rasmussen, \emph{The superpolynomial for knot
homologies}, Experiment. Math. \textbf{15} (2006), 129--159,
Equations~(69) and~(71), p.~25; the equality of those displays, asserted there
without the intervening algebra, has been verified term by term in a symbolic
derivation on file.
\item Axiom~\ref{ax:gausscode}: J.~S.~Carter, \emph{Classifying immersed
curves}, Proc.\ Amer.\ Math.\ Soc.\ \textbf{111} (1991) 281--287, which
classifies generic immersed curves in the plane by their Gauss data up to
homeomorphism of the sphere; the use of Gauss codes as a complete presentation
of a classical link, with the classical codes exactly the realizable ones, is
the framing of L.~H.~Kauffman, \emph{Virtual knot theory}, European J.\
Combin.\ \textbf{20} (1999) 663--690. For diagrams satisfying the paper's
standing Rule~1, C.~H. Dowker and M.~B. Thistlethwaite, \emph{Classification
of knot projections}, Topology Appl. \textbf{16} (1983), 19--31,
Corollary~1.2, p.~24, supplies the orientation refinement; this citation is
only for that prime/reduced subclass. The universal signed-oriented
Gauss-data route uses M.~Goussarov, M.~Polyak and O.~Viro, \emph{Finite type
invariants of classical and virtual knots}, Topology \textbf{39} (2000),
1045--1068, Theorem~1.A, p.~4, together with G.~Kuperberg, \emph{What is a
virtual link?}, Algebr. Geom. Topol. \textbf{3} (2003), 587--591, Theorem~1,
p.~588, and the oriented-link conclusion on p.~591. For the isotopy step see
S.~Smale, \emph{Diffeomorphisms of the 2-sphere}, Proc. Amer. Math. Soc.
\textbf{10} (1959), 621--626, Theorem~6; the topological case is attributed
there to H.~Kneser, \emph{Die Deformationss\"atze der einfach
zusammenh\"angenden Fl\"achen}, Math. Z. \textbf{25} (1926), 362--372.
\end{itemize}


\end{document}
